\quellseite{692}
\part{IV. Auslegungen über das fünfte Buch Mosis}

\chapter[A. Luthers Anmerkungen zu dem fünften Buche Mosis]{A. Luthers Anmerkungen zu dem fünften Buche Mosis.\footnote{Diese Schrift ist eine Sammlung der Vorträge über das Deuteronomium, welche Luther seinen Klosterbrüdern im Augustinerkloster hielt vom 23. Februar 1523 an (Buchwald, Andreas Poachs Sammlung ungedruckter Predigten D. M. Luthers, S. XVI) bis in die ersten Monate des Jahres 1525 (Rößlin, Martin Luther, 3. Aufl., Bd. I, S. 615); herausgegeben wurde sie, versehen mit einer Widmung Luthers an Georg von Polenz, Bischof von Samland, der dem Evangelio hold war und für dessen Ausbreitung Sorge trug, spätestens im April 1525, denn schon im Mai erschien der erste Nachdruck in Straßburg. Die erste Ausgabe erschien unter dem Titel: Deuteronomion Mose cum annotationibus. Mart. Luther. Wittembergae. 1525. Nach der Widmung an Polenz folgt ein anderer Titel: Deuteronomion Mose ex Ebraeo castigatus cum annotat. Martini Lutheri. Wittembergae 1524. Im Eink.: Wittembergae, apud Johannem Lufft. Anno Domini 1525. Noch in demselben Jahre erschienen Nachdrucke in Straßburg bei Johannes Knobloch (Mense Majo) und zu Basel bei Adam Petri (Calendis Juliis). Auch eine deutsche Uebersetzung kam schon in demselben Jahre zu Straßburg heraus (ohne Angabe des Druckers) unter dem Titel: „Verzeichnungen über das fünfft Buch Mose, Deuteronomium genant“ etc. In den Sammlungen lateinisch: Wittenberger (1549), Tom. III, fol. 14b, ohne die Widmung; Jenauer (1603), Tom. III, fol. 74b und Erlanger exeget. opp., Tom. XIII, p. 1. Die Anordnung des Textes in der Wittenberger und der Erlanger ist so, daß erst der nach dem Hebräischen verbesserte lateinische Text des ganzen fünfften Buches Mosis vorgeht, und dann die Anmerkungen folgen; in der Jenauer dagegen schließen sich die Anmerkungen an den Text jedes einzelnen Capitels. Die Widmung findet sich auch noch lateinisch bei Aurifaber, Bd. II, Bl. 285 und bei De Wette, Bd. II, S. 647. In deutscher Uebersetzung findet sich die Schrift in der Wittenberger (1556), Bd. VIII, Bl. 149, ohne die Widmung; in der Altenburger, Bd. V, S. 1031; in der Leipziger, Bd. IV, 92 und bei Walch. In den deutschen Ausgaben ist natürlich der Abdruck des Textes des ganzen fünfften Buches Mosis weggelassen. Wir geben eine neue Uebersetzung dieser Schrift nach der Erlanger Ausgabe unter Vergleichung der Wittenberger und der Jenauer.}}

\begin{center}
Seinen Klosterbrüdern vorgetragen vom 23. Februar 1523 bis zu Anfang 1525. Herausgegeben vor Mai 1525.

Aus dem Lateinischen neu übersetzt.
\end{center}

\section{[Luthers Widmungsschreiben.]}

Dem ehrwürdigen Vater und Herrn in Christo, dem Herrn Georg von Polenz, dem rechtschaffenen Bischofe der Kirche zu Samland, seinem in dem HErrn Christo hochzuverehrenden Vater Gnade und Frieden.

Ich habe mir vorgenommen, hochverehrter Vater, das fünfte Buch Mosis, welches ich für meine Brüder in zwanglosem Gespräche zu Hause behandelt habe, und welches auf ihre Bitte öffentlich ausgehen soll, deinem Namen zu widmen und es deiner Ehrwürd(en) darzubieten. Dies habe ich schon lange im Sinne gehabt, bin aber durch Geschäfte verhindert worden; jetzt endlich thue ich es auch durch diesen Brief, da Gott es gnädiglich gestattet hat; nicht der Meinung, als ob alles das, was meine Wenigkeit an der vornehmsten und hauptsächlichsten Quelle (auctore) aller heiligen Bücher versucht hat, eines so großen Mannes würdig wäre, denn das Ansehen und die Majestät derselben sollte mit Recht mich und einen jeglichen Menschen von der Auslegung abschrecken, wie geschrieben steht [Ps. 50, 16. Vulg.]: „Zum Sünder hat GOtt gesagt: Warum erklärst du meine Rechte?“ Es hat mich aber die Noth und die Rücksicht auf die Seligkeit gezwungen, daß ich mich dessen unterstanden habe, um Gottseligkeit und Gottesdienst (religione) zu lernen und zu lehren.

Aber um endlich einmal meines Herzens Gesinnung gegen dich an den Tag zu geben, habe ich diese Gelegenheit ergriffen, damit ich auch wenigstens in etwas zu denen gehören möchte, welche Zeugniß ablegen von deinem aufrichtigen Glauben an Christum und von deiner thätigen Liebe gegen seine Brüder. So glauben wir, wenn etwa der HErr geruhen sollte, solches zu
\quellseite{693}
geben, daß, durch dein so neues und außerordentliches (singulari) Beispiel angezeigt, auch etliche andere Bischöfe und Fürsten die Kirche GOttes mit einem neuen Wunder erfreuen möchten, die Satzungen und menschlichen Meinungen verwerfen, und das reine Wort GOttes annehmen, ehren und ausbreiten, wie es ihr Beruf und Amt erfordert. Denn, indem wir die Gaben GOttes an dir loben, loben wir nicht dich, sondern preisen das herrliche Wunder der Gnade Gottes, welche, wie wir sehen, in dir stark ist, herrscht und triumphirt, und hören es mit Freuden, daß, wie Amos redet [Cap. 3, 12.]: „Gleichwie ein Hirte dem Wolfe zwei Kniee oder ein Ohrläpplein aus dem Maule reißt“, so der HErr dich einzig und allein unter allen Bischöfen auf Erden erwählt und befreit hat aus dem Rachen des Teufels, den er weit aufgesperrt hat wie die Hölle, und alle verschlingt. Denn wir sehen an den anderen Bischöfen (wiewohl ich hoffe, daß noch etliche Leute wie Nicodemus unter ihnen sind) nichts Anderes, als daß sie, nachdem der Kaiser und Könige und Fürsten verkehrt sind (subversis), wüthen und rasen wider das wiederkehrende oder vielmehr aufgehende Evangelium, auf daß sie von neuem das Wort des zweiten Psalms [V. 2.] erfüllen: „Die Könige im Lande lehnen sich auf, und die Herren rathschlagen mit einander wider den HErrn und seinen Gesalbten“, wodurch sie zugleich veranlassen, daß auch GOtt an ihnen das erfülle, was da folgt [V. 4.]: „Aber der im Himmel wohnet, lachet ihrer, und der HErr spottet ihrer. Er wird einst mit ihnen reden in seinem Zorn, und mit seinem Grimm wird er sie schrecken.“

Dir aber ist die ganz besondere und wunderbare Gnade geschenkt, daß du nicht allein das Wort angenommen hast und geglaubt, sondern dasselbe auch frei öffentlich bekannt, gelehrt und nach deinem bischöflichen Ansehen Sorge dafür getragen hast, daß es in deinem ganzen Sprengel (dioecesin) gelehrt werde, hast auch diejenigen, welche am Worte arbeiten, reichlich versorgt: da die Gottlosen dieses sehen, zürnen sie und knirschen sie mit den Zähnen und härmen sich vergeblich darüber ab. Diese Gnade nun, welche GOtt dir erwiesen hat, durfte ich nicht mit Schweigen übergehen, sondern mußte sie, wie es Paulus [1 Cor. 8, 1. ff.] mit den Macedoniern gegen die Corinthier thut, auf jegliche Weise ausbreiten, preisen und davon singen zur Ehre GOttes, zur Ausbreitung des Evangelii, zum Wachsthum des Glaubens, zum Trost der Schwachen und derer, welche bei uns an verschiedenen Orten mancherlei leiden, zur Aufrichtung der Kleinmüthigen; endlich auch zum Schrecken und zur Schande der Feinde des Worts und der Götzen, die an der Bischöfe Statt sitzen, die nichts Anderes wissen noch können, als daß sie Tyrannei üben an dem elenden Volke. Ich habe mich auch davor nicht gescheut, daß ich dich durch dieses Lob belade mit dem Haß und Neid der anderen Bischöfe, Könige und Fürsten, oder dich gar in Gefahr des Lebens bringe, da ja heutzutage keine Schandthat so groß ist als dies Eine, und dem höchsten Schänden des Heiligthums gleichgeachtet wird, nämlich daß man das Evangelium Gottes bekennt, und ihm dessenwillen schon viele mit der größten Beständigkeit ihr Blut vergossen haben. Unter diesen leuchten besonders hervor Heinrich und Johannes, die von den Löwenschen Sophisten zu Brüssel verbrannt worden sind, und nach ihnen unser Heinrich von Zütphen, der Evangelist von Bremen, der durch die Ditmarschen Bestien nicht bloß mit Einem Feuer oder Einer Art des Todes getödtet worden ist. Was unser noch warte, das weiß GOtt, wenn wir anders würdig sein sollten, die Genossen jener zu werden, und auch Schmach zu leiden um des Namens des HErrn willen. Ich sage, ich habe mich nicht gescheut, dich dieser Gefahr auszusetzen, da ich ganz fest davon überzeugt bin, daß der, welcher dich geehrt hat mit dem Worte des Kreuzes, dich auch stärken wird mit dem Geiste des Kreuzes, und auch in deinem Fleische triumphiren wird, nicht allein über die Lästerungen deines Namens und den Haß und das Wüthen der Mächtigen in der Welt, sondern auch über die Kräfte und Gewalt (so es vonnöthen sein sollte, dieselben zu erleiden) des Satans mit seinem ganzen Reiche. Denn wir wissen, daß, der in uns ist, größer ist, denn der in der Welt ist [1 Joh. 4, 4.].

Ueberdies, damit diese Gnade und Gabe GOttes noch überschwänglicher wäre, hat der HErr noch dieses hinzugefügt und durch seine Güte verschafft, daß du einen Landesfürsten hast, der das Evangelium von Herzen liebt, jenen berühmten Helden Adelbert, aus dem Stamme der Markgrafen zu Brandenburg, den Großmeister
\quellseite{694}
Preußens, welchem in demselben Geiste gegeben ist, um mit den Worten des Propheten zu reden [Jes. 32, 8.] fürstliche Gedanken zu haben. So geschieht es, daß unter eurer beider christlichem Schutz das Volk Preußens, welches das Evangelium vielleicht nie anders als verdunkelt oder verfälscht gehabt hat, es nun, da es weithin läuft und Frucht schaffet, ganz rein und klar hört durch die unaussprechliche Wohlthat GOttes. Denn ich glaube, daß bis zu dieser Zeit das Evangelium nach Deutschland niemals ganz klar und in vollem Lichte leuchtend gekommen sei, wiewohl ich nicht zweifle, daß der Herr immer gesucht hat, die Seinen überall zu erretten, auch in der Mitte von Sodom und Gomorra. Und siehe das wunderliche Geschick: nach Preußen läuft das Evangelium in vollem Laufe und mit vollen Segeln, wohin es nicht gerufen und wo es auch nicht gesucht wurde, aber in Ober- und Niederdeutschland, wohin es aus freien Stücken kam und anlangte, wird es mit aller Muth und Unfrömmigkeit gelästert, zurückgestoßen und vertrieben, so daß man hierin das bemerkenswerte (egregiam) Geschick des Evangelii erfüllt sehen kann, von dem Paulus Röm. 10, 20. f. [Jes. 65, 1. 2.] sagt: „Ich bin erfunden von denen, die mich nicht gesucht haben, und bin erschienen denen, die nicht nach mir gefragt haben.“ Von Israel aber sagt er: „Den ganzen Tag habe ich meine Hände ausgestreckt zu dem Volk, das sich nicht sagen läßt und widerspricht, welches auf Wegen wandelt, die nicht gut sind.“

Doch ich will zu dir zurückkehren. Laß daher zu, lieber Bischof (praesul), daß dies Büchlein dir zugeschrieben werde, damit durch dich und in dir bei dieser Gelegenheit Christus gepreiset werde, zugleich auch unter deinem Namen denen gedient werde, welche etwa vorher das fünfte Buch Mosis nicht so wohl verstanden haben, als es mir gegeben ist. Denn es sind viele, und unter diesen etliche, welche sich haben dünken lassen, daß sie gewiß die Meister von allen wären, welche Mosen und das ganze alte Testament für gar gering achten, da sie meinen (velut), man solle mit dem Evangelio zufrieden sein; von deren Meinung soll ein jeder Christenmensch sich völlig fern halten. Denn es ist gewiß, daß, gleichwie die Weisen der Welt sagen, Homer sei der Vater aller Dichter, die Quelle, ja, ein Meer aller Gelehrsamkeit und Weisheit und Beredsamkeit, so unser Moses die Quelle und der Vater aller Propheten und heiligen Bücher ist, das heißt, der himmlischen Weisheit und Beredsamkeit. Da aber zu dieser Zeit (hoc saeculo) alles anfängt, wiederhergestellt zu werden, als ob der Tag der Wiederbringung aller Dinge bevorstehe, so bin ich darauf gekommen, daß ich dies unternehme, ob nicht vielleicht auch Mosis wiederhergestellt werden könnte, und ich die Büchlein wieder zu der Quelle zurückleiten könnte. Denn es ist sehr lieblich und, sowohl für das Verständnis als auch um das Gedächtnis zu unterstützen, sehr bequem und nützlich, wenn man die Fußtapfen der späteren Propheten in Mose sieht, wie sie ihn gelesen, gelernt, gelehrt und Tag und Nacht in Händen gehabt haben. Aus seinem Füllhorn (copiae cornu) sind alle ihre Schätze genommen, so daß sie auch selbst von ihm sagen können: Von seiner Fülle und Ueberfluß haben wir alle empfangen, wie er es selbst zuvor gesehen und vorhergesagt hat, daß es geschehen werde, in seinem Gesange [5 Mos. 32, 2.]: „Meine Lehre triefe wie der Regen, und meine Rede fließe wie Thau“ etc.

Ich habe mir aber Mühe gegeben, daß ich erstlich alles auf das einfältigste behandelte, und habe mich nicht zu den (sogenannten) heimlichen Auslegungen hinreißen lassen, wenn bisweilen, wie es etlichen scheint, ungereimte und thörichte Gesetze vorkamen. Denn erstlich lehrt Moses die rechte Göttlichkeit, indem er gar ausführlich und reichlich den Glauben predigt, mit angehängten sehr schönen Ceremonien, durch welche das unverständige Volk gefangen und gehalten werden muß, damit es nicht eigene erdichte, welche GOtt haßt. Sodann hat er damit zu schaffen, daß das bürgerliche Regiment geordnet, und die Liebe unter einander erhalten werde, indem er alles mit sehr passenden und billigen Gesetzen einrichtet und darüber Verfügungen trifft. Da ist nichts Thörichtes oder Nichtiges, sondern alles ist nothwendig und nützlich, was der leicht verstehen wird, welcher weiß, was es auf sich habe, das obrigkeitliche Amt zu bekleiden bei einem Volke, das sehr ungebunden (paulo liberiore) und in seinen eigenen Augen weise ist, wie dieses jüdische Volk war, wo die Noth nicht selten zwingt, solche Dinge zu ordnen und zu thun und zu gestatten, welche sonst mit gutem Grunde verlacht und verspottet würden.

Ich glaube auch nicht, daß GOtt aus einer
\quellseite{695}
andern Ursache dies Volk durch solche Ceremonien habe bilden wollen, als weil er sah, daß der große Haufe der Menschen durch diese Larven und äußeren Schein am meisten bewegt und eingenommen werde. Damit es nun nicht leere Larven und ein bloßes Schauspiel wäre, hat er sein Wort hinzugethan, gleichsam als einen gewichtigen Inhalt und Wesen der Larven, dadurch sie ernst und bedeutend (graves) würden, und sie [die Juden] wüßten, daß Gott selbst das geredet, was sie [die Ceremonien] enthielten, damit nicht, wenn sie ohne das Wort selbst andere aufrichteten, bloße Spielerei und Possen daraus würden. Derartig sind unsere Ceremonien im Reiche des Papsts, die nach dem Exempel der Mosaischen erdacht sind. Weil diese keinen Befehl in GOttes Worte haben, wodurch sie ernst und wichtig werden, sind sie nichts Anderes als ganz bedeutungslose und lächerliche Nachäfferei der Juden.

Ich habe aber hernach auch kurze Allegorieen hinzugefügt, fast bei jedem einzelnen Capitel, nicht als ob mir viel daran liege, sondern um dem thörichten Bestreben etlicher zuvorzukommen, die Lust haben mit Allegorieen umzugehen, wie wir sehen bei Hieronymus, Origenes und anderen alten Schriftstellern, daß sie eine gar wenig glückliche und heilsame Weise hatten, Allegorieen zu machen, da sie alles auf das Leben und die Werke ziehen, während vielmehr alles auf das Wort und den Glauben gezogen werden sollte, und sie sind in Wahrheit nur auf Allegorieen, das heißt, auf Reden über fernliegende Sachen (alieniloquia) bedacht gewesen. Damit nun die Leser nicht mit einem falschen Schein in Allegorieen betrogen werden möchten, habe ich es der Mühe werth gehalten, ihnen zu zeigen, daß sie dann recht auf Allegorieen bedacht sind, wenn sie, soviel als möglich, überall den Dienst des Worts oder den Lauf des Evangelii und des Glaubens in ihnen finden. Denn darauf gehen alle Figuren und Bedeutungen im Gesetz und in dem Volke Mosis, so viele derselben auch immer sein mögen.

Der HErr aber, der da wirkt alles in allem, der auch in dir sein gutes Werk angefangen hat, erhalte dich und lasse dich zunehmen, daß du auch in diesem Leben ein wahrhaft großer Bischof in dem Worte GOttes werdest, und in dem zukünftigen Leben, wenn der Erzhirte der Hirten und Bischöfe kommen wird, die unverwelkliche Krone empfängest [1 Petr. 5, 4.]. Dem, bitte ich, befehl ich mich in deinem heiligen Gebete, und in seiner Gnade gehab dich wohl. Amen. Wittenberg 1525.

Dein D. Martin Luther.

\section{Anmerkungen M. Luthers zum fünften Buch Mosis}

Dies Buch wird Deuteronomium genannt, das heißt, das andere Gesetz. Denn so nennt es Moses selbst 5 Mos. 17, 18. f., wo er dem, der König in Israel sein würde, befiehlt, daß er dies andere Gesetz von den Leviten nehmen und alle Tage darin lesen soll, und 5 Mos. 31, 10. ff. befiehlt er, daß es je über sieben Jahr, wenn das Erlaßjahr gefeiert würde, vor dem ganzen Israel gelesen werden sollte, am Fest der Laubhütten, an dem Orte, den der HErr erwählt hätte.

Es wird aber das andere genannt, nicht weil es ein anderes ist als das, welches auf dem Berge Sinai gegeben ist, sondern daß es von neuem in einem neuen Bunde durch den Dienst Mosis vor denen wiederholt worden ist, die das frühere auf dem Berge Sinai nicht gehört hatten, da ja diejenigen, welche es von dem HErrn selbst gehört hatten, alle in der Wüste umgekommen waren.

Es ist daher dieses Buch ein kurzer Inbegriff [der Bücher] Mosis und eine Summe des ganzen Gesetzes und der Weisheit des Volkes Israel, in welchem, mit Uebergehung dessen, was die Leviten und die Priester angeht, nur das gelehrt wird, was dem Volke und den gemeinen Leuten zu wissen vonnöthen ist. Und wenn du dies Büchlein recht eigentlich zu unserem Gebrauche benennen willst, so kannst du es richtig eine überaus reiche und ganz vortreffliche Auslegung der heiligen zehn Gebote nennen. Wenn du diese kennst, wird dir nichts mangeln, was zum Verständniß
\quellseite{696}
der zehn Gebote nöthig ist. Denn es lehrt dieses Volk nach den zehn Geboten recht zu leben, sowohl geistlicher als leiblicher Weise, und ordnet das innerliche Reich des Gewissens, zugleich auch das weltliche Regiment in Betreff der Güter, sodann die äußerliche Weise der Ceremonien mit ganz göttlicher Billigkeit und Weisheit, so daß auf dem ganzen Gebiete des Lebens nichts ist, was hier nicht auf das weislichste und billigste geordnet wäre.

\section{Das erste Capitel}

\begin{schriftzitat}
V. 1. 2. Das sind die Worte etc.
\end{schriftzitat}

Dies ist eine Beschreibung des Landes. Die Hebräer gebrauchen ein und dasselbe Wort „jenseit“ [\hebraeisch{עבר}, trans], wenn sie auf jener (ultra) oder auf dieser (citra) Seite des Jordans sagen wollen, so daß man nur aus der Stellung, welche der Redende und der Hörende einnimmt, beurtheilen kann, ob man „jenseit“ oder „diesseit“ sagen müsse.

Ferner, eine Wüste wird bei ihnen nicht deshalb so genannt, weil daselbst keine Menschen oder Dinge sind, da in der Schrift viele Städte, die in der Wüste gelegen sind, beschrieben werden, sondern weil es ein waldiger Ort ist. Das hebräische \hebraeisch{ערבה} ist daher richtig mit „Gefilde“ übersetzt, denn es bedeutet eigentlich eine Ebene, wo ohne Wälder und Berge flache Gegenden (campi) sich weithin erstrecken, auf deutsch: ein Feld. Und siehe, daß Moses [4 Mos. 22, 1. 33, 48.] hier „die Gefilde Moab“ macht in der Wüste, indem er gleichsam die Gefilde der Wüste gegenüberstellt, als wenn es ein Gefilde wäre in waldiger Gegend (silvestribus).

„Das rothe Meer“ wird es von den Griechen genannt, aber auf hebräisch: das Meer \hebraeisch{סוף}, das heißt, das voller Niethgras oder Binien ist. Es scheint, daß Hieronymus [in der Vulgata das Wort \hebraeisch{סוף}] falsch übersetzt habe: „Wo sehr viel Goldes ist“, da es wahrscheinlicher\footnote{Statt vero similius in den Ausgaben haben wir verisimilius angenommen.} ist, daß „Dilahab“ der Eigenname eines Ortes ist, welcher Ueberfluß an Gold bedeutet. Denn auch anderswo macht ebenderselbe Mann aus Eigennamen Gattungsnamen und umgekehrt. Dieser Fehler kann in der hebräischen Sprache leicht vorkommen.

Aber die folgende Beschreibung des Landes, da er sagt: „gegen dem Schilfmeer, zwischen Paran und Thophel, Laban, Hazeroth und Di-Zahab, elf Tagereisen von Horeb, durch den Weg des Gebirges Seir bis gen Kades-Barnea“, wollen wir den Einwohnern jener Orte überlassen, damit wir uns mit ungewissen Dingen nicht vergeblich abmühen.

\begin{schriftzitat}
V. 3–5. Und es geschah im vierzigsten Jahr etc.
\end{schriftzitat}

Dies ist die Beschreibung der Zeit. Das vierzigste Jahr hier ist der Auszug aus Egypten; im ersten Monat, am ersten Tage des Monats, nämlich am Ende des Lebens Moses, der noch fast zwei Monate bei ihnen sein sollte, in welchen er sein Testament macht und ihnen Valet sagt.

Merke aber dies: „Er redete zu den Kindern Israel alles, wie ihm der HErr an sie geboten hatte.“ Er redet nichts, als was ihm der HErr geboten hat, denn er wiederholt es und schärft es ein bis zum Ueberdrus, damit er lehre, daß man im Volke GOttes nichts reden soll als das, von dem man gewiß ist, daß es GOttes Wort sei, ja, auch das ist nicht genug, daß man gewiß sei, es sei GOttes Wort, sondern es muß durch das Gebot GOttes ein jeglicher gezwungen werden, das Wort GOttes zu verkündigen, wie er hier sagt: „Alles, was ihm der HErr geboten hatte.“ Er sagt nicht: „Was ihm der HErr eingegeben hat, wie sich die gottlosen Menschen verwegen selbst gefallen\footnote{Es scheint uns, daß besser placent statt placeant gelesen werde. Darnach haben wir übersetzt.} wegen der Eingebung (afflatu) des Heiligen Geistes, da sie sich dünken lassen, daß sie in ihren Aufsätzen und Rathschlägen nicht irren. Denn was hier gesagt wird: „Es fing Moses an, auszulegen das Gesetz“, gehört auch dazu, daß er nichts rede, als was Gott geboten hat.

\quellseite{697}
\begin{schriftzitat}
V. 6. Der HErr, unser GOtt, redete mit uns am Berge Horeb, und sprach: Ihr seid lange genug an diesem Berge gewesen etc.
\end{schriftzitat}

Hier fängt die Erzählung an, in welcher er die ganze Geschichte gar schön wiederholt und kurz zusammenfaßt, indem er die Wohlthaten und Wunder GOttes und zugleich auch die Werke der gottlosen Menschen anführt, damit er den Ruhm und die Herrlichkeit GOttes preise, um dadurch das Volk, durch die Erfahrung belehrt, dazu zu locken, daß sie der Güte GOttes vertrauen möchten und seinen Zorn fürchten, auf daß sie in solcher Weise geschickt werden möchten, sein Gesetz von Herzen aufzunehmen. Denn dies ist die allerbeste Vorbereitung auf das Hören des Gesetzes, und zur Anregung des Zuhörers, welche geschieht durch das evangelische Preisen der Barmherzigkeit und des Zornes GOttes.

„Horeb“ und Sinai sind dasselbe. Ferner ist Hieronymus hier etwas dunkel, da er sagt [V. 6.]: Sufficit vobis [Es genügt euch], desgleichen [V. 7.]: Revertimini [Kehret zurück]. Denn er befiehlt nicht, daß sie zurückkehren sollen, sondern, nachdem sie am Sinai fast ein ganzes Jahr geblieben waren, in welchem sie das Gesetz empfangen hatten, die Stiftshütte erbaut, und das Priestertum und die Obrigkeit geordnet waren zu rechter Regierung des Volks, desgleichen das Volk in Fähnlein und Heereshaufen abgetheilt und gezählt worden war, und alles in die rechte Ordnung gebracht, befahl er, daß sie im zweiten Jahre, im zweiten Monat, aufbrechen sollten und so geordnet und gerüstet auf geradem Wege nach dem Lande Canaan ziehen, wohin sie auf dem Wege durch das Gebirge Seir in elf Tagen kommen könnten. Das ist es, was er hier sagt: „Ihr seid lange genug an diesem Berge gewesen; wendet euch, und ziehet hin, daß ihr zu dem Gebirge der Amoriter kommet.“ Und so sind sie aufgebrochen und kamen an ihre dreiundzwanzigste\footnote{Nach der Zählung der Lagerstätten im 33. Capitel des 4. Buchs Mosis ist zu Kades die einundzwanzigste Lagerstätte gewesen nach dem Auszug von Sinai [V. 16.].} Lagerstätte zu Pharan oder Kades, welches auch Zin genannt wird, wie das 4. Buch Mosis [Cap. 20], 1. und 33, 36.] lehrt.

\begin{schriftzitat}
V. 7. f. In dem Berge des Amorites etc. (Ad montem Amorraei).
\end{schriftzitat}

Es sei hier ein für allemal gesagt für ungelehrte Leute, daß Berg (montem) hier und an ähnlichen Stellen, in welchen er Länder und Völker anzeigt, für gebirgige Gegenden genommen werde, oder vielmehr für ein ganzes Land, in welchem verschiedene Berge sind. Daher lesen wir öfter: Auf dem Berge Ephraim, auf dem Berge Juda. So auch hier: „Der Berg des Amoriter’s.“ Sodann steht der Singular Amorraeus für den Plural oder für den Namen des Volks, nicht für eine einzelne Person.

\begin{schriftzitat}
V. 9–12. Da sprach ich zu derselben Zeit zu euch etc.
\end{schriftzitat}

Nämlich da ihr noch am Berge Sinai waret (davon er [V. 6.] gesagt hatte: „Der HErr redete mit uns am Berge Horeb“). Er sagte dies aber 2 Mos. 18, 14. auf Eingeben Jethro’s, ehe sie das Gesetz empfangen hatten. Darum ist dieser Gegenstand im Erzählen später gesetzt, welcher doch in der That eher geschehen ist als das Vorige.

Lieber, beachte hier, daß man sich eher nach Personen und Obrigkeit umsieht, als die Gesetze gegeben werden, denn vergebens werden Gesetze gemacht, wenn die Leute nicht da sind, welche sie verwalten und vollziehen; und zugleich, daß die Obrigkeit ein Amt sei, welches Arbeit, Lasten und Hader des Volkes tragen muß, daß sie sich als einen Knecht erkenne, sich nicht für einen Herrn halte, wie es denn die Liebe erfordert. Diejenigen aber, welche ihren Nutzen, Müßiggang und Ruhe im obrigkeitlichen Amte suchen, bringen es in verkehrter Weise\footnote{Wir haben mit der Wittenberger und der Jenaer perverse angenommen statt perversi in der Erlanger.} dahin, daß das Volk ihnen vielmehr dienen muß, und gebrauchen des Volkes Nachtheils zu ihrem Muthwillen.

„Mühe“ [V. 12.] nennt er das beständige Beschäftigkeit und die fortwährende (paratam) Bekümmerniß für das Volk, um es zu regieren, wodurch sie [die obrigkeitlichen Personen] von Erholung und Ruhe ausgeschlossen werden.

„Last“ nennt er die überaus große Geduld im Tragen der Schwachheit und ungeschlachten Sitten des Volks. Wenn eine obrigkeitliche Person damit nicht begabt ist, und alles nach der Schärfe des Rechts erzwingen will, so wird sie eher das Volk verderben als erheben. Man muß das Sprichwort des Kaisers Friedrichs des Dritten wissen: Wer nicht durch die Finger sehen kann, der kann nicht regieren.

„Hader“ sind die Händel, die das Volk unter
\quellseite{698}
sich hat; um diese zu schlichten und beizulegen ist die Obrigkeit notwendig. Es kommt vom Volke an die Obrigkeit nichts als Mühe, Kümmerniß, Geduld, Beschwerde und Arbeit, wie jener Bias richtig gesagt hat: Regiment weist aus, was einer für ein Mann ist, ja, wie die Schrift sagt, sie werden Götter genannt [Psalm 82, 6.], weil sie in ganz göttlichen Tugenden erglänzen sollen, weit über das Maß der Menschen, weil sie das schwache und mannigfach geartete gemeine Volk (varium corpus vulgi) unter dem Reiche des Satans in feindlichem Lande seliglich versorgen sollen.

Die eingeschalteten Worte [V. 11.]: „Der HErr mache euer noch viel tausendmal mehr“ etc., sind bei uns auch im Gebrauche, und zwar gar christlich, wo man sagt von Kindern oder von Vieh, um die es wohl steht: GOtt behüte dich, oder GOtt segne dich. Nämlich dies ist von Alters her überliefert oder auch von diesem Exempel Mosis hergenommen, wo er sofort, nachdem er gesagt hatte, das Volk sei gemehret worden wie die Sterne des Himmels, hinzufügt: „GOtt mache euer noch viel mehr“, gleichsam als fürchte er einen Zauber, wie unsere Weiblein zu argwöhnen pflegen, wenn die kleinen Kinder von jemandem ins Angesicht gelobt werden.

\begin{schriftzitat}
V. 13–16. Schaffet aus eurer Mitte weise Leute etc.
\end{schriftzitat}

Die Thiere werden durch Gewalt und Kunst regiert. Die Menschen müssen durch Weisheit und Verstand regiert werden, weil der Mensch seine Stärke hat (viget) in der Vernunft, welche man nicht erreichen kann dadurch, daß man ihm einen Strick anlegt, oder ein Schwert in ihn stößt, sondern durch das Wort, welches durch das Ohr zu ihm gelangt; wenn aber die Vernunft gefangengenommen ist durch das Wort, so wird der ganze Mensch gezogen und regiert [und dahin gebracht], wohin du nur willst.

Und hier siehst du, daß die Obrigkeit durch die Stimmen des Volks erwählt werden soll, was auch die Vernunft vorschreibt. Daher wird auch hier das Volk auf diese Weise durch das Wort Mosis gefangen, gehorsam willig und lobt Mosen. Denn daß die Obrigkeit dem Volke wider seinen Willen aufgedrängt wird, ist gefährlich oder schädlich. Er nennt sie „wohlbekannte Leute“, weil das Volk sie kennen soll, viel mehr aber, weil sie solche Leute sein sollen, die erfahren und verständig seien in Sachen, so daß du verstehen mußt, die „wohlbekannten“ und die „kundigen“ (scientes) Leute seien dieselben. „Weise“ sind diejenigen, welche göttliche und menschliche Dinge verstehen, nämlich, welche die Gesetze und Rechte und alles, was notwendig ist für das Leben der Menschen, wissen. „Verständige“ werden im Hebräischen diejenigen genannt, welche aufmerksam sind und mit scharfem Urtheil beobachten, oder die einen wachsamen Geist haben, so daß sie über alle Dinge ein richtiges und gesundes Urtheil fällen können, und verstehen, was die Billigkeit erfordert. Denn viele Dinge führen die Weisen irre, wenn sie nicht auch verständig sind, und viele Dinge täuschen sie, wenn sie nicht erfahren und kundig sind.

Du siehst also, daß im göttlichen Rechte keine Rücksicht genommen werde auf die Reichen, Gewaltigen, Edlen, Starken, Freunde, bei der Verwaltung eines obrigkeitlichen Amtes, wie die Welt zu thun pflegt, sondern auf die Weisen, Verständigen und Erfahrenen, wenn sie auch arm, verachtet und schwach sind etc.

Und nun hat mit Recht (sane) die Frage aufgeworfen, ob es besser sei, wenn ein Fürst gut sei, aber unverständig (imprudens), oder wenn er verständig sei, jedoch zugleich auch böse. Sicherlich erfordert Moses hier beides [daß er gut sei und auch verständig], jedoch, wenn man beides nicht haben kann, ist es besser, wenn er klug und nicht gut ist, als wenn er gut und nicht verständig ist, weil ein guter [Fürst] durchaus in nichts regiert, sondern sich nur regieren läßt, und zwar nur von den allerschlimmsten Leuten. Wenngleich ein kluger Fürst den Guten auch Schaden thut, so regiert er doch zugleich die Bösen, was für die Welt mehr vonnöthen und geeignet ist, da sie nichts Anderes ist als ein großer Haufe von bösen Leuten. Diese Meinung scheint auch Moses hier zu billigen, da er zuerst handelt von der Weisheit der Fürsten, darnach erst davon, ob sie gut seien.

\begin{schriftzitat}
V. 17. 18. Keine Person sollt ihr ansehen; desgleichen: auch vor niemandes Person sollt ihr euch scheuen.
\end{schriftzitat}

Dies ist die höchste und die schwierigste Tugend der Fürsten, nämlich Billigkeit und Unparteilichkeit (simplicitas) im Gerichte. Denn es ist leicht zu richten über Arme und Üble,
\quellseite{699}
aber Mächtige, Reiche, Freunde zu verurtheilen, Blutsverwandtschaft, Ehre, Furcht, Gunst, Gewinn hintanzusetzen, und allein auf die Sache zu sehen, das ist göttliche Tugend. Dies thut auch kein Fürst, der nicht von dem Heiligen Geist durch den Glauben an GOtt muthig gemacht ist. Darum ist die Welt voller Fürsten, aber wer will finden Einen, der rechtschaffen fromm sei?“ sagt Salomo [Spr. 20, 6.]. Wie oft wird doch dieser Ausspruch in den Propheten wiederholt, wo sie der Unterdrückung und Ungerechtigkeit beschuldigt werden gegen die Armen, Waisen und Wittwen.

\subsection{Denn das Gerichtamt ist GOttes}

Auch diese Stelle ist stark für die bürgerliche Obrigkeit. Denn gleichwie er hernach Cap. 32, 35. sagt, daß die Rache sein sei, so sagt er auch hier, daß das Gericht seine Ordnung sei. Dies sagt er aber zum Trost und zur Ermahnung der Obrigkeit, damit sie wisse, daß sie nicht ihre eigene Sache führe, sondern GOttes, und sie gewiß seien, daß sie GOttes Diener seien und GOttes Werk ausrichten und in einem seligen (salutis) Stande seien, auf daß sie sich dessen desto mehr erfreisten, recht und schlecht (simpliciter) zu richten.

Siehe auch das an, daß er, ehe das Gesetz gegeben wurde, Richter eingesetzt hat, um uns die Billigkeit ans Herz zu legen. Denn die Richter sind lebendige Gesetze oder die Seele des Gesetzes. Darum sind sie eher und höher als die Gesetze und auch als das Schwert. Deshalb, wiewohl sie nach den Gesetzen das Urtheil fällen und dies mit dem Schwert vollstrecken sollen, so sollen sie sowohl das Gesetz als auch das Schwert in ihrer Gewalt haben, damit sie die Gesetze nicht wider die Gesetze gebrauchen, wie jene thörichten Leute, welche unverwandtes Auges auf das Gesetz starren, die Umstände der Sache nicht ansehen, und allein dies im Munde haben: So steht es geschrieben, so geschehe es. Doch hiervon anderswo.

\begin{schriftzitat}
V. 19–26. [V. 20.] Ihr seid an das Gebirge der Amoriter gekommen.
\end{schriftzitat}

Kades-Barnea ist nahe bei dem Lande Juda gegen Mittag. Darum spricht Moses ihnen Muth ein, als sie schon die Grenzen des Landes der Verheißung erreicht hatten, und facht ihren Glauben an durch Wiederholung der Verheißungen, damit sie ohne Furcht hineingehen möchten. Aber hier wankt der Glaube, und der Unglaube überwältigt das Volk, [uns] zu einem jämmerlichen Exempel, so daß sie, als sie schon am Eingange des Landes waren, welches sie einnehmen sollten, aus Strafe hiefür wieder zurückgehen mußten, und achtunddreißig Jahre lang in der Wüste umherschweifen, bis daß sie alle umkamen und niemand hineingelangte, mit Ausnahme von Josua und Caleb; nur zwei von einer so großen Zahl.

Es fing aber dies Fallen in Unglauben an mit einem geringen und schon anziehenden Zweifel, da sie vorwendeten, man müsse Kundschafter schicken, welche den Weg erforschen sollten; und sie vertrauen sich nicht dem Worte GOttes an, der ihnen die Verheißung gegeben hatte, der sie auch durch das rothe Meer geführt hatte, wo ein Weg unmöglich war; als ob er sie nicht auch in einem offenen und mit Wegen wohl versehenen Lande führen könnte oder wollte. So mißtrauen die GOtt in kleinen Dingen, die ihm in großen Dingen geglaubt hatten, auf daß du wissest, der Glaube sei nicht ein Werk des freien Willens, sondern allein der Gnade GOttes.

Es ist aber auch dies Werk GOttes wunderbar, wie denn alle seine Werke sind. Denn das jüdische Volk mußte fallen, da es den Amoriten schon auf dem Halse saß, damit die Langmuth GOttes gepriesen würde, welche jenen Heiden, die bereits umgebracht werden sollten, Raum zur Buße gäbe, und sie bekehrt werden möchten. So trägt er, mit Schaden und Gefahr für sein eigenes Volk, auch für die Heiden Sorge, indem er jenes schlägt, dieser verschont, auf daß er zeige, er sei nicht allein der Juden GOtt, sondern auch der Heiden [Röm. 3, 29.], wiewohl auch diese Güte und Langmuth GOttes den Heiden durch ihre Blindheit zu einem Strick und Aergerniß wird, so daß sie sich verhärten und Unbußfertigkeit hinzuthun, daß der Zorn und das gerechte Gericht GOttes mit Haufen über sie kommen muß [Röm. 2, 5.].

Denn wer zweifelt, daß die Amoriter, als sie sahen, daß der Feind zurückwich und in der Wüste umkam, und sie nun sicher und frei wären, diese Ueberzeugung angenommen haben, daß sie GOtt für sich hätten als einen Streiter wider ihre Feinde? Und es wird ihnen dieses Werk GOttes eine Falle, daß sie vermessen werden wegen ihrer Verdienste, und den Zorn GOttes weit
\quellseite{700}
von sich abschieben auf andere Völker, so daß sie nun ohne Furcht in ihren Sünden fortfahren. Das ist es, was der Psalm [Ps. 66, 5. Ausg.] sagt: „Er ist erschrecklich in seinen Rathschlägen über die Menschenkinder“, und wir sehen hier ein Exempel nach beiden Seiten hin, was da in der Welt der Unglaube und die Blindheit der Menschen thue.

Im vierten Buch Mosis schreibt Moses, Cap. 13, 2. f., daß die Kundschafter auf Befehl des HErrn gesandt worden seien. Wie sagt er denn nun hier, daß sie auf das Dringen des Volkes gesandt worden seien? So steht es um die ganze Schrift, daß sie sich selbst auslegen will, indem die Stellen von überall her mit einander verglichen werden, und daß sie verstanden werden will, indem sie allein die Meisterin (magistra) ist. Und das ist die allersicherste Weise, den Sinn\footnote{Wir haben mit der Wittenberger und der Jenaer sensum angenommen statt sensus in der Erlanger.} der Schrift zu erforschen, wenn man durch Vergleichung der Stellen und genaue Beobachtung derselben nach dem rechten Sinne trachtet. So hat Moses oben gesagt, daß er gleichsam aus eigenem Antriebe gesagt habe: „Schaffet aus eurer Mitte weise Leute“, während er nach 2 Mos. 18, 17. dies auf den Vorschlag des Jethro gesagt hat. So ist es hier gewiß, daß die Kundschafter gesandt worden seien, weil das Volk in seinen Unglauben darauf drang, und daß der HErr ihrer Schwachheit nachgegeben und nachgegeben habe, und so zu Mose gesagt, 4 Mos. 13, 3.: „Sende“ etc.

Aber das Volk sündigte durch diese Bangigkeit, und sie fingen an, sich auf menschlichen Rath zu gründen, nicht allein auf das Wort GOttes. Darum ist ein so großer Fall gefolgt auf diesen kleinen Zweifel, und, wie Paulus sagt [1 Cor. 5, 6.], „ein wenig Sauerteig versäuerte den ganzen Teig“. So hat die Schlange die Eva, welche anfangs nur ein wenig zweifelte, bald ganz niedergeschlagen durch den Unglauben, mit dem sie das ganze menschliche Geschlecht verderbt hat. So gar ist nicht zu scherzen mit dem Worte GOttes und dem Glauben.

\begin{schriftzitat}
V. 27–33. Der HErr ist uns gram etc.
\end{schriftzitat}

Siehe, welch einen großen Fall zieht endlich jener ganz geringe Anfang des Zweifelns nach sich und das Vertrauen auf menschlichen Rath.

Denn nun haben sie aller Verheißungen und Wunder GOttes gänzlich vergessen, sind allein in ihrem Rathe ersoffen, und werth geworden, daß sie falsche Boten und Prediger hörten, und, wie Paulus sagt 2 Thess. 2, 11.: „GOtt sendet denen kräftige Irrthümer, welche der Wahrheit nicht geglaubt haben.“ Jene predigen größere Dinge, als die Sache an sich ist, und diese glauben es, und anstatt auf das rechte Wort GOttes gründen sie sich auf die Lügen der Menschen.

Jetzt erst sind die Mauern bis an den Himmel aufgeführt, jetzt ist das Volk groß und hoch, jetzt sind Enakin da, jetzt frißt das Land seine Einwohner [4 Mos. 13, 33.]. Nachdem das Wort GOttes verloren ist, wüthet der Unglaube, und vergeblich wird ihnen von Josua und Caleb die Verheißung und die Wahrheit GOttes eingeschärft. Die Gottlosigkeit siegt, und die Lüge bekommt die Oberhand, die Wahrheit wird zu Boden geschlagen. Hier ist die Furcht der menschlichen Klugheit und Tugend in göttlichen Dingen.

Endlich schreiten sie zur Gotteslästerung und sagen: „Der HErr ist uns gram.“ Das ist nun die Dankbarkeit für so viele Wunder, die ihnen erwiesen worden sind. Und wer kann doch dieses überaus böse Exempel des Unglaubens gebührend groß machen? Merke also: der Unglaube macht der Gefahren mehr und dieselben größer, als sie sind, das Wort GOttes aber achtet er für nichts. Dagegen die Gottseligkeit hält alle, auch die größten Gefahren, für nichts, und hält das Wort GOttes für eine Kraft GOttes, Röm. 1, 16., wie auch Moses hier sagt [V. 29.]: „Fürchtet euch nicht“ etc., indem er GOttes Kraft verheißt und die vergangenen Exempel wiederholt, um ihren Glauben aufzurichten, und kräftiglich durch das Wort GOttes streitet; aber er richtet nichts aus.

Merke: wenn dem Moses, der doch mit so vielen Worten und Zeichen GOttes gerüstet war, nicht geglaubt wird, sondern er gezwungen wird, das ganze Volk dem Unglauben zu überlassen: was wundern wir uns denn, daß heutzutage wenige uns glauben, und die ganze Welt im Unglauben wüthet? Aus der so großen Zahl hängen nur zwei dem Moses an; dennoch läßt er auch da noch nicht ab, oder verläßt das Amt des Wortes und hört auf ihr Führer zu sein (ducatus), sondern predigt vergebens den Ungläubigen. Dies ist eine große Historie und ein treffliches Exempel des göttlichen Wortes
\quellseite{701}
und menschlichen Rathes, in welchem man sehen kann, wie es mit der allmächtigen Kraft des Wortes in der Welt stehe.

\begin{schriftzitat}
V. 34–39. Als aber der HErr euer Geschrei hörte, ward er zornig und schwur etc.
\end{schriftzitat}

Aus dieser Stelle ist der Vers im 95. Psalm genommen [V. 11.]: „Ich schwur ihnen in meinem Zorn“ etc., und das Exempel dieser Historie wird herrlich gehandelt in der Epistel an die Hebräer. Aber was ist es, daß Mose [V. 37.] sagt, der HErr sei auch über ihn zornig geworden um des Volkes willen, da er dies doch durch seine eigene Sünde verdient hatte, wie 4 Mos. 20, 12. geschrieben steht? Er deutet hier freilich ein Geheimniß an. Aber nach der Historie verstehe ich es so, daß Moses, veranlaßt dadurch, daß er ein so widerpenstiges und ungläubiges Volk regieren mußte, endlich in Unglauben gefallen sei, und das bezeichnet er hier so, daß der Zorn des HErrn um ihretwillen ihn getroffen habe.

Doch dies ist uns wiederum zur Lehre und zum Troste geschrieben, damit wir nicht verzweifeln, wenn wir gesündigt haben, da in dieser Anfechtung nicht allein viele vom Volke, ohne Zweifel große und heilige Männer, gefallen sind, sondern auch der treffliche Mann und überaus große Führer des Volks, Moses, sammt seinem heiligen Bruder Aaron, damit wir den HErrn fürchten und uns mißtrauen, da wir alles, was wir sind, allein durch seine Kraft sind.

\begin{schriftzitat}
V. 40. Ihr aber wendet euch und ziehet etc.
\end{schriftzitat}

Das ist die Klage, daß sie gezwungen werden, von dem Eingange in das Land wieder zurückzugehen zu dem rothen Meer; zwar nicht nach Egypten, sondern in die Wüste; damit sie daselbst allmälig in achtunddreißig Jahren aufgerieben würden.

\begin{schriftzitat}
V. 41–46. Wir haben an dem HErrn gesündigt, wir wollen hinauf etc.
\end{schriftzitat}

Diese Eine Stelle wäre genugsam, um den freien Willen zu widerlegen, daß wir wissen könnten, daß nichts Anderes etwas tauge, als was GOtt in uns wirkt, wie er hier sagt [V. 42.]: „Denn ich bin nicht unter euch.“ Denn was unterlassen sie hier, das menschlicher Rath und Kraft vermag? Sie haben Reue, es thut ihnen leid, sie sind bereit, sie nehmen die Waffen, sie ziehen hinauf, und es mangelt hier durchaus an nichts, was der Mensch thun kann. Und dennoch taugt es nicht allein nichts, richtet auch nichts aus, sondern wird auch verboten und verdammt. Desgleichen, auch nach der Strafe, und nachdem sie die Niederlage erlitten haben und weinen und klagen, nimmt er sie nicht an, wiewohl sie dies doch sonder Zweiffel ohne Heuchelei thaten aus allen Kräften des freien Willens. Aber weil sie sich vermaßen, durch sich und ihr Werk GOtte zu dienen und zu gefallen, und die Gnade zu verdienen nach Billigkeit oder nach Recht (de congruo aut condigno), und die Werke thaten, die von ihnen erwählet, nicht von GOtt befohlen waren, so thaten sie alles vergeblich.

Aber auch dies ist geschehen, um die Amoriter zu verstöcken, damit sie durch Mehrung ihrer Sünden ihren Untergang beschleunigten, da sie als Sieger hervorgingen und glaubten, daß GOtt ihnen günstig sei. Und siehe die unterschiedlichen Gerichte GOttes: seinem Volke, welches auf seine eigenen Kräfte vermessen ist, läßt er nichts glücklich hinausgehen, er läßt es besiegt und geschlagen werden, als ob er nicht ihr GOtt wäre; aber die Feinde, die Amoriter, welche auf ihre Kräfte vertrauen, läßt er Glück haben und siegen, auf daß du wissest, es sei mehr zu fürchten und ein augenfälliges Zeichen des beständigen Zornes GOttes, wenn es dem Gottlosen auf seinem Wege wohlgeht und in den Dingen den Sieg erlangt, die GOtt weder geboten hat noch wirkt, sondern die er ausrichtet nach seinem eigenen Rathe und in seiner eigenen Kraft, wie GOtt bisher denen zugelassen hat, die Glück gehabt haben in ihrer Trügerei mit menschlichen Satzungen.

\subsection{Von den heimlichen Deutungen (mysteriis) in diesem Capitel}

Was ich anderswo oft erinnert habe, wiederhole ich hier und will wiederum erinnern, daß der christliche Leser vor allen Dingen seine Mühe dadurch gut anwende, daß er den sogenannten buchstäblichen Sinn suche, welcher allein das ganze Wesen des Glaubens und der christlichen Theologie ist, welcher allein in Trübsal und Anfechtung Stand hält, und die Pforten der Hölle mit Sünde und Tod überwindet, und triumphirt zum Lob und Ruhme GOttes. Die
\quellseite{702}
Allegorie ist aber öfters ungewiß und nicht zuverlässig und durchaus nicht sicher, um den Glauben zu stützen, da sie öfters von menschlicher Muthmaßung und Meinung abhängt. Wenn sich jemand darauf verläßt, so stützt er sich auf den Rohrstab Egypten. Darum soll man sich hüten vor Hieronymus, Origenes und ähnlichen Vätern, oder ihre Schriften mit freiem Urtheile lesen, ja auch vor der ganzen Alexandrinischen Schule, von welcher der Jude Philo erwähnt (wie Eusebius und Hieronymus bezeugen), daß sie sich in solchen Bemühungen um Allegorieen einst hervorgethan habe. Denn das Beispiel dieser, welches allzusehr gelobt und angenommen worden ist, haben spätere Schriftsteller in unglücklicher Weise nachgeahmt, und nach ihrem Belieben alles aus der Schrift gemacht und gelehrt, bis daß etliche die Worte GOttes auf die ungereimtesten Seltsamkeiten (monstra) gezogen haben, und, wie Hieronymus sich auch über seine Zeit beklagt, die Schrift mit sich selbst in Widerspruch bringen durch unpassende (incongrua) Zeugnisse; welcher Uebelthat er doch auch selbst schuldig war.

Solche Leute sind diejenigen, welche heutzutage fast die ganze Schrift von der heiligen Jungfrau auslegen, wo sie nur ein Wörtlein finden, das weibliches Geschlechts ist; desgleichen, welche aus dem Städtchen\footnote{Luc. 10, 38. steht das Wort castellum, das ist, ein befestigtes Städtchen. Im Griechischen κώμη, welches in unserer Bibel durch „Märkt“ übersetzt ist.} der Martha die Klöster erbauen, aus den Städten in Israel die magistros nostros auf den hohen Schulen machen, und ähnliche, unzählige Albernheiten (portenta), bis sich endlich einer gefunden hat, der die ganzen Metamorphosen Ovids auf Christum angewendet hat, worüber Hieronymus mit Recht entrüstet ist in der Epistel an den Paulinus, und sie Gaukler (circulatores) nennt.

Darum muß man hier die Regel Pauli halten, daß man die Allegorieen als etwas Untergeordnetes (secundo loco) ansehe, und sie anwende, um die Lehre des Glaubens zu stärken, schmücken und auszustreichen (locupletanda), oder, wie er 1 Cor. 3, 11. f. sagt, daß sie nicht der Grund seien, sondern auf den Grund gebaut werden, nicht als Heu, Holz und Stoppeln, sondern als Silber, Gold und Edelsteine. Dies geschieht dann, wenn die Weissagung, wie Röm. 12, 7. vorgeschrieben ist, dem Glauben ähnlich ist, nämlich, daß du einen gewissen Ausspruch vor dich nehmest, der irgendwo in der Schrift steht, der nach dem buchstäblichen Sinne gehandelt ist, und dann erst deine Allegorie, welche auf dasselbe hinauskommt (idem sonantem), anwendest, nicht als ob die Allegorie den Spruch beweisen oder stützen solle, sondern daß sie aus dem Spruche bewiesen und gestützt werde, wie ein Gebäude nicht den Grund trägt, sondern von der Grundlage getragen wird.

Damit wir nun hier sicher allegorisiren können, ist uns Paulus 2 Cor. 3, 7. ff. vorangegangen, wo er den ganzen Moses in einer sehr reichen Allegorie handelt, indem er die Klarheit des Angesichtes Mosis mit der Klarheit des Angesichtes Christi vergleicht. Nach diesem Beispiel kannst du nun in rechter Weise eine Allegorie zu Stande bringen, indem du aus Mose das Amt des Gesetzes und der Sünde und des Todes, gleichwie aus Christo das Amt des Evangelii, der Gnade und des Lebens machst, wie Paulus dies im vierten und achten Capitel des Briefes an die Römer schön lehrt. Ferner macht Paulus aus dem Einen Angesichte Mosis\footnote{Statt Moses in der Erlanger und der Jenaer ist mit der Wittenberger Mosi zu lesen.} ein solches, welches auf zwiefache Weise angesehen werden kann, ein verhülltes und ein offenbares. Das verhüllte schreibt er denen zu, welche durch den Dienst des Gesetzes nur zu Werken und zur Heuchelei gebracht werden. Diese verstehen nie die Kraft des Gesetzes, sondern ihr Herz bleibt blind beim Lesen des Alten Testamentes, daß sie die Herrlichkeit Christi weder erkennen noch darnach seufzen. Das offenbare schreibt er denen zu, welche durch das Gesetz nicht Werke thun, sondern die Sünde erkennen und getödtet werden, so daß sie nach der Gnade verlangen.

Die Beschreibung des Orts nun und die Beschreibung der Zeit bedeutet den Ort selbst, wo das Gesetz gelehrt wird, das heißt, das Volk des Gesetzes, nämlich die Synagoge, wozu alle Wörter sich schicken, welche Mosis hier setzt, nämlich: 1. „Jenseits des Jordans“, das heißt, da man noch nicht vom Gesetze zur Gnade übergegangen ist, das Evangelium auch noch nicht öffentlich gelehrt wird. 2. „In den Gefilden Moab“, weil da das Volk ist, welches sich rühmt, daß es von den Vätern herkomme. Denn „Moab“ heißt: vom Vater. „Wir sind Abrahams Same“, sagen sie. Und „die Gefilde
\quellseite{703}
„Moab“ sind das, was sie von den Vätern haben, nämlich der Name, die Gebräuche, die Güter, das Geschlecht und die Verheißung der zukünftigen Dinge. 3. „In der Wüste“, weil es ein unfruchtbares und unbebautes Volk ist, wenn es verglichen wird mit der Offenbarung der Gnade. 4. „Gegen das rothe Meer“, weil dies Volk einen Abscheu hat gegen das Kreuz und die Tödung des alten Menschen, welcher in jenem Meere ersäuft ist. 5. „Paran“ heißt Prahlerei und Ruhmredigkeit, und kommt her von dem hebräischen \hebraeisch{נז}, denn solches pflegt das Geschlecht der Heuchler zu thun. 6. „Tophel“ heißt Fall, weil sie von Tag zu Tage mehr dahin fallen und es ärger mit ihnen wird, gleichwie die Christen von Tag zu Tage erneuert werden. 7. „Laban“ bedeutet weiß, und zeigt ihren schönen Schein der Göttlichkeit an. 8. „Hazeroth“ bedeutet Vorhöfe, weil die Gerechtigkeit des Gesetzes auswendig im Fleische ist, und niemals zugelassen wird zu dem Heiligthum des Geistes. 9. „Disahab“ ist Ueberfluß an Gold; die Gerechtigkeit aus den Werken verschlingt die Reichthümer der Welt. 10. „Elf Tagereisen“ sind nur um Eins entfernt von zwölf; denn das Amt des Gesetzes steht dem apostolischen Amte des Evangeliums sehr nahe, aber dennoch entfernt es sich von demselben. 11. „Der Weg des Gebirges Seir“ ist der Weg der Werktreiber des Gesetzes, welcher schwer, bergig und rauh, und „Seir“, das ist, voll Gestrüpp ist; der Weg Christi aber ist lieblich, leicht und angenehm im Geist der Gnade. 12. „Die Könige der Amoriter“, welche jenseits des Jordans getödtet sind, bedeuten, daß durch das Amt des Gesetzes GOttes zwar die Bemühungen der Menschen äußerlich verworfen werden, und daß sie sich allein des Gesetzes GOttes rühmen, aber dennoch an dasselbe, dem Geiste nach, nicht hinankommen. Denn „Amoriter“ bedeutet einen beredten oder geschwätzigen Menschen, nämlich die Philosophie und die menschliche Theologie, deren erster König „Sihon“ ist, das heißt ein Verwüster, weil diese Art der Lehrer die Schrift GOttes in unglaublicher Weise verwüstet mit mancherlei Deutung (aequivocationibus) und Verkehrung der Worte auf ihre Meinung. Die Stadt „Hesbon“, was im Hebräischen Künste und Vernunft bedeutet, oder alles, was durch Verstand und Kunst erdacht wird, bezeichnet ohne Zweifel die vernunftgemäße (rationalem) Weisheit und (wie man sagt) wortreiche (verbosam) Theologie, welche den Worten GOttes angedichtet ist; die gottlose Amoriter, das ist wortreiche Verführer macht, und alles verwüstet, wie damals die Pharisäer, und wie jetzt die hohen Schulen und das Papstthum thun. Der andere ist „Og“, das ist ein Koch, der für den Bauch sorgt. „Basan“ heißt Fettigkeit, weil nämlich das Reich der Amoriter mit seiner Gottlosigkeit und Geschwätzigkeit Reichthum und Ueberfluß für den Bauch verdient, wie denn die Exempel vor Augen sind.

„Ihr seid lange genug“ etc. [V. 6.]. Daß hier der HErr befiehlt, daß sie von Kades-Barnea ziehen sollten in das Land der Verheißung, bedeutet, daß eine andere Lehre sein werde, nämlich die des Evangelii, von welcher das Gesetz und die Propheten Zeugniß ablegen, zu welcher die Juden eingehen sollten zur Zeit des neuen Testamentes. Denn „Kades“ bedeutet etwas Heiliges, „Barnea“ einen unverschweifenden unbeständigen Sohn, wie das Volk ist, welches auf ein weltliches Heiligthum und auf ungewisse Werke vertraut, aber niemals Ruhe hat in seinem Gewissen. Daher ist dieses Wort damals erfüllt, als Johannes der Täufer und Christus anfingen zu lehren, daß sie aus der Gerechtigkeit und Heiligkeit des Fleisches in das Reich GOttes ziehen sollten, indem sie sagten [Matth. 3, 2.]: „Thut Buße, das Himmelreich ist nahe herbei gekommen“, das heißt, das rechte Land der Verheißung.

Hier fängt die elende Abtrünnigkeit der ungläubigen Synagoge an, und es sterben unter dem Zorne GOttes alle streitbaren Männer, wie dies die Epistel an die Hebräer in würdiger Weise behandelt, indem sie dies Exempel des Unglaubens derer anführt, welche nicht in die Ruhe GOttes eingehen wollten, indem sie gleichsam mit Fingern auf Kades-Barnea, die unruhige Heiligkeit, zeigt. Doch haben die Juden gerade dies auch in der ganzen Zeit des Gesetzes gethan, so oft sie den heiligen Propheten nicht glaubten, die ihnen befahlen, von Kades-Barnea auszugehen zu der wahren Gerechtigkeit, zu der Gerechtigkeit des Glaubens; besonders aber ist dies unter Christo geschehen.

„Die Kundschaften“, die ausgeschickt wurden, sind die Meister und Lehrer dieses Volkes, welche die Herrlichkeit und den Reichthum der Verheißung sahen, aber durch Mißtrauen und aus Furcht vor dem Kreuze und Unwilligkeit das
\quellseite{704}
Fleisch zu tödten wieder zurückgingen, und das Volk durch gottlose Lehren vom Glauben abwendig machten. Auch während der ganzen Zeit der Synagoge gingen sie selbst nicht hinein, und wehrten denen, die hineingehen wollten: Da sind Riesen, Enakim, befestigte und große Städte, ein großes und hochgewachsenes Volk. So sagten die Juden [Joh. 11, 48.]: „Damit nicht die Römer kommen und nehmen uns Land und Leute“, denn sie fürchten die Macht der Welt. Was „Enakim“ und andere seien, wollen wir hernach sehen.

\section{Das zweite Capitel}

\begin{schriftzitat}
V. 1–6. Und umzogen das Gebirge Seir eine lange Zeit etc.
\end{schriftzitat}

Dies ist (wie er selbst im Texte sagt) achtunddreißig Jahre, damit die gottlosen Murrer allmälig umkämen und den Heiden Raum zur Buße gegeben würde, und den ungläubigen Amoritern ein Aergerniß, daß sie verstockt würden durch diese Schwachheit des Volkes GOttes. Es war aber diese Reise im steinigen Arabien (denn Idumäa ist der größte Theil desselben) von KadesBarnea an der Grenze des gelobten Landes im Süden zu dem rothen Meer; darnach wurde nach einer langen Reise wieder umgewandt nach Osten zu den Moabitern. Dann folgt:

„Ihr habt das Gebirge nun genug umzogen, wendet euch gegen Mitternacht“ [V. 3.], nämlich indem die Edomiter zur Linken gelassen wurden, und die Moabiter und Ammoniter (welche östlich vom gelobten Lande wohnen) zur Rechten, so daß sie gerades Weges gegen Norden sich wandten zu den Amoritern, darnach aber, nachdem die Amoriter besiegt waren, sich gegen Westen kehrten, und über den Jordan in das Land Canaan gingen, wie man solches auf der Karte (tabula) gut sehen kann.

Moses aber schreibt an dieser Stelle, daß die drei Völker, die Edomiter, die Moabiter und die Ammoniter, von dem HErrn ausgenommen worden seien, während er doch 4 Mos. 20, 21. sagt, daß die Edomiter Israel mit Gewalt den Durchzug verwehrt haben. Aber dies hat sich an verschiedenen Orten zugetragen. Das, was im vierten Buche Mosis geschrieben steht, ist in Kades geschehen, an der Grenze Edoms und der Amoriter, wo sie auf einem kurzen Wege zu den Amoritern hinüber gehen konnten. Aber der HErr verschaffte es, daß die Edomiter sich weigerten. An dieser Stelle aber schreibt er von der Grenze der Edomiter und Moabiter, nach dem sehr langen Umwege im Gebirge Edom, durch den Weg des rothen Meeres, wo jene Völker, schon erschrocken durch die Wunder und die Kraft GOttes, die in Israel erwiesen worden waren, des Schutzes GOttes bedurften. Der HErr aber nimmt sich ihrer an wegen dessen, was ihre Väter Lot und Isaak verdient hatten, die Gott geglaubt hatten, darum wollte er nicht, daß ihr Name ausgetilgt werden sollte. Aber lieblich und freundlich befiehlt er, daß sie die Edomiter nicht überfallen und plündern sollen, indem er ihnen die empfangenen Wohlthaten und Verheißungen GOttes einschärft und sagt:

\begin{schriftzitat}
V. 7–23. Denn der HErr, dein GOtt, hat dich gesegnet etc.
\end{schriftzitat}

Als ob er sagen wollte: Es ist nicht vonnöthen, daß du voll Sorge oder geizig seiest, und laß es dir nicht beschwerlich sein, von ihnen das zu kaufen, was dir nothwendig ist, denn du hast überflüssig, daß du kaufen kannst etc.

Merke aber: wiewohl die Geschichte dieses Volkes, wenn sie im Einzelnen (in specie) angesehen wird, den heidnischen Historien gleich zu sein scheint, denn in gleicher Weise wechselt Sieg und Niederlage, bald ist Ueberfluß bald Mangel da, und wenn die Wunder auch noch so häufig erwiesen werden, so verbirgt sich doch GOtt so sehr, daß es scheint, als ob alles durch menschliche Klugheit und Zufall regiert werde. Ohne Zweifel ist es allen den gottlosen Heiden so vorgekommen, darum haben sie sich auch wider sie gesetzt und so gegen sie gehandelt, als ob kein
\quellseite{705}
GOtt in Israel wäre. Darum\footnote{Hier beginnt der Nachsatz zu dem im Anfange dieses Absatzes begonnenen Vordersatze: „wiewohl die Geschichte dieses Volkes“ etc. Hier sollte nun statt „Darum“ (Ideo) „Dennoch“ (tamen) gesetzt sein.} müssen die Historien dieses Volkes von allen Historien der Heiden so weit geschieden werden, als der Himmel von der Erde ist. In den Historien der Heiden kann man entweder die Größe oder die Kleinheit der Werke sehen, in dieser aber ist nur dies Eine zu bewundern und zu verehren, nämlich das Wort GOttes, durch dessen Leitung und nach dessen Willen alles ausgeführt wird und geschieht. Ja, sie werden in Wahrheit heilige Historien genannt, nicht, weil jene Dinge von heiligen Leuten gethan sind, sondern weil sie nach dem heiligen Worte GOttes gethan sind, welches alle Dinge heiligt, und in seinem heiligen Namen und an GOttes Statt.

Deshalb, obgleich die Thaten aller Heiden ebenfalls nichts als Wunder und Werke GOttes sind, so haben sie doch nicht das Zeugniß des Wortes und des Wohlgefallens GOttes. Darum sind ihre Sachen und Historien nur Zeugnisse (testes) des Zornes [GOttes], sehr furchtbare Kundgebungen des erschrecklichen Gerichts GOttes. Darum sind die Historien der ganzen Welt, wenn man sie auf Einen Haufen nimmt, unvergleichlich geringer als eine einzige und zwar die allergeringste Historie dieses Volks. Und daher kommt es, daß die hochmüthigen und fleischlichen Leute durch die Einfachheit der Dinge in der Schrift betrogen werden, weil sie keine Rücksicht nehmen auf GOttes Wort und nur die Dinge in Anschlag bringen.

„Er kennt deinen Weg“ [V. 7. Zulg.] ist eine der hebräischen Sprache eigenthümliche Redeweise (idiotismus), wie Ps. 1, 6.: „Der HErr kennt den Weg der Gerechten“; Matth. 25, 12.: „Wahrlich, ich sage euch, ich kenne euer nicht.“ Dies ist ein Wort eines vortrefflichen Trostes, durch welches der Glaube angefacht wird, wie wir im Deutschen sagen: „Er nahm sichs an“, er erkannte es und ließ es sich angelegen sein, so daß es mehr das Herz als den Verstand bezeichnet, wie es die Wortfügung und die Worte dieser heiligen Sprache mit sich bringen. Denn wenn du mit der That zeigst, daß du jemanden kennst oder nicht kennst, dann kann man in Wahrheit von dir sagen, du kennest ihn, oder du kennest ihn nicht.

„Elath“ und „Ezengaber“ [V. 8.] sind am rothen Meer gelegene Städte des Königreichs Edom, die auch sonst in der heiligen Schrift gar berühmt sind. „Ar“ [V. 9.] ist die Hauptstadt der Moabiter, wie Rabba die der Ammoniter, und Sela, das ist Stein, die der Edomiter, von der das steinige Arabien seinen Namen hat. Ferner macht Moses hier einen Unterschied unter den Riesen, welche im Hebräischen mit einem allgemeinen Namen Rephaim (\hebraeisch{רפאים}) genannt werden. Diejenigen, welche von den Edomitern aus dem Gebirge Seir vertrieben wurden, nennt er „Horiter“ [V. 12.]; die Moabiter nannten die ihrigen, die sie vertrieben hatten, „Emim“ [V. 10. f.], die Ammoniter die ihrigen „Samsummim“ [V. 20.], die Cananiter die ihrigen „Enakim“ [V. 10.]. Daher sind dieselben Tyrannen bei verschiedenen Leuten mit verschiedenen Namen genannt worden. Damit stimmt 1 Mos. 14, 5., wo er erwähnt, daß von dem Könige der Elamiter die Riesen zu Ashtaroth (welches später zum Königreiche Ogs in Basan gehörte) geschlagen worden seien, und die Susin und Emim und Horiter auf ihrem Gebirge Seir. Desgleichen 1 Mos. 6, 4. nennt er die Riesen vor der Sündfluth Nephilim.

Es haben aber die Nephaim ihren Namen vom Heilen (\hebraeisch{רפא}), nämlich weil jene Vornehmen und Helden des Landes auch Väter des Vaterlandes und Heilande genannt werden wollten. Emim kommt von \hebraeisch{אימה}, welches Schrecken bedeutet; sie werden schreckliche und furchtbare Leute genannt, weil sie mit Gewalt und Waffen wüthen. Horiter (\hebraeisch{חורים}) heißen Edle, Vornehme, Glänzende (candidi), die Besten (optimates), von der weißen Farbe (candore), entweder weil die Fürsten im Morgenlande weiße Kleider tragen, wie die Lateiner das Staatskleid (trabea) und Purpur, oder weil sie vor anderen ansehnlich sind. Sammumsumim kommt her von \hebraeisch{זמם}, welches Verbrechen oder Schandthat bedeutet, nämlich weil sie hervorragende Mörder und Räuber waren, und sich alles unterstanden, wie es sie gelüstete. Enakim kommt her vom Tragen einer Halskette, als ob sie mit einer Halskette versehen gewesen wären (torquati). Nephilim kommt von fallen (\hebraeisch{נפל}), weil sie einfallen, niederwerfen und mit Gewalt unterdrücken.

Von dem, daß Moses schreibt, die „Avim“ [V. 23.] seien von den Cappadociern aus Hazerim vertrieben, weiß man nicht gewiß, was er
\quellseite{706}
damit meine, es sei denn, daß er ein Exempel anführt, durch welches er zeigen will, daß, wenn der HErr gebietet und die Reiche verändert, wie er will, ein Volk das andere vertreibt, auf daß er nicht allein nur der Juden, sondern auch der Heiden GOtt sei, wie er hier sagt [V. 9. 19.]: „Denn ich habe es den Kindern Lot zu besitzen gegeben“, und [V. 5.]: „Sihon habe ich das Gebirge Seïr gegeben.“ Denn, daß fast alle hier aus „Caphtorim“ die Cappadocier machen, weiß ich nicht, ob sie sich durch die Aehnlichkeit des Namens leiten lassen, da die Cappadocier in Asien sind; es sei denn, daß sie wollen, es sei nur eine kleine Mannschaft der Cappadocier wider die Awim aufgebrochen. Ferner die, welche Hieronymus Awim (Evaeos) nennt, sind nicht die Heviter (Hevaei) im Lande Canaan, welche den Kindern Israel übergeben waren, denn diese heißen im Hebräischen Hiphim (\hebraeisch{חוי}), mit n geschrieben, jene aber Awim (\hebraeisch{עוים}), mit g geschrieben. Mit diesem Namen wird Jos. 18, 23. auch eine Stadt im Stamme Benjamin benannt. „Hazerim“ aber scheint dasselbe zu sein als Hazeroth, eine von den Lagerstätten, 4 Mos. 33, 17., davon er auch oben Cap. 1, 1. gesagt hat. „Gaza“ ist sicherlich die Stadt in Palästina. Also scheint es, daß die Awim gewohnt haben zwischen dem Lande Canaan und Egypten, am Rücken des Gebirges Seïr oder Arabiens, angrenzend an die Midianiter.

\begin{schriftzitat}
V. 24. Gehet über den Bach Arnon.
\end{schriftzitat}

Ueberall lauten die Worte, als ob „Arnon“ der Name eines Baches wäre, wiewohl Hieronymus sagt, es sei der Name eines hoch aufsteigenden Felsens zwischen den Moabitern und den Amoriten. Wenn dies wahr ist, so muß der Bach in dem Thale unter dem Felsen gewesen sein, auf den Moses auch 4 Mos. 21, 28. anspielt, wo er erzählt, daß man im Volksliede gesungen habe: „Das Feuer hat gefressen die Bürger der Höhe Arnon“, und am Ende dieses Capitels [V. 36.] scheint er „den Bach“ von „Arnon“ zu trennen, da er von der Stadt sagt, die am Bache liegt, und nicht hinzufügt „Arnon“.

\begin{schriftzitat}
V. 24. 25. Siehe, ich habe Sihon in deine Hände gegeben etc.
\end{schriftzitat}

Es wird die Verheißung erneuert, und GOttes Wort wird vorausgeschickt, nach dessen Befehl und Leitung die Sache im Glauben ausgeführt werden sollte, nicht durch Vermessenheit oder menschlichen Rath, wie wir gesagt haben. Es ist aber eine herrliche Verheißung, daß der König Sihon in ihre Hände gegeben sei, damit sie des Sieges gewiß wären, und nicht allein dies, sondern daß er vorhersagt, daß auch der Schrecken des Namens Israel über alle Völker kommen werde, so daß sie überall sicher sein können, es würden weder von der Seite noch vom Rücken her Feinde da sein, die den Amoritern zu Hülfe kämen, und die Israeliten sich nichts angelegen sein lassen sollten, als wie sie die Amoriter umbrächten.

\begin{schriftzitat}
V. 26–29. Da sandte ich Boten aus der Wüste Kedemoth.
\end{schriftzitat}

„Kedemoth“ wird Jos. 13, 18. unter den Städten des Stammes Ruben aufgezählt, wiewohl andere aus dieser Stelle [entnehmen] wollen, daß es auch der Name einer Wüste sei. Ich aber glaube, daß Moses so sagen wolle: Ich habe Boten gesandt aus der Wüste, die zu der Stadt Kedemoth gehört, so daß es allein der Name der Stadt ist, nämlich daß Moses will, er sei nicht in der Stadt gewesen, sondern in der Wüste bei der Stadt, nämlich im Anfang des Gebietes der Amoriter.

Du siehst aber hier, daß auch den Feinden zuerst Friede angeboten werde, von denen doch 5 Mos. 20, 16. geboten wird, daß sie völlig vertilgt werden sollen, so daß alles getödtet werde, weshalb auch die ganze Gemeine (vulgus) im Buche Josua Cap. 9, 18. wider die Obersten murrte, weil sie einen Bund mit den Gibeoniten gemacht hatten. Dies muß daher so verstanden werden, daß sie zuerst allen den Frieden anbieten mußten; wer den ausschlug, sollte umgebracht werden; diejenigen aber, welche sich ergaben, sollten zinsbar werden [5 Mos. 20, 10. ff.].

\begin{schriftzitat}
V. 30–35. Der HErr, dein GOtt, verhärtete seinen Muth.
\end{schriftzitat}

Hier wird der Ruhm und die Kraft des freien Willens gepriesen, nämlich, daß er nichts Anderes ist als nichts, wie es auch Röm. 1, 24. heißt: „GOtt hat sie dahingegeben in ihrer Herzen Gelüste“, und Cap. 9, 18.: „Er verstockt, welchen er will.“ Wenn er also irgend etwas vermag, wohlan, so richte er das aus, daß er vom HErrn nicht verstockt werde. Verstocken aber heißt, zulassen, daß die Herzen steif, stolz
\quellseite{707}
und sicher werden, aber nicht erschrecken, daß sie sich aller Dinge vermessen und weder Zorn noch Gericht fürchten.

\begin{schriftzitat}
V. 36. 37. Von Arver an, die am Ufer des Bachs bei Arnon liegt.
\end{schriftzitat}

Er erwähnt zwei Städte, Arver, die höher gelegen ist als das Ufer des Baches (super ripam), die andere am Bache. Dies verstehe ich so, daß Arver auf dem Gipfel des Berges Arnon gewesen sei, wie Hieronymus sagt; dieser Berg hat einen Bach, der unten im Thal daran hinfloß, an welchem die andere Stadt lag. Arver aber gab er [Moses] dem Stamme Gad, Jos. 13, 24.

Beachte hier, wie es scheint, als ob in Mose eine große Lüge vorkomme. Der König Sihon hatte nur ein Reich, welches so groß war, daß zwei Stämme, Ruben und Gad, es einnehmen, das heißt, ein Volk, welches etwa sechsmal so klein war, als das Mosis. Was ist das also für ein herrlicher Sieg, daß das so zahlreiche Volk Israel ein so kleines Volk überwunden hat? Was bedurfte es hier der Verheißungen GOttes? oder wo ist hier der so schöne Schein eines großen Wunders? Aber, wie ich gesagt habe, so hatte er diese vierzig Jahre lang sein Volk in der Wüste gedemüthigt, und gleichsam als ein fliehendes regiert; darnach hatte er zugelassen, daß es von den Amalekitern und Cananitern bisweilen geschlagen wurde. Darum wurden die Feinde, die Amoriter, aufgeblasen und vermaßen sich, sie mit geringer Mannschaft zu besiegen. Inzwischen aber hatte das Volk GOttes gelernt, sich allein auf die Kraft GOttes zu verlassen. So regiert der HErr immer sein Volk, daß es den Schein hat, als sei es völlig besiegt, und die Feinde desselben hätten den Sieg gänzlich erlangt, und gerade dadurch erfüllt er seine Wundertaten.

Die heimlichen Deutungen der Könige und der Namen haben wir im ersten Capitel gesehen. Und das Tödt(en) bedeutet sicherlich das Tödt(en) der Sünde durch das Wort GOttes.

\section{Das dritte Capitel}

\subsection[{(V. 1–17. [V. 2.] Fürchte dich nicht etc.)}]{(V. 1–17. [V. 2.] Fürchte dich nicht etc.)\footnote{Von uns eingefügt.}}

Wiederum wird eine Verheißung vorausgeschickt, als sie mit Og, dem Könige von Basan, kämpfen sollen, damit die Historie durch das Wort GOttes geheiligt sei, und die Sache allein durch die Weisheit GOttes ausgerichtet werde. Alles andere ist klar, denn die Namen der Orte und der Berge, Argob, Basan, Gilead, Hermon, welcher auch Sirion und Senir genannt wird [V. 9.], und die übrigen kommen auch in den folgenden Büchern sehr oft vor.

Es muß der König Og ein besonders großer Riese gewesen sein, welcher es allein, auch nachdem er gehört hatte, daß der König Sihon geschlagen war, wagte, mit den Israeliten zu kämpfen, ist auch dem Könige Sihon nicht zu Hülfe gekommen, nämlich weil er vertraute auf die Größe seiner Stärke, wie hier geschrieben steht, daß er neun Ellen\footnote{Eine Elle ist das Maß vom Ellbogen eines Mannes bis zu der Spitze seines Mittelfingers, welches anderthalb Fuß lang ist.} lang gewesen sei. Und ohne Zweifel war nicht allein er so groß, sondern sein Reich ist durch solche Riesen regiert worden, deren jeglicher sich vermaß, viele aus Israel zu tödten, wie sie 4 Mos. 13, 34. von „Enaks Kindern“ (Enakim) sagen: „Im Vergleich zu ihnen erschienen wir als die Heuschrecken.“ Wenn der Einwohner viele waren, so ist es glaublich, daß auch der Riesen des Königs Og viele waren, damit das Werk des Glaubens in dem Volke Israel und die Wunder der Verheißung GOttes desto augenfälliger würden. Denn unter Saul flohen vor Einem Riesen alle Israeliten, was Og hier auch hoffte, und es würde geschehen sein, wenn nicht der Glaube des Volkes und die wahrhaftige Verheißung GOttes ein Wunder gewirkt hätte.

Das, was wir verdolmetscht haben: Sub scatebris Pisga [unter den Quellen Pisga], hat Hieronymus: Ad radices Phasga [am Fuße des Phasga] und lautet im Hebräischen: Asedoth Pisga. Und es ist noch nicht gewiß, ob es ein Eigenname oder ein Gattungsname sei, wiewohl ich dazu geneigt bin, daß ich fast glauben
\quellseite{708}
möchte, es sei der Name einer Stadt, begegnen durch Jos. 13, 20. [Bulga.], wo unter den Städten des Stammes Ruben Asedoth Pisga aufgezählt wird, das heißt, Asedoth am Fuße des Berges Pisga. Aber auch Hieronymus sagt in dem Buche „von den hebräischen Ortschaften“ (locis), daß es der Name einer Stadt sei. Sonst bedeutet es hervorsprudelndes Wasser, Quellen, Ergiebung.

\begin{schriftzitat}
V. 18–20. So ziehet nun gerüstet vor euren Brüdern, den Kindern Israel, her.
\end{schriftzitat}

Dieses sagt er zu denen vom Stamme Ruben, Gad und Manasse. Aber dies lautet, als ob er wollte, daß sie aus ihren Kräften den Krieg führen sollten. Aber dies geschieht, damit Raum sei für den Glauben und die göttliche Verheißung, und damit Gott nicht kärglich versucht werde. Denn wiewohl GOtt alle Menschenkinder schafft, ernährt, erhält und regiert, will er dennoch nicht, daß irgend jemand müßig sei, sondern er hat die Glieder gegeben, durch welche einer den andern zeugen, ernähren und regieren soll, durch welche Gottes Gelegenheit gegeben wird zu schaffen, zu nähren und zu regieren, damit so durch den Glauben wahrgenommen werde, daß unter dem menschlichen Werke die verborgene Kraft GOttes alles wirke, und die Ungläubigen betrogen werden, welche allein das Werk des Menschen ansehen und die Kraft GOttes nicht kennen.

So wird ihnen auch hier geboten, daß sie der Waffen gebrauchen, und doch wird ihnen verboten, auf die Waffen zu vertrauen. Sie gehen in GOttes Werke einher mit ihren Kräften, und dennoch sind sie nicht vermessen auf ihre Kräfte. Was sind also da, wo der Glaube und das Wort herrscht, die menschlichen Kräfte anders als eine Art Larven GOttes, unter welchen verborgen er seine Wunder wirkt, während er durch die Schwachheit derselben die Hoffärtigen, Starken, Weisen und Heiligen gegen sich herausfordert? und wenn sie angelaufen sind, macht er sie bald zunichte und triumphirt über sie und trägt sie öffentlich zur Schau (Col. 2, 15.). Wohl denn, dessen Glieder und Waffen Gottes so dienen! Ja, wer da erkennt, daß GOttes Macht und Weisheit eine solche ist, der vertraut gänzlich, nicht auf die Larve GOttes, sondern auf das Wort, welches unter der Larve mächtig ist, und vermag und thut große Dinge, ja, alles in dem HErrn. So wird es geschehen, daß er frei und sicher wandeln kann, gleichweit von (in medio) Vermessenheit und Verzweiflung, daß er GOtt nicht versuche, da er das hat, dessen er gebrauchen kann, wiederum auch nicht verzweifeln, wenn ihm das mangelt, dessen er gebrauchen könnte. Denn GOtt, welcher mit ihm wirkt unter den Dingen, die er hat, wird mit ihm auch wirken ohne diese Dinge, wenn er nichts hat, so daß er in solcher Weise weder Ueberfluß noch mehr hat, wenn alles vorhanden ist, noch Mangel leidet oder weniger hat, wenn ihm alles gebricht. Denn er hat GOtt auf beiden Seiten, wie Paulus [Phil. 4, 12.] sagt: „Ich habe gelernt, Ueberfluß zu haben, ich habe gelernt, Mangel zu leiden“ etc. Die Gottlosen aber werden aufgeblasen durch die vorhandenen Güter; wenn sie fehlen, werden sie niedergeschlagen, weil sie allein die Larve der Güter ergreifen; die übrigen, welche das, was ihnen gegeben ist, nicht gebrauchen wollen, versuchen GOtt.

\begin{schriftzitat}
V. 21–29. Und Josua gebot ich etc.
\end{schriftzitat}

Er wiederholt die Verheißungen der künftigen Siege und die Einsetzung eines Führers, und zugleich schärft er das Vergangene ein, um ihren Glauben zu stärken sowohl durch die Verheißung künftiger Dinge als auch durch das Vorhalten der vergangenen, da er selbst nicht mit ihnen über den Jordan gehen sollte, dies auch von dem HErrn durch seine Bitten nicht erlangen konnte. Und so hat er in diesem Capitel die angefangene Erzählung beschlossen, in welcher er die ganze Historie kurz wiederholt hat. Im folgenden Capitel fängt er an zu ermahnen und das Gesetz auszulegen, und es beginnt das Deuteronomium.

Aber warum wird Moses nicht erhört, da es doch wahrscheinlich ist, daß er im Gebete gebetet habe? Dies ist uns zum Exempel und zum Troste geschrieben. Denn obgleich der HErr ihn nicht erhört, wodurch Moses inne ward, daß er erzürnt auf ihn sei, wie er hier sagt, so verläßt er ihn doch nicht, befiehlt ihm, auf den Berg zu steigen und das Land anzusehen, und dem Josua den Befehl zu geben. So sollen auch wir uns nicht wundern, wenn wir nicht erhört werden, da wir nicht wissen, wie wir beten sollen, zugleich aber auch sollen wir auf
\quellseite{709}
keine Weise zweifeln, daß wir GOtte angenehm und theuer seien, und sollen seine Güte unter dem Zorn ergreifen, damit wir nicht verzagt werden.

\subsection{Heimliche Rettung}

Daß Moses nicht über den Jordan geht in das gelobte Land, sondern ihm befohlen wird, dies einem anderen Führer zu überlassen, das ist das große Geheimniß, daß, wie es im Briefe an die Hebräer [Cap. 7, 19.] heißt, daß Gesetz nichts vollkommen machen konnte. Denn das Gesetz gibt auch nicht den Geist der Gnade, darum führt es nicht in das Reich GOttes ein, wo das rechte Erbe des HErrn ist. Moses tödtet zwar zwei Könige jenseit des Jordans, das heißt, das Gesetz demüthigt die Sünder und zeigt, daß sie getödtet sind, und läßt nichts lebendig, das heißt, kein Vertrauen auf irgend eine Sache, denn es führt in die Hölle und ist ein Amt des Todes, läßt aber das Vieh und die Güter unverletzt, welche zur Beute werden. Denn das Gesetz tödtet den Menschen nicht leiblicher Weise, sondern nur die Zuversicht des Herzens; wenn diese getödtet ist, so kommt alles, was der Mensch ist und hat, in den Dienst des Gesetzes als eine Beute, und es lebt der elende Mensch mit allen seinen Kräften als ein Gefangener des Gesetzes. Aber weil es die Vergebung der Sünden nicht predigt, welche es durch seinen Dienst rege gemacht hat, so muß es im Lande Moab sterben außerhalb des Reiches GOttes. Denn es gebührt sich nicht, daß man in dem Volke der Freiheit den Zwang des tödtenden Gesetzes lehre, sondern die Gabe des Geistes, der da lebendig macht. Darum fällt hier das ganze Gesetz, so daß nichts von ihm übrig ist, daß auch der Mensch nicht weiß, wo es begraben sei. Josua aber, welcher sowohl dem Namen als auch der That nach Christum bedeutet, wiewohl er ein Diener Moses gewesen ist, führt er dennoch, nachdem sein Herr gestorben ist, das Volk hinein und theilt das Erbe des HErrn aus. So hat Christus, der zuerst unter das Gesetz gethan ist, demselben für uns gedient; endlich aber, da es ein Ende nahm, hat er ein anderes Amt, das des Evangelii, eingesetzt, durch welches wir von ihm in das geistliche Reich eines fröhlichen und sicheren Gewissens gegen GOtt geführt werden, wo wir ewiglich herrschen.

\section{Das vierte Capitel}

\begin{schriftzitat}
V. 1. Und nun höre, Israel, die Gebote und Rechte \&c.
\end{schriftzitat}

Nachdem nun die Erzählung und kurze Wiederholung beendet ist, hat er in diesem Capitel eine Ermahnung an die Hörer des Gesetzes gerichtet, welches er im folgenden Capitel anfangen wird. Es ist aber eine gar reiche Ermahnung und, wie denn die Art dieser Sprache ist, auch in vielen Worten. Denn er wiederholt hier, wie auch anderswo, dasselbe so oft, daß, wenn ein fleischlicher Leser dabei wäre, der nicht versteht, um was es sich handelt, derselbe geärgert werden müßte. Aber uns kommt es zu, unsern Ueberdruß niederzuhalten, weil wir wissen, daß es GOttes Worte sind, die wir hören. Aber GOtt hören ist Seligkeit, wenn er auch fort und fort nur eine und dieselbe Silbe spräche, wie viel mehr, wenn er so mancherlei und so fleißig redet? Darum gebühret es uns, nicht bloß auf die Worte, sondern viel mehr und vor allen Dingen auf die Person dessen, der da redet, nämlich GOttes, mit Furcht Acht zu haben.

Es dient aber diese häufige Wiederholung dazu, daß wir, die wir wegen unseres Fleisches immer fallen, immer kalt sind und faul zum Guten, immer geneigt und aufgelegt zum Bösen von Jugend auf, unablässig einen Zuchtmeister und Treiber haben, der die Sünde im Zaume halte. Zugleich wird hierin die eigentliche Art (virtus) des Gesetzes gezeigt, dem es eigen ist, daß es anhalte, dringe und, bis daß die Gnade zu Hülfe kommt, nicht ablasse mit Müdemachen, Anklagen, Verdammnen, wie der Apostel sagt [Röm. 4, 15.]: „Das Gesetz richtet nur Zorn an.“

Im Hebräischen sind ungefähr sechs oder sieben Wörter für Gesetz: \hebraeisch{תּוֹרָה}, \hebraeisch{מִצְוָה}, \hebraeisch{חֻקָּה}, \hebraeisch{פִּקּוּד}, \hebraeisch{עֵדוּת}, \hebraeisch{מִשְׁפָּט}, \hebraeisch{מִשְׁמֶרֶת}.

\quellseite{710}
Erstlich \hebraeisch{תורה} bezeichnet eigentlich das Gesetz, und hat seinen Namen von lehren (\hebraeisch{ירה}), daher wird es Lehre\footnote{ab instruendo . . . vocatur instrumentum.} genannt.

\hebraeisch{מצוה} ist ein Gebot, wie in Mose öfter gelesen wird: „Wie mir der HErr geboten hat.“

\hebraeisch{חקה} bezieht man auf Ceremonien, und das gefällt mir wohl. Denn eigentlich bedeutet es den Brauch und die Sitten in heiligen Dingen, welche ich Gebote (statuta) nenne, wie im 2. Psalm, V. 7.: „Ich will das Gebot des HErrn erzählen“, wo es im Hebräischen heißt: „Ich will predigen von dem Brauch und von der Satzung“, nämlich von der neuen Weise und der rechten Ceremonie, die da ist, an Christum, den Sohn GOttes, glauben, wie ebendaselbst folgt, in welchem überhaupt alle (omnes omnium) Ceremonien erfüllt werden, und auf den sie sich beziehen.

\hebraeisch{פקד}, das ist eigentlich, was die Lateiner mandatum nennen, auf deutsch: „befehlen“. Denn das Verbum \hebraeisch{פקד} bedeutet besuchen, einsetzen, ordnen, urtheilen (censere) und überhaupt gebieten. Daher kommt im 109. Psalm, V. 8., \hebraeisch{פקדה}, was man griechisch ἐπισκοπή oder διοίκησις nennt: „Und sein Bisthum wird ein anderer empfangen“ [Apost. 1, 20.], das heißt, das ihm befohlene Amt oder Regierung. Unser Uebersetzer gibt es fast überall durch Rechtfertigung wieder.

\hebraeisch{עדות} ist Zeugniß, daß GOtt durch sein Gesetz und Wort bezeugt, daß er zugegen sei und für uns sorge, damit wir ein Bewußtsein (conscientiam) von GOtt haben, den wir sonst weder sehen, noch fühlen, noch erkennen.

\hebraeisch{משפט} übersetzen sie beständig durch Gericht, und nicht übel, nur daß es vonnöthen ist zu wissen, wie weit sich dies Wort im Hebräischen erstreckt. Denn es bedeutet ebenfalls das Amt oder die Verwaltung des Rechts, wie im ersten Psalm, V. 5.: „Die Gottlosen werden nicht bestehen im Gericht“, wo einige sagen würden: im Regimente. Desgleichen im 72. Psalm, V. 1.: „GOtt, gib dein Gericht dem Könige.“ Daher hat auch das Buch der Richter seinen Namen, welches im Lateinischen das Buch der Fürsten, der Regenten, der Obersten (consulum) oder der obrigkeitlichen Personen genannt werden würde. Desgleichen Ps. 110, 6.: „Er wird richten unter den Völkern“, das heißt, er wird herrschen. Zweitens bezeichnet es das Werk selbst und die Vollstreckung des Urtheils. Daher findet sich in der Schrift überall das Wort: Gericht und Gerechtigkeit schaffen [Ps. 103, 6.]: „Der HErr schaffet Gericht allen, die Unrecht leiden“, und [Ps. 17, 2.]: „Von deinem Angesichte komme mir mein Recht“ und: „Er hat Gerichte mit ihm geredet.“

Endlich bedeutet es eigentlich das, was im Lateinischen jura heißt, nämlich die Gesetze, durch welche die Menschen unter einander im bürgerlichen Leben in weltlichen Angelegenheiten zu Schlichtung ihrer Händel kommen (disponuntur), daher es auch öfter für Gewohnheit und Brauch genommen wird, als, für das Völkerrecht, oder das natürliche Recht, wie 2 Mos. 22, 17.: „Nach dem Recht der Töchter soll er mit ihr handeln“,\footnote{In der Jenaer steht: facit statt: faciet.} und Ps. 119, 132.: „Wie du pflegst zu thun (secundum judicium) denen, die deinen Namen lieben.“ Diese weite Bedeutung des Wortes erzeugt überall in der Schrift Dunkelheit und Zweideutigkeit. Deshalb ist vonnöthen, daß man wohl Acht habe.

\hebraeisch{משמרת} heißt etwas beobachten (observantiam) oder Acht darauf geben (custodiam), nämlich ein allgemeiner Name; der über alle zuvor genannten geht, ja, wenn wir es genau betrachten, so werden nur zwei oder drei Wörter sein, welche eigentlich das Gesetz bezeichnen, nämlich \hebraeisch{חקה} und \hebraeisch{משפט}, Gebot (statutum) und Recht, deren eines das lehrt, wodurch wir zu einem rechten Verhalten gegen GOtt angewiesen werden (ordinamur), nämlich \hebraeisch{חקה}, das andere das, wodurch wir zu einem rechten Verhalten gegen die Menschen angewiesen werden, nämlich \hebraeisch{משפט}. Und diese zwei sind genugsam, und dieser Unterschied scheint mir bewährt und beständig zu sein. Denn so beginnt auch Moses in diesem Capitel: „Höre die Gebote und Rechte, die ich euch lehre“, indem er mit diesen beiden Worten gleichsam alles begreift, als ob er sagen wollte: Ich will euch lehren, wie ihr recht leben sollt vor GOtt und Menschen.

Die andern Wörter sind nun alle allgemein, und zeigen die Bedeutung (vim) und das Wesen der Gebote und der Rechte an, nämlich \hebraeisch{תורה}, daß sie uns belehren und unterweisen,\footnote{Statt instruant und erudiant in der Erlanger ist mit der Wittenberger und der Jenaer instruant und erudiant zu lesen.} damit
\quellseite{711}
wir nicht Thiere seien und in Finsterniß handeln. \hebraeisch{מִצְוָה}, daß sie uns binden und zwingen, damit wir nicht frech (libere) das verachten, was uns gelehrt worden ist. \hebraeisch{פִּקּוּדִים}, daß sie uns geboten werden zum Halten, damit wir etwas haben, womit wir umgehen, auf daß wir nicht müßig seien. \hebraeisch{עֵדֹת}, daß sie (wie ich gesagt habe) uns zum Zeugniß seien, daß GOtt sich unser annehme, und wir uns sicher darauf verlassen, daß er überall und immer bei uns gegenwärtig sei. \hebraeisch{מִשְׁמֶרֶת}, daß sie erfordern, daß man sorgfältig und achtsam sei ihnen gegenüber, damit wir nicht träge und schläfrig seien, sie zu thun.

Außer diesen sind auch diese Ausdrücke: \hebraeisch{דָּבָר}, \hebraeisch{נְאֻם}, \hebraeisch{אִמְרָה}, Wort, Rede, Ausspruch; desgleichen \hebraeisch{דָּתָן}\footnote{Statt \hebraeisch{דָּתָן} sollte wohl \hebraeisch{דַּעַת} gelesen werden, welches sich Ps. 119, 66. findet. Es kommt wohl \hebraeisch{דָּת}, aber nicht \hebraeisch{דָּתָן} vor.} und \hebraeisch{טַעַם}, aber diese beiden sind selten im Gebrauch, als Ps. 119, 66.: „Gutes, \hebraeisch{טעם} und Erkenntniß, worfür wir [in der Vulgata] bonitatem et disciplinam haben, und bezeichnet fast das Verhalten (gestum), welches wir Gebärden und gesittetes Betragen (habitus corporis) nennen, wie es im 34. Psalm heißt: „David verstellte seine Gebärde vor dem Könige Achis“, wie der Titel dieses Psalms sagt.

Es sind also nur diese zwei übrig: \hebraeisch{חֻקָּה} und \hebraeisch{מִשְׁפָּט}, (Gebote und Rechte, die in der Schrift sehr häufig sind, von denen die Richter und die richterliche Gewalt herkommt. Aber von jenem Worte [\hebraeisch{חֻקָּה}] kommt der Meister, zu regieren in geistlichen Dingen, den die Schrift 1 Mos. 49, 10. \hebraeisch{רעה} nennt: „Noch ein Fürst von seiner Hüfte“ [Vulg.], wo der hebräische Text hat: „Noch ein Lehrer von seinen Füßen“, das heißt, es wird in Juda immer ein solcher sein, der die göttlichen Dinge lehrt und Schüler zu seinen Füßen hat, und wird von ihnen nicht weggenommen werden, bis daß Silo komme.

\begin{schriftzitat}
V. 2. Ihr sollt nichts dazu thun, das ich euch gebiete, und sollt auch nichts davon thun.
\end{schriftzitat}

Ehe Moses irgend etwas gebietet oder lehrt, trägt er vor allen Dingen dafür Sorge, daß es unverletzt und unverändert bewahrt werde, damit nicht irgend etwas dazu oder davon gethan werde. Und diese Stelle verdammt kurzweg alle menschlichen Gesetze, und entscheidet, daß in Sachen des Gewissens nichts gelte als GOttes Gesetz und GOttes Wort. Und hierin sollte der Wille GOttes genugsam sein, welcher so urtheilt, wiewohl auch die Vernunft und die Noth ebendaselbe erfordert. Denn da wir durch Adams Sünde in Blindheit ersoffen sind, so daß wir GOtt ganz und gar nicht kennen in allen seinem Willen und Rathe, so ist es nicht nur thöricht, sondern auch unmöglich, daß wir aus uns selbst das Licht und den Weg zurichten, auf dem wir zu GOtt gelangen und das ewige Leben können, was er will, daß wir thun sollen, wie der [Schreiber] im Buch der Weisheit Cap. 9, 14. 13. sagt: „Die Gedanken der sterblichen Menschen sind mißlich und ungewiß. Denn welcher Mensch kann wissen, was GOtt will?“ Wenn daher GOtt nicht sein Gesetz offenbart, durch welches er uns seinen Willen bekannt macht, so bleibt nur diese Auslage über uns (nostrum elogium) übrig, Ps. 116, 11.: „Alle Menschen sind Lügner“, [Ps. 39, 6.:] „alle Menschen sind gar nichts“; und Jes. 55, 9.: „So viel der Himmel höher ist denn die Erde, so sind auch meine Wege höher denn eure Wege.“ Darum ist es ein überaus thörichter Beweisgrund der Sophisten für den freien Willen: Ich habe die Freiheit, eine Kuh zu regieren, Geld wegzuwerfen, also habe ich die Freiheit zu thun, was GOtte gefällt, und ihm zu dienen. Das ist dasselbe, als wenn du sagen würdest: Ich kann mit den Füßen die Erde treten und auf der Erde einhergehen, also kann ich auch im Himmel gehen und auf den Wolken wandeln.

Es hat zwar der Mensch von der ersten Einsetzung [1 Mos. 1, 26.] her die Kenntniß und das Vermögen, über das, was geringer ist als er, zu herrschen und frei damit zu handeln, aber sich selbst zu beherrschen, und zu thun, was dem gefalle, der über ihm ist, das weiß er nicht, vermag es auch nicht. Da hört der freie Wille auf; [da herrscht er nicht,] sondern ist nothwendiger Weise blind, unvermögend, ja, todt und verdammt. Daher darf er sich hier keines Dinges vermessen, oder etwas erwählen, was er thun wolle, sondern muß sich einzig und allein auf das Wort GOttes gründen, auch nichts dazu oder davon thun, so daß er, wenn er nicht ganz gewiß ist, daß GOttes Wort da sei, durchaus nicht darnachgehe, wo er aber gewiß ist, daß GOttes Wort da sei, nicht davon abgehe. O eine nothwendige Ermahnung!

Aber du wirst sagen: Warum sind denn so
\quellseite{712}
viele Bücher der Propheten hinzugethan? Und David verordnete, daß die Hütte Mosis abgethan und ein Tempel gebaut werden sollte, und setzte die Ordnungen der Leviten ein, und Josua that vieles und befahl es dem Volke, und so kann man auch von vielen anderen Dingen sagen. Ich antworte: Das heißt nicht zu dem Worte Mosis hinzu oder davon thun, wenn jemand eben dasselbe mit anderen oder mehr Worten lehrt. So hat Moses die Propheten, nicht sie selbst sich gelehrt und ausgelegt, und nach seinem Gesetze haben sie geurtheilt und die Könige und das Volk gestraft.

Sodann sagt Moses hier nicht von GOtte, sondern vom Volke: „Du sollst nichts dazu thun“ etc. Denn wer zweifelt, daß GOtt selbst nach den Umständen der Zeit dazu oder davon thun könne? Denn er, mag er nun dazu oder davon thun, bleibt immer wahrhaftig, und sein Wort ist immer wahrhaftig, und leitet und behütet uns treulich. So bezeugt David im ersten Buch der Chronika [Cap. 28, 6. ff.], daß er nicht nach seinem Willen (denn das hatte GOtt verworfen, wiewohl er von dem Propheten Nathan gelobt ward [1 Chron. 18, 2. ff.]), sondern auf den Befehl GOttes geboten habe, daß der Tempel gebaut werden solle, und die Leviten verordnet habe. So hatte Josua den Befehl GOttes, daß er gebieten sollte, was zu gebieten war. Wenn nun in solcher Weise alle Propheten irgend etwas Anderes lehrten, so hat ihnen allen der HErr gerade dies offenbart, ebenso wie dem Moses, oder wie Petrus [2. Ep. 1, 21.] sagt, sie sind getrieben worden (inspirati sunt) von dem Heiligen Geist, daß sie reden mußten.

\begin{quote}
Auf daß ihr bewahren möget die Gebote des HErrn, eures GOttes.
\end{quote}

Siehe hier Mosen an, welcher geltend macht, daß man zu den Worten GOttes nichts hinzu noch davon thun müsse, damit nicht das daraus erfolge, daß die Gebote GOttes nicht bewahrt werden. Durch diesen Ausspruch zeigt er, daß die Satzungen der Menschen in heiligen Dingen wider die Worte GOttes streiten, da sie sich dawider setzen, und das Aergerniß veranlassen, daß die Gebote GOttes nicht gehalten werden. Deshalb macht ihnen auch Paulus Tit. 1, 14. einen üblen Namen, indem er lehrt, daß sie die Wahrheit verkehren, und Röm. 16, 18., daß sie die unschuldigen Herzen verführen, und Christus befiehlt [Matth. 16, 6.], daß wir uns hüten sollen vor dem Sauerteig der Pharisäer und [Matth. 7, 15.] vor den Wölfen in Schafskleidern.

\begin{schriftzitat}
V. 3. 4. Eure Augen haben gesehen x.
\end{schriftzitat}

Moses hat seine Ermahnung aus mancherlei Dingen zusammengesetzt, damit sie die Zusätze und Fälschungen der Menschen verwerfen, und dem reinen Worte GOttes, und zwar allein, anhangen möchten. Zuerst [nimmt er die Ermahnung her] von dem ganz nahe liegenden Exempel des Baal-Peor; da er [GOtt] wegen dessen Verehrung, weil sie über (ultra) das Wort GOttes von Menschen erfunden war, so viele Kinder Israel verderbte [4 Mos. 25, 9.], so zeigt er genugsam an, wie er durchaus nichts zu seinen Geboten hinzu gethan wissen will.

Es gilt auch nicht, wenn jemand für die Satzungen der Menschen geltend machen will: Die Moabiter seien Götzendiener gewesen und hätten den Teufeln geopfert, wir aber dienen mit unseren Ceremonien dem wahren GOtte, und sie seien erfunden in gottloser Meinung und aus einem guten Eifer, deßhalb bei unsern Vornehmen (studia) dem jener Leute nicht ähnlich.

Ich antworte: Die Moabiter und andere Heiden haben den Teufeln nicht so gedient, daß sie wußten, es wären Teufel, sondern glaubten, daß sie dem wahren GOtte dienten, nicht weniger als die götzenbieterischen Juden, ja, als alle unsere Papisten, auch die heiligsten und geistlichsten, sondern ihre GOttlosigkeit war das, daß sie den rechten Namen des wahren GOttes annahmen, und ihn verehrten mit Gebräuchen, die GOtt nicht geboten, sondern welche sie erfunden hatten.

Denn diejenigen, welche solches thun, denken in ihrem Herzen, der wahre GOtt sei so beschaffen: er wolle verehrt werden, und es sei ihm angenehm, was die Menschen, auch die noch gottlos sind, ohne GOttes Wort vornehmen. Aber er ist nicht ein solcher GOtt, dem gottlos Wesen gefällt, wie der fünfte Psalm, V. 5., sagt, sondern er ist der GOtt der Gerechten, und will von gottseeligen und gerechten Leuten verehrt werden, die durch sein Wort und seine Gnade, nicht durch ihre eigenen Kräfte gerechtfertigt sind. Daher sind die Gedanken jener Leute von GOtt nothwendiger Weise eine bloße Erdichtung und Lüge, und der Gott, den sie sich in solcher Weise erdichten und bilden, ist nicht der wahre GOtt,
\quellseite{713}
sondern ein Götze ihres Herzens, unter welchem sie den Teufel, den Meister und Vater dieser Lüge, anbeten. Und so verehren sie denn in Wahrheit unter dem rechten Namen des wahren GOttes in falscher Einbildung Götzen und Teufel.

Darum muß man in der Schrift „andere (alienos) Götter“ nicht so verstehen, als ob deren Verehrer den Namen des rechten GOttes gänzlich verleugnet hätten, ja, sie nahmen denselben für sich aufs hartnäckigste in Anspruch, wie wir in den Propheten lesen, sondern weil sie mit Beibehaltung des Namens des wahren GOttes ohne rechte Erkenntniß GOttes (welche das Fleisch und der Mensch aus sich nicht haben kann) sich Gedanken von GOtt machten, welche ihnen recht zu sein schienen, als ob GOtt sich eine Gestalt geben ließe und veränderlich wäre, nach den Einbildungen und Dünkel (visiones) ihres Herzens, während er doch nicht gestaltet werden kann (informabilis) und unveränderlich ist. Wo daher das Wort GOttes nicht ist, da ist auch keine rechte Erkenntniß GOttes; wo die Erkenntniß GOttes nicht ist, da ist gottlose Unwissenheit, Einbildung, Wahn vom rechten Gotte, als ob er Gefallen habe an diesem und jenem, was wir nach unserem eigenen Sinne (proprie) erwählen. Aber dies alles ist gottlose Verehrung, durch welche der wahre GOtt nicht erreicht wird, sondern an seiner Statt und unter seinem Namen die Vorstellungen (phantasmata) und Götzen des Herzens. Dies aber sind andere Götter, nämlich die Teufel, welche diesen Wahn eingeben.

Dann geht man auch nach außen vor und errichtet nach diesem Wahn Ceremonien, Lehre und Regiment, denen sie verschiedene Namen geben; und sie nennen GOtt bald Baal, bald Astaroth, bald Dagon etc., wie auch wir gethan haben, indem wir den Benedictus, Franciscus, Dominicus und vor allen den Papst einen Gott nennen, weil wir uns dünken lassen, durch alle diese Gebräuche und Lehren Gotte zu gefallen. Dieser Wahn ist der lügenhafte Götze des Herzens, unter welchem an der Statt und im Namen GOttes der Satan da sitzt und angebetet wird.

Das ist es, was der Apostel Röm. 1, 21. berichtet: „Da sie wußten, daß ein GOtt ist, und haben ihn nicht gepriesen als einen GOtt, noch gedankt, sondern sind in ihrem Dichten eitel geworden, und ihr unverständiges Herz ist verfinstert“ etc. Ich sage, sie haben GOtt erkannt, das heißt, sie haben Kunde gehabt, daß irgend eine Gottheit wäre, was ihre Bemühungen beweisen, mit denen sie die Götzen verehrt haben an der Statt GOttes. Denn wenn sie keine Erkenntniß von GOtt gehabt hätten, so hätten sie den Götzen die Gottheit (eam) nicht beilegen können, noch den Namen GOttes [für sich] in Anspruch nehmen. Aber das ist ihre Unverständigkeit gewesen, daß sie ihn nicht gepriesen haben als einen GOtt, das heißt, sie haben sich ihn als einen solchen Gott erdichtet, der verehrt werden möchte durch ihre Meinungen und Bemühungen, so daß sie in solcher Weise sich selbst rühmen möchten, [als Leute] die Götte aus ihren eigenen Kräften etwas leisten könnten, das ihm angenehm wäre, und seine Gnade verdienen. Das aber heißt, GOtt verunehren, und Gott nicht danken für das, was man von ihm empfangen hat, sondern hochmüthig sein und vermessen wegen dessen, was man Gotte dargebracht hat. Dies ist nicht Anderes als wollen, daß sich GOtt von uns gestalten lasse, und daß wir uns von GOtt nicht gestalten lassen wollen; wollen, daß GOtt seinen Rath und seine Gedanken unseren Rathschlägen und Gedanken anpasse, wie er spricht in Jesaia [Cap. 40, 18. 25.]: „Wem wollt ihr denn GOtt nachbilden?“ und Cap. 55, 9.: „So viel der Himmel höher ist denn die Erde, so sind auch meine Wege höher denn eure Wege.“

Daß aber die Moabiter ihrem BaalPeor mit Fressen und Unkeuschheit gedient haben, das beweist nicht, daß sie um deswillen die Teufel verehrt haben. Denn das ist das Schicksal (fortuna) der menschlichen Satzungen, daß sie nur Bauchdiener (ventris animalia) groß ziehen, so daß Paulus sich untersteht [Phil. 3, 19.], den Bauch ihren Gott zu nennen, und Röm. 16, 18. sagt, daß sie ihrem Bauche dienen. Und alle Propheten beschuldigen sie, daß sie das Volk GOttes fressen. Dies bezeugen auch die Haine und die Höhen in der Synagoge, wo sie unter den grünen Bäumen ihren Muthwillen trieben. So sehen wir auch in unserem Reiche des Papstes, wie sie öffentlich und unverschämt dem Wohlleben, dem Müßiggang, dem Fressen und Saufen und der Unkeuschheit dienen, zugleich aber auch ihre Stundengebete (horas) murmeln und die Messen kreuzigen.

BaalPeor heißt ein Mann, der das Maul
\quellseite{714}
auffert,\footnote{vir hiatus. Vgl. St. Louiser Ausgabe, Bd. XIX, 688, § 49.} von dem viele meinen, es sei der Priapus gewesen, vielleicht dadurch beweget, daß die Moabiter in seinem Dienste sich der Unkeuschheit hingaben. Aber dies beweist (movet) nichts, weil die Unkeuschheit (wie ich gesagt habe) ein Merkmal ist, welches alle menschlichen Satzungen und Gottlosigkeiten haben; sodann scheint auch diese Geschichte Israels älter zu sein und eher geschehen, als Priapus in Griechenland aufkam. Deshalb glaube ich, daß der wahre GOtt von den Moabitern unter diesem lugenhaften Namen und falschen Gottesdienste geehrt worden sei, wie er von den Israeliten Baalim genannt wurde, wie der Prophet Hosea bezeugt. Sodann sehe ich, daß Baal bei den Hebräern in zusammengesetzten Wörtern die Vorzüglichkeit (primatum) bedeutet. Denn so sagen sie Balsam (balsamum = \hebraeisch{בעל}), das heißt, das vorzüglichste Öl (denn Öl nennen sie \hebraeisch{שמן}); desgleichen Beelzebub [\hebraeisch{בעל זבוב}], die vornehmste Fliege, wie wir sagen: Erzengel, Erzkoch, Erzherzog, Erzbischof. So haben die Moabiter den Baal-Peor etwa den vornehmsten Schlund (hiatum) genannt, aus einer mir unbekannten Ursache, es sei denn, daß du das als eine Vermuthung annehmen wollest (divinare), daß GOtt furchtbar und ein Erzrächer sei, bereit alle Gottlosen und Widersacher seines Volkes zu verschlingen, wenn er nicht auf diese Weise verehrt werde, wie bei uns die Papisten erdichtet haben, GOtt sei ein ganz geschwinder Rächer, wenn jemand aus Eingeben des Teufels einem Geistlichen Schaden gethan hätte. So haben sie auch aus St. Antonius, Sebastian, Valentin und Rochius eine Art Baal-Peor gemacht, um ihren Bauch zu nähren.

\begin{schriftzitat}
V. 5. Wie mir der HErr, mein GOtt, geboten hat etc.
\end{schriftzitat}

Diese Ursache allein sollte genugsam sein, daß sie alles gerne thäten, daß er GOttes Gebote lehrt, nicht seine eigenen, wenn auch keine Strafe oder Belohnung darauf folgte. Und dies ist das andere Stück (locus) der Ermahnung, welches in richtiger Ordnung folgt. Denn das erste handelt von dem Schrecken und dem Zorn GOttes, womit ein Gesetzgeber anfangen muß, sodann aber die Herrlichkeit und die Liebe GOttes daranfügen, damit er die Erschreckten locke zur Liebe.

\begin{schriftzitat}
V. 6. Denn das ist eure Weisheit etc.
\end{schriftzitat}

Das dritte Stück ist hergenommen von dem Vorzug (excellentia) und der Würde, die sie über alle Völker haben werden, welche hieraus folgt. „Wer mich ehret (sagt er [1 Sam. 2, 30.]), den will ich auch ehren.“ Wenn sie GOtt in solcher Weise verherrlichen würden, daß sie jene Gebote und Rechte hielten, so sollte es geschehen, daß der HErr sie wiederum verherrlichte, indem in der ganzen Welt das Gerücht von der Heiligkeit der Gesetze und des Volkes, welches sie hielt, ausgebreitet würde. Durch dieses Gerücht würden viele zu dem wahren GOtte gezogen werden. Und so würde auch durch diese Ehre, mit der GOtt sie unter den Völkern verherrlichte, GOtt selbst mehr verherrlicht werden. Denn so dienen auch alle Werke GOttes, mit denen er seine Heiligen verherrlicht, den Gottlosen dazu, daß auch sie bekehrt werden. Dies ist die Liebe, die alles, was unser ist, zum Nutzen des Nächsten wendet (facit) und annimmt.

„Weisheit“ und „Verstand“ (intelligentia), glaube ich, unterscheiden sich so, daß Weisheit gehört zu der Kenntniß der Dinge, Verstand zum Urtheil; so daß der weise ist, welcher das weiß, was zum Glauben und der Ehrfurcht (religionis) gegen GOtt gehört, und zum Recht und zur Liebe gegen den Nächsten, daß aber der verständig ist, der in beiden Stücken [gegen GOtt und gegen den Nächsten] scharfsinnig, klug und vorsichtig ist, so daß er recht darin wandelt und sich durch keine Trügerei verführen lasse; auf deutsch: „weise und klug“.

Aber er ermahne das noch viel mehr, mit wie wenigen Worten Mosis die gesamte Weisheit und das Wissen (scientiam) der Menschen verwirft und verdammt, indem er zeigt, daß die Weisheit und Verstand des Volkes GOttes nicht bestehe in den spitzfindigen Forschungen (speculationes) der Sophisten, auch nicht in den herrlichen Verordnungen (sanctiones) der Fürsten und klugen Leute, sondern in der Kenntniß und Beobachtung des Gesetzes GOttes. „Das ist (sagt er) eure Weisheit und Verstand.“ Mag daher jemand immerhin weise sein im Regieren des Wesens dieser Welt (wie sie denn ja auch weise sind, Böses zu thun), so soll doch das Volk GOttes nirgendes nach Weisheit trachten, noch irgend eine andere wissen als aus dem Gesetze seines GOttes; da wird es reichlich und seliglich finden, wie es sich gegen GOtt und Men-.
\quellseite{715}
ischen, wie in Glück und Unglück, wie in Frieden und Krieg verhalten soll. Denn anderswoher geholte Weisheit ist nichts als Thorheit vor GOtt. Daher macht Moses durch dieses Wort alle Menschen vor GOtt zu Thoren, in aller ihrer Weisheit und Klugheit.

\begin{schriftzitat}
V. 7. 8. Zu dem Götter also nahe sich thun etc.
\end{schriftzitat}

Die Ehre und der Ruhm dieses Volkes ist größer als die aller Völker der Erde, weil ihr GOtt nahe bei ihnen ist, sondern auch, weil er gerechte Rechte und Gebote hat; wenn sie dieselbigen hielten, so sollten die Völker sich über sie verwundern und Heil von ihnen empfangen. Das ist ein gar schöner und lieblicher Trost und eine zuverlässige Verheißung, daß GOtt nahe sei. Diese Lehre (locum) hat David sehr oft in seinen Psalmen angeführt und die anderen Propheten überall. Aber zugleich schließt er das Kreuz des Volkes GOttes mit ein und zeigt es an. Denn daß er denen nahe ist, die ihn anrufen, beweist, daß die Anrufenden in Trübsal sind, welche sie zwingt anzurufen, wie denn auch dieses Volk in der That immer [in Trübsal] war, und alle, die gottselig leben wollen, Verfolgung leiden [2 Tim. 3, 12.], damit der anrufende Glaube Raum habe und eine Gelegenheit da sei, GOtte nahe zu kommen.

Die Gottesdienste anderer Völker aber, wie sie gottlose Satzungen wider GOtt haben, und unbillige Rechte gegen den Nächsten, so haben sie auch Götter, die ferne von ihnen sind, welche nicht selig machen können, ja, welche sich verfolgen und vertheidigen lassen von ihren Verehrern.

\begin{schriftzitat}
V. 9–11. Und sollst deinen Kindern und Kindeskindern kund thun den Tag etc.
\end{schriftzitat}

Diesen Tag der ἐκκλησία oder Versammlung schätzt Moses oft ein, weil an demselben das Volk das Gesetz GOttes gehört hatte an dem Berge Sinai unter Ungewitter und erschrecklichem Donner, wo nicht ein Mensch redete, sondern die höchste Majestät durch den Dienst der Engel, und von diesem Tage das Ansehen Mosis begann, wie im zweiten Buche Mosis, Cap. 19, gesagt wird, und die ganze Ordnung (institutio) des Volks.

\begin{schriftzitat}
V. 12–19. Die Stimme seiner Worte höret ihr, aber kein Gleichniß sahet ihr außer der Stimme etc.
\end{schriftzitat}

Diese Stelle verwirft nicht allein die äußerlichen Bilder von GOtt, welche er dem rohen und kindischen Volke verboten wissen wollte, sondern vor allem und am allermeisten die innerlichen, welche (wie wir oben gesagt haben) der Wahn und die Gedanken sind, die wir in uns, ohne die Stimme GOttes, erdichtet haben. Hier wird sicherlich allein die Stimme seiner Worte gepriesen, und alles, was über GOtt geredet oder gedacht wird, das nicht diese Stimme seiner Worte ist, ist durchaus gottlos und verdammt. Denn er hat uns allein durch sein Wort seinen Willen, seine Rathschläge vorgebildet wissen wollen, nicht durch unsere Vorstellungen oder Einbildungen. Es gefällt oder mißfällt Gotte also nicht das, was du meinst, daß ihm gefalle oder mißfalle, wie heilig und gottselig es dir auch vorkommt, was du denkst (wie die Urheber und Lehrer der geistlichen Orden [religionum] und der Secten gethan haben), sondern das, wovon er selbst durch die Stimme seines Wortes erklärt, daß es ihm gefalle oder\footnote{In den Ausgaben: et statt: aut.} mißfalle. Denn niemand kann GOttes Willen anzeigen oder kund thun, als GOtt selbst, ja, jedermann irrt in gottloser Weise, wenn er sich dessen untersteht, da niemand die Tiefen GOttes kennt als der Heilige Geist, der in GOtt ist. Darum ist es unmöglich, daß die Menschen ohne ein gewisses Wort GOttes richtig von GOtt halten, denken, predigen oder ihn recht verehren. Der Ausspruch steht fest: „Ihr habt seine Stimme gehört.“

\begin{schriftzitat}
V. 20. Aus dem eisernen Ofen.
\end{schriftzitat}

Mit einem starken Ausdruck nennt er die Knechtschaft in Egypten einen eisernen Ofen, nicht allein um der Erschrecklichkeit der Geschichte willen, sondern auch wegen des Elendes der Gewissen, welche unterdrückt waren unter dem Zorn des Gesetzes und der Tyrannei der Satzungen, deren Vorbild das Volk Israel war in seinem egyptischen Ofen.

\begin{schriftzitat}
V. 24. 25. Denn der HErr, dein GOtt, ist ein verzehrend Feuer und ein eifriger GOtt.
\end{schriftzitat}

Dieser Ausspruch wird oft gebraucht auch in der Epistel an die Hebräer, und es ist würdig, daß er oft gebraucht werde, um die Gottlosen zu schrecken und zu demüthigen. Es sind aber zwei Dinge: „ein verzehrendes Feuer“ und: „ein
\quellseite{716}
eifriger GOtt“. Ein Feuer wird er genannt, weil er gänzlich verzehrt und den Gottlosen nichts übrig läßt, auch nichts ist, was seinem Zorne widerstehen könne; ein Eiferer, weil er ein solches Gemüth hat, daß er nicht schonen will. Wer sollte nun diesen nicht fürchten, von dem man weiß, daß er nicht verschonen will und eine solche Weise (virtutem) hat, daß er unversöhnlich und ohne Aufhören Rache übt? Ferner, diesen Zorn übt er in zwiefacher Weise aus; erstlich geistlicher Weise und plötzlich, indem er den Gottlosen seinen Geist entzieht, so daß der Teufel alles verwüstet; sodann auch leiblich und langsam, indem er auch die Leiber tödtet, die Länder verheert und alles zunichte macht.

Und merke sonderlich, daß GOtt von Mose beschrieben werde als ein Feuer und ein Eiferer hinsichtlich des ersten Gebots, in welchem er am meisten verachtet wird, wo er doch am meisten gefürchtet und verehrt werden sollte, damit du sehest, auf welches Gebot er selbst vorzüglich dringe, nämlich daß wir ihn für unsern GOtt halten sollen, damit wir kein Vertrauen auf die Werke aufrichten, sondern uns allein auf seine Barmherzigkeit gründen, wiewohl er in seiner Sache mehr durch die Finger zu sehen scheint, als in dieser Gottlosigkeit, da er die Gottlosen überhand nehmen und ungestraft herrschen und sehr viele verführen läßt, so daß weder Feuer noch Eifer, sondern Kälte und Trägheit bei GOtt zu sein scheint, wie auch viele Psalmen klagen, daß er schlafe, und ihn bitten, daß er aufstehe.

\begin{schriftzitat}
V. 26–29. So rufe ich heutiges Tages zu Zeugen etc.
\end{schriftzitat}

Das letzte Stück dieser Ermahnung ist die Drohung. Denn so pflegen wir zu ermahnen, indem wir unsere Gründe hernehmen von dem, was nützlich ist, was nothwendig, was ehrbar, was gottesfürchtig ist, endlich von der Gefahr etc., was Moses hier auch thut.

\begin{schriftzitat}
V. 30–40. In den letzten Tagen etc.
\end{schriftzitat}

Dies „in den letzten Tagen“ ist allgemein geredet, wofür wir sagen: später. Denn Moses will solchen Zorn GOttes beschreiben, daß er denen, die nachher wieder zurechtkommen, immer Raum zur Buße lasse. Dies haben auch die Ereignisse in der Geschichte bewiesen, wo GOtt Israel so oft gestraft hat, und es doch, wenn es wiederkehrte, hernach immer wieder angenommen hat. Er redet also nicht allein von den letzten Tagen vor dem Tage des Gerichts, in welchen, wie man meint, die Juden noch bekehrt werden sollen, sondern allgemein von aller Zeit (wie ich gesagt habe), in der sie nach der Sünde Buße gethan haben und die Strafe erkannt.

Und um zu dieser Buße zu reizen und die Verzweiflung zu verbannen, fügt er die lieblichste und freundlichste Verheißung hinzu, auf welche sich gar schön auch das Gebet Manasse, des Königs Juda, gründet, nämlich, daß GOtt, dein HErr, barmherzig ist, der dich nicht verlassen noch verderben wird, noch des Bundes vergessen, den er deinen Vätern geschworen hat [5 Mos. 4, 31.]. Dadurch zeigt er zugleich an, was die wahre Buße sei, nämlich ein brennender Trieb nach Barmherzigkeit in der Angst des Gewissens. Denn allein diesen wird die süße Barmherzigkeit vorgehalten, nicht aber den Heuchlern.

\begin{schriftzitat}
V. 41–43. Da sonderte Moses drei Städte aus etc.
\end{schriftzitat}

Diese Sache menget Moses mit in die Rede, indem er die Ordnung der Historie hält. Denn diese Reden des Deuteronomiums hat er sicherlich nicht in Einem Tage zu Ende gebracht, sondern vielleicht kaum in zehn Tagen, oder an jedem einzelnen Tage kaum zwei oder drei Capitel; und inzwischen hat er mit Sachen im Volke zu thun gehabt, wie sie eben vorfielen, wie diese Aussonderung der Städte zu Zufluchtsörtern.

\begin{schriftzitat}
V. 44. ff. Das ist das Gesetz etc.
\end{schriftzitat}

Dies ist der Schluß der Rede, wie Moses Weise ist. Ferner, der Berg Hermon, welcher hier [3. 48.] Sion genannt wird, ist ein anderer als der in Jerusalem. Denn dieser wird mit einem Hade geschrieben, Zion, dieser mit einem Sin, Sion. Eben Cap. 3 haben wir gesagt, daß „die Quellen Pisga“ für den Eigennamen einer Stadt zu halten seien (dici), [nämlich] Albedoth, welches Hieronymus übersetzt hat: „am Fuße“ [des Berges],\footnote{In der alten Uebersetzung lautet dies: „daß Hieronymus muß, wurzeln [radices], verdolmetschet hat.“} aber in dem Buche „von den hebräischen Ortschaften“ macht er den Namen einer Stadt daraus.

\quellseite{717}
\section{Das fünfte Capitel}

\begin{schriftzitat}
V. 1. 2. Und Moses rief dem ganzen Israel etc.
\end{schriftzitat}

Hier hast du es, daß Moses, um das Deuteronomium zu vollenden, verschiedene Predigten hält, und das Volk öfters zusammenberufen hat. Daher wiederholt er auch immer einiges am folgenden Tage, was er am vorhergehenden gesagt hatte, wie hier, da er den Bund GOttes und die Stimme des Herrn etc. einschärft. Hier fängt nun die hauptsächlichste Predigt von den zehn Geboten an.

\begin{schriftzitat}
V. 3–5. Und hat nicht mit unsern Vätern diesen Bund gemacht etc.
\end{schriftzitat}

Moses zeigt hier den Unterschied des neuen und des alten Testamentes an. Das neue Testament ist das älteste, welches verheißen ist von Anfang der Welt, ja, „vor den Zeiten der Welt“, wie Paulus im Briefe an den Titus [Cap. 1, 2.] redet, aber nur unter Christo erfüllt. Das alte Testament, welches unter Mose verheißen worden ist, ist unter Josua erfüllt worden. Es ist aber dies der Unterschied zwischen beiden, daß das neue sich allein auf die Verheißung des erbarmenden und treuen GOttes stützt, ohne unsere Werke, das alte aber auch auf unsere Werke. Daher gehen die Verheißungen Mosis nicht weiter, als sofern sie die Gebote und Rechte halten. Dies war die Ursache, daß es selbst endlich veralten und abgethan werden mußte und ein Vorbild sein des neuen und ewigen Testamentes, welches vor der Welt angefangen hat und nach der Welt bleiben wird.\footnote{Statt duravit in der Jenaer und in der Erlanger haben wir mit der Wittenberger durabit angenommen.} Jenes aber hat in der Zeit angefangen und nach einiger Zeit wieder aufgehört.

\begin{schriftzitat}
V. 6. ff. Ich bin der HErr, dein Gott etc.
\end{schriftzitat}

Es scheint dies Gebot bloß äußerliche Satzungen (ceremonias) über die Bilder anzuordnen für das rohe und kindische Volk, dennoch ist es in Wahrheit ein geistliches Gesetz, welches den inwendigen Dienst des Geistes erfordert, auch was die äußerlichen Bilder betrifft. Uebrigens kann niemand diese Gebote des Dekalogs besser auslegen, als Moses es in diesem Deuteronomium thut, wie wir sehen werden. Daher wollen wir hier nicht mehr davon reden.

Es zeigt aber Moses die Kraft des Gesetzes, indem er sagt, daß sie die Stimme GOttes nicht hätten ertragen können, nämlich weil das Gesetz, wenn es im Geiste verkündigt und gehört wird, gänzlich tödtet und den Menschen zu einer ihm unerträglichen Erkenntniß seiner Sünde und zur Furcht vor dem Tode bringt, so daß der Mensch nach einem Mittler seufzt und ein lieblicheres Wort begehrt, das ist, das Evangelium der Gnade, wie hier das erschreckte Volk Moses als Mittler erbittet, der gelinder reden werde, damit es nicht sterbe von der Stimme GOttes. Denn was würde ein Mensch doch leiden und thun, damit er das Gesetz nur nicht hören müßte?

\begin{schriftzitat}
V. 22. ff. Und thut nichts dazu etc.
\end{schriftzitat}

Und nachher [V. 32.] sagt er: „Weichet nicht, weder zur Rechten noch zur Linken“, nämlich er wiederholt wiederum und schärft ein, daß zu dem Worte GOttes nicht irgend etwas hinzu oder davon gethan werde, worüber wir oben genug geredet haben.

\section{Das sechste Capitel}

\begin{schriftzitat}
V. 2. 3. Daß du den HErrn, deinen Gott, fürchtest etc.
\end{schriftzitat}

Es scheint hier [hebräischen] Sprache eigen zu sein, daß sie „die Furcht GOttes“ sagt, wo wir Gottesdienst oder Gottesverehrung oder griechisch θεοσέβεια sagen, wie Moses hier die beiden Stücke verbindet, die Furcht und das Halten der Gebote.

\quellseite{718}
Von dieser Stelle an wollen wir diese Auslegung, welche Mosis über seinen Dekalog selbst verfaßt hat (denn was ist, von dieser Stelle an, das ganze übrige Buch anders als eine überaus reiche und zwar sehr klare Erklärung der zehn Gebote, wie wir sehen werden), in zwei Theile theilen. Der erste Theil soll die Erklärung der drei Gebote der ersten Tafel sein, welche fast bis zum zwanzigsten Capitel fortgeht; der zweite Theil aber die Erklärung der übrigen Gebote der zweiten Tafel bis zum Ende des Buches. Es haben freilich viele Leute Auslegungen über die Gebote (praeceptoria) und ganze Sammelwerke (summas) geschrieben, aber bloßen Dreck, wenn es mit dieser Auslegung Mosis verglichen wird. Denn Moses wird dich lehren, was das heiße, nicht andere Götter haben, was da sei, GOttes Namen nicht unnützlich führen, was der Feiertag sei etc. Er sagt nun:

\begin{schriftzitat}
V. 4. Höre, Israel, der HErr, unser GOtt, ist ein einiger HErr.
\end{schriftzitat}

Du siehst, daß er das erste Gebot selbst in bejahender Weise (affirmative) auslegt, nämlich, daß Ein HErr sei. Denn es wird hier der aus vier Buchstaben bestehende Name (tetragrammaton [das ist \hebraeisch{יהוה}]) gesetzt, welcher Gotte allein eigen ist. Er handelt aber diese Einheit GOttes im Geiste, das heißt, er hat nicht sowohl das im Auge, daß GOtt Einer sei, als vielmehr, daß er von uns für einen einigen gehalten werde, da uns ja daraus, daß Ein GOtt ist, kein Nutzen entspringt; wenn wir ihn aber für Einen und für unsern GOtt halten (wie er hier sagt), das ist das Leben und die Seligkeit und die Erfüllung aller Gebote.

So sagt Jacob 1 Mos. 28, 21.: „Der HErr soll mein GOtt sein.“ Wie soll der HErr sein GOtt sein? Ist er es denn zuvor nicht gewesen? Dies sagt er nur aus dem Grunde, daß er sich vorbehielte, allein den HErrn als seinen GOtt zu haben mit gewissen Ceremonien (ritu) und Einer Art der Gottesverehrung. So wird GOtt und ändert sich, je nachdem sich unser Herz gegen ihn ändert, wie es Ps. 18, 26. 27. heißt: „Bei den Heiligen bist du heilig, und bei den Verkehrten bist du verkehrt.“ So erhebt sich der Antichrist über jeden Gott, das heißt, er wird einen eigenen Gottesdienst einsetzen, welchen er über allen wahren und falschen Gottesdienst aller Götter erhöhen wird, weil niemandes Wort mehr gefürchtet und angebetet wird [als das seine].

Es ist also die erste Erklärung Mosis vom ersten Gebot diese, daß der HErr, unser GOtt, für Einen HErrn zu halten sei, das heißt, daß er nicht mit immer anderen Gottesdiensten verehrt werden soll, die von uns erfunden sind, sondern allein mit dem, welchen er eingesetzt hat.

\begin{schriftzitat}
V. 5. Du sollst den HErrn, deinen GOtt, lieb haben von ganzem Herzen etc.
\end{schriftzitat}

Die zweite Erklärung des ersten Gebotes; denn die erste, eben gegebene, betrifft den Glauben. Denn niemand kann Einen GOtt haben, wenn er nicht allein an ihn hanget und ihm allein vertraut; sonst wird er zu mancherlei Werken hingerissen und sich verschiedene Götter erdichten. Die zweite, welche aus der ersten folgt, betrifft die Liebe. Denn wenn wir alles Vertrauen auf ihn setzen, an ihm hangen, und erkennen, daß alles von ihm herfließe, und daß er sich unser sorgfältig annehme, dann folgt nothwendiger Weise eine süße Liebe gegen ihn. Darum gebraucht er der verneinenden (negativa) Redeweise in dem Gebote: „Du sollst nicht andere Götter haben“ etc., als ob er sagen wollte: Es ist vonnöthen, daß du dich selbst verleugnest und an dir verzweifelest, damit du nicht viele Götter machest, daß du Einen GOtt habest, denn die Natur kann nicht anders, als Abgötterei treiben.

Wenn er daher sagt: Der HErr, dein GOtt, ist ein einiger HErr, so nimmt er alles [andere] Vertrauen hinweg. Wenn er sagt: Du sollst den HErrn lieb haben, so regt er an zu einem fröhlichen und freien Dienst gegen GOtt. Denn wenn ich GOtt liebe, so will ich in Wahrheit alles, was GOtt will, so ist mir auch nichts lieber, als das zu hören und zu thun, was GOtt will, wie auch die fleischliche Liebe thut gegen den geliebten Gegenstand. So empfangen wir durch die Einheit mit GOtt im Glauben alles umsonst von GOtt, durch die Liebe thun wir gegen GOtt alles umsonst.

Aber daß er hinzufügt: „Von ganzem Herzen, von ganzer Seele, von allem Vermögen“, das ist so schwer, daß, wenn nicht Gott vorgäbe, es kein Heiliger erfüllen könnte. Ja, wer ist da, der nicht an beiden Stücken Mangel hätte, sowohl in dem, daß er Einen GOtt haben soll, als auch in dem, daß er ihn lieben soll? denn
\quellseite{719}
es ist niemand, der nicht bisweilen im Glauben wankte, und andere Dinge zugleich mit GOtte liebe. Das beweisen die Früchte. Denn wenn ich GOtt von ganzem Herzen liebte, so würde mich nichts mehr verletzen, als daß die Gebote GOttes verachtet werden, wie Paulus dies in den Briefen an die Corinther und an die Galater beklagte, da er sah, daß die Ehre GOttes mit Füßen getreten werde. Aber wo sind die Leute, welche beweinen, daß der Name GOttes so in der ganzen Welt mit Füßen getreten wird?

„Von ganzem Herzen“ bedeutet: von Grund des Gemüthes und mit ganzem Gemüthe; „von ganzer Seele“: mit dem ganzen leiblichen (animali) Leben; „von allem Vermögen“, das ist, mit allen Kräften und Gliedern, wovon ich anderswo genug gesagt habe. Nicht daß wir keine anderen Dinge lieben sollten, da alles, was GOtt gemacht hat, sehr gut und liebenswerth ist, sondern daß man in der Liebe GOttes, und den Dingen, die GOttes sind, nichts gleichstellen oder vorziehen soll, und die Liebe aller Dinge dahin gedrungen werde, daß die Liebe gegen GOtt völlig werde. Hieraus merke, wie verkehrt diejenigen sind, welche ihre Satzungen und Aussprüche so streng erzwingen, aber zulassen, daß GOttes Gebote so vernachlässigt werden.

\begin{schriftzitat}
V. 6–9. Und diese Worte, die ich dir heute gebiete, sollst du zu Herzen nehmen etc.
\end{schriftzitat}

Nicht bloß in dem Buche, nicht bloß in Gedanken, sondern im innersten Gemüthe, so daß sie dir der höchste Schatz seien, denn wo dein Schatz ist, da wird auch dein Herz sein. Es soll also in deinem Herzen nichts regieren als der Glaube und die Liebe gegen GOtt. Darüber soll dein Herz Tag und Nacht nachdenken, Ps. 1, 2. Denn wenn sie in solcher Weise zuerst im Herzen sind, dann kann die rechte Folge nicht ausbleiben, daß sie auch in deinem Munde seien.

„Und sollst sie deinen Kindern schärfen“, das heißt, täglich sollst du sie wiederholen und einbläuen, damit sie nicht matt und kalt werden und vom Rost verzehrt werden. Dann wird folgen, daß du überall und immer von ihnen redest. Endlich sollst du sie auch als ein Denkmal auf deine Hand binden und vor die Augen. Zuletzt sollst du sie an die Pfosten schreiben.

Siehe die Ordnung, wie man GOttes Wort handeln soll: erstlich soll man es im Herzen bedenken; zweitens den Kindern treulich und anhaltend mit dem Munde einschärfen; drittens öffentlich und überall reden; viertens zeichnen auf die Hand und vor Augen malen; zum fünften und letzten schreiben, und zwar an die Pfosten und Thore des Hauses, nicht in Bücher, weil Moses selbst sie schon in ein Buch geschrieben hat. Hier übergehe ich diese heimlichen Deutungen, daß das Zeichnen der Gebote auf die Hand und das Stellen derselben vor Augen bedeute, daß man sie im Werke erfülle und sie betrachte. Er will einfach, daß uns diese Worte überall entgegentreten und im Gedächtniß sein sollen.

Du wirst wahrnehmen, daß Moses diese Ermahnung nicht zu den anderen Geboten setzt, sondern nur zu dem ersten und zu seiner Erklärung, daß sie im Herzen sein sollen, daß man sie schärfen, reden, zeichnen und schreiben soll, damit du wissest, daß das erste Gebot das Maß und die Regel für alle anderen sei, denn sie weichen und gehorsam sein müssen; als, wenn es dem Glauben und der Liebe gemäß ist, sollst du tödten wider das fünfte Gebot, gleichwie Abraham 1 Mos. 14, 15. die Könige tödtete, und der König Ahab sündigte, daß er den König von Syrien nicht tödtete, 1 Kön. 20, 34. ff. So steht es auch mit dem Diebstahl, so mit Kinderlist und Trug wider die Feinde GOttes, so mit dem Raube, Gütern, Weibern, Töchtern, Söhnen, Knechten der Feinde. So sollst du Vater und Mutter hassen, daß du den HErrn liebest.

Kurz, wo irgend etwas wider den Glauben und die Liebe wäre, da sollst du nichts davon wissen wollen, daß es noch irgend ein anderes Gebot gebe, weder von GOtt noch von Menschen. Wo es aber dem Glauben und der Liebe gemäß wäre, sollst du wissen, daß die Gebote in allen Dingen und überall da sind, denn der Ausspruch steht fest: „Diese Worte sollen in deinem Herzen sein“, da sollen sie regieren. Ferner, wenn sie nicht auch im Herzen wären, so könnte sicherlich niemand diese verschiedene Anwendung je nach den Umständen (hanc epikian) verstehen noch befolgen, noch könnte er jemals der Gesetze heilsamlich und sicher und recht gebrauchen, wie auch Paulus sagt 1 Tim. 1, 9., „daß dem Gerechten kein Gesetz gegeben ist“, darum, weil [1 Tim. 1, 5.] die Erfüllung des Gesetzes ist Liebe aus einem guten Herzen und ungefärbtem Glauben, welche des Gesetzes recht gebraucht, indem sie keine Gesetze hat, und
\quellseite{720}
alle Gesetze hat; keine, weil keine [Gesetze] binden, wenn sie nicht dem Glauben und der Liebe dienen; alle, weil alle binden, wenn sie dem Glauben und der Liebe dienen.

Es ist also die Meinung Mosis an dieser Stelle: Wenn du das erste Gebot recht verstehen willst, und in Wahrheit keine anderen Götter haben, so halte dich dermaßen, daß du den einigen GOtt glaubest und liebest, dich selbst verleugnest, alles umsonst empfängest, und alles umsonst thuest.

\begin{schriftzitat}
V. 10–12. Wenn dich nun der HErr, dein GOtt, in das Land bringen wird etc.
\end{schriftzitat}

Er hält an mit der Auslegung des ersten Gebots und beginnt zu zeigen, auf welcherlei Weise gegen das erste Gebot gesündiget werde, und lehrt, daß man den Anlaß, wider dasselbe zu sündigen, meiden soll, und nimmt zuerst vor sich den Reichthum und den Ueberfluß, nämlich den Mammon und den Geiz, von dem Baruch [Cap. 3, 17.] schreibt: „Das Gold, darauf die Menschen ihr Vertrauen setzen.“ Und Paulus sagt, daß der Geiz Abgötterei sei [Col. 3, 5.] und eine Wurzel alles Uebels [1 Tim. 6, 10.]. Hüte dich, sagt Moses [V. 1, 12.], wenn du satt wirst, daß du nicht des HErrn, deines GOttes, vergessest. Denn das menschliche Herz vertraut auf das vorhandene Gut; wenn kein Gut da ist, so ist es voll Mißtrauens, wie man gemeiniglich sagt: Gut macht Muth, Armuth wehe thut. Aber das Vertrauen auf Reichthum und der Glaube und die Liebe können nicht zugleich im Herzen regieren. Und das nennt er hier „vergessen des HErrn, deines GOttes“. Denn das heißt nicht des HErrn gedenken, wenn du seinen Namen im Munde führst, sondern wenn du im Herzen mit beständigem Glauben an ihn hangst und ihn liebst.

Und siehe du selbst, wie Moses das erste Gebot zuerst geistlich handelt wider die Abgötterei des Herzens, wider das Vertrauen auf die Güter, um Vertrauen zu GOtt zu erwecken, ehe er von den Götzen sagt, damit du nicht meinest, daß Moses nichts Anderes im Auge gehabt habe als die Götzen. Es ist also dies der Sinn: Du sollst Einen GOtt haben, daß du ihm allein vertraust. Wenn dir gleich Reichthum zufällt, so sollst du nicht darauf vertrauen; wenn er nicht da ist, so sollst du nicht kleinmüthig werden, sondern immer auf den HErrn vertrauen, wie auch der Psalm [Ps. 62, 11. Vulg.] sagt: „Fällt euch Reichthum zu, so hänget das Herz nicht daran“, und: „Trachtet nicht nach Raube, werdet auch nicht eitel.“

\begin{schriftzitat}
V. 13. Sondern sollst den HErrn, deinen GOtt, fürchten und ihm allein dienen.
\end{schriftzitat}

„Du sollst fürchten“, das heißt, du sollst verehren, und aus Ehrbietung gegen ihn sollst du dich scheuen, auf den Mammon zu vertrauen, das heißt, wenn die Güter reichlich vorhanden sind, da die Menschen in Sicherheit zu leben pflegen und GOtt zu verachten und ihren Lüsten zu dienen, sollst du mich fürchten, und herrschen über die Güter, sollst mir dienen, und sollst dich durch sie nicht hinreißen lassen nach deinen Lüsten. Auch dies ist eine herrliche Erklärung des ersten Gebots, nämlich, daß wir dann uns am meisten fürchten sollen, wenn alles überflüssig vorhanden ist, wohl verwahrt\footnote{Statt tota in der Jenaer und der Erlanger haben wir mit der Wittenberger tuta angenommen.} und sicher, und kein Anlaß zur Furcht da ist. Denn es ist dies die Art (virtus) des Glaubens und das Wesen (vis) des ersten Gebots, daß wir im Wohlergehen furchtsam seien und GOtt fürchten, wiederum in Widerwärtigkeiten (wie folgt) sicher und unbesorgt und vertrauensvoll gegen GOtt und, mag es nun wohl oder übel gehen (utroque tempore), nur an GOtt hangen.

Hier haben die Sophisten mancherlei Träume erdichtet von der δουλεία, λατρεία und ὑπερδουλεία. Der Hebräer bezeichnet mit einem und demselben Worte den Dienst gegen GOtt und Menschen, so daß ihre Unterscheidung ohne Nutzen ist. Er will aber dies sagen: „Du sollst ihm allein dienen“, das heißt, alles, was du thust und lebst, sei es unter der Dienstbarkeit der Menschen oder in der Regierung der Dinge, das sollst du mit Bezug auf mich ausrichten, und es soll in keiner anderen Weise (nomine) von dir geschehen, als daß du gewiß seiest im Glauben, daß darin mir allein gedient werde, und wenn du nicht wüßtest, daß mir in der Sache gedient werde, solltest du bald ablassen, so daß du nichts thätest, den Menschen zu gefallen, wider mich, selbst wenn du ihr Gefangener wärest, und sollst deiner Begierde oder Ehre nichts nachgeben, wenn dir auch alle Dinge reichlich zufallen. So lehrt auch Paulus, daß die
\quellseite{721}
Knechte ihren Herren unterthan sein sollen, aber als dem HErrn [Eph. 6, 5.], und daß die Weiber ihren Männern unterthan sein sollen, aber in dem HErrn [Eph. 5. 22.], und er selbst will des Philemon genießen, aber in dem HErrn [Philem. V. 20.], und er will, daß eine Jungfrau freie [1 Cor. 7, 36.], aber in dem HErrn [V. 39.], und [1 Cor. 11, 12.] daß der Mann durch das Weib komme, aber in dem HErrn. Er träumt nichts von der δουλεία und ὑπερδουλεία, sondern erhebt die eine und einzige Dienstbarkeit gegen GOtt höher als alle Dinge, und breitet sie über alle Dinge aus. So sagt Petrus in der Apostelgeschichte [Cap. 13, 36.], daß David dem Willen GOttes gedient habe.

Ja, wenn die δουλεία von der λατρεία getrennt wird, dann ist es schon weder λατρεία noch δουλεία, sondern ein Götzendienst (idololatria), da du hier hörst: „Du sollst ihm allein dienen.“ Also wird entweder dein Dienst (dulia) ein Gottesdienst (latria) sein, oder ein Götzendienst (idololatria), wenn es getrennt wird. Aber dies haben die Sophisten erdichtet, indem sie den äußeren Schein (larvam) der Werke und die Secten ansehen, und die Verehrung GOttes maßen nach der Verschiedenheit der Werke, obgleich GOtt durch dieses Wort die Verschiedenheit der Werke aufhebt und sie zusammenfaßt in die Einigkeit (unitas) des Glaubens im Herzen, damit du allein das ins Auge fassest, daß du irgend welche Werke, welcher Art sie auch sein und wo sie auch immer geschehen mögen, in der Furcht GOttes und in seinem Namen thuest, und ihm allein durch dieselben dienst. Es fällt also auch hier das Papstthum und jedes Reich menschlicher Satzungen, durch welche solche Werke gelehrt werden, dadurch gottlosen Menschen gedient wird, die entweder von GOtt abwenden, oder zwingen, GOtt zu verleugnen. In diesem Sinne hat auch Christus dem Satan, der ihn aufforderte zu Werken, die nicht allein GOtte geleistet werden sollten, dieses Wort [Matth. 4, 10.] geantwortet: „Hebe dich weg von mir, Satan, denn es stehet geschrieben: Du sollst anbeten GOtt, deinen HErrn, und ihm allein dienen.“

\subsection{Und sollst bei seinem Namen schwören}

Hier sind zwei Dinge zu beachten. Erstlich, daß Christus Matth. 5, 34. schlechthin verbietet zu schwören; hier wird geboten zu schwören. Aber auch anderswo haben wir gesagt, daß der Brauch des Eides ein zwiefacher ist; der eine, den wir nach unserer Leichtfertigkeit für uns selbst ohne Ursache schwören; dies verbietet Christus ganz und gar; der andere, da wir für die Ehre GOttes und das Heil des Nächsten aus Glauben und Liebe schwören, um die Wahrheit zu bestätigen. Dieses gebietet Moses, in solcher Weise, daß er nicht befiehlt zu schwören, sondern darauf dringt: Wenn geschworen werden muß, so sollst du nicht anders schwören als bei dem Namen GOttes.

Aber wie reimt sich das zu dem Vorhergehenden? Er hatte gesagt: „Du sollst ihm allein dienen“; bei diesem Dienste hat er nicht die religiösen (religiosa) und ceremonialen Werke der Verehrung GOttes von den weltlichen (profanis) unterschieden, sondern hat durchaus alle in eins zusammengefaßt, mögen sie nun GOtte oder den Menschen geschehen, daß sie alle von Herzen allein GOtte erwiesen werden sollen. So ist auch hier von dem Namen GOttes zu urtheilen. Er will nicht, daß man allein den Namen „GOtt“ in den Mund nehme, wenn man schwören muß, da Paulus auch sagt [1 Cor. 15, 31. Vulg.]: „Bei eurem Ruhm, ich sterbe täglich“ etc. Und Christus spricht [Joh. 16, 23.]: „Wahrlich, wahrlich, ich sage euch“, und Paulus [Röm. 9, 1.]: „Deß gibt mir Zeugniß mein Gewissen.“

Daher schwörst du bei dem Namen GOttes, wenn du das, bei dem du schwörst, auf GOtt beziehst und im Namen GOttes ergreiffst; sonst würdest du nicht schwören, wenn du wüßtest, es mißfalle ihm; auf jegliche Weise, wodurch du GOtte allein dienst, wenn du den Menschen im Namen GOttes dienst; sonst würdest du nicht dienen. Durch diesen Eid wirst du bewahrt, daß du GOtte allein dienst, und nicht gezogen werdest zu einem gottlosen Werke oder Eide. So erklärt Christus auch Matth. 23, 16–22., daß der bei GOtt schwöre, welcher bei dem Tempel und dem Altar und dem Himmel schwöre. Und Matth. 5, 34–36. verbietet er, bei Jerusalem, bei dem Haupte, bei dem Himmel, bei allen Dingen zu schwören, weil in allem diesem bei GOtt geschworen wird. Wenn man aber bei GOtt leichtfertig und unnütz schwört, das ist den Namen GOttes unnützlich führen. Daß er also will, man solle bei dem Namen GOttes schwören und bei keinem anderen, dafür ist nicht allein das der Grund, daß für die
\quellseite{722}
Wahrheit (welche GOtt ist) die Bestätigung von niemand anders, als von GOtt selbst beigebracht werde, sondern, daß der Mensch allein in GOttes Dienst bleibe, und lerne, ihm alles heimzutragen, und in seinem Namen alles zu thun, zu haben, zu gebrauchen und zu leiden, damit sie nicht, wenn sie sich eines anderen Namens bedienten, von GOtt abgewandt würden, und sich gewöhnten zu schwören, als ob GOtt das nicht anginge, und sie so endlich durch bösen Gebrauch einen Unterschied machten zwischen den Werken, mit welchen GOtt gedient wird, und den Werken, mit denen GOtt nicht gedient wird, da er doch will, daß man ihm in allen Dingen diene, und daß alles in seiner Furcht geschehe, indem er gegenwärtig ist, darauf sieht und darüber urteilt.

Daher soll man des Eides gebrauchen wie auch des Schwertes und des geschlechtlichen Verkehrs (copula sexus). Es ist verboten, das Schwert zu nehmen, da Christus sagt [Matth. 26, 52.]: „Wer das Schwert nimmt, der soll durch das Schwert umkommen“, weil er es ohne Befehl und aus eigener Willkür nimmt. Aber es ist das Gebot da, das Schwert zu führen, und es ist ein Dienst GOttes, Röm. 13, 4., wenn es von GOtt oder durch Menschen verordnet wird. Denn dann wird es im Namen des HErrn geführt zum Besten des Nächsten, wie Paulus sagt: „Sie [die Obrigkeit] ist GOttes Dienerin, dir zugut.“ So ist es verboten, des Fleisches zum geschlechtlichen Umgang zu gebrauchen, weil es eine unordentliche (vaga), böse Lust ist. Aber wo durch die Ehe das Geschlecht dir zugestellt ist, da ist es eine Pflicht, des Fleisches zu gebrauchen, welche geleistet werden muß, und nun magst du im Namen des HErrn des Fleisches gebrauchen, und thust deine Schuldigkeit gegen das göttliche Gesetz, das heißt, gegen die Liebe. Auf dieselbe Weise soll man des Eides gebrauchen, daß du nicht für dich schwörest, sondern GOtte oder dem Nächsten im Namen des HErrn. So wirst du immer in dem Dienste GOttes allein bleiben.

\begin{schriftzitat}
V. 16. Ihr sollt den HErrn, euren GOtt, nicht versuchen.
\end{schriftzitat}

Wie er in dem vorhergehenden Worte von der Furcht GOttes gelehrt hat, daß wir uns im Wohlergehen recht verhalten, damit wir nicht sicher und unbesorgt, und gewiß seien, daß GOtt sich unser annehme, der uns nicht verlasse, sondern nahe bei uns sei und gegenwärtig in allen unseren Nöthen, was die Ungläubigen und Gottlosen nicht thun, welche an den [sichtlichen] Dingen hangen.

Es wird aber GOtt in zwiefacher Weise versucht. Erstlich, wenn wir der nothwendigen Dinge, die zur Hand sind, nicht gebrauchen, sondern andere suchen, welche nicht da sind, wie der Teufel Christum versuchte, indem er ihn aufforderte, sich von der Zinne des Tempels hinabzustürzen, während doch Treppen da waren, auf denen er hinabsteigen konnte. So würde der [GOtt] versuchen, welcher bei kaltem Wetter das Kleid nicht gebrauchen wollte, sondern ein Zeichen vom Himmel erwartete, damit er nicht friere, gleichwie die Juden die Zeichen verachteten, welche da waren, und ein anderes begehrten vom Himmel. Gleicherweise [versucht derjenige GOtt], welcher schnarcht und nicht arbeiten will, und vermeint, daß GOtt ihn ohne Arbeit ernähren müsse, während GOtt ihm zugesagt hat, daß er ihn durch die Arbeit versorgen wolle, wie es heißt Spr. 10, 4.: „Die Hand der Fleißigen macht reich, aber die lässige Hand wird hungern.“

Solcher Art ist der gemeine ehelose Stand (coelibatus), da GOtt durch Erschaffung und Zugesellung des Weibes Rath geschaffen hat für die Sünde und die Schwachheit des Fleisches, und die thörichten Leute dies fahren lassen, und sich vermessen, Keuschheit zu halten durch ein Wunder vom Himmel. Es ist auch oben gesagt worden, wie GOtt unter der Larve der vorhandenen Dinge seine Werke erzeige, und wolle, daß wir derselben gebrauchen, aber nicht auf dieselben vertrauen. Denn wiewohl es wahr ist, daß eine fleißige Hand reich macht, so ist es doch zugleich auch wahr, daß allein der Segen des HErrn reich macht, wie derselbe Salomo [Spr. 10, 22.] sagt, nämlich durch eine fleißige Hand. Denn wenn eine fleißige Hand mit Gewalt gehindert würde, so würde dennoch der Segen des HErrn reich machen. So gibt GOtt allein durch das Schwert Errettung, und dennoch ist die Errettung (salus), die durch einen Menschen geschieht, eitel (wie David [Ps. 44, 7.] sagt): „Und mein Schwert kann mir nicht helfen.“ Es wird aber GOtt helfen durch das Schwert, wenn es da ist, und ohne Schwert, wenn man es nicht haben kann. Man muß also
\quellseite{723}
der Dinge gebrauchen, aber nicht darauf vertrauen. Aber allein auf GOtt muß man vertrauen, und zwar ebensowohl, wenn das da ist, dessen du gebrauchen mögest, als wenn es nicht da ist.

Zweitens wird GOtt versucht, wenn nichts vorhanden ist, dessen man bedarf, als das bloße und alleinige Wort GOttes. Von dieser Versuchung redet Moses hier eigentlich, da er hinzufügt: „Wie ihr ihn versuchtet zu Massa“, wo sie zankten und sagten [2 Mos. 17, 2. 7.]: „Ist der HErr unter uns oder nicht?“ Denn hier sind die gottlosen Leute mit dem Worte nicht zufrieden; wenn GOtt das, was er verheißen hat, nicht thut zu der Zeit, an dem Orte, auf die Weise, welche sie vorschreiben, so weichen sie von ihm und glauben nicht. Aber wenn man GOtte den Ort, die Zeit oder die Weise vorschreiben will, so ist das recht eigentlich ihn versuchen, und gleichsam mit Händen greifen wollen, ob er da sei. Das ist aber nichts Anderes als GOtt einschränken und ihn unserem Willen unterwerfen wollen, und ihm die Gottheit gänzlich nehmen, welche frei und unbeschränkt und unbegrenzt sein muß, und vielmehr uns Ort, Weise und Zeit vorschreiben.

Daher ist beiderlei Versuchung wider das erste Gebot, sowohl die, welche bloß aus Willkür und Vorwitz geschieht, wenn die Dinge reichlich vorhanden sind, als auch die, welche eintritt, wenn der Mangel drängt, und die Schwachheit des Glaubens dazu überredet. Und hier siehst du, daß das erste, allergeistlichste Gebot durch Moses auf das geistlichste und vollkommenste ausgelegt wird. Denn er ist noch nicht zu den Bildern gekommen, sondern mit dem Verdamm(en) der Neigung zur Gottlosigkeit, welche die Wurzel des äußeren Götzendienstes und der Bilder ist, beschäftigt er sich an erster Stelle.

\begin{schriftzitat}
V. 20–24. Wenn dich nun dein Sohn heute oder morgen fragen wird.
\end{schriftzitat}

Du siehst, wie sehr es sich Moses angelegen sein läßt, das erste Gebot nach seinem geistlichen Verstande einzuschärfen, und wir haben gezeigt, daß er auf diese Weise am meisten auf das Gebot gedrungen wissen will, nämlich daß man wisse, daß, wenn nicht vor allen Dingen (primo loco) diese Meinung und Gesinnung richtig gestellt ist, es nichts sei, wenn sie auch alle Bilder zerstören und alles thun; wenn aber hierin das Herz recht steht, so wird alles andere in rechter Weise geschehen. Bei den anderen Geboten hält er sich nicht so lange auf, gibt sich nicht so viel Mühe, macht auch nicht so viele Worte.

\begin{schriftzitat}
V. 25. Und es wird unsere Gerechtigkeit sein vor dem HErrn, unserm GOtt.
\end{schriftzitat}

Was ist dies? Wird denn die Gerechtigkeit vor dem HErrn durch Werke zuwege gebracht? Das sei ferne, sondern (wie ich gesagt habe) er hat mit dem ersten Gebote und seinem hauptsächlichen Werke zu schaffen, nämlich daß er den Glauben lehre, die Liebe und die Furcht GOttes von ganzem Herzen. Denn durch diese wird das Gesetz erfüllt und GOtte genuggethan. Darum wird hier an dieser Stelle die Gerechtigkeit richtig gesetzt und gerühmt, welche vor GOtt gilt. Denn dann haben wir in Wahrheit die Gerechtigkeit, welche uns von GOtt zugerechnet wird (welche eigentlich das Wort \hebraeisch{צדקה} bedeutet Ps. 33, 5., indem es die zurechnende Barmherzigkeit einschließt), wenn wir glauben, lieben, GOtt fürchten, und zwar von ganzem Herzen. Denn der Gerechte wird seines Glaubens leben. Deshalb setzt Moses diesen Ausspruch nicht, wo er von anderen Werken des Gesetzes redet, weil sie nicht rechtfertigen, sondern vielmehr gerechtfertigt werden durch das Herz, welches bereits durch den Glauben gerechtfertigt ist. Denn er rühmt allein das, daß er Gebote und Rechte etc. mache, damit es ihm [Israel] wohlgehe etc. [Cap. 6, 1. 3.].

\quellseite{724}
\section{Das siebente Capitel}

\begin{schriftzitat}
V. 1. f. Wenn dich der HErr, dein GOtt, ins Land bringet etc., wirst du sie schlagen etc.
\end{schriftzitat}

Moses fährt fort, auch andere Gelegenheiten zur Gottlosigkeit wider das erste Gebot anzuzeigen, und geht in das äußerliche Werck, in einer gar richtigen Ordnung; da er zuerst das Hertz in die rechte Stellung gebracht hat (correxerit), die Quelle der Werke, kann er nachher auch die Hand und die anderen Glieder richtig stellen, indem er lehrt, daß sie die Völker in dem Lande tödten und verbannen sollen, auch keine Ehen mit ihnen schließen noch Bündnisse mit ihnen eingehen sollen, ferner die Altäre, die Säulen, die Haine und die Bildwerke zerstören. Er hält auch in diesem Werke die Ordnung, daß er zuerst befiehlt, die Urheber der Bildnisse zu verderben, darnach auch die Bilder, weil die Bilder vergebens abgethan werden, wenn man die Verfertiger und die Lehrer derselben als ihre Verehrer bleiben läßt.

Hier ist zu beachten, daß er nicht schlechthin das Urtheil spricht, daß man diese Völker ausrotten solle, sondern, wenn sie in ihrer Verstocktheit beharren. Sonst hätte man ihnen Frieden anbieten müssen und sie tragen, wenn sie sich zu Israel bekehrten, wie den Gibeoniten und der Hure Rahab geschah. Sodann hat er dies Werck befohlen, nicht, weil er wollte, daß dieses für sein Volk zu allen Zeiten vonnöthen gewesen wäre, sondern weil er beschlossen hatte, diese Völker wegen ihrer Sünden gänzlich auszurotten; zu diesem Werke wollte er sein Volk gebrauchen, während er Sodom ohne ein anderes Volk umkehrte, und sonst ein Volk durch das andere zu strafen pflegt. Darum darf man diesen buchstäblichen und fleischlichen Verstand des ersten Gebotes nicht auf die Christen ziehen, denen es zukommt, mit dem Schwerte des Geistes die Heiden zu tödten und die Bilder abzuthun. Denn dies Werk ist insonderheit diesem Volke eine Zeitlang aufgelegt, wie auch alles andere, was diesem Volke befohlen wird, desgleichen von der Ehe, von Bündnissen und allen äußerlichen Ceremonien insgesamt.

Hier muß ich [ein wenig] abschweifen und von unseren neuen Propheten reden, welche sich rühmen, daß sie durch das erste Gebot gedrungen worden (wiewohl sie nicht Juden, sondern Christen sein wollen), die Bilder mit Gewalt und mit der Hand zu zerstören. Und hier müthen sie außerordentlich und rühmen sich, sie seien ganz voll des Geistes, und verdammen alle als schuldig des ersten Gebotes und der Majestätsbeleidigung. Ich aber behaupte, daß sie Bluthunde, Aufrührer und Mörder sind, welche nichts Anderes im Sinne haben als Mord und Todtschlag. Dies kann ich leicht aus dieser Stelle beweisen, auf diese Weise: Wenn sie wollen, daß man im ersten Gebote nichts habe, was Ceremonien betreffe (ceremoniale), sondern behaupten, daß alles nothwendiger Weise gethan werden müsse, so werden sie durch diese Stelle gezwungen werden, die ganze Welt zu tödten, welche Bilder hat, wenngleich sie dieselben nicht anbetet. Denn es wird hier ebenwohl geboten, daß sie die Völker tödten sollen, als daß sie die Bilder zerstören, und wenn eins behauptet wird, werden sie das Uebrige mit Nothwendigkeit zugestehen müssen; ja, er befiehlt hier, zuerst die Heiden zu tödten, als das nothwendigere Werk, ehe man die Bilder zertrümmere.

Und dies ist es, was ich kommen sah, wenn ihre Lehre durchdringen würde, daß das Volk nicht zurückgehalten werden könnte vom Morden, da es sich verließe auf dieses Gebot und ihre Lehre. Daher habe ich gerathen (censui), daß man sie aus dem Lande jagen solle als rechte Bluthunde und Aufrührer, die in ihrem ganzen Leben nichts Anderes zu schaffen hätten, als daß sie Blut vergössen. Da es nun feststeht, daß bei den Christen die Gottlosen nicht mit dem leiblichen Schwerte getödtet werden müssen, und dieser Theil des Gesetzes als ein bloß ceremoniale aufgehoben ist, da er nur für eine Zeitlang galt bis auf Christum, so folgt zugleich, daß auch die Bilder nicht von Noth wegen abgethan werden müssen, da sie zu diesem Theile des Gesetzes gehören. Darum soll niemand zweifeln, der die Bilderstürmer (imaginicidas) so müthen sieht wider Holz und Stein, daß da ein mörderischer Geist verborgen sei, nicht ein lebendigmachender, und daß sie, wenn die Gelegenheit
\quellseite{725}
sich böte, auch die Menschen tödten würden, wie etliche derselben angefangen haben zu lehren. Sie werden aber durch eine nothwendige Folgerung gezwungen, so zu lehren und zu thun, indem sie das Gesetz Mosis dazu bringt, wiewohl Andere, listige Leute, diesen mordgierigen Geist gar fein verbergen.

Es ist auch noch ein anderes Wüthen dieses Geistes, daß sie die Vollstreckung dieses Gesetzes, sogar wenn es uns betreffen möchte, nicht richtig [auf die Personen, denen es zukommt,] vertheilen. Denn Moses hat sowohl in diesem Buche als auch anderswo zuerst die obrigkeitlichen Personen eingesetzt, welche die Gesetze ausüben sollen. Diese Schwärmgeister befehlen dies aber dem gemeinen Volke mit Verachtung der obrigkeitlichen Personen. Denn GOtt hat nicht bloß an einer Stelle befohlen, daß man die Sünden mit öffentlichem Gerichte, Zeugen und Urtheil handeln solle, sie aber fallen darüber her mit eigenem (privato) Wüthen. Denn man liest kein Beispiel, daß das Volk ohne einen Führer oder eine Obrigkeit, sei es nun die ordentliche, oder eine von GOtt gegebene, die Bilder abgethan habe, wie man an Gideon [Richt. 6, 25.], Hiskia [2 Kön. 18, 4.], Josia [2 Kön. 23, 4.] und Jehu\footnote{Erlanger: Ahab statt Jehu.} [2 Kön. 10, 27. ff.] sehen kann.

Ihr drittes Wüthen ist, daß sie durchaus alle Bilder zerstören, wiewohl Moses dies allein von denen geboten hat, welchen man dient und denen man vertraut. Dies zeigt nicht allein der Text selbst an, sondern auch die Absicht des ersten Gebotes, da er sagt, man solle kein Gleichniß GOttes machen, um es anzubeten. Aber andere Bilder als GOttes, zumal die nicht angebetet werden, verbietet GOtt nirgends. Ja, er selbst hat die eherne Schlange in ebendemselben Volke aufgerichtet und geduldet, bis daß man anfing, sie zu verehren. Ja, auch die Rubeniter errichteten einen Altar am Jordan [Jos. 22, 10.], wovon die anderen [Stämme] urtheilten, daß es gottlos sei, aber da sie vernahmen, er sei nur zu einem Zeichen und Denkmal erbaut, nicht um darauf zu opfern und Gottesdienst zu verrichten, ließen sie ihn unangetastet. Ueberdies heißt es 3 Mos. 26, 1. klarlich, daß man keine Bilder machen solle, um sie anzubeten.

Darum wollen wir diese Bluthunde\footnote{viros sanguinum. Wegen der Uebersetzung vergleiche 2 Sam. 16, 7.} meiden und uns nicht ins Judenthum ziehen lassen.

Zu uns sagt Paulus [1 Cor. 8, 4.]: „Wir wissen, daß ein Götze nichts in der Welt sei“, und daß alle diese äußerlichen Dinge frei sind, wenngleich es Bilder sind, die zu irgend einer göttlichen Verehrung bestimmt sind. Denn wir sollen sie durch das Wort abthun, oder durch die öffentliche (communi) Einwilligung der Obrigkeit und derer, unter deren Gewalt (ditione) sie sind, beseitigen. Die wir aber nur zu einem Zeichen und Denkmal haben, mögen wir frei haben, damit wir nicht endlich auch in den Geist des Mordes und des Aufruhrs verfallen, wenn wir leiden, daß aus der Freiheit eine Nothwendigkeit gemacht werde. Denn man möchte diese wüthenden Leute immerhin tragen, wenn sie nur die Bilder zerstörten, und nicht auch das Gewissen machten, es sei ein nothwendiges Werk, und uns unter den Zorn des Gesetzes schickten und uns der Freiheit beraubten. Jetzt aber, da wir in der Freiheit bestehen müssen, die GOtt uns geschenkt hat, sollen wir ihnen sagen, daß Mosis in allen seinen Gesetzen uns nichts angehe, sondern nur die Juden, außer wo er mit dem natürlichen Gesetze übereinkommt, von dem Paulus Röm. 2, 15. lehrt, daß es in den Herzen der Heiden beschrieben sei. Alles, was da nicht geschrieben ist, sollen wir zählen zu den Ceremonien, welche dem Volke Mosis vonnöthen sind, aber frei für uns, wie auch der Sabbath ist, wie Paulus, Col. 2, 16., und Jesaias, Cap. 66, 23., bezeugen.

Man muß sich aber wundern, warum diese Feinde der Bilder so fromm und gütig sind gegen die Bilder auf goldenen und silbernen Münzen, desgleichen gegen die, welche auf silbernen Gefäßen angebracht sind. Warum lieben sie nur diese Bilder und verbrennen sie nicht auch, oder werfen sie weg? Können wir denn nicht hier die Bosheit des Satans verstehen, der in ihren Herzen durch den größten Geiz und die höchste Heuchelei regiert? Ferner, warum zerfleischen sie denn nicht auch ihre Herzen, da sie nicht ohne ein Bild sein können, so oft sie hören, daß entweder Christus, der Gekreuzigte, gepredigt wird, oder sie selbst an ihn denken oder an andere Heilige? Oder ist ein Bild in den Augen, das außer uns ist, schädlicher als das, welches im Herzen in uns ist? Es ist ein Wüthen und eine Unsinnigkeit, wodurch sie nur Ehre suchen, als ob es ein treffliches Werk wäre. Uns ist es
\quellseite{726}
genug, zu wissen, daß ein Götze nichts in der Welt sei. Wenn er nichts ist, so wird er nicht schaden, er stehe oder falle. Doch auch ich liebe die Bilder nicht gar sehr, und wollte, daß sie in der Kirche nicht aufgestellt würden, nicht, weil ich allein das im Auge habe, daß sie angebetet werden, was, wie ich glaube, sehr selten geschieht, sondern weil ein Vertrauen auf das Werk gesetzt wird, da sie so viel kosten und so schönen Schein haben, als ob durch dieses Werk GOtte ein Dienst geleistet werde, während doch unterdessen die Kosten und alles verloren ist, was darauf verwendet wird, das zu einem besseren Brauche für die Noth der Brüder verwendet werden sollte. Sonst kann ich es nicht verdammen, daß man in einem Privathause schön gemalte (levi pictura) Bilder habe. Aber weil andere diese Ursache ungezeigt haben, und die rasenden Propheten sehen, daß sie darin keinen Ruhm erlangen können, so erdichten sie eine Nothwendigkeit des Gesetzes wider die Freiheit des Geistes, was auf keine Weise zu dulden ist.

\begin{schriftzitat}
V. 2. So sollst du sie verbannen.
\end{schriftzitat}

Dieses Wort [verbannen] wird häufig gebraucht in den Büchern Mose, Josua und der Richter, welches hebräisch \hebraeisch{חרם}, lateinisch excommunicare heißt. Daher kommt \hebraeisch{חרם} oder \hebraeisch{חרם}, Bann, griechisch ἀνάθεμα. Daher sagt Paulus 1 Cor. 16, 22.: „So jemand den HErrn JEsum Christum nicht lieb hat, der sei Anathema, Maranatha.“ So wurden die Heiden durch das Schwert der Kinder Israel verbannt und wurden ein Fluch, indem sie ganz und gar zu nichts gemacht und unwiederbringlich zu Grunde gerichtet wurden. Denn GOtt wollte sie durch das beschlossene Urtheil ganz und gar von der Erde vertilgen, auch mit einem Fluche. Dies nennt Moses eigentlich Haram oder Herem, so daß keine Hoffnung des Lebens oder der Errettung übrig wäre, gleichwie das, was verbannt wird, vernichtet und ein Fluch ist, 3 Mos. 27. Aber obgleich diese Völker des Todes würdig waren, so wäre es doch niemandem, auch Israel nicht, erlaubt gewesen, sie zu tödten, wenn sie nicht durch ein bestimmtes und klares Gebot und Wort GOttes gezwungen gewesen wären, damit der Ausspruch feststehe [Matth. 26, 52.]: „Wer das Schwert nimmt, der soll durchs Schwert umkommen“, und wiederum [Röm. 12, 19.]: „Die Rache ist mein, ich will vergelten“, desgleichen [Hebr. 10, 30.]: „Der HErr wird sein Volk richten.“ Denn er, der das Leben gegeben hat, kann es mit Recht wegnehmen, gegen den sie auch allein gesündigt haben. Daher hat er hier der Israeliten zu seinem Dienst gebraucht, daß er seinen Zorn durch ihre Hände ausrichtete, gleichwie er von Anfang an immer ein Volk durch das andere zu verderben pflegte, und ein Reich durch das andere [an einen anderen Besitzer] zu übertragen, wo sie gesündigt hatten, wie er in den Propheten Daniel [Cap. 2, 44.] und Amos [Cap. 9, 8.] spricht: „Siehe, ich komme über ein sündiges Königreich, daß ich es zermalme.“ Aber darin sind die Schwerter anderer Völker unterschieden von den Schwertern Israels, daß GOtt sich der Wuth jener bedient aus einem geheimen Gerichte, und die Gottlosen durch die Gottlosen zermalmt, aber das Schwert dieser geheiligt wird durch ein offenbares und gewisses Gebot GOttes, so daß mit heiligem und reinem Gewissen die Gottlosen die Gottlosen verderben, und das Blut vergießen in heiligem Gottesdienste (religione).

Das ist es, weshalb er durch ein Urtheil entscheidet, daß die Amoriter und die Cananiter getödtet werden sollen, und die Israeliten nicht gegen irgend welche Leute wüthen läßt, wie es sie gelüstet. Denn er hat die Moabiter, Ammoniter und Edomiter ausgenommen, welche nicht weniger gottlose Völker waren als die Amoriter und Cananiter, um anzuzeigen, daß man auch noch so große Sünder dennoch nicht mit gutem Gewissen tödten dürfe, wenn sie nicht durch den Befehl und Willen GOttes überantwortet und mit Namen genannt sind. Denn wer ist unter den Menschen, der vor GOtt nicht des Todes werth sei, ausgenommen derjenige, den GOtt, indem er mit seiner Barmherzigkeit zuvorkam, zuerst durch sein Wort lebendig gemacht und ihm dann befohlen habe, andere zu tödten? So daß die Juden keine Ursache haben, sich zu rühmen, als ob sie allein unter den Menschen dazu geboren wären, die Heiden zu tödten, während sie von GOtt nur dazu erwählt waren, daß GOtt seinen Zorn an den Cananitern durch ihr Schwert ausübte.

Und dies zeigt Moses hier gar schön an, da er sagt [V. 6.]: „Denn dich hat der HErr erwählt
\quellseite{727}
, bist auch nicht durch dein Verdienst dazu gekommen, daß du diese Völker tödtest, sondern dadurch, daß GOTT dich dazu erwählt, berufen und es dir befohlen hat. Ja, damit du wissest, daß kein Unterschied sei vor GOTT zwischen dir und den Heiden, so nimm das zu Herzen, daß, wenn du das Wort GOttes nicht erfüllst und diese Heiden nicht tödtest, sondern mit ihnen Bündnisse und Ehen schließest, dich auch derselbe Zorn erwartet. Er spricht [3. 4.]: „So wird dann des HErrn Zorn ergrimmen über dich, und dich bald vertilgen.“ Dieselbe Strafe zeigt an, daß sie vor GOTT derselben Gottlosigkeit schuldig sein werden, so daß sie wissen sollen, daß sie allein durch das Wort GOttes einen Vorzug haben vor jenen Heiden; wo sie das nicht hätten, so könnten keine Heiden so gottlos sein, daß sie ihnen nicht gleich sein würden sowohl der Schuld als auch der Strafe nach.

\begin{schriftzitat}
V. 7–11. Nicht hat euch der HErr angenommen und euch erwählt, daß euer mehr wäre, denn alle Völker.
\end{schriftzitat}

Siehe, wie fleißig Moses das erste Gebot handelt. Ohne hatte er den wahren Gottesdienst im Glauben, in der Furcht und in der Liebe des Geistes gelehrt. Hier aber, nachdem er zum äußerlichen Werke gekommen ist, damit sie nicht das, was sie im Geiste angefangen hatten, im Fleische vollenden und wegen ihres Werkes aufgeblasen werden möchten, und da sie äußerlich die Gottlosen und die Götzen vertilgt hätten, selbst ärgere Abgöttische würden und schlimmere Götzen in ihren Herzen aufrichteten, nämlich Vertrauen auf das Werk und Ruhm von demselben (wie das Fleisch zu thun pflegt), und eben dadurch, aus Anlaß eines so scheinbaren Werkes, schwerer wider das erste Gebot sündigen möchten, als die Heiden, welche sie tödteten: so kommt er ihnen heilsamlich zuvor und erhält sie in dem rechten Brauch des Gesetzes, indem er den Ruhm des Werkes und das Vertrauen darauf gänzlich aufhebt und sie allein zum Vertrauen auf die Gnade beruft, indem er sagt: „Nicht, daß euer mehr wäret, denn alle Völker“ etc. Als ob er sagen wollte: Daß GOTT eures Schwertes gebraucht, geschieht nicht, weil ihm eure Kraft vonnöthen wäre, oder er es nicht ohne euch vermöchte, da euer sehr wenige sind. Der Ruhm des Werks ist nicht euer, sondern dessen, der sich eurer geringen Zahl bedient und

dadurch eine so große Menge verderben wird. Sonst, wenn er durch die Menge hätte siegen wollen, so hätte er euch nicht erwählt, sondern andere Völker, deren mehr ist, denn ihr seid.

Was bleibt also übrig, dessen ihr euch in diesem Werke rühmen könntet? Nichts ist euer, „sondern, daß der HErr euch geliebt hat (sagt er), und daß er seinen Eid hält“ etc. Du siehst, daß in einem jeglichen Werke dem menschlichen Vertrauen nichts vorgehalten wird als die unverdiente Liebe GOttes, durch die er bewogen wird, daß er uns zuvorkomme mit seinem Worte und mit seiner Verheißung, auch ehe wir geboren werden. So viel fehlt daran, daß er irgend etwas wiedervergelte, nachdem wir geboren sind oder irgend ein Werk gethan haben. Und dies ist der lautere und reine Verstand des ersten Gebots, daß wir erkennen, daß wir nichts durch unsere Verdienste haben, sondern alles allein durch seine Barmherzigkeit und Liebe haben, erlangen, vermögen und thun zu seiner Ehre. Wie er diese Barmherzigkeit zuerst durch sein Wort zusagt, so bestätiget er dieselbe darnach auch durch das Werk, welches er durch uns thut, wie mit einem Zeichen, wie er hier den Auszug aus Egypten und die Vertilgung der Canaaniter anführt.

Du mußt also merken, daß es das Werk eines großen Geistes ist, daß man die Dinge, welche einen großen Schein haben durch das Vorwenden von GOttes Namen und GOttes Dienst, vertilgen könne, als da sind die Altäre und Gottesdienste der Heiden. Denn wen sollte das nicht abschrecken und betrügen, wenn man ihm GOttes Namen entgegenhält? ja, wen sollte jene erheuchelte Frömmigkeit nicht aufblasen? Sicherlich verehrten diese Heiden, wie ich gesagt habe, nur den wahren GOtt, aber mit falschem und selbst erfundenem Dienste und Meinung, die von GOtt nicht geboten waren. Aber ein viel größerer Geist gehört dazu, daß man sich dessen nicht rühme, daß man einen solchen Gottesdienst vernichtet habe, und anstatt des äußeren Götzen nicht einen inwendigen Götzen eines ganz nichtigen Vertrauens aufrichtete. Das ist eine so große Sache, daß man den wahren und einzigen GOtt rein erkenne, aufrichtig fürchte, verehre, liebe und glaube, daß Moses so viele Worte darauf verwenden muß, um das erste Gebot zu erklären. Daher wiederholt er auch hier die Worte des ersten Gebotes und sagt [V. 9.]:

\quellseite{728}
„So sollst du nun wissen, daß der HErr, dein GOtt, ein treuer GOtt ist, der den Bund und Barmherzigkeit hält“ \&c., damit er zeige, daß dies zum Verständniß des ersten Gebotes gehöre, was er von dem Vertrauen auf die Werke gesagt hatte.

\subsection{A. 12–16. Und wenn ihr diese Rechte hört, und haltet sie \&c}

Was ist dies? Bisher hat er gelehrt, daß das Volk geliebet sei, nicht durch sein Verdienst, sondern durch die Erwählung GOttes, der ihnen zuvorkommt. Jetzt verheißt er den Werken und Verdiensten die Liebe, den Segen, ja, was noch mehr ist, die Zuverlässigkeit (fidem) des Bundes und der Verheißung schreibt er ihren Werken zu, indem er spricht: „Wenn ihr haltet, so wird der HErr auch halten den Bund und Barmherzigkeit“, was er hier durch einen langen Text von dem zeitlichen Wohlergehen auslegt. Ich antworte: Es ist die von uns schon so oft gegebene Unterscheidung der vorhergehenden und der nachfolgenden Barmherzigkeit bekannt. Die vorhergehende (praeveniens) ist die, durch welche wir erwählt, berufen und gerechtfertigt werden vor allen unseren Werken. Die nachfolgende ist, da wir uns selbst durch die Werke die Barmherzigkeit gewiß machen und die vorhergehende Barmherzigkeit empfinden, von der gesagt ist [Mal. 3, 7.]: „Bekehrt euch zu mir (das geschieht durch die vorhergehende Barmherzigkeit), so will ich mich zu euch auch kehren“, das ist die nachfolgende Barmherzigkeit, welche den schon Gerechtfertigten Frieden, Sicherheit und alles Gute verleiht, so daß man sagen kann, die eine sei die Gnade, die andere Friede, oder, die eine sei die Gabe und die Sache, die andere sei die Gewißheit der empfangenen Gabe und der Sache, die man besitzt. So will er hier sagen: Wenn du diese Gebote hältst (was sie doch ohne die Gnade nicht vermochten), so wirst du dadurch, daß du die Dinge erfährst, inne werden, daß der HErr treu ist, daß er das verheißene Land schenke, und, nachdem er es dir gegeben hat, auch erhalte.

\begin{schriftzitat}
V. 17–20. Wirst du aber in deinem Herzen sagen: Dieses Volk ist mehr, denn ich bin \&c.
\end{schriftzitat}

Moses nimmt einen anderen Anlaß zur Uebertreibung des ersten Gebots hinweg, nämlich den Mangel an Vertrauen, welcher entsteht aus dem Unvermögen und der Schwachheit unserer Kräfte, wenn dieselbe verglichen wird mit der Sache und dem gebotenen Werke. Dadurch sündigten die Kundschafter, welche das Volk vom Eingehen in das Land abschreckten, da sie viel Wesens davon machten, daß Riesen, Enakim, und befestigte Städte in demselben wären, wie wir oben gesehen haben.

Es lehrt aber Moses hier wiederum, wie GOttes Gebote erfüllt werden müssen, nämlich nicht durch unsere Kräfte, sondern durch die uns verheißene göttliche Kraft, indem er mit Einem Worte zugleich sowohl das Vertrauen und die Vermessenheit auf unsere Kraft auslöscht, als auch das Mißtrauen und die Verzweiflung wegen unserer Schwachheit; die Vermessenheit dadurch, daß er anzeigt, es werde Größeres geboten, als wir vermögen, wie er hier zugibt, daß der Heiden, welche er zu vertilgen gebietet, mehr seien und stärker als Israel, ihr Zerstörer; die Verzweiflung dadurch, daß der HErr durch sie alles thun werde was er ihnen gebietet.

Wenn daher die Kinder Israel auf ihre Kräfte gesehen und sie verglichen hätten mit den Kräften der Heiden, welche vernichtet werden sollten, so hätten sie ganz und gar daran verzweifelt und das Gebot GOttes unbeachtet gelassen. Jetzt aber verheißt GOtt, daß er ihnen beistehen werde, damit sie alles im Glauben erfüllen möchten, und er fügt zu der Verheißung das Exempel seiner vorigen Barmherzigkeit, daß er sie aus den Händen Pharao’s und der Egypter befreit hat, damit er sie so durch Wort und That zum Glauben reize, in welchem sie auf Befehl GOttes die Heiden vernichten sollten. Ja, er sagt ihnen auch durch eine neue Verheißung zu, daß er Hornisse [unter die Feinde] senden werde, damit du sehen, wie viel daran gelegen sei, daß der Glaube gestärkt werde. Durch den Glauben dient man also GOtte, durch den Glauben werden die Gebote GOttes erfüllt, durch den Glauben verdienen wir, daß die göttliche Kraft uns beistehe in allen unseren Werken, so daß Christus mit Recht gesagt hat [Marc. 9, 23.]: „Alle Dinge sind möglich dem, der da glaubet.“

\begin{schriftzitat}
V. 21. Laß dir nicht grauen vor ihnen; denn der HErr, dein GOtt, ist unter dir, der große und schreckliche GOtt.
\end{schriftzitat}

Lieber, wen sollte dieses herrliche Wort nicht aufgeblasen machen und zwingen, stolz zu sein, nicht allein wider jene Heiden, die ausgerottet
\quellseite{729}
werden sollten, sondern auch wider alle Pforten der Hölle? Was ist es Wunder, wenn Einer tausend in die Flucht schlägt, und wie David sagt [Ps. 18, 30.]: „Mit meinem GOTT kann ich über die Mauern springen“? Was kann Stolzeres und Hoffährtigeres geredet werden: „Der große und schreckliche GOTT ist mitten unter uns“? Ps. 46, 6.: „GOTT ist bei ihr darinnen, darum wird sie wohl bleiben.“ So heißt es auch 1 Joh. 4, 4.: „Der in uns ist, ist größer, denn der in der Welt ist.“ Es mögen also jene Heiden Enakim, Riesen, genannt sein, was liegt daran? Wer ist in ihrer Mitte? Nichts als das eitle Vertrauen auf ihren Arm und auf das Schwert des Fleisches, hier aber der große GOTT der Geister. Das wird nun treulich genug gesagt. Wohl dem, der es fest glaubt, denn es wird geschehen, wie er geglaubt hat.

\begin{schriftzitat}
V. 22. Du kannst sie nicht eilend vertilgen, auf daß sich nicht wider dich mehren die Thiere auf dem Felde.
\end{schriftzitat}

Wie lieblich und väterlich zeigt er sie zum Glauben, so daß er auch den künftigen schwachen Gedanken des Glaubens zuvorkommt, damit sie nicht beständig denken möchten, wenn nun die Verheißungen GOttes anfingen erfüllt zu werden, und die Heiden größtentheils geschlagen wären, aber doch noch einige übrig blieben, sie seien verlassen, oder sie seien mit der Verheißung betrogen, sondern es geschehe alles ihnen zum Besten, damit sie das Land desto fester und völliger einnehmen möchten, und nicht gezwungen würden, anstatt der gottlosen vertilgten Menschen noch grausamere wilde Thiere zu erdulden, wiewohl gerade dieser Aufschub auch den Heiden zu ihrer Bekehrung gegeben wird, damit diejenigen, welche wollten, wieder zurechtkommen möchten.

Jetzt wollen wir die heimlichen Deutungen dieses Capitels ein wenig behandeln.

Unser Schwert ist das Wort GOttes, mit welchem das geistliche Volk kämpft und zuerst die Heiden selbst tödtet, das heißt, sie von dem Irrthum bekehrt, in welchem sie lebten. Jes. 11, 4.: „Und er wird mit dem Stabe seines Mundes die Erde schlagen, und mit dem Odem seiner Lippen den Gottlosen tödten.“ Die Bildwerke sind die gottlosen Lehren, welche aus gottlosen Meinungen hervorgehen, durch welche man meint GOTT zu verehren durch Werke, ohne den Glauben. Die Altäre sind die Bemühungen und Uebungen, mit welchen wir uns nach diesen Lehren anstrengen und vergeblich peinigen und völlig zu Grunde richten (mactamus). Die Haine sind die Lockungen des Fleisches, durch welches wir uns in diesen Werken wohlgefallen und rühmen, als ob sie vor GOTT angenehm wären. Dies alles muß durch das Wort des Glaubens zu nichts gemacht werden, damit GOTT allein unter uns sei, wie er hier sagt.

Daß Hornisse vor uns und unter die Heiden geschickt werden, das ist, daß die Gottlosen durch das Schrecken des Gesetzes und den tödtlichen Stachel des Gewissens getrieben werden, so daß sie sich nicht verbergen können, bis daß sie zu Grunde gehen, und ihre Gottlosigkeit zu bekennen gezwungen werden. Ferner, wie es nicht in der Hand der Israeliten stand, Hornisse vor ihnen herzuwenden, so steht es auch nicht bei uns, irgend einen Gottlosen zu erschrecken, sondern GOTT ist es, der ihre Herzen in Angst setzt und furchtsam macht, so daß sie erschreckt mit Paulus sprechen [Apost. 9, 6.]: „HErr, was willst du, daß ich thun soll?“ und mit denen in der Apostelgeschichte, Cap. 2, 37.: „Lieben Brüder, was sollen wir thun?“

Daß die Heiden nicht eilend, sondern allmählig vertilgt werden sollen, damit die wilden Thiere im Lande sich nicht mehren, das ist, daß, wenn die Gottlosigkeit des Herzens überwunden ist, der Gottlose noch nicht ganz und gar gereinigt ist, sondern noch der Kampf des Geistes und des Fleisches übrig sei, wie Paulus Röm. 7, 14. ff. und Gal. 5, 17. ff. lehrt. Dieser Kampf ist in diesem Leben dazu nothwendig, damit wir nicht hoffärtig und faul seien und ärger werden, wie Paulus 2 Cor. 12, 7. spricht: „Auf daß ich mich nicht der hohen Offenbarung überhebe, ist mir gegeben ein Pfahl ins Fleisch.“ So wird in unseren Grenzen der Heviter, der Samariter und die Philister übrig gelassen, das heißt, die Ueberbleibsel der Sünde, welche uns zu schaffen machen und in den geistlichen Waffen üben sollen. Sonst, eine wie grausame und wilde Bestie ist doch die sichere Vermessenheit, die Eitelkeit, der Hochmuth, die Trägheit, indem wir unserer Schwachheit vergessen, und unseren Kräften das zuschreiben, was allein der Gnade und der Barmherzigkeit zugehört, welche über uns herrscht und uns unsere Gebrechen nicht zurechnet!

Von dem Golde der Götzenbilder soll nichts
\quellseite{730}
in das Haus gebracht werden, weder die Ueberbleibsel gottloser Worte noch Werke sollen im Gewissen bewahrt werden, sondern die Herzen sollen mit dem lauteren Worte GOttes, und damit allein, erfüllt und bewahrt werden, damit das gottlose Wesen uns nicht endlich auch anstecke. Zum Beispiel, das gottlose Wort „freier Wille“ und alles, was darüber an Lehren ausgebreitet ist, soll man ganz und gar verabscheuen und nicht zulassen, als ob es in gutem Brauche oder mit irgend einer Glosse geduldet werden könnte. So sind auch die Klostergelübde und Regeln für ein verbanntes Gold (aurum anathematis) zu halten, und man soll nur deshalb davon handeln, damit man sie strafe und verwünsche, nicht, damit unseren Gewissen dadurch geraten werde, mit welchem Glößlein man auch immer meinen möchte, daß es geschehen könne. Das Haus des Gewissens muß von ihnen völlig rein sein, weil es nichts Zarteres und Weicheres gibt als das Gewissen und den Glauben, wie man spricht: Zucht und Ehre, Glaube und Auge kein Scherz leidet [man hat alle drei bald verderbt].

\section{Das achte Capitel}

\begin{schriftzitat}
V. 1. Alle Gebote, die ich dir \&c.
\end{schriftzitat}

Dies ist eine neue und eine andere Predigt, wie ich gesagt habe, daß Moses dieses Buch in vielen Tagen und mit vielen Predigten vollendet habe. Er handelt aber in diesem Capitel von dem Anlaß zur Uebertretung des ersten Gebotes, den das Wohlgehen und der Ueberfluß gibt, der das Herz viel stärker abwendet, als Widerwärtigkeit und Mangel, wie er [Moses] in seinem Gesange sagt [5 Mos. 32, 15.]: „Er ist fett und dick und stark worden, und ward geil.“ Und Spr. 1, 32.: „Der Ruchlosen Glück bringt sie um“, wie man auch im deutschen Sprüchworte sagt: Es müssen starke Beine sein, die da sollen gute Tage tragen. Denn ein Mensch erträgt leichter Böses als Gutes, wie der Dichter sagt: Luxuria saevior hostis incubuit etc. [Der Ueberfluß ist als ein überaus grausamer Feind hereingebrochen.]

\begin{schriftzitat}
V. 2. Und gedenkest alle des Weges \&c.
\end{schriftzitat}

Zuerst sagt er, daß das Volk von GOtt in der Wüste vierzig Jahre lang geschlagen und versucht worden sei, auf daß es kund würde, was in seinem Herzen sei, nämlich, ob es GOtt mit aufrichtigem Gemüthe liebe oder um des Bauches willen. Denn auch wir selbst erkennen unser Herz (welches Gotte immer offenbar ist) nie gewisser, ob es Gotte diene um des Bauches willen oder nicht, als wenn wir von Mangel und Unglück angefochten werden. Ein gläubiges Herz liebt, glaubt und dient ebensowohl, wenn es Mangel und böse Tage leidet, als wenn es mitten im Reichthum und guten Tagen ist. Aber ein ungläubiges Herz glaubt, so lange das Gute da ist, weicht aber zurück, wenn Unglück hereinbricht, wie es im 78. Psalm [V. 37. 36.] heißt: „Ihr Herz war nicht fest an ihm, und sie logen ihm mit ihrer Zunge.“

\begin{schriftzitat}
V. 3–16. Er demüthigte dich, und ließ dich hungern, und speiste dich mit Manna, auf daß er dir kund thäte, daß der Mensch nicht lebe vom Brod allein \&c.
\end{schriftzitat}

Was ist das für eine Folgerung, daß durch das Geben von Manna kund gethan werde, daß der Mensch von dem Worte GOttes lebe? Ist denn das Manna das Wort GOttes? Keineswegs, sondern er erklärt vielmehr das, was er vom Mangel gesagt hatte, nämlich auf diese Weise: Ehe er dir Manna gab, ließ er dich zuvor Hunger leiden. Dies that er deshalb, damit er dir zeigte, daß, wenn auch niemals das Manna käme, er dich doch durch sein Wort erhalten könnte, durch welches er zugesagt hatte, er werde dir ein GOtt sein, welcher dich nicht verlassen würde, wie er dich auch erhalten hat. Denn der Glaube an das Wort GOttes speist, auch mitten im Hunger, nicht allein die Seele, sondern auch den Leib, das heißt, wie er hier sagt, den ganzen Menschen, da es unmöglich ist, daß der, welcher mit dem Glauben an dem Worte hängt, umkomme, wenn er auch nichts Leibliches essen oder trinken sollte, obgleich GOtt
\quellseite{731}
in Wahrheit nicht verläßt, ja, endlich, wenn seine Zeit gekommen ist, auch leiblich nährt, wie den Elias durch die Raben und die Witwe zu Sarepta, und hier die Väter durch das Manna.

Dieses Kundthun also, daß der Mensch nicht allein vom Brode lebe, bezieht sich nicht auf das Manna, sondern auf beides, nämlich auf die Trübsal des Hungers, und auf die Menge des Manna, so daß der Sinn ist: Diese seine Güte wollte GOtt dir durch deine Erfahrung zeigen, damit du im Hunger lerntest, auf das Wort des Verheißenden zu vertrauen, durch welches du auch im Hunger erhalten werden möchtest und nicht umkämest, sodann aber endlich durch diesen Glauben auch leibliche Sättigung empfingest. Denn dies alles thut er um deswillen, daß er dir den Bauch als Götzen nähme, und du lernest, daß daran nicht das Leben eines Menschen gelegen sei, wenn der Bauch mit seinem Brode versorgt wäre, sondern vielmehr [daran, daß] beides, die Seele und der Bauch, [mit dem Worte versorgt werde,]\footnote{Von uns hinzugefügt.} was du niemals lernen würdest, wenn der Bauch immer satt wäre, und du niemals durch das Wort lerntest, GOtt zu vertrauen auch zu der Zeit, da der Bauch Mangel leidet.

Auf diese Weise führt Christus Matth. 4, 4. diese Stelle an, wo der Teufel ihm nur die Sorge für den Bauch vorhält; er aber sagt: „Der Mensch lebt nicht vom Brod allein“, indem er nicht leugnet, daß der Mensch vom Brode lebe, aber nicht allein. Denn wenn das Brod nicht da ist, dann lebt er von dem Worte. Wenn er dies durch den Glauben besitzt, so muß nothwendiger Weise auch endlich das Brod kommen, wenn es auch aus den Steinen hervorgebracht, oder, wie hier, vom Himmel hernieder gesandt werden müßte. Da er nun sagt: „Er demüthigte dich, und ließ dich hungern, und versuchte dich“, zieht er von dem Bauche und der Bauchssorge ab, damit er lehre, mitten im Hunger dem Worte zu vertrauen, und [durch das Wort] zu leben. Da er sagt: „Er speisete dich mit Man“, zeigt er an, daß es denen an Brod nicht fehlen werde, welche durch das Wort im Hunger leben. Du siehst also, daß durch diese Worte nichts gelehrt werde als der Glaube, durch den wir festhalten und glauben, daß wir einen GOtt haben nach dem Sinne des ersten Gebotes.

Er gebietet aber, daß sie dieser Thatsache eingedenk sein sollen, wie sie in der Wüste sowohl durch das Wort als auch durch das Brod ernährt worden seien, indem er gleichsam anzeigt und prophezeit, daß sie später einmal durch ebensolche Anfechtung des Hungers versucht werden würden, damit sie dann durch den Trost dieses Exempels sich zum Glauben aufrichten sollten, und wüßten, daß sie ernährt werden würden, wie sehr auch der Hunger wüthe, wenn sie nur dem Worte der Verheißung im ersten Gebote glauben, in welchem er zusagt, daß er ihr GOtt sein will, wie es im 37. Psalm, V. 19., heißt: „Und in der Theurung werden sie genug haben.“ Denn es ist unmöglich, daß der Gerechte verlassen werde, oder sein Saame nach Brod gehe. Ihn wird zwar hungern, aber er wird nicht Hungers sterben, denn der Hunger übt seinen Glauben, der am Worte hängt; aber der Glaube erlangt auch Brod für den Leib.

Es steht also das herrliche Wort Mosis fest, daß GOtt mit den Seinen so handle, daß er durch Hunger versucht, und übt durch sein Wort, sodann den Gläubigen speist, sollte es auch mitten aus dem Himmel sein, wenn es nicht anders geschehen kann, damit sie so durch Erfahrung lernen, daß man für den Bauch nicht bestimmt sein müsse, daß auch das Leben nicht stehe in den Dingen, welche wir besitzen, oder im Brod, sondern im Worte, durch welches wir in GOtt reich werden, wie im Evangelio [Luc. 12, 21.] steht. Denn da wir durch das Wort im Herzen leben, so zwingen wir GOtt gleichsam, daß er auch den Bauch nähre.

Wie er aber den Gläubigen die Nahrung zusagt, so auch die Kleidung, sodann auch gute Gesundheit des Leibes, wie er hier bald nach dem Manna anfügt [V. 4.]: „Deine Kleider sind nicht veraltet an dir, und deine Füße sind nicht geschwollen diese vierzig Jahr“, damit er durch diese Worte lehre, daß denen nichts mangeln werde, die durch das Wort leben und glauben, sondern daß GOtt sich unser annimmt in allen Dingen, und durch alle Dinge, wie Petrus spricht [1. Ep. 5, 7.]: „Alle eure Sorge werfet auf ihn, denn er sorget für euch“, und Ps. 34, 11.: „Die den HErrn suchen, haben keinen Mangel an irgend einem Gut.“

Wenn man diese und ähnliche Verheißungen GOttes, die gar herrlich und gewiß sind, recht versteht, so heißt das die Verheißung des ersten
\quellseite{732}
Gebotes recht verstehen, in welchem er sagt: „Ich bin der HErr, dein GOtt.“ Dein, dein, der ich mich dir als GOtt zeigen und erweisen will, und ich will dich auch nicht verlassen, wenn du es nur glauben willst. Denn alle solche Verheißungen hängen ab von dem ersten Gebote und fließen aus demselben. Dagegen, wenn man ihnen nicht glaubt, so heißt das in Wahrheit, das erste Gebot nicht verstehen, sondern andere Götter haben.

Es wendet aber Moses diese Lehre des Glaubens an auf den künftigen Gebrauch inmitten des Ueberflusses an Gütern, wie ich im Anfang des Capitels gesagt habe, nämlich, daß sie dann daran gedenken sollten, wie sie einst mitten im Mangel durch das Wort mit Manna ernährt worden seien, und dies Exempel und Lehre überdenken, und lernen, daß sie nicht um desswillen mehr auf GOtt vertrauen, weil sie Ueberfluß haben und satt seien, indem ihr Bauch wohl versehen ist. Ein solches Vertrauen ist nicht ein Vertrauen auf GOtt, sondern vielmehr auf die empfangenen Güter und Gaben. Daher kommt es, daß sie sowohl des rechten Vertrauens als auch des Wortes GOttes vergessen, und niemals lernen, im Mangel Gotte zu vertrauen. Und (wie ich gesagt habe), es ist eine große Sache und das Werk eines reichen (multi) Geistes, wenn alles wohl steht, GOttes nicht zu vergessen, und sich so zu verhalten, wie Paulus sagt [1 Cor. 7, 30. 31.], als heißest du nichts, der Welt zu brauchen, als ob du derselben nicht brauchtest, daß man wisse, übrig zu haben und Mangel zu leiden, wisse, niedrig zu sein und hoch zu sein [Phil. 4, 12.], und mit dem Propheten [Psalm 62, 11.] das Herz nicht an Reichthum zu hängen, wenn er zufällt, und nicht eitel zu werden, sondern allein an Gotte zu hangen.

Denn das ist es, womit Moses in diesem Capitel zu schaffen hat, indem er gegenüberstellt und einführt mitten in den Ueberfluß das Exempel des Manna, welches mitten im Mangel gegeben wurde, damit er sie von dem Bauche zu dem Worte zurückbringe. Daher wiederholt er auch am Ende, nachdem er die künftigen Reichthümer abgemalt hatte, und spricht [V. 15–17.]: „Er ließ dir Wasser aus dem harten Felsen gehen, und speisete dich mit Manna in der Wüste, von welchem deine Väter nichts gewußt haben, auf daß er dich demüthigte und versuchte, daß er dir hernach wohl thäte. Du möchtest sonst sagen in deinem Herzen“ etc. Ist es nicht ein stärkeres Exempel, daß er aus dem harten Felsen Trank, und Manna in der Wüste gegeben hat, als daß Christus Matth. 6, 28, lehrt von den Lilien des Feldes und den Vögeln unter dem Himmel, daß GOtt sie versorgt? Denn was könntest du weniger hoffen, als daß Wasser und Trank aus einem Felsen kommen sollte? Was weniger in einer Wüste, als Brod und Nahrung? Warum hat er das Wasser nicht aus einem grünen Baume oder einer Quelle des Himmels gegeben? Warum nicht die Nahrung aus Zweigen, oder Wurzeln, oder Kräutern? Warum hat er dies gethan?

Freilich, um die über alle Maßen große Sorge GOttes, die er für uns trägt, zu preisen, der uns auch dann ein GOtt sei und alles reichlich mittheile, wenn alles ganz verzweifelt steht, da er ja auch einen Felsen in deinen Trank verwandeln will und kann, die Wüste in deine Speise, die Blöße in schöne Kleidung, die Armuth in Reichthum, den Tod in Leben, die Schande in Ehre, das Böse in Gutes, die Feinde in Freunde. Und was sollte er nicht können? Er kann alles aus allem, alles in allem, alles aus nichts, alles in nichts geben und verwandeln. Nur das lasse dir angelegen sein, daß du glaubest, und wenn Mangel da ist, nicht kleinmüthig und sorgenvoll seiest, wiederum, wenn Ueberfluß da ist, nicht aufgeblasen seiest und sicher GOttes vergessest, sondern immer mit demselben Herzen deinem Gotte anhängest, gleichviel ob alles da ist oder alles fehlt, indem du weißt, daß der, welcher Ueberfluß gibt, auch geben wird, wenn Mangel drückt. Siehe also, wie geeignet und passend Moses sich dieses Wunders GOttes bedient, um den Sinn des ersten Gebotes darzulegen.

\begin{schriftzitat}
V. 17. 18. Du möchtest sonst sagen in deinem Herzen: Meine Kräfte, und meiner Hände Stärke etc.
\end{schriftzitat}

Dies fügt er hinzu, um zu beschreiben, was das Vergessen GOttes sei, und wie es entstehe. Denn so steht das gottlose Herz des Menschen, wenn ihm Reichthum zufällt: „Ich habe dies durch meine Arbeit zuwege gebracht, und hat nicht Acht darauf, daß es lauter Segnungen GOttes sind, bisweilen durch unsere Arbeit, bisweilen ohne unsere Arbeit, die aber nie von (ex) unseren Arbeiten herkommen, sondern uns
\quellseite{733}
immer aus seiner Barmherzigkeit umsonst geschenkt sind. Denn, wie wir oben gelehrt haben, er gebraucht unserer Arbeit gleichwie einer Art Larve, unter welcher er uns segnet und das schenkt, was sein ist, so daß Raum für den Glauben da sei, damit wir nicht meinen, es sei durch unsere Kräfte oder Arbeit zuwege gebracht, was wir haben, oder, wie er hier sagt, damit wir nicht dafür halten, daß wir durch unsere Kräfte und die Stärke unserer Hände dies Vermögen besitzen, sondern des HErrn, unseres GOttes, gedenken, daß er es sei, der die Kräfte gibt, dies Vermögen zu erwerben, nicht durch unsere Verdienste, sondern weil er es zuvor verheißen hatte.

Dies ist daher die besonders ausgezeichnete (insignis) Stelle, aus der viele treffliche Sprüche geflossen sind, als das Wort in den Sprüchen Salomonis [Cap. 10, 22.]: „Der Segen des HErrn macht reich“, desgleichen im Psalm [Ps. 136, 25.]: „Der allein Fleische speise gibt“, und der ganze Psalm Salomonis [Ps. 127.] „Wo der HErr nicht das Haus bauet“, wo ganz und gar in Abrede genommen wird, daß irgend etwas durch unsere Kräfte zuwege gebracht oder erhalten werde, wie auch hier Moses sagt, daß durch unsre Kräfte oder die Stärke unserer Hände der Reichthum nicht erworben werde, sondern durch GOtt, der da segnet und die Kräfte gibt, daß wir ein „Vermögen ausrichten“ [5 Mos. 8, 17.], das heißt, Reichthum und Hab und Gut (substantiam), denn das bedeutet das hebräische Wort \hebraeisch{חיל}. Der Glaube, welcher sich nicht überhebt im Wohlergehen und kleinmüthig wird in Widerwärtigkeit, ist also der kurze Inbegriff,\footnote{Wir haben mit der Wittenberger nach summa rei ein Komma angenommen.} die Auslegung und der wahre Verstand und die Erfüllung des ersten Gebotes.

\begin{schriftzitat}
V. 19. 20. Wirst du aber des HErrn, deines GOttes, vergessen, und andern Göttern nachfolgen etc.
\end{schriftzitat}

Gleichwie den Gläubigen die Verheißung der Gnade, so wird den Ungläubigen die Drohung des Zornes vorgelegt, damit er von beiden Seiten zum Glauben ziehe und dränge. Der Unglaube macht aus GOtt und dem Vater einen Richter und Feind, wie er hier sagt: „Ich bezeuge euch, daß ihr umkommen werdet.“ Der Glaube macht aus dem Feinde und Richter einen GOtt und einen Vater, wie er kurz zuvor [Cap. 1, 31.] gesagt hat: „Wie ein Mann seinen Sohn trägt, so hat dich der HErr getragen.“ Und wiederum sollst du hier sonderlich merken, daß durchaus kein Unterschied gemacht werde zwischen dem Volke Israel und den Heiden, die ausgerottet werden sollten, es sei denn durch den Glauben, wie er hier sagt, daß sie ebensowohl umkommen sollen als jene Heiden, welche der HErr vor ihrem Angesichte umbringen wird. Sie haben also nichts, dessen sie sich wider die Heiden rühmen möchten, weder vom Gebet, noch von der Gerechtigkeit der Werke, noch von dem Geblüt der Väter, noch von den Wundern GOttes, noch von den göttlichen Reden, noch vom Priestertum, noch vom Königreiche, noch von irgend etwas Anderem. Es steht das Urtheil fest: Wenn sie GOttes vergessen und anderen Göttern dienen sollten, so sollen sie umkommen, als ob dies alles nichts wäre, und sie auch Heiden wären. Denn wer nicht glaubt, der wird verdammt. Es ist aber genugsam gesagt worden, daß „andere Götter ehren“ heißt ungläubig sein und wider den wahren GOtt handeln.

Die heimlichen Deutungen oder Allegorieen dieses Capitels sind leicht zu verstehen. Manna, welches auf hebräisch etwas Bereitetes oder ein Geschenk heißt, weil es ohne unsere Hände bereitet und geschenkt wird, bedeutet das Evangelium GOttes, wie es auch Christus Joh. 6, 32. auslegt: „Moses hat euch nicht Brod vom Himmel gegeben, sondern mein Vater gibt euch das rechte Brod vom Himmel.“ Denn dieses himmlische Wort ist nicht durch unsere Kräfte oder Verdienste der Welt gegeben worden. Mangel und Trübsal in der Wüste vierzig Jahre lang bedeutet das Elend eines Gewissens, welches sich abmüht mit der Sünde, wo da Hunger und Durst, sodann auch feurige Schlangen und Scorpione sind, nämlich das Beißen und der Stachel der Sünde, welches uns grausam peinigt. Es kommt das Manna und speist uns, es kommt der Bach, der aus dem Felsen geschlagen ist, das ist der Geist, der uns durch Christum, den Gekreuzigten, erlangt ist, und uns erquickt, so daß wir niemals, ewiglich nicht dürsten werden, Joh. 4, 14. Ferner, so ungereimt es war, daß aus einem Steine, den man schlägt, Wasser fließen sollte, so ist es wider alles Bedenken (inopinatum), daß aus dem gestorbenen Christo der Geist des Lebens komme, so daß auch Moses
\quellseite{734}
hier zu Falle kam. Denn wer sollte den Geist, der da lebendig macht, erhoffen aus einem sterbenden und begrabenen Fleische? Ueberfluß und Glückseligkeit in dem gelobten Lande sind die Reichthümer an Gaben des Geistes, die Schätze der Weisheit und der Erkenntniß, die uns in Christo geschenkt sind, da wir sicher wohnen und reich an allem himmlischen Segen, 1 Cor. 1, 4. 5. Hier ist es vonnöthen, daß wir nicht aufgeblasen werden, wenn wir gesättigt sind, sondern in Dankbarkeit den HErrn preisen, der uns dies alles nicht allein ohne unser Verdienst gegeben hat, sondern auch, da wir unter dem Fluche und verloren waren unter der Noth der Sünde. Dies sei für diese Gelegenheit genug.

\section{Das neunte Capitel}

In diesem Capitel schneidet Moses eine andere Gelegenheit ab, das erste Gebot zu übertreten, die man den geistlichen Hochmuth nennt, welcher sich der Gerechtigkeit und der Verdienste rühmt. Dies ist das Vertrauen auf eigene Werke, welches das schädlichste Verderben und der größte Widerfacher des Glaubens oder des Vertrauens auf die Barmherzigkeit GOttes ist. Darum zerstört Moses dasselbe mit vielen Worten in diesem ganzen Capitel. Denn dies ist ein Raub an GOttes Ehre, da neben einander zugleich nicht bestehen kann, daß wir uns GOttes rühmen wegen der Barmherzigkeit, und Ruhm suchen an uns selbst\footnote{Wir haben die Lesart der Wittenberger angenommen: gloriari in Deo de misericordia, et gloriari in nobis ipsis. In der Jenaer und in der Erlanger sind die Worte et gloriari fälschlich vor de misericordia gestellt.} wegen unserer Gerechtigkeit und unserer Werke. Er widerlegt aber diese Gerechtigkeit mit drei starken Gründen.

Erstlich, daß jene Heiden verdient hätten, daß sie ausgetrieben würden wegen ihrer Gottlosigkeit. Dies sagt er [V. 5.]: „Denn du kommst nicht herein, ihr Land einzunehmen, um deiner Gerechtigkeit und deines aufrichtigen Herzens willen, sondern der HErr, dein GOtt, vertreibt diese Heiden um ihres gottlosen Wesens willen.“ Durch dieses Wort sollen auch wir uns unterweisen lassen, wenn wir sehen, daß andere geschlagen werden, entweder durch uns oder durch andere. Denn es folgt nicht: Du oder andere werden jenen schlagen, also bist du und andere gerecht vor jenen. Sonst würde der Thurm in Siloah, der viele in Jerusalem erschlug, Luc. 13, 4., auch für gerecht erklärt werden müssen. Sondern GOtt ist es, der die Gottlosen schlägt; mag er dies nun thun durch einen Thurm, durch Feuer, Wasser, wilde Thiere, durch deine oder eines anderen Menschen Hand, daran liegt nichts; es geht da nichts Anderes vor, als daß der gerechte GOtt den Gottlosen schlägt.

Uebrigens sollst auch du dies Exempel fürchten, denn dir zum Schrecken wird der Gottlose geschlagen, wie es Christus beibehält, Luc. 13, 5., auslegt, indem er sagt: „Ihr alle werdet auch also umkommen“, und Paulus sagt Röm. 11, 21. 22., indem er mit demselben Grunde verbietet, daß sich die Heiden nicht erheben sollen wegen des Falles der Juden: „Siehe zu, daß er vielleicht deiner auch nicht verschone“, sonst wirst du auch abgehauen werden.“ Und Röm. 2, 3.: „Denkst du aber, o Mensch, der du richtest die, so solches thun, und thust auch dasselbe, daß du dem Urtheil GOttes entrinnen werdest?“ Es ist daher ein trefflicher Ausspruch des heiligen Gregorius, der da sagt: Wenn wir irgend jemand sündigen sehen, so sollen wir vor allen Dingen uns selbst in ihrem Unglück hinein, weil wir in gleicher Weise entweder auch schon gefallen sind, oder fallen können, was jemand in diesen Vers zusammengefaßt hat:

\begin{quote}
Aut sumus, aut fuimus, aut possumus esse, quod hic est [entweder wir sind, oder waren, oder können sein, was dieser ist].
\end{quote}

Und in den Lebensbeschreibungen der Väter sagt er, als er hörte, daß ein Bruder gefallen war: Gestern jener, heute ich. Darum, wenn ein anderer geschlagen wird, so muß man daran denken:

\begin{quote}
Tunc tua res agitur, paries cum proximus ardet [Das geht dich an, wenn die nächste Wand brennt].
\end{quote}

\quellseite{735}
Ja, mit einem Herzen voll Barmherzigkeit und Mitleid sollen die Gottlosen von uns verderbt oder geschlagen werden, indem wir bedenken, daß wir das Werkzeug GOttes sind, die vielleicht auch selbst wie eine Ruthe ins Feuer geworfen werden sollen, nachdem die Uebelthäter durch uns gestraft worden sind.

Zweitens [widerlegt Moses die eigene Gerechtigkeit]\footnote{Von uns hinzugefügt.} durch das Ansehen der göttlichen Verheißung. [V. 5.:] „Der HErr hat sie vertrieben, um seinen Eid und sein Wort zu erfüllen, welches er deinen Vätern zugesagt hat.“ Es konnte nichts Stärkeres gesagt werden wider das Vertrauen auf die eigene Gerechtigkeit. Denn wo waren die Kinder Israel, als GOtt ihrem Vater Abraham das Land Canaan verhieß, da er noch unfruchtbar war und keine Hoffnung auf Nachkommenschaft hatte? Wenn sie durch die Verheißung GOttes das Land empfangen und besitzen, so haben sie es also nicht durch ihre eigenen Verdienste oder ihre eigene Gerechtigkeit, sondern aus lauter Gnade und Güte GOttes, die ausgeschüttet ist über Unwürdige, die selbst noch nicht geboren waren. Denn warum gibt er die Verheißung? Thut er es etwa darum, daß es in solcher Weise diejenigen verdienen sollten, die erst nach 430 Jahren kommen sollten? Das sei ferne, sondern der es aus lauter Güte verheißen hat, der hat auch aus bloßer Barmherzigkeit sein Wort erfüllt. Dieses Grundes bedient sich auch Paulus in den Briefen an die Galater [Cap. 2, 16. ff. 3, 2. ff.] und an die Römer [Cap. 4, 1. ff.], da er aufs gewaltigste erweist, daß die Gerechtigkeit nicht herkomme aus den Verdiensten, sondern aus der Barmherzigkeit GOttes, die er einst verheißen hat.

Drittens [widerlegt er die eigene Gerechtigkeit] aus ihrer eigenen Erfahrung, da er sagt [V. 6.]: „Sintemal du ein halsstarrig Volk bist“ etc. Durch diese Halsstarrigkeit haben sie das Land so gar nicht verdient, daß es nahe daran war, weil GOtt zornig war, daß er sie in der Wüste eher vertilgte als die Heiden, wie der Psalm [Ps. 106, 26.] sagt: „Und er hob auf seine Hand wider sie, daß er sie niederschlüge in der Wüste“, und seine Verheißung anderstwohin gemandt hätte, nämlich auf die Nachkommen Moses, 4 Mos. 14, 12. Ja, er schlug alle bis auf den letzten Mann nieder in der Wüste, ausgenommen Josua und Caleb, so daß auch weder Moses noch Aaron in das Land kamen. So viel fehlt also daran, daß sie durch ihr Verdienst in das Land eingehen sollten, daß ihnen auch weitaus das Gegentheil hätte begegnen müssen, wie es denn auch wirklich geschah, wenn mit ihnen nach ihren Verdiensten hätte gehandelt werden sollen. Dieses Grundes bedient sich auch Paulus Röm. 4, 14. f., wo er sagt: „Denn wo die vom Gesetz Erben sind, so ist der Glaube nichts und die Verheißung ist ab, sintemal das Gesetz richtet nur Zorn an“, damit du sehest, daß Moses und Paulus sich derselben Schlußkunft (dialectica) des Geistes bedienen wider die Gerechtigkeit aus den Werken, für die Gnade und Barmherzigkeit GOttes.

Deshalb ist dies alles auch uns zur Lehre geschrieben, damit wir lernen, daß das Unglück, welches über die Gottlosen hereinbricht, zwar verdienter Weise über sie komme, übrigens, wenn es über uns nicht hereinbreche, sondern wir guter Tage genießen, wir solches nicht unserer Gerechtigkeit, sondern der Güte GOttes zu danken haben, durch welche es uns von Ewigkeit zuvor verziehen und verheißen ist, da auch wir weitaus das Gegentheil verdienen. Denn es steht das Urtheil fest, daß uns nicht wegen unserer Gerechtigkeit irgend etwas Gutes gegeben werde, sondern damit GOtt sein Wort erfülle, daß er von Ewigkeit gewollt hat, daß wir nicht aufgeblasen sein sollen, und einen Götzen aus unserer Gerechtigkeit machen, sondern wissen sollen, daß wir einen GOtt haben, von dem wir alles umsonst empfangen aus lauter Güte, die über uns Unwürdige ausgeschüttet ist, wie auch der Patriarch Jakob bekennt, indem er spricht, 1 Mos. 32, 10.: „Ich bin zu geringe aller Barmherzigkeit, die du an deinem Knechte gethan hast“, das heißt, nicht allein Einer Barmherzigkeit, die auch groß sein möchte, sondern irgend welcher, wie klein und gering sie auch immer sein mag.

Ferner wendet Moses ein heilsames Gegenmittel gegen diese schädliche Seuche an (denn so gewaltig ist dies Ungeheuer der eigenen Gerechtigkeit, daß man nicht genug dagegen thun kann). Denn er legt ihnen alle vergangenen Uebelthaten des Volks vor Augen, und befiehlt ihnen, daß sie derselben gedenken sollen, und die Sünde, daß sie am Berge Sinai das Kalb angebetet haben, rückt er ihnen vor andern mit gewaltigen Worten
\quellseite{736}
auf. Denn was kann diese Krankheit des Stolzes gewinder und besser (praesentius) heilen, als daß man an die vergangenen Gottlosigkeiten und Schandthaten gedenkt? Denn was hat das Volk anders als solche Dinge, wegen deren es sich schämen muß, die Augen gen Himmel zu erheben? wie Paulus zu seinen Römern sagt [Cap. 6, 21.]: „Welcher ihr euch jetzt schämet.“ Denn darum hat GOtt auch den David, Moses, Petrus und Paulus und andere große Männer fallen lassen, damit sie gedemüthigt würden und sich vor GOtt schämten, und allein auf seine Güte sich gründeten. Daher straft Petrus, 2 Petr. 1, 9., diejenigen hart, welche die Vergebung der vorigen Sünden in Vergessenheit kommen lassen, und sicher und kalt werden, darnach auch halsstarrig und hoffährtig.

Endlich beschließt er [V. 24.]: „Ihr seid dem HErrn immer ungehorsam gewesen, so lange ich euch gekannt habe.“ O welch ein feines Lob ist das, welche Verdienste, welche Gerechtigkeit des heiligen Volkes, nämlich, daß sie der Stimme GOttes ungehorsam gewesen sind! Jetzt gehe hin und rühme dich, prahle, blase dich auf damit, daß dir um deiner Gerechtigkeit willen das Land gegeben sei, da du doch nichts Anderes zu hören verdienst, als daß du ein halsstarrig Volk seiest und immer ungehorsam gegen den HErrn: freilich, diesen Verdiensten gebührt jener königliche Reichthum, und nicht vielmehr tausendfältiger Tod und Kreuz! Was kann Erschrecklicheres gesagt werden, als daß man dem HErrn ungehorsam sei? Und doch wird diesem Ungehorsam das gepriesene (laudata) und gesegnete Land gegeben. Was hat hier Israel übrig, darauf es hoffährtig sein könnte, und nicht vielmehr seinen Mund in den Staub lege und spreche: Ich bin ärger als alle Heiden, und groß und unverdient ist deine Gnade, daß ich dieses Land empfange. Ferner, wenn dieses heilige Israel und das auserwählte (peculiaris) Volk so beschaffen ist vor GOtte, was sind denn wir Heiden und Sünder? Siehe, mit wie gewaltigen Worten hier Moses das Volk anklagt selbst in dem Gebete, in welchem er für sie bittet, da er spricht [V. 27. 28.]: „Siehe nicht an die Härtigkeit, und das gottlose Wesen und Sünde dieses Volks, damit nicht etwa die Einwohner des Landes sagen“ etc. Hier wird von dem Volke nichts erzählt, als solche Dinge, wegen deren sie den Tod verdient haben. Es war nur das Eine da, was ihnen zur Rettung diente, daß dem Namen des HErrn Gefahr drohte, wenn sie ausgerottet wurden. Deshalb, damit der Name GOttes nicht verlästert werde, wird ihnen Gnade (venia) gegeben, die ihnen sonst nicht geschenkt sein würde, sondern sie hätten gänzlich ausgerottet werden müssen, wenn der Name GOttes, nach dem sie genannt waren [Jer. 15, 16.], hätte unverletzt bleiben können. Daher nahmen auch viele andere, besonders David im Psalter, diesen gar sicheren und zuverlässigen Vorwand, daß sie sagten [Ps. 25, 11.]: „Um deines Namens willen, HErr, sei gnädig meiner Missethat“ etc., und Josua [Cap. 7, 9.]: „Was willst du denn bei deinem großen Namen thun“, nach dem wir genannt sind?

\section{Das zehnte Capitel}

Was Moses hier über die zweiten Tafeln sagt, gehört zu dem vorhergehenden Capitel, um die Sünde, daß sie das Kalb angebetet hatten, groß zu machen. Denn die Capitel sind nicht recht abgetheilt; wie häufig an anderen Stellen, so auch hier. Was aber gesagt wird von den Reisen und den Lagerstätten Israels, vom Tode Aarons und der Einsetzung seines Sohnes Eleasar an seine Statt, von der Verordnung, daß die Leviten die Lade tragen und im Dienste stehen und den Namen des HErrn loben sollten, wird erwähnt, entweder weil dies geschehen ist, als dies vorging, oder weil Moses an diesen Lagerstätten das ausgerichtet hat mit den Leviten, was auf dem Berge Sinai befohlen war, oder, was ich eher glaube, weil nach dem Tode Aarons, da ein neuer Priester, Eleasar, eingesetzt worden war, auch der Dienst der Leviten von neuem und noch einmal gebessert und eingerichtet ist, wie es in menschlichen Dingen zu geschehen pflegt, die täglich dahinfällen und sich ändern. Danach, nachdem dies eingefügt ist, kehrt Moses
\quellseite{737}
wieder zurück zu der Sünde, die sie mit dem Kalbe begangen haben, und nachdem er Gnade erlangt hatte [für das Volk], empfängt er den Befehl, das Volk in das verheißene Land zu führen.

Die heimliche Deutung des Kalbes und der Tafeln gehört auch zu dem vorhergehenden Capitel. Ueber die Tafeln belehrt uns Paulus genugsam 2 Cor. 3, 3., nämlich, daß die steinernen Tafeln die harten Herzen des Gesetzesvolks sind, während das Volk der Gnade die fleischernen und weichen Tafeln des Herzens hat. Es sind aber die Gesetze auf steinerne Tafeln geschrieben, weil des Gesetzes Werk beschrieben ist in aller Herzen, sintemal ihr Gewissen sie bezeuget, Röm. 2, 15. Aber des Herzens Härtigkeit macht, daß es mit dem Gemüthe nicht ergriffen wird, nicht geliebet, nicht gehalten, auch nicht im Werke erfüllt wird. Wie die steinernen Tafeln nur die geschriebenen Buchstaben bewahren, aber nichts thun, so haben die harten Herzen, ohne den Geist, das Gesetz nur zum Zeugniß wider sich selbst. Besonders das jüdische Volk ist ein solcher Wirth (hospes) des Gesetzes und des ganzen Moses, zum Zeugniß wider sich selbst.

Das Kalb, welches aus Gold verfertigt ist, ist die Werklehre, die aus Worten der Schrift und dem verfälschten Gesetze verfertigt ist, und verdreht durch den Dienst und die Kunst der Priester, die dem Volke vorstehen sollen. Diese [Werklehre] wird angebetet, wenn man darin Ehre sucht und die Zuversicht des Gewissens (die allein dem wahren Gotte gebührt) darauf gesetzt wird, wie die Werkheiligen (justitiarios) nothwendiger Weise thun müssen. Moses aber nimmt es, verbrennt es mit Feuer, zerbricht und zermalmt es, und wirft es in den Bach, der vom Berge fließt, das heißt, der wahre Dienst des Gesetzes lehrt die Erkenntniß der Sünde, und daß die Werke sammt ihren Lehrern nichts seien, sondern daß die Gnade GOttes vonnöthen sei, um zu rechtfertigen. Daher wird diese ganze Lehre zermalmt und hinweggeworfen und verschlungen durch den Bach des Evangelii, welcher aus dem Berge, Christo, fließt, der die Erde erfüllt. Aus demselben gibt Moses dem Volke zu trinken, indem er zeigt, daß das Evangelium nothwendig sei, und zu denselben hindränget durch die Erkenntniß der Sünde.

Dann kommt Moses wieder darauf zurück, daß er zum ersten Gebote vermahnt und dasselbe preist, indem er sagt [V. 12.]: „Nun, Israel, was fordert der HErr, dein GOtt, von dir?“ als ob er sagen wollte: Dies ist die Summa des ersten Gebotes, von dem ich alles dies gesagt habe, und um dessentwillen alles Zuvorgesagte geschehen ist, daß du den HErrn, deinen GOtt, fürchtest, und [in seinen Wegen] wandelst, und ihn liebest, und ihm dienest, daß du die Gebote haltest, die ich dir heute gebiete. Ich habe oben gesagt, daß bei den Hebräern „die Furcht GOttes“ dasselbe bedeutet als bei uns die Verehrung GOttes. Und siehe hier die Ordnung: fürchten, wandeln, lieben, dienen von ganzem Herzen etc. setzt Moses vor das, was da folgt [V. 13.]: „Daß du haltest die Gebote und Rechte“ etc.,\footnote{In der Jenaer fehlen die Worte: praemittit huic, quod sequitur: Ut serves praecepta et statuta etc. — In der Erlanger ist mit „Praemittit“ (fälschlich) ein neuer Satz begonnen.} damit du wissest, daß von den äußerlichen Geboten nichts gehalten werden könne, daß es Gott gefalle, wenn es nicht hervorgeht aus einem Herzen, das da fürchtet, liebt, gehorsam ist, und dient, so daß die Kraft des ersten Gebotes, welche der Glaube ist, in allen Geboten und Werken regiere und wirke; ohne denselben ist alles nur Schein und Larve.

\begin{schriftzitat}
V. 14. 15. Siehe, der Himmel und aller Himmel Himmel ist des HErrn, deines Gottes etc.
\end{schriftzitat}

Wiederum, nach der Weise eines guten Gesetzgebers schärft er ein und wiederholt, wie groß und wie gut GOtt sei, der da gebietet. Siehe, sagt er, es sind verschiedene Himmel und verschiedene Götter, aber dein GOtt ist Einer über alle (wie Paulus 1 Cor. 8, 5. f. auch fast redet), und dennoch hat er aus allen Völkern, unter allen Himmeln, deine Väter und ihren Samen geliebt und erwählt. Es ist bekannt, daß in der Schrift „die Himmel“ nicht unterschieden werden nach den Kreisen (sphaeras), wie die Mathematiker zu thun pflegen, sondern nach Theilen der Erde, wie auch die Erde [in verschiedene Theile] eingetheilt wird, wie es anderswo heißt [5 Mos. 28, 23.]: „Dein Himmel, der über deinem Haupt ist, wird eher sein“, nämlich, daß sie [die Schrift] einen anderen Himmel anzeigt für dieses Volk, einen anderen Himmel für ein anderes Volk. Aber dieser GOtt ist ein GOtt aller Himmel, das heißt, er herrscht
\quellseite{738}
über alle Himmel im Gesichtskreise (horizonte) aller Menschen, und nicht bloß der Juden GOtt, sondern aller Völker. So heißt es Ps. 115, 16.: „Der Himmel allenthalben ist des HErrn, aber die Erde hat er den Menschenkindern gegeben.“

\begin{schriftzitat}
V. 16. 17. So beschneidet nun eures Herzens Vorhaut etc.
\end{schriftzitat}

Dies ist ein klarer Spruch von der zwiefachen Beschneidung, wo zugleich das Vertrauen und die Vermessenheit auf die fleischliche Beschneidung verdammt wird. Sonst, wenn sie genugsam wäre, was wäre es denn vonnöthen, die Beschneidung des Herzens zu fordern? So sagt auch Jeremias [Cap. 9, 26.], daß die Heiden zwar unbeschnitten seien am Fleisch, aber Israel sei unbeschnitten am Herzen. Mit diesem Worte wirft Moses jenen ganzen Ruhm dieses Volkes zu Boden, mit welchem er in dem vorhergehenden Worte, wie es schien, sie aufgeblasen hatte, da er sagte, sie seien vor allen Völkern, unter allen Himmeln, auserwählt, damit sie nicht hoffärtig wären gegen die übrigen Völker. Ja, es folgt bald, daß der HErr der HErr sei über alle Herren, und groß, der weder die Person ansieht noch Geschenke annimmt [V. 17.]. Was [achtet er] denn? Nichts als den Gehorsam des Herzens oder die Beschneidung des Geistes, dessen Lob aus GOtt ist, Röm. 2, 29. Da er nun die Person nicht ansieht, die Beschneidung des Fleisches aber und die Zeugung nach dem Fleische (sanguinis) aus den Vätern, und alles, außer der Beschneidung des Herzens, persönliche Dinge sind\footnote{Statt personaliae in der Erlanger ist personalia zu lesen.} und Ansehen der Person (personas) machen, so folgt, daß sie dadurch Gott nicht gefallen können, sondern allein durch die Beschneidung des Herzens. Wenn die Heiden diese haben, wie das zweite Capitel des Briefes an die Römer [V. 26. 27.] darlegt, so erfüllen sie das Gesetz, und richten die Juden, die nach dem Fleische beschnitten sind.

Es zerstört also Moses durch diesen Ausspruch die Gelegenheit zur Sünde wider das erste Gebot, welche entsteht aus dem Halten der Beschneidung und allen Werken des Gesetzes. Ein gar wunderlicher Gesetzgeber, welcher zugleich die Werke gebietet und sie dennoch verdammt, wenn sie gethan sind. Aber dies thut er, damit sie wissen sollen, daß Gesetz werde nicht erfüllt durch die Werke des Gesetzes, sondern durch die Beschneidung des Herzens, und erst diese Werke seien wahrhaft gute Werke, welche nicht aus dem Zwang des Gesetzes, sondern aus einem Herzen hervorgehen, welches zuvor beschnitten ist. Denn das fördert der, welcher die Personen nicht ansieht; sonst machen die Werke nichts als Schein (personas). Und dieses Lehrstück Moses ist dem ähnlich, welches Paulus im Briefe an die Römer behandelt.

\begin{schriftzitat}
V. 18. Und hat die Fremdlinge lieb, daß er ihnen Speise und Kleider gebe.
\end{schriftzitat}

Diese sehr lieblichen Verheißungen GOttes sind wohl zu merken: „Der keine Person achtet“, desgleichen: „Er nimmt kein Geschenk“, desgleichen: „Er schafft den Witwen und Waisen Recht“, desgleichen: „Er hat die Fremdlinge lieb, daß er ihnen Speise und Kleider gebe.“ Denn mit diesen Worten tröstet Gott klarlich alle Elenden, Verlassenen und Betrübten in der ganzen Welt, und läßt ihr Herz zu sich, daß sie sich ihm anvertrauen und sich auf seine Güte verlassen, wie es sich gebührt gegen den wahren, Einen und alleinigen GOtt. Denn wenn er seiner Person achtet, so habe ich schon, daß ich mich rühmen kann und mit aufgeblasenem Herzen stolz sein kann wider alle Könige, Fürsten, Reichen, Starken der Welt, gewiß, daß ich, der ich von ihnen als ein Armer, Unedler, Schwacher verachtet und verlassen werde, doch von ihm, welcher der GOtt der Götter und der HErr aller Herren ist (wie Moses hier auch mit Worten prangt und hoffärtig ist), mächtig und schrecklich über alle meine Verächter, angesehen, aufgenommen und behütet werde. Was sollte mir also mangeln, wenn mich irgend ein Herr über wenige Knechte verachtet, während mich der HErr über alle Herren und der GOtt aller Götter seiner Sorge würdigt?

So auch, wenn ich gleich ein Narr sein muß, ein Sünder, ein verzweifelter und verlorener Mensch, dazu, wenn ich gleich verurtheilt, verdammt und in den Bann gethan bin als ein Sünder, ein Ketzer, ein Verfluchter, ein Gotteslästerer, ein Knecht des Teufels: was ist es denn? Jene sind heilig, gerecht, geehrt, gepriesen, unter die Kinder GOttes gerechnet und unter die Sterne gezählt: was macht das aus? Hier ist der GOtt aller Götter, welcher diesen
\quellseite{739}
Schein (personas) nicht ansieht, sich auch nicht um ihre Gaben (dona) kümmert; er sorgt für mich Sünder und Narren ebensowohl, als jene sich vermessen, daß er für sie sorge. Lieber, wenn man diese Verheißungen im Glauben festhielte, was für ein Ereigniß (res) könnte dann wohl eintreten, das mich betrüben könnte? Welche Sünde könnte mich drücken, welche Sünde könnte mich in Verzweiflung bringen? Ja, welche Höhe, welche Tiefe, welches Gegenwärtige, welche Creatur könnte uns abgeschieden machen oder demüthigen? Laß nur dein Herz an diesen Worten hangen, und dir wird nichts mangeln, wenn auch alles dich verlassen sollte; aber nichts wird da sein, wenn dir auch alles günstig wäre (arriserint). Nur glaube, daß der nicht lüge, der dies redet, wie er ja nicht lügen kann, und du wirst in allen Dingen und durch alles ein sicheres und fröhliches Herz haben.

Siehe, das heißt das erste Gebot erklären, hier ist die Auslegung von Moses selbst. So lehrt er verstehen, was da heiße, daß ein GOtt sei, was es sei, einen GOtt haben, was da sei, das erste Gebot halten. Lieber, was für große Quellen sind sogar nur diese Worte allein für die Propheten gewesen, welche hieraus alles genommen haben, was sie von der Sorge GOttes für die Schwachen, Geringen, Armen, Sünder, Wittwen, Waisen, Verurtheilten, Verdammten, Betrübten und Elenden schreien; wiederum [alles], was sie wider die Reichen, Tyrannen, Mächtigen und Richter, wider die Gewalttäthigen, Harten und Hoffärtigen donnern vom Zorn und von der Rache GOttes. Denn alles dies fließt aus dem großen Ocean des ersten Gebotes und fließt wieder in denselben zurück, so daß nie ein Wort gehört worden ist, welches fruchtbarer und reicher wäre an Trost, noch je gehört werden wird, wiederum auch kein härteres und strengeres als dies Wort des ersten Gebotes: „Ich bin der HErr, dein GOtt.“

Und, um nun weiter zu gehen zu dem Uebrigen: Die Wittwen und Waisen leiden vieles, wer sorgt denn für sie? Niemand, niemand. Aber sie sollen hieher sehen mit gewissem Herzen, da gesagt wird zu ihrer Freude: Der GOtt aller Götter, der HErr über alle Herren, der große GOtt, mächtig und schrecklich, der schaffet den Wittwen und Waisen Recht. Daher ist auch das geflossen Ps. 68, 6.: „Der ein Vater ist der Waisen und ein Richter der Wittwen“, der diesen Titel der Ehre so gerühmt wissen will, damit man sehe, daß er aus lauter Barmherzigkeit die Tyrannen ermahne, daß sie nicht etwas begehen, wodurch er gezwungen werde, diesem Titel genugzuthun, sondern sie treulich schreckt, damit sie sich fürchten, den Wittwen und Waisen Schaden zuzufügen, ja vielmehr nach seinem Exempel ihnen wohlthun und sie lieben. Und das ist noch viel herrlicher, daß er, mögen nun die Wittwen und Waisen glauben oder nicht, ihnen nichtsdestoweniger Recht schafft, und die Tyrannen richtet, wiewohl der Unglaube der Wittwen und Waisen dies nicht verdient in diesem Leben, oder es zu erfahren, so lange der Unglaube dauert. Denn diese Verheißung geht fort und wird erfüllt wider die Gewalttäthigen und die Harten, wenngleich niemand hier glauben sollte; sie würde aber noch viel gewaltiger und schneller erfüllt werden, wenn auch die Wittwen und Waisen glaubten.

Er liebt die Fremdlinge so, daß er sie speist und kleidet. Ich bitte dich, wer wollte sich nun verlassen auf sein Haus oder sein väterliches Erbe? wiewohl dies auch nicht zu verachten ist. Aber wenn es sich etwa begäbe, daß man ein Fremdling sein müßte, so soll niemand verzweifeln oder weniger glauben, als wenn er des väterlichen Hauses genießen könnte, sondern soll mit dem Psalmisten singen [Ps. 27, 10.]: „Mein Vater und meine Mutter haben mich verlassen, aber der HErr hat mich aufgenommen.“ Das Buch der Weisheit [Cap. 10, 10.] sagt von dem Patriarchen Jakob: „Die Weisheit leitete den Gerechten, so vor seines Bruders Bruder Zorn flüchtig sein mußte, straß Weges, und zeigte ihm das Reich GOttes und gab ihm zu erkennen, was heilig ist, und half ihm in seiner Arbeit, daß er wohl zunahm und viel Gutes an seiner Arbeit gewann“ etc. Wer zweifelt, daß daher es geschehe, daß in der Fremde mehr Leute reich, hoch- und berühmt werden, als derer, die zu Hause bei den Eltern bleiben und da aufgebracht werden? Denn GOtt ist ein Vater aller Menschen, der den Fremdling lieb hat, daß er ihm Speise und Kleider gebe.

\begin{schriftzitat}
V. 21. Er ist dein Ruhm und dein GOtt.
\end{schriftzitat}

Das ist eine hebräische Redeweise. „Ruhm“ (laus) scheint mir an dieser Stelle eigentlich das zu sein, was Paulus im Griechischen καύχησιν oder καύχημα nennt, gloriationem, auf deutsch:

\quellseite{740}
„Ruhm“, so daß der Sinn ist: Wegen keiner Sache, weder wegen deiner Person noch wegen deiner Werke sollst du dich aufblasen oder rühmen, sondern rühme dich deines GOttes, in ihm blaße dich auf, in ihm sei hoffährtig, da du weißt, daß er es sei, der die Person nicht ansieht etc., der bei dir solche großen und schrecklichen Dinge gethan hat, die deine Augen gesehen haben, als wollte er sagen: Durch die Erfahrung hast du gelernt, daß er dir, da du ein Fremdling warst, Speise und Kleidung gegeben hat, daß er die Person nicht ansieht, daß er den Waisen und Wittwen Recht schafft etc. Du hast also, dadurch du aufgeblasen und hoffährtig sein kannst wider alle Tyrannen, ja, wider alles Uebel, nämlich deinen GOtt, der dich mit allem Guten versorget. Du siehst also, daß das erste Erfüllung des ersten Gebotes sei, wenn GOtt der Ruhm, der Preis, das Prangen, die Hoffahrt unseres Herzens ist, beides zur Zeit des Wohlergehens und der Widerwärtigkeit. Denn das heißt, GOtt in Wahrheit haben.

Die heimliche Deutung von den zerbrochenen Tafeln und von den unverfertigten ist dies, daß das Amt des Gesetzes in zwiefachem Brauche ist: Erstlich, wenn nur Werke gethan werden, ohne daß das Herz dabei ist, das heißt in Wahrheit ein Kalb aufrichten und anbeten, das ist, aus dem Gesetze eine fleischliche Werklehre machen, und in derselben sich rühmen, was dasselbe ist als das Gesetz übertreten. Dies zeigt Moses an, da er die Tafeln zerbricht, nachdem er das Kalb gesehen hatte, das heißt, er thut kund, daß durch die Lehre und den Wahn von den Werken das Gesetz GOttes übertreten worden sei. Daher kommt auf solche Weise das Gesetz nicht zu dem Volke, weil es von dem fleischlichen Volke nicht verstanden wird. Zweitens, wenn es zu dem Volke gelangt, aber wegen des Glanzes des Angesichts Mosis nicht ertragen wird, und die Anbeter des Kalbes getödtet werden, das heißt, daß das geistliche Gesetz Zorn anrichtet [Röm. 4, 15.], indem es die Sünde anzeigt, Röm. 7, 7., und man kann die Stärke (cornu) seines Glanzes nicht ertragen, weil es etwas ganz Anderes erfordert als das Werk. Darum so wird hier zugleich beschrieben die Reise der Kinder Israel, und der Tod Aarons, und die Reformation der Leviten, nämlich, daß der zweite Brauch des Gesetzes nach dem ersten folgt, und daß die Priester und Meister (magistri) des Kalbes abgethan werden müssen, das ist, die Lehre der Werke, damit das Volk zu einer andern Gesinnung komme, und die Leviten, das ist, die Diener des Wortes, reformirt werden, damit sie im Dienste des HErrn, nicht der Werke, stehen, und den Namen des HErrn loben [V. 8.], das heißt, das Wort der Gnade predigen, so daß allein der HErr ihr Theil sei in reinem Glauben. Dann wird das erste Gebot erfüllt.

\section{Das elfte Capitel}

Dies Capitel ist eine Art Schluß (epilogus) aller vorhergehenden Ermahnungen über das erste Gebot. Denn er wiederholt und schärft die Wohlthaten ein, die sie in Egypten und in der Wüste empfangen haben, und er verheißt, daß sie das Land und auch in Zukunft Wohlthaten empfangen sollen, wenn sie die Gebote GOttes halten und keinen anderen Göttern dienen. Denn dieses erste Gebot hat er bisher gehandelt und getrieben, darum ist er auch wegen desselben so besorgt, daß er befiehlt, es zu Herzen zu nehmen, es als ein Zeichen auf die Hände und vor die Augen zu binden, und daß die Kinder unterrichtet werden sollen, wie er auch oben [Cap. 6, 6. ff.] gethan hat, weil in diesem Gebote der ganze Inbegriff und die Erfüllung aller folgenden enthalten ist. So sehen wir, daß Moses nichts unterlassen hat, was zum Verständniß des ersten Gebotes dienen möchte, wie er alles reichlich behandelt hat, was dem Glauben fördern kann, und alles, was ihn hindert. In dem Folgenden wird er die Gebräuche und die Ceremonien der Werke dieses Gebotes behandeln.

Er mischt aber auch etliche Verheißungen ein, nämlich, wenn sie dem HErrn anhangen, werde er die Völker vertreiben, die stärker seien als sie, desgleichen, alle Oerter, darauf sie ihren Fuß setzen, sollten ihnen gehören, desgleichen, daß er sagt [V. 25.]: „Niemand wird euch widerstehen
\quellseite{741}
mögen“, und: „Eure Furcht und Schrecken wird der HErr, euer GOtt, über das ganze Land kommen lassen, in welches ihr hineingehen werdet“; sodann neben den Verheißungen auch den Fluch. Ja, beides zugleich, Segen und Fluch, befiehlt er auf den Bergen Grisim und Ebal zu verkündigen, davon wir nachher sagen werden, was alles zu der letzten Schlußrede gehört. Er fügt auch noch hinzu, wie das gelobte Land beschaffen sei, als eine Art Verheißung, nämlich, daß es nicht sei wie das Land Egypten, sondern der HErr nehme sich desselben sonderlich an, denn seine Augen sähen darauf von Anfang des Jahres bis ans Ende, nämlich, damit er sie durch den Glauben an GOtt hänge, und sie wissen sollten, daß Regen auf dieses Land fallen werde, wenn sie glauben, durch GOttes Gnade; wenn sie aber nicht glauben, werde es an Regen mangeln.

Es ist bekannt, daß Egypten nicht durch Regen befeuchtet wird, sondern durch die Ueberschwemmung des Nil, die jährlich im Sommer stattfindet. Dies Wunder GOttes ist geringgeachtet worden, weil es so häufig geschah, wie es auch mit allen anderen Dingen geht. Denn dadurch hat Egypten einen wunderbaren Unterschied gegen andere Länder. Aber dieses gelobte Land hat Berge und Thäler, daher wird es nicht durch die Ueberschwemmung eines Flusses, sondern durch befeuchtende Regen vom Himmel befruchtet, nicht daß andere Länder nicht auch auf dieselbe Weise befeuchtet und versorgt würden, und die Augen des HErrn nicht auch auf sie Acht hätten von Anfang des Jahres bis ans Ende, wie in diesem Lande, da er allem fleischliche Speise gibt [Ps. 136, 25.], und, wie es im vorigen Capitel [V. 18.] heißt, „den Fremdlingen Speise und Kleider gibt“, und „die Herzen der Menschenkinder mit Speise und Freude erfüllt“, Apost. 14, 17., sondern daß kein Volk hierüber eine Verheißung GOttes hat. GOtt gibt zwar allen alles, aber diesem sonderlichen Volke gibt er noch das Wort der Verheißung dazu, daß sie nicht allein vom Brod leben sollen, wie andere Völker, sondern auch von dem Worte, und auf der Erde nicht allein Sorge für den Bauch tragen sollen, sondern viel mehr auch für den Geist, damit sie nicht meinen möchten, ihnen werde das Land gegeben, sich darin zu mästen wie die Schweine, und nicht vielmehr glauben, daß sie durch das Wort GOttes ernährt würden und alles durch das Wort GOttes empfingen, das heißt, GOtte dienten. Also nicht um des Landes selbst willen, sondern um des Volkes willen, das darin ist, trägt GOtt für dasselbe Sorge und sieht darauf, damit er sie durch den Glauben in demselben regiere, wie es anderswo heißt: Er hat das Volk nicht um des Ortes willen, sondern den Ort um des Volkes willen erwählt [2 Mos. 20, 24.]. Denn er hat Abraham und seinen Samen nicht um des Landes willen erwählt, welches er ihm verheißen hatte, sondern erwählte das Land um Abrahams willen, um es ihm zu schenken.

Es bedeutet aber das Land das Reich Christi, welches zuvor von gottlosen Teufeln und Lehrern und Werkstreibern besessen, durch das Wort des Evangelii befreit wird von Sünden und Irrthümern. In diesem Lande lebt man so durch den Glauben, daß wir mit getrostem und zuversichtlichem Gewissen empfinden, daß GOtt sich unser annehme und fort und fort die Augen seiner Gnade über uns offen halte. Darnach wird es mit Regen der himmlischen Lehre befruchtet. Denn nicht durch das Werk oder Wort eines Menschen, wie in Egypten, werden die befeuchtenden Gewässer dahingeleitet, sondern es wird allein durch die Rede GOttes befeuchtet vom Himmel herab etc.

Vom Frühregen und Spätregen ist auch anderswo gesagt, daß der Frühregen der erste Regen sei, der Spätregen der nachkommende, denn so gibt es fröhliche Saaten und alle Dinge die Fülle, wenn nach dem Regen heiteres Wetter kommt und Sonnenwärme, darnach aber, wenn die Hitze stark wird, wiederum Regen folgt. Denn beständiger Regen verderbt ebensowohl alles, als beständige Wärme und Hitze. So muß auch die Lehre nicht allein im Volke getrieben werden, sondern nach der Lehre muß auch das Werk des Glaubens geübt werden. Wo aber die Kraft durch Wirken und Leiden erschlafft ist, da muß mit der Lehre das Herz wieder aufgerichtet, gestärkt und getröstet werden. So wächst der Mensch in der Erkenntniß GOttes. Sonst, wenn man nur lehrt und nichts thut, ist es ebenso, als wenn es beständig regnete, so daß alles erstickt würde und verdürbe. Wiederum, wenn man thut und leidet, und nicht lehrt, ist es ebenso, als wenn die Hitze beständig glühte, daß auch so die Kraft des Geistes erschlaffen und ausgelöscht werden würde. Es ist also beides vonnöthen, Frühregen und Spätregen, nämlich lehren und ermahnen, Röm. 12, 7. 8.: „Lehret jemand, so warte er der Lehre; ermahnet jemand, so warte er des Ermahnens.“

\quellseite{742}
\section{Das zwölfte Capitel}

\begin{schriftzitat}
V. 1. ff. Das sind die Gebote und Rechte, die ihr halten sollt etc.
\end{schriftzitat}

Nachdem nun die geistliche Erklärung des ersten Gebots beendet und das Herz unterwiesen worden ist, geht Moses über zu den Werken und äußerlichen Gottesdienst im Werk und Ceremonien, und macht zuerst eine Bestimmung über den Ort des Gottesdienstes, wo er fast fünf- oder sechsmal wiederholt, daß sie nicht an einem jeglichen Orte, der ihnen beliebt, sondern allein an dem, welchen der HErr erwählt hat, opfern und GOtt verehren sollen. Daher gebietet er auch, daß der Heiden Orte, Altäre und Säulen zerstört und ihr Gedächtnis vertilgt werden soll. Es war aber der Ort, den der HErr erwählt hatte, überall wo die Bundeslade und die Hütte war, wenn geredet wird von dem allgemeinen und beständigen Orte für das stete Opfer des ganzen Volkes. Sonst erwählte er öfter auf Einen Tag oder Stunde einen andern Ort, wie man an Samuel, Elias, Gideon, Manoah und anderen erkennt, aber auch dies geschah nicht ohne sonderliche Offenbarung durch einen Engel oder durch einen Propheten.

Dies alles gebietet er um deswillen, damit das Volk sich in der Verehrung GOttes nicht leiten ließe durch seine eigene Neigung, so heilig und gut sie auch immer sein möchte, sondern durch das Wort regiert werde; denn wenn der Mensch ohne das Wort nicht leben kann, dem Bauche nach, wie viel weniger wird er ohne das Wort im Werke GOttes und im Geiste leben können? Denn GOtt will, daß unser Gewissen getrost und sicher sei, daß es Gotte gefalle. Dies ist nicht möglich, wenn es durch sein Gutdünken geleitet wird, sondern nur, wenn es sich auf GOttes Wort gründet. Wenn sie deshalb GOtt an einem Orte verehren würden, den sie selbst erwählet hatten, so konnten sie doch nicht gewiß sein, daß sie Gotte gefielen, wenngleich sie sich selbst darin gefielen. Sie waren aber gewiß, daß sie Gotte gefielen, wenn sie an dem Orte opferten, der durch das Wort GOttes bestimmt worden war.

Es entsteht in diesem Capitel eine gewisse Dunkelheit und Schwierigkeit dadurch, daß Moses dasselbe so oft wiederholt, sodann auch dadurch, daß er wider sich selbst streitet, da er zweimal nachläßt [V. 15. 21.], daß sie in ihren Städten opfern (sacrificent) und Fleisch essen mögen nach aller Lust ihrer Seele, einmal, wenn der Ort des HErrn nahe dabei ist, das andere Mal, wenn dieser Ort ferne sein sollte, und doch verbietet er zugleich auch, daß sie nicht an einem jeglichen Orte essen oder opfern sollen (sacrificent), worauf er (wie ich gesagt habe) in diesem ganzen Capitel dringt. Die andere Schwierigkeit ist: Wie kann er zugeben, daß sie von den Zehnten und Opfern essen [V. 18.], was doch den Priestern und Leviten gebühre?

Auf das erste wollen wir so antworten: daß Moses in diesem Capitel nur darauf dringe, daß sie nicht anderswo opfern sollen als an der Stätte, die GOtt erwählt hat. Darum hat er dies im Auge, und hat nicht darauf Acht, wie er unterschiedlich rede, sowohl von den Opfern als auch von den Speisen, nur, daß er Eine Stätte bestimme. Darum will er so sagen: Ich gestatte dir zwar, daß du zu Hause allerlei Fleisch schlachtest (sacrifices) und essest, und an deinem Orte, beide der Reine und der Unreine [mögen’s essen, V. 15. 22.]. Aber alles, was du Gotte opfern (offere) sollst oder willst, von dem will ich durchaus nicht, daß du es anderswo opferst oder essest als an der Stätte GOttes. Wenn dieser Ort auch fern sein sollte, so lasse ich es dir dennoch zu, daß du allerlei Fleisch schlachtest und essest; aber was du Gotte opfern willst, das sollst du dahin bringen, oder wenn du es nicht dahin bringen kannst, so sollst du es zu Gelde machen, und kaufe an dem Orte GOttes anderes, wie er Cap. 14, 24–26. sagen wird, so daß der Sinn ist: das Fleisch, welches man nicht schuldig ist zu opfern, könne man frei essen an jeglichem Orte, aber die Opfer und das, was Gotte gebührt, könne man nur an der Stätte GOttes schlachten und essen.

Hieraus ist offenbar, daß Moses das Wort „opfern“ (sacrificandi) für das Wort „schlachten“ (mactandi) gebraucht, da nichts geopfert wird, es werde denn geschlachtet, wie auch Lucas in der Apostelgeschichte Cap. 10, 13. sagt:

\quellseite{743}
„Stehe auf, Petre, schlachte (sacrifica) und iß“, das heißt: Tödte und iß. Daß es sich so verhalte, erzwingen die Umstände des Textes, denn er verordnet auf das strengste, daß sie nur an dem von GOtt erwählten Orte opfern sollen (sacrificent), und doch erlaubt er ihnen Fleisch zu schlachten (sacrificari) und zu essen an jeglichem Orte. Denn da er dies von dem Fleisch sagt, welches Gotte nicht geopfert werden soll, so sagt er ohne Gefahr und Aergerniß, daß es entweder geschlachtet oder geopfert werde, da jedermann wohl versteht, daß da nicht vom Opfer gehandelt werde. Aber wenn er von dem Fleisch redete, welches Gotte geopfert werden sollte (offerendis), würde er nicht ohne Aergerniß zulassen, daß es an einem jeglichen Orte geopfert werden könne (sacrificandas). Doch, er zeigt zugleich unter diesem Worte heimlicher Weise das gemeinsame und geistliche Priestertum an, durch welches wir alle geistlicher Weise (mystice) uns selbst opfern (sacrificamus), Röm. 12, 1.

Zu der zweiten Schwierigkeit sagen wir dasselbe, daß Moses, da er sein Augenmerk auf die Bestimmung des Ortes gerichtet hat, nicht unterschiedlich redet vom Essen, sondern das ganze Volk mit den Priestern und Leviten zusammenfaßt, so daß der Sinn ist: Die Zehnten und die Opfer sollt ihr an der Stätte GOttes opfern und essen in Fröhlichkeit, aber ein jeglicher sein Theil, nämlich den Zehnten die Priester und Leviten, von den Opfern auch das gemeine Volk. Denn Moses sieht (wie ich gesagt habe) auf das ganze Volk, welches vor GOtt versammelt ist, wo einige die Zehnten, andere die Opfer, wieder andere nur Theile von den Opfern empfangen und essen, und alle fröhlich sind an einem und demselben Orte. Darum folgt nicht, daß das gemeine Volk die Zehnten gegessen habe, sondern, daß unter dem ganzen Volke Leute gewesen sind, welche sie gegessen haben, nämlich die Priester und die Leviten, während unterdessen das gemeine Volk sein Theil von den Opfern aß.

Bemerkenswerthe Aussprüche in diesem Capitel sind diese zwei [V. 8.]: „Ihr sollt der keins thun, das wir heute allhie thun, ein jeglicher, was ihm recht dünket“, und [V. 32.]: „Alles, was ich euch gebiete, das sollt ihr halten, daß ihr darnach thut. Ihr sollt nichts dazu thun, noch davon thun.“ Der erste Spruch verdammt den freien Willen ganz und gar, denn er hält das, was wir vornehmen (nostra studia), gegen das Wort GOttes, indem er in einem allgemeinen Satze den Willen ausspricht, daß von uns nichts erwählt oder gethan werden soll, ohne daß GOttes Wort vorhergeht und leuchtet; so soll auch nicht irgend ein Ort erwählt werden etc. Daraus folgt, daß alles, was wir ohne das Wort vornehmen, lauter Finsterniß und Irrthum sei. Denn wenn dies nicht wäre, so wäre es nicht vonnöthen gewesen, uns zu warnen, daß wir nicht thun sollten, was uns recht dünkt; wir würden auch des Wortes nicht bedürfen.

Der andere Spruch nimmt die Vermessenheit weg, daß wir ja nicht etwas, was GOtt geordnet hat, besser machen sollen, und beschenkt uns zugleich auch mit Freiheit, indem er uns losspricht von allen Werken, Bemühungen, Gesetzen und Satzungen der Menschen, und unsere Gewissen allein an das Wort GOttes bindet, wovon anderswo sehr vieles gesagt ist.

Merke aber, daß an dieser Stelle das Volk von Moses als ein gottloses beschrieben werde, da er sagt, daß sie thun, was ihnen, nicht was Gotte recht dünkt. Hieher gehört das Wort Apost. 7, 42., wo Stephanus aus Amos [Cap. 5, 25. f.] anführt, daß sie Götzendiener gewesen sind, und Paulus, Apost. 13, 18., sagt, daß GOtt ihre Weise geduldet habe, und der 94. Psalm [V. 10.], daß GOtt vierzig Jahre Mühe gehabt habe mit diesem Volk. Wie kann es denn nun bestehen, daß das Volk im vierten Buche Mosis Cap. 24, 6. ein gesegnetes Volk genannt wird, und Bileam es so herrlich lobt? Dies kann so bestehen, daß an beiden Stellen von der Sache geredet wird nach der [Redefigur der] Synecdoche,\footnote{Synecdoche = ein Theil für das Ganze oder das Ganze für einen Theil.} daß der größere Theil gottlos gewesen ist und der kleinere Theil gottselig. Aber um des Wortes GOttes willen, welches unter ihnen war in den Gottseligen, wird das ganze Volk gelobt, wie Paulus im Briefe an die Römer Cap. 3, 1. die Beschneidung, daß ist das ganze Volk, lobt, daß ihnen vertraut sei, was GOtt geredet hat, wiewohl viele nicht glaubten.

Die heimliche Deutung (allegoria) der Einen Stätte, die GOtt erwählt hat, ist der einige Sinn des Glaubens, in welchem alle Heiligen übereinkommen und darin Gemeinschaft haben.

\quellseite{744}
Hier wird also mit der größten Strenge geboten, daß wir in der Einigkeit (unitate) des Glaubens bleiben sollen, und nicht in Secten und unsere Meinungen zertrennt werden sollen, welche die Hügel, die Thäler, die Haine und überhaupt alle die Stätten sind, die uns gefallen; alles, was wir in diesen opfern, mißfällt [G]otte]. Aber denen, welche in demselben Sinn des Glaubens bleiben, schadet nichts, wenn sie sonst auch noch Sünder sind (peccent), sondern alles, was sie opfern, das gefällt [G]otte].

Aber anderswo als an der Stätte GOttes darf man allerlei Fleisch essen etc. Das heißt, allerlei äußerliche Dinge, welche nicht den Glauben unbetreffen, sind frei, das Gewissen kann ihrer gebrauchen, und nicht gebrauchen, als da sind solche fleischlichen Dinge als Kleider, Stätten, Speisen, Zeit, Personen, Geberden. Aber was des Glaubens ist, als unsern Leib zum Opfer begeben, und die Früchte des Geistes sind nicht frei, sondern müssen im Glauben geopfert werden, und man soll fröhlich sein vor Gotte wegen aller Güter, die man von GOtt empfangen hat, besonders wegen des Wortes der Gnade.

Daß man das Blut auf die Erde gießen soll, wie Wasser, und nicht essen mit dem Fleisch [V. 16. 23. f.], sei es bei dem Opfer oder bei anderem Essen, bedeutet, daß man sowohl in Sachen des Glaubens als auch der christlichen Freiheit nichts lehren noch dem folgen soll, was den alten Menschen, Fleisch und Blut an sich hat (sapit). Denn es ist ebenso gottlos, wenn man die gläubigen Gewissen mit fleischlichen Lehren der Werke gefangen nimmt, oder durch Werke gerecht machen will, als wenn man sie wiederum so frei machen will, daß sie ganz und gar nichts thun und durch die Freiheit dem Fleische Raum geben. Derjenige nun ißt nicht das Blut, sondern gießt es auf die Erde, wie Wasser, welcher, gerechtfertigt durch den Glauben, die gerecht machenden Werke verachtet, und dennoch den Müßiggang und die Freiheit des Fleisches verdammt. Wenn wir das thun, so wird es uns wohl gehen, wie hier [V. 25.] Moses sagt, und wir thun, was dem HErrn gefällt. Denn „das ist der gute, der wohlgefällige und der vollkommene GOttes Wille“, Röm. 12, 2.

\section{Das dreizehnte Capitel}

\begin{schriftzitat}
V. 1. ff. Wenn ein Prophet oder Träumer unter euch wird aufstehen, und gibt dir ein Zeichen oder Wunder etc.
\end{schriftzitat}

Zur Bestätigung dessen, was er am Ende des vorigen Capitels gesagt hat: „Ihr sollt nichts dazu thun“ etc., setzt Moses dieses ganze Capitel, indem er will, daß man dem Worte GOttes mit so großer Ehrfurcht anhange, daß wir uns weder durch Personen noch Zeichen bewegen lassen, mögen sie auch so gelehrt und heilig sein als die Propheten, oder gut und liebenswürdig sein wie Brüder, Kinder und Freunde, oder groß und viel wie Städte und Gewaltige. Man muß sich ganz und gar allein auf das Wort gründen und alles aus den Augen und Sinnen setzen, weil man GOtt verliert, wenn man das Wort verliert. Es ist aber besser, daß man Freunde, Brüder, Heilige, Gewaltige verliere, als GOtt.

Und hier siehst du, daß einem jeglichen das Recht gegeben werde, auch über die Lehren der rechten Propheten zu urtheilen, wie auch Christus befiehlt, Matth. 7, 15.: „Sehet euch vor vor den falschen Propheten“, wiewohl niemandem erlaubt ist, das äußerliche Recht zu vollstrecken und einen solchen Propheten zu tödten, sondern der Obrigkeit, wie uns im folgenden Capitel gelehrt wird. Denn ein jeglicher soll sein Gewissen bewahren, darum ist auch für jeden Menschen das Recht, die Geister und die Propheten zu richten, nothwendig, aber niemand hat das Recht, aus sich selbst das Schwert zu nehmen. Denn wozu bedürfte man sonst der öffentlichen Obrigkeit?

Hier erhebt sich die Frage: Wie verbietet er denn hier, einem Propheten zu glauben, der solche Zeichen thut und vorherjagt, welche kommen, während er doch im 18. Capitel, V. 22., sagt, daß man daran merken könne, daß es GOttes Wort sei, wenn das kommt, was der Prophet geredet habe? Ich antworte: In diesem
\quellseite{745}
Capitel redet er von dem schon vorhandenen (praeterito) Worte GOttes, welches schon angenommen und mit den nöthigen (suis) Zeichen bekräftigt ist, wie er sagt: „Ihr sollt nichts dazu thun noch davon thun“, und wiederum [Cap. 12, 1.]: „Haltet die Gebote GOttes“ und [Cap. 28, 1.]: „Gehorchet seiner Stimme.“ Wider ein solches Wort soll man keine Propheten zulassen, wenn sie auch Zeichen und Wunder regneten, nicht einmal einen Engel vom Himmel, wie Paulus im Briefe an die Galater [Cap. 1, 8.] sagt. Hier aber wird von einem Propheten geschrieben, der darum Zeichen thun werde, um andere Götter einzuführen. Dies nun ist wider das empfangene Wort, daß man nur Einen GOtt verehren soll; deshalb soll man ihn nicht hören. Im achtenzehnten Capitel aber redet er von einem neuen Worte, welches noch nicht empfangen worden ist, und welches auch nicht wider das empfangene Wort streitet; daher sagt er dort [V. 20.]: „Wenn ein Prophet reden sollte in meinem Namen, das ich ihm nicht geboten habe“ etc. Einen solchen Propheten soll man nicht zulassen, wenn er nicht Zeichen thut, wie auch Christus sein Evangelium bekräftigte mit Zeichen, als über Mosen hinaus gepredigt werden mußte. Denn GOtt offenbart nicht irgend ein neues Wort, ohne es mit Zeichen zu bestätigen. Er läßt es zwar zu, daß ein neues Wort gelehrt werde, damit er uns versuche, aber er ist getreulich uns zur Seite, daß nicht Zeichen geschehen, oder das nicht komme, was sie vorherverkündigt haben, wie an Elias gezeigt wird, da er mit den Baalspfaffen zu thun hatte, wie auch Paulus sagt [1 Cor. 11, 19.]: „Denn es müssen Rotten unter euch sein, auf daß die, so rechtschaffen sind, offenbar unter euch werden.“ So läßt er auch zu, daß Zeichen geschehen wider das angenommene Wort, wiederum, um uns zu versuchen (wie Moses hier sagt, V. 3.), ob wir ihn von ganzem Herzen lieben. So hat uns bisher der Teufel mit lügenhaften Zeichen und Wundern betrogen, und uns in kräftige Irrthümer gebracht [2 Thess. 2, 11.], wie Paulus vorhergesagt hat, indem wir Wallfahrten, Geistererscheinungen und etliche Heilungen bei gewissen Gräbern bewundert haben, wider das Evangelium, welches wir empfangen haben, so daß auch heilige Leute hier zu Falle gekommen sind, wie Augustinus, Bernhard, Hieronymus und viele andere, welche Orden und Regeln auf Grund der Werke gestiftet haben wider den reinen Glauben, die gewißlich verdammt worden wären (wie Wittek sagt), wenn sie nicht Buße gethan hätten, und, durch die Fülle des Glaubens bei ihrem unerkannten Irrthum, erhalten worden wären.

Wir haben es oft gesagt und sagen wiederum, daß andere Götter nicht bloß ein äußerlicher Götze sei, sondern vielmehr ein Wahn oder ein irrige Gewissen, welches man sich von dem wahren GOtte erdichtet; denn wie das Gewissen ist, so ist auch GOtt. Denn wenn du glaubst, man diene GOtt mit diesen oder jenen Opfern, an diesem oder jenem Orte, und zwar ohne das Wort GOttes, dann hast du schon den wahren GOtt verloren, und jener Wahn von solchem Opfer ist dein Gott, welchen [Wahn] du unter dem Namen des wahren GOttes festhältst. Wenn du glaubst, daß durch die Kappe und die Platte, durch Armuth, Gehorsam, Fasten, Speise und Trank GOtt verehrt werde (da du hier kein Wort GOttes hast), so ist schon die Kappe und die Platte dein Gott, oder jener Wahn von der Kappe und Platte. Darum, wie du inwendig den Wahn von der Kappe an der Statt GOttes hast, so schreibest du auch fort nach außen und richtest die Kappe als einen äußerlichen Götzen auf nach dem Bilde des inwendigen Wahns, du trägst sie, du ehrst sie, du dienst ihr und hältst sie in großem Werthe.

Siehe, das heißt andere Götter machen und ihnen folgen, das heißt den Göttern dienen, die du nicht kennst, denn du fühlst und verstehst auch nicht, daß du an der Statt des wahren GOttes einen Wahn und einen Götzen verehrst, hast auch nicht Acht darauf, wie ungewiß du seiest in diesem Dienste, und wie gar nichts du gedenkest oder haltest von dem wahren GOtte, weil du ohne sein Wort dir über ihn Gedanken machst. Er aber kann nicht anders weder gehalten noch gedacht werden als durch sein Wort. So siehst du, daß jede Weise, fremde Götter zu erfinden oder zu verehren, keine andere ist, als der gottlose Wahn, durch den wir ohne das Wort GOttes erwählen und glauben, daß wir durch dieses oder jenes Werk, durch diese oder jene Stätte, durch diese oder jene Ceremonie (ritu) GOtte gefallen, während er doch nicht ein solcher [Gott] ist, und dennoch unter seinem Namen ein anderer [Gott] im Herzen lügenhaft erdichtet wird. Daher folgen auch äußerlich verschiedene Namen, so wie
\quellseite{746}
verschiedene Götzen [aufkommen], so daß der eine Baal, der andere Astharoth, ein anderer Dagon, ein anderer Moloch, Peor, Camos genannt wird, und so noch mit viel mehr Namen, wie unsere Mönche, der eine vom weißen, der andere vom schwarzen Kleide, und ein jeglicher mit seinem Namen und nach seinem Werke, welches äußerlich von anderen unterschieden ist, genannt wird. Alle diese sind Propheten, welche durch ihre Träume betrogen werden, und sprechen: Lasset uns gehen und anderen Göttern dienen, das heißt, lasset uns neue Gebräuche (ritus) erwählen ohne das Wort, unter dem Namen des wahren GOttes.

\section{Das vierzehnte Capitel}

\begin{schriftzitat}
V. 1. Ihr seid Kinder des HErrn, eures GOttes; ihr sollt euch nicht zerschneiden.
\end{schriftzitat}

Nachdem der Ort des Gottesdienstes festgesetzt ist, bestimmt er hier die Personen oder die Gebräuche, welche die Person betreffen, nämlich, wie sie an ihrem Leibe beschaffen sein sollen. Er verbietet aber drei Dinge: das Zerschneiden, das Bescheeren wegen eines Todten, und das Essen unreiner Dinge. Unter „Zerschneiden“ wird hier nicht das Beibringen (impressio) von Malen verstanden, wovon er nachher reden wird, sondern ein gewisser Gebrauch bei Leichenbestattungen. Mit diesem Zerschneiden zerfetzten und zerrissen vielleicht die Egypter und andere Heiden ihr Angesicht; dann schoren sie auch die Augenbrauen ab zum Zeichen der Betrübniß und der Trauer. Denn er redet hier nicht von dem Kahlmachen des Haupts, sondern des Gesichtes zwischen den Augen, wie im hebräischen Text steht. Mit diesem Gebrauche will er nichts Anderes, als daß sein Volk solche Sitten haben soll, die denen anderer Völker ganz unmöglich wären, wie er sagt: „Ihr seid Kinder des HErrn, eures GOttes.“ Die heimliche Deutung hiervon ist dies, daß wir, da Christus von den Todten auferweckt ist, und nun durch den Glauben der Tod und die Sünde überwunden ist, kein Zeichen der Traurigkeit oder des Todes an uns tragen sollen, wie die anderen, welche durch das Bewußtsein der Sünde ohne Hoffnung des Lebens sind. Darum sind an ihnen keine Zeichen noch Worte des Lebens zu sehen, denn der Tod herrscht in ihnen durch die Kraft der Sünde, wir aber, die wir Kinder GOttes sind, leben mit ihm.

Hierauf lehrt Moses, welche Thiere für sie rein und unrein sein sollen, und es ist wahrscheinlich, daß diese Thiere buchstäblich (ad literam) alle unrein seien, das heißt, schädlich und unzuträglich zur Speise und für den menschlichen Leib.

Aber ein stärkerer Grund ist dies, daß dies Volk durch die Menge und Verschiedenheit der Gesetze nicht allein von anderen Völkern unterschieden werde, daraus ihnen Haß und Verfolgung wegen des Gesetzes GOttes erwüchse, sondern auch, daß sie zu schaffen hätten, damit sie nicht aus eigener Wahl (per sese) Gebräuche und Werke und Gesetze erfinden möchten, indem sie müßig wären und nichts zu thun hätten; sodann auch, damit dies Volk in Unterdrückung gehalten werde, auf daß durch die Menge der Sünden sein Hochmuth gebrochen würde, und es desto brünstiger nach Christo, dem ihm verheißenen Heiland, seufzte; zugleich auch, damit sie durch dieses äußerliche Exempel zur Reinigkeit des Geistes bewegt und zur Liebe und Sanftmüthigkeit gebildet würden, indem sie sähen, daß ihnen nur unreine, reißende und grausame Thiere verboten würden. Dies sei von dem buchstäblichen Verstande gesagt.

Nun folgt die Allegorie, von der wir gesagt haben, daß sie am sichersten auf das Amt des Wortes gezogen werde, da sie unkräftig (frigida) ist, und vielmehr nicht ein verborgener Verstand; sondern der historische Verstand ist es, der auf Sitten und Werke gezogen wird. So bedeutet auch hier die Unreinigkeit der Thiere die Lehre von den Werken, welche die Lehrer des Gesetzes und der Satzungen lehren, deren Herz und Gewissen niemals gereinigt wird, ja, „den Unreinen ist nichts rein, sondern unrein ist beide ihr Sinn und Gewissen“ [Tit. 1, 15.]. Denn je mehr sie sich reinigen durch Werke, desto mehr
\quellseite{747}
beflecken sie sich durch [falsches] Vertrauen und Heuchelei, durch Unglauben und solch gottloses Wesen. Daher vergleicht auch Petrus [2. Ep. 2. 22.] die Gottlosen, welche vom Glauben in die Werke zurückgefallen sind, mit einer Sau, die sich nach der Schwemme wieder im Koth wälzt. Sie sind aber reißend und grausam, weil ihre Lehre nicht allein mehr und mehr verunreinigt, sondern auch zerfleischt und verschlingt, weil sie nicht ein ruhiges und fröhliches Gewissen machen, sondern ein zerfleischtes, zerrissenes, wüstes und elend verschlungenes Gewissen; außerdem verführen sie auch sehr viele, und reißen sie von der lauteren Lehre des Glaubens. Daher nennt Christus sie „reißende Wölfe“ [Matth. 7, 15.], und Paulus „greuliche Wölfe“, die der Heerde nicht verschonen [Apost. 20, 29.]. Deshalb sind die gottlosen Lehrer der Werke unrein und grausam, indem sie die Herzen mehr verunreinigen, verschlingen und von Christo abreißen.

Hier überlasse ich müßigen Leuten die Aufgabe, daß sie durch alle diese Thiere erforschen, welches die Sitten der gottlosen Lehrer seien. Bekannt ist die Thorheit des Straußes, welcher meint, er sei ganz verdeckt, wenn er seinen Hals mit irgend einem Zweiglein bedeckt hat. So ergreift ein gottloser Lehrer irgend einen Ausspruch der Schrift und meint, daß sein Wahn schön sei, indem er nicht sieht, daß er von allen Seiten nackt und ohne Waffen seine Sache behauptet. Das Käuzlein ist gut für den Vogelfänger, weil der Schein der Frömmigkeit wunderbar anlockt. Die Nachteule sieht nicht, wenn es hell ist, bei Nacht sieht sie, denn der Gottlose gehorcht der Wahrheit nicht, kehrt sich aber zu seinen Fabeln. Der Kukuk säuft die Eier der Grasmücke aus und legt die seinigen in das Nest. Wenn nachher seine Jungen größer geworden sind, verschlingen sie ihre Mutter, die Grasmücke, im Nest, denn der gottlose Lehrer stiehlt die Kinder der Kirche und schiebt dafür seine Nachkommenschaft ein, welche die Kirche schließlich verderbt, von der sie ernährt worden war. So besudelt der Wiedehopf sein Nest, denn die gottlosen Lehrer verunreinigen (wie ich gesagt habe) sich selbst und die Kirche, wie Petrus [2. Ep. 2, 13. Vulg.] sagt: „Sie sind Schmutz und Flecken.“ Die geschwätzige Schwalbe thut sich zu den Menschen, aber sie ist kein nütze, und wenn sie verletzt wird, so nimmt sie schwere Rache am Vieh.

Ein gottloser Lehrer macht viele Worte und ist schmeichlerisch, hat aber keinen Nutzen, und wenn er verletzt wird, ist er unversöhnlich [und ruht nicht], bis daß er sich gerächt hat. So mag es jemand, den es gelüstet, mit dem Adler und anderen Thieren machen, und er wird sehen, daß alles gar gut auf die gottlosen Lehren passe.

Solche Dinge essen, das ist ihnen glauben und folgen, wie das Fleisch Christi essen an ihn glauben heißt, Joh. 6, 54. 56. Nicht essen ist, diese Greuel vermeiden. Aber die Klauen spalten ist (wie Paulus den Timotheus lehrt), das Wort GOttes recht theilen [2 Tim. 2, 15.], das heißt, so lehren, daß du das Evangelium und das Gesetz recht anwendest, mit dem Evangelium das Gewissen aufrichtest, lebendig und frei machest, und nicht mit dem Gesetze oder Werken und Sünden drückest und beladest; wiederum, das Fleisch nicht frei machest durch das Evangelium, sondern es mit dem Gesetze und Werken drückest und tödtest, wie es sich denn für den alten Menschen und den Leib der Sünde gebührt, daß er vernichtet werde. Wiederkäuen aber ist, mit dem Herzen das Wort aufnehmen, und es mit dem höchsten Fleiße bedenken, so daß man es nicht (wie das Sprichwort sagt) zu einem Ohre einfallen und zu dem anderen hinausfallen lasse, sondern es beharrlich im Herzen behalte, und es wiederhole, und zu seinem Fleisch und Blut (viscera) mache.

Wer daher entweder allein das Evangelium, oder allein das Gesetz über den zwiefachen (utrumque) Menschen, nämlich den innerlichen und den äußerlichen, setzt, der spaltet die Klauen nicht (das heißt, er übt das Amt des Lehrens nicht recht aus), weil er allein durch das Evangelium die ganze Freiheit lehrt, das ist, die Freiheit sowohl des Geistes als auch des Fleisches, so daß zugleich das Gewissen sich freut und befreit wird von Werken, und auch der Leib des Dienstes der Liebe müßig geht; wiederum allein durch das Gesetz die ganze Gefangenschaft lehrt, das ist, sowohl des Geistes als auch des Fleisches, so daß zugleich das Gewissen geplagt werde mit Bedenken (scrupulis) und Unruhe, und der Leib mit Mühe und Arbeit. Beides ist unrein und befleckt, und zeigt an, daß der Lehrer nicht wahrhaft gerechtfertigt ist durch den Glauben. Es werde also das Herz frei durch das Wort der Gnade, und der Leib ein Knecht
\quellseite{748}
durch das Gesetz der Liebe, so wird die Klaue recht gespalten sein.

Wiederum gibt es Thiere, welche die Klaue spalten und nicht wiederkäuen, wie die Sau oder das Schwein. Diese sind es, die beides recht lehren, aber während sie anderen predigen, selbst verwerflich werden [1 Cor. 9, 27.], weil sie ihren Leib nicht züchtigen, auch nicht in Knechtschaft bringen, weil sie todt sind mit der Welt, die in Wollüsten lebt [1 Tim. 5, 6.], und unter den Dornen und Disteln das Wort ersticken, daß es nicht Frucht bringe [Matth. 13, 22.]. Sie schöpfen das Wort, aber bald lassen sie es fallen, merken auch nicht mit ganzem Herzen darauf oder denken darüber nach. Diese sind Schweine in dem Pfuhl der Sünde. Andere aber käuen wieder, doch spalten sie die Klauen nicht. Diese sind es, die entweder allein das Evangelium oder allein das Gesetz mit ganzem Herzen annehmen, beständig daran hangen und festhalten, aber entweder allein das Gewissen quälen, oder allein den Leib frei machen, und diese sind also unrein wie das Kameel, der Hase und das Kaninchen.

Dasselbe ist von den Wassertieren zu halten, welche ohne Floßfedern und Schuppen sind. Denn sie haben die Floßfedern anstatt der Flügel, mit denen sie ihren Lauf lenken. Durch dieselben wird die Lehre des Glaubens bedeutet, welche durch das Wort GOttes die Herzen in dieser Welt regiert und führt. Die Schuppen aber sind gleichsam die Waffen und der Schmuck ihres Leibes, ohne welche die anderen schwimmenden Thiere gleichsam nackt sind. Sie bedeuten aber, daß der Glaube nicht müßig sei, sondern thätig durch den Geist der Liebe. Wenn daher Etliche Schuppen haben, aber ohne Floßfedern sind, so bedeuten sie die Werktreiber, ohne Glauben und Wort; als die Schlangen. Wenn sie aber Floßfedern haben, ohne Schuppen zu haben, so bedeuten sie diejenigen, welche das Wort haben, ohne Frucht und Werk; als, die Aale und die Muränen. Andere haben weder Floßfedern noch Schuppen, bei denen steht es weder recht mit dem Glauben noch mit den Werken; als, die Würmer, Blutegel \&c. Diese alle darf man nicht essen, noch sich damit gemein machen (communicari). So sind unter den Vögeln diejenigen, welche kriechen und Federn haben, unrein, und bedeuten die Heuchler, welche die Lehre des Worts nur zum Schein haben, als ob sie fliegen könnten; dennoch kriechen sie ganz und gar, und sind irdisch gesinnt.

Das Aas soll nicht gegessen werden; anderswo sagt er, daß man es auch nicht anrühren soll, sondern dem Fremdlinge verkaufen oder geben. Das heißt, lasset die Todten ihre Todten begraben [Matth. 8, 22.]. Wir sollen auch keine Gemeinschaft haben mit denen, die ohne das Wort des Lebens sind. Und wenn sie es inne gehabt haben, aber, nachdem sie das Wort verloren, gestorben sind, wie die Ketzer, sollen wir sie den Fremdlingen geben, das heißt, sie sollen uns sein wie Heiden. Wir sollen sie sich zu denen gesellen lassen und sie von uns thun. „Denn du bist ein heilig Volk dem HErrn, deinem GOtt“, das heißt, du bist durch sein Wort geheiligt, darum sollst du dich fremder Sünden nicht theilhaftig machen.

\begin{schriftzitat}
V. 21. Du sollst das Böcklein nicht kochen, weil es noch seine Mutter säuget.
\end{schriftzitat}

Gleichwie er an den unreinen Thieren die bürgerliche Reinigkeit verordnet, so auch an dem Böcklein die bürgerliche Barmherzigkeit, damit sie sich durch diese äußerliche Sanftmuth und Lindigkeit gewöhnten, sich auch der Menschen zu erbarmen und ihrer zu schonen, auch ihr Recht fahren lassen. Sodann, da er hier vom Gottesdienst handelt, will er nicht, daß ein unzeitiges (immaturum) Böcklein geopfert werde, sondern ein Jährling (wie er anderswo [2 Mos. 12, 5.] gebietet), damit sie nicht untüchtige und unnütze Dinge opferten. Es bedeutet dies aber das Wort des Apostels Röm. 14, 1.: „Den Schwachen im Glauben nehmet auf“, damit wir nicht diejenigen, welche der Milch bedürfen, mit fester Speise verderben, und sie nicht kochen oder tödten durch Märtyrerthum oder andere starke Versuchungen, ehe sie Festigkeit im Leben und Gewißheit in der Lehre erlangt haben. Denn GOtt läßt uns nicht versuchen über unser Vermögen. Deshalb müssen auch wir so glauben und lehren. Denn im Reiche Christi ist nur Ein Gesetz: „Einer trage des andern Last“, Gal. 6, 2.

\begin{schriftzitat}
V. 22. ff. Du sollst alle Jahre den Zehnten absondern \&c.
\end{schriftzitat}

Gott hat diesem Volke dreifache Zehnten auferlegt. Erstens den, welchen sie jedes Jahr den Leviten gaben von allen Früchten des Landes. Zweitens den, welchen sie jedes dritte Jahr, über
\quellseite{749}
den ehegenannten, von allen denselben Früchten des Landes absonderten, der gleichsam zu einem gemeinen Kasten und zu einem öffentlichen Almosen verwendet werden sollte, sowohl für die Leviten als auch für die Fremdlinge, auch für die Waisen und Wittwen, weil er nicht wollte, daß Arme im Lande wären, wie er im folgenden Capitel sagt. Von diesen beiden Zehnten redet er in diesem Capitel. Der dritte Zehnte war den Leviten aufgelegt, daß sie von ihrem Zehnten den Priestern auch den Zehnten gäben; davon ist im dritten Buche Mosis gesagt. Diesen Priestern gab GOtt auch alle Erstgeburt, so daß die Erstgeburt der Menschen mit Geld gelöst wurde. Es waren also die ordentlichen Einkünfte der Priester die Erstgeburt oder die Erstlinge und die Zehnten, die sie von den Leviten empfingen, wozu noch die Opfer und die Geschenke des ganzen Volkes hinzukamen. Die Einkünfte der Leviten aber waren die Zehnten, die sie von dem Volke erhielten und ein Theil der Opfer, wenn sie den Priestern dienten. Und wollte doch GOtt, daß alle anderen Schatzungen abgethan würden, welche heutzutage Lande und Leute verzehren, die auch nothwendiger Weise unbillig sein müssen, und [an deren Statt] die Zehnten eingesetzt würden; wenn Einer nicht ausreichte, so könnten noch drei, vier oder fünf hinzugethan werden. So könnte ein Volk unter seinen Herren leben. Dies wäre die ehrbarste und gerechteste Weise, um Einkünfte (reditum) zu beschaffen, welche allein von der Güte GOttes abhängen würde. Denn wenn GOtt in einem Jahre reichen Segen gäbe, so hätte das Volk reiche Einnahme, und die Obrigkeit reiche Zehnten. Wenn aber GOtt wenig Segen gäbe, so trüge die Obrigkeit gleiche Last mit den Unterthanen und empfinge weniger. Jetzt aber, da die jährlichen Steuern festgesetzt und gewiß sind, und die Jahre des Segens selten fallen, so werden die Unterthanen gezwungen, die vollen Steuern zu bezahlen, wenngleich sie in zehn Jahren nacheinander oder länger kein Jahr des Segens haben, ja, was noch elender ist, sie werden gezwungen, die vollen Steuern zu bezahlen, wenngleich sie in dem Jahre entweder kaum so viel oder nicht so viel von dem Lande empfangen haben, als sie bezahlen müssen, nämlich mit Gelde, welches sie anderswo geborgt haben, davon sie entweder die Steuern bezahlen oder davon sie in dem Jahre leben.

Ist dies nicht mehr als Barbarei? ja, eine Art thierische und grausame Raubgier? Könnte man also sagen, daß diese Steuern gerecht seien? daß diese Gesetze billig seien? Was ist es Anderes, wenn das Volk ausgezogen und aus allen Orten verjagt wird? Was ist heutzutage die Herrschaft in Deutschland anders als eine bloße Tyrannei? Von auswärtigen Reichen rede ich nichts, ich weiß nichts davon; wegen der Unseren bin ich bekümmert. Wer von den Fürsten oder obrigkeitlichen Personen, meinst du, kann heutzutage selig werden, wenn er dies wissentlich oder williglich leidet oder thut, wenn er der Sache helfen (wie es sich für einen Fürsten gebührt) könnte? Du siehst, wie sehr hier der Grund (ratio) der Seligkeit dazu bringt, daß die Fürsten mit einander rathschlagten, zusammen dahin arbeiteten und ins Werk setzten, damit durchaus alle anderen Schatzungen abgethan und in Zehnten (so viel ihrer noch sein möchten) verwandelt würden, und sie so das Volk erhielten und die Welt mit anderen Gesetzen regierten. Ich gestehe, daß das gemeine Volk in Ueppigkeit und Leichtfertigkeit versunken ist, wenn es Ueberfluß hat, aber diesem Uebelstande (rebus) muß man anders abhelfen als durch ungerechte Schatzungen, und nicht Sünde mit Sünde heilen. Sodann liegt auch die Schuld davon, daß der Pöbel so verderbt ist, bei den Fürsten, welche ohne Gesetz, ohne Sorgfalt, ohne Strenge regieren, und nur darauf beflissen sind, daß sie Schatzungen auflegen und Vergnügen haben.\footnote{Dieser ganze Absatz ist in der lateinischen Wittenberger Ausgabe weggelassen worden und statt dessen von den Herausgebern eingefügt (Witt. Tom. III [1549], fol. 57 b): „Aber über diese Dinge wird heutzutage vergeblich disputirt, denn wenn man den Zehnten zugestände zu nothwendigem Gebrauche, so würde der Geiz leicht den neunten, achten, siebenten, fünften oder dritten Theil der Früchte erfinden. Daher beklagen sich die Kinder Israel, daß sie unter Salomo bedrückt worden seien, und bitten den Rehabeam um Erleichterung, 1 Kön. 12, 4.“}

Aber wir wollen zu Mose zurückkehren. Siehe, eine wie schöne Ordnung darin ist. Zuerst handelt er von den jährlichen Zehnten und den täglichen Erstgeburten, durch welche für die Priester und die Leviten gesorgt wird, dann erst von den Zehnten in jedem dritten Jahre, durch die irgend welchen Armen geholfen werden sollte, damit wir erkennen, daß an erster Stelle die Diener des Wortes versorgt werden sollen, durch welche nicht der Bauch, sondern das Herz und der Geist
\quellseite{750}
gemeidet wird; darnach soll auch der Bauch der Armen versorgt werden. Denn der Glaube und das Wort ist eher als die Liebe und ein gutes Werk, darum muß man auch eher dafür sorgen, daß der Glaube und das Wort versehen werde als die Liebe und das gute Werk. Es wird aber nicht anders versehen, als daß die ernährt werden, welche lehren sollen. Ja, wenn der Glaube und das Wort untergingen, so würde nichts übrig bleiben von der Liebe und den guten Werken. Aber Gesetz ist Gesetz; es wird geredet, aber nicht gethan. Denn gleichwie das Volk Mosis dieses Gesetz anstehen ließ und seine Priester und Leviten nicht ernährte, daher sie gezwungen wurden, entweder das Wort und den Dienst GOttes zu verlassen, sich mit ihrer Hände Arbeit zu ernähren, oder neue Lehren und Abgötterei aufzurichten; hier nährte das Volk seine Verführer reichlich, welches die rechten Hirten vernachlässigte, weil auch sie die Ohren von der Wahrheit wandten und sich zu den Fabeln kehrten: so geschieht es auch jetzt; die Diener des Wortes ernährt niemand, so daß sie durch Mangel an Brod gezwungen werden, das Predigtamt zu verlassen, Ackerbau zu treiben und ein Handwerk zu erlernen, während sie doch die gottlosen Verführer bisher nicht allein ernährt haben, sondern auch die Welt mit Bisthümern und Klöstern erfüllt haben, welche den Reichthümern der Könige und Fürsten gleich sind. Dies ist geschehen aus gerechtem Gericht GOttes, daß diejenigen, welche einem Diener Gottes und Boten des Heils auch nicht Ein Brod gegeben haben noch geben, an die Diener des Teufels und die Boten des Todes ganze Königreiche und Fürstenthümer verschwenden.

Du siehst, mit welcher Sorgfalt Paulus 1 Cor. 9, 7. f. 1 Tim. 5, 17. f. Gal. 6, 6. f. und anderswo darauf dringt, daß man die Diener des Worts nähren soll, so daß es schmachvoll ist, daß man in dieser ganz geringen Sache in einem so großen und herrlichen Volke von einem so großen Apostel so viele Worte hören muß, wo dagegen die Liebe so brennen sollte, daß es vonnöthen wäre, das Volk vom Geben abzuhalten, wie 2 Mos. 36, 5. f., da wir nach dem Exempel der Galater [Gal. 4, 15.] (wenn es möglich wäre) auch unsere Augen ausreißen und solchen Dienern geben sollten. Daher fügt auch hier Moses hinzu, nachdem er von den Zehnten und der Erstgeburt geredet hatte [V. 23.]: „Auf daß du lernest fürchten den HErrn, deinen GOtt, dein Lebenlang“, nämlich, daß das Ernähren der Diener des Worts die erste und höchste Erweisung (usus) dessen ist, daß man GOtt ehre und ihn fürchte. Denn wer [die Diener des Worts] nicht ernährt, wie sollte der sich um GOttes Wort kümmern? Wer sich aber um GOttes Wort nicht kümmert, wie sollte der sich um GOtt kümmern? Wer sich um GOtt nicht kümmert, wie sollte der ihn fürchten? Also, wenn man den Diener des Wortes vernachlässigt und verachtet, ist es dasselbe, als ob man GOtt und sein Wort verachte. „Wer euch höret, der höret mich, und wer euch verachtet, der verachtet mich“, das ist es, was Moses hier und anderswo so sorgfältig einschärft, damit sie die Leviten nicht verlassen, welche kein anderes Erbe hatten.

Das, was er von dem allzuweiten Wege und der Entfernung der Stätte GOttes sagt, von dem Verkauf der Erstgeburt und der Zehnten, und daß anderes für dies Geld gekauft werden solle, ist droben in einem vorhergehenden Capitel [Cap. 12] erklärt, daß es nicht erlaubt war, in anderen Städten das, was geopfert werden mußte, zu schlachten und zu essen, oder wenn es gegessen wurde, so mußte man anderes an dessen Statt für Geld kaufen und es wieder ersetzen.

Die heimliche Deutung der Erstlinge kann man aus dem geistlichen Verstand (mysterio) der Tötung der Erstgeburt in Egypten verstehen. Ich halte dafür, daß diese Erstgeburt die erste Frucht des Gesetzes bedeute, das heißt, die Gerechtigkeit aus den Werken. Denn das Gesetz zwingt zuerst zu den Werken, und wenn es in fleischlichem Sinne genommen wird, so erzeugt es scheinende Heuchler, welche sich dünken lassen, sie seien die Ersten von allen, und ihnen gebühre alles. Diese Gerechtigkeit hat Christus durch seinen Durchgang, das ist, durch seinen Tod getödtet, durch den er gelehrt hat, daß durch des Gesetzes Werke kein Fleisch gerecht sein mag vor Gott, Röm. 3, 20. Aber die Erstgeburt Israels, welche dem HErrn geheiligt ist, ist die Gerechtigkeit des Glaubens, welche man ganz und gar dem verdankt, der die Gerechtigkeit der Werke verdammt hat. Daher, daß man die Erstgeburt und die Erstlinge opfert, das ist, daß man bekenne, die empfangene Gerechtigkeit des Glaubens sei durch seine Gnade, nicht durch unsere Kräfte zuwege gebracht. Denn die erste Frucht des Evangelii ist der Glaube im
\quellseite{751}
Geist, welcher auch uns zu Erstgebornen macht und uns nach zwiefachem Rechte zu Erben einsetzt, sowohl des Königreichs als auch des Priesterthums; die Zehnten aber sind, gleichsam als die letzte Frucht, die Güter der Liebe und des Geistes, die wir ebenfalls opfern und entrichten, indem wir bekennen, daß sie uns durch seine Gnade mittelst des Glaubens gegeben sind, daß wir sie nicht durch unsere Verdienste empfangen haben.

\section{Das fünfzehnte Capitel}

\begin{schriftzitat}
V. 1–3. Ueber sieben Jahre sollst du ein Erlaßjahr halten etc.
\end{schriftzitat}

Weil er im vorigen Capitel angefangen hatte zu handeln von den Zehnten im jedem dritten Jahre zur Versorgung der Armen, bei Gelegenheit der jährlichen Zehnten, die den Leviten gebührten für den Gottesdienst, bleibt er nun bei dem, was er angefangen hatte, und beendigt nun gänzlich dies Lehrstück (locum) von den Armen, indem er [von der Hauptsache] abschweift, in einem solchen Grade (adeo), daß er auch dies bürgerliche Gesetz vom Erlassen in die Behandlung des ersten Gebotes einmengt. Aber die Summa dieses Capitels ist, daß sie sich das nicht zu Schulden kommen lassen sollen, daß die Armen im Lande verlassen seien. Er sagt aber, daß immer Arme da sein würden, so daß sie Leute haben, denen sie geben können, entweder von den Zehnten im dritten Jahre, oder mit Borgen, oder auf andere Weise. Von den Zehnten ist geredet worden, nun wollen wir von dem Borgen sagen.

Er sagt, daß Geborgte könne vor dem siebenten Jahre wiedergefordert werden, aber im siebenten Jahre selbst solle alles, was dann noch als Geborgtes sich finden sollte, dem Bruder frei und geschenkt sein, und zwar nach göttlichem Rechte. Daß aber dieses siebente Jahr ein feststehendes und gewiß geordnetes war, wie das Jubeljahr war, in welchem das, was gekauft worden war, wieder zurückgegeben und freigelassen wurde, ergibt sich mit Nothwendigkeit daraus (necesse est), daß es nicht von dem Tage anfing, an welchem das Leihen stattfand, sondern es war ein allgemeines über das ganze Land und Volk, wie die Feste, das Pascha, die Pfingsten und andere waren. Es ist wahrlich ein sehr schönes und sehr billiges Gesetz; wollte doch GOtt, daß die Regenten der Welt demselben auch heutzutage nachfolgten, dann hätten sie weniger Fragen und Beschwerungen, da sie wüßten, daß die Processe, Hadersachen, Schulden, Handlungen (acta), Verträge, Rechte, Siegel, Briefe endlich im siebenten Jahre, möchte es nun nahe bevorstehen oder weit entfernt sein, auf einmal alle aufgehoben und ausgelöscht werden müßten, und daß sie nicht in ewige und endlose Untersuchungen (quaestiones) hinausgeschoben und beibehalten werden könnten. Zugleich wären sie gezwungen, vorsichtig zu sein, daß sie nicht eine so große Summe ausliehen, daß sie keine Hoffnung haben könnten, daß sie vor dem siebenten Jahre wiedergegeben werden könne, und so könnten sich auch Verschwendender und Schurken (perditis) nicht auf fremde Güter verlassen, die sie durch Schulden und Verträge zusammengebracht haben.

Aber was willst du zu Christo sagen, welcher Matth. 5,\footnote{In den Ausgaben: Matth. 6.} 40. ff. und Luc. 6, 34. f. verbietet, daß man das Geliehene wiederfordere, und gebietet, daß man leihe, ohne zu hoffen, daß man Gleiches wiedernehme? Ich antworte: Christus redet zu den Christen, welche über alle Gesetze sind, und mehr thun, als die Gesetze vorschreiben. Moses aber hat das Volk, welches der Obrigkeit und dem Schwerdte unterworfen ist, insofern dessen Angehörige Bürger sind (civilem populum), mit Gesetzen versehen, damit die Bösen im Zaum gehalten würden, und der öffentliche Friede bestände. Hier ist nun das Gesetz so zu verwalten, daß derjenige, der geborgt hat, wiedergebe, wiewohl ein Christ, wenn ein solches Gesetz ihm nicht zur Hülfe käme, ihm auch das Geborgte nicht wiedergegeben würde, dies mit Gleichmuth ertragen soll. So auch, wiewohl das Gesetz zwingt, daß man niemanden
\quellseite{752}
verlebe, und Rache nimmt an denen, die gewaltthätig sind, so trägt es doch ein Christ, wenn man ihm Unrecht thut, und rächt sich nicht, begehrt auch nicht Rache, wiewohl er die Strenge des rächenden Schwertes nicht verhindert, weil er weiß, daß das Schwert eingesetzt ist zur Rache über die Uebelthäter, wie Petrus 1 Petr. 2, 14. sagt.

Sodann, was ist das, daß er zuläßt, von einem Fremden das Geliehene wiederzufordern, auch im siebenten Jahre, das heißt, immer, aber nicht von einem Bruder? Ist denn nicht auch die Gerechtigkeit und die Liebe gegen einen Fremden zu beobachten? Darauf ist zu antworten, daß auch dies geschehe aus einer billigen Ursache des Weltregiments (publicae politiae), daß die Bürger vor den Auswärtigen und Fremdlingen mit einem Vorrechte geehrt werden, damit nicht alles durch einander und gleich sei; so geschah es auch bei den Römern, da sie einigen Städten das Recht einer Colonie gaben, anderen das Recht einer römischen Stadt. Denn diese Unterschiede (formas), wiewohl sie einen Schein der Unbilligkeit zu haben scheinen, muß die Welt in ihrem Regimente nothwendiger Weise haben, wie denn auch der Knechte und der Mä gde Stand ist, desgleichen der Handwerker (operariorum) und Arbeiter. Denn nicht alle können in gleicher Weise Könige, Fürsten, Rathsherren, Reiche, freie Leute sein, da die Welt ohne mancherlei und verschiedene Personen nicht bestehen kann. So sehr nun auch vor GOtt kein Ansehen der Personen ist, sondern alle gleich sind, so ist doch für die Welt das Ansehen der Personen und die Ungleichheit nothwendig, und zwar darum, damit die Bösen im Zaume gehalten werden, und der öffentliche Friede bestehe, welcher unter der Gleichheit und Unterschiedslosigkeit der Personen nicht bestehen kann.

Doch das Volk der Juden hatte ein weitergehendes und höheres Recht, nicht allein das Geliehene wiederzufordern, sondern, wie er hier [V. 6.] sagt, auch an den Heiden zu wuchern und Zinsen von ihnen zu nehmen, nämlich aus göttlicher Gewalt, welche selbst dies ordnet und zuläßt. Denn er ist GOtt und der HErr über alle Dinge, der nicht bloß Geld und Sachen, sondern auch Königreiche und Herrschaften (imperia) wegnimmt, wenn er will, auf welche Weise er immer will, und gibt, wem er will. Wenn also GOtt, um Rache an den Heiden zu nehmen, sie durch Zins und Wucher bestrafen will, und dies den Juden aufträgt, daß sie es thun sollen, so thun die Juden wohl daran, daß sie sich Gott gehorsamlich zum Werkzeug hergeben, und seinen Zorn an den Heiden durch Wucher und Zinsen ausrichten, nicht anders, als da er ihnen gebot, daß sie die Amoriter und Cananiter vertreiben sollten. So, wenn ein Ehemann Gott dessen werth zu sein schiene, daß ihm sein Weib oder seine Kinder genommen werden sollten, und mir dies durch sein Wort aufgetragen würde, so wäre ich nicht ein Ehebrecher oder ein Menschenräuber, wenn ich ihm sein Weib und seine Kinder raubte, sondern eine gehorsame Ruthe GOttes über den gottlosen Mann.

Hier wird die Frage aufgelöst, wie es den Juden erlaubt gewesen sei zu wuchern. Denn es ist zu antworten, daß es ihnen nicht freigestanden habe oder freistehe wegen ihrer Verdienste oder nach dem gemeinen Rechte, sondern durch den Zorn GOttes über die Heiden, den er durch die Juden als die Werkzeuge seines Zornes ausrichten will. Dennoch hätten sie sich auch dieser Freiheit nicht gebrauchen dürfen, wenn es ihnen nicht durch ein gewisses und offenbares Wort GOttes befohlen, und sie zu solchen Werkzeugen erwählt worden wären. Denn sie waren auch nicht besser als irgend welche Heiden, wie ich oben gesagt habe, nur daß GOtt sie allein aus Barmherzigkeit erwählt und angenommen hat, so daß, wenn du die Sache recht betrachtest, nicht die Juden selbst es sind, welche wuchern, sondern GOtt ist es, der durch den Wucher der Juden die Heiden verfolgt. Das ist genugsam bewiesen worden, da er die Juden, die ihm ungehorsam waren und sündigten, wiederum den Heiden übergab, nicht allein daß sie mit Wucher beschwert werden sollten, sondern daß sie mit aller möglichen Schmach gequält werden sollten, ja auch bedeutend grausamer, als er die Heiden jemals den Juden übergeben hat; wie er auch im 28. Capitel [V. 12.] dieses Buches vorhergesagt und in diesem Capitel [V. 6.] anfügt, daß sie vielen Völkern leihen werden, jedoch nur, wenn sie der Stimme des HErrn gehorchen [V. 5.], als ob er sagen wollte: wenn sie nicht gehorchen würden, sollten sie nicht allein nicht leihen, sondern den Heiden gleich, oder noch elender sein, was ihnen auch widerfahren ist.

Jedoch heutzutage, da die Juden aufgehört haben das Volk GOttes zu sein, ist das Gesetz aufgehoben, und sie haben mit ihrer Gottlosigkeit
\quellseite{753}
und Gotteslästerung den Zorn GOttes verdient, und es ist ihnen der Wucher nicht zu gestatten, sondern sie sollen unter die Gesetze der Heiden gezwungen werden, unter denen sie leben; wiewohl, wenn du den Text wohl ansiehst, so gebietet er nicht zu wuchern, es sind auch nicht Worte des Gesetzes, sondern vielmehr Worte der Verheißung, da er sagt: „Wenn du der Stimme des HErrn gehorchest, so wirst du vielen Heiden leihen“ (foenerabis). Aber die Worte der Verheißung zu erfüllen, steht nicht bei Menschen, wie die Worte des Gesetzes, sondern allein bei GOtt, der dies verheißt, so daß dies der Sinn sein muß: Wenn du der Stimme des HErrn gehorchst, so werden durch GOttes Wirkung die Heiden in ein solches Elend geraten, daß sie auch, wenn du nicht daran denkst noch es suchst, Wucher bezahlen müssen, und dir in allen Dingen unterworfen werden, und du über sie mit allen ihren Gütern herrschen wirst, so daß du nehmen kannst (rapias), einfordern (exigas), wuchern, wie es dir gefällt. Der HErr ist es, der dir die Heiden so unterwerfen und sie durch dich demütigen wird.

\begin{schriftzitat}
V. 4. (Es soll allerdings kein Bettler unter euch sein.
\end{schriftzitat}

Eine sehr schöne weltliche Ordnung, die aber niemals gehalten worden ist, darum bleibt auch dies Gesetz Mosis, soweit es das ganze Volk betrifft, nur in Worten. Wenn in diesem Volke die Bettelei verboten ist, mit welchem Rechte wird sie unter den Christen, als ob sie heilig wäre, auch durch Gesetze aufgerichtet? Die Armuth wird gepriesen, aber in solcher Weise, daß man ihr zu Hülfe kommen solle. Sodann wird die Armuth des Geistes gelobt, aber es wird geboten, daß man der äußeren Armuth abhelfe, nicht anders als anderen Widerwärtigkeiten des Nächsten. Und es ist zu verwundern, warum diese Lobredner der äußeren Armuth nicht auch Wunden, Krankheiten, Gefängniß, Blöße, Verbannung, Hunger, Durst, Schwert, Gefahren, Tod, Sünden, den Teufel und alle anderen Uebel auf sich nehmen und sich dazu verbindlich machen (profiteantur), indem hierzu neue Gelübde eingesetzt werden, wie sie es mit der Armuth gethan haben, so daß einer eine Krankheit leide, ein anderer Gefängniß, ein anderer Hunger, ein anderer Sünden oder den Teufel, da Christus befiehlt, daß man für diese Dinge bei dem Nächsten ebensowohl Sorge tragen und sie abwenden soll als die Armuth, da er sagt [Matth. 25, 43.]: „Ich bin krank gewesen, und ihr habt mich nicht besucht“ etc.

Aber unsere Lobredner der Armuth haben anstatt der Krankheiten und Wunden eine wohlgepflegte (obesam) Haut und ein gemästetes Fleisch, mehr als Schlemmer und Huren, anstatt der Verbannung haben sie Häuser, die prächtiger sind als die Paläste der Könige, anstatt des Hungers verschlingen sie die Vorrathskammern aller Leute, für den Durst haben sie volle Keller, anstatt des Todes haben sie ein ganz angenehmes und sicheres Leben. Dann singen sie uns von der Herrlichkeit der Armuth, die man tragen müsse. Aber GOtt hat geboten, daß man sie abthun solle, damit wir seien, wie in der Apostelgeschichte auch für das christliche Volk ein Exempel geschrieben steht [Apost. 4, 34.]: „Es war auch keiner unter ihnen, der Mangel hatte.“

Es soll also in dem Volke GOttes keine Armuth und Bettelei sein, sondern Fürsorge und Aufsehen, daß da keine Armuth und Bettelei sei, auf daß du wissest, daß die Bettelorden und alle, welche sich zur äußeren Armuth verbindlich machen (professores) und dieselbe rühmen, Jünger und Diener des Teufels sind, die geradezu wider den HErrn und seinen Gesalbten (Christum) wüthen, gleichwie in dem Volke GOttes auch nicht Krankheit, Hunger, Durst, Verbannung, Tod, Sünde und Teufel sein soll, sondern Fürsorge und Aufsehen, wenn etwas derartiges unter ihnen eintreten sollte, daß es so bald als möglich beseitigt werde, und Sorge getragen, daß es nicht unter ihnen sei. Ich sage, man soll die Armuth nicht beschwören, nicht erwählen, nicht lehren, denn sie ist von selbst überall genugsam vorhanden, wie er hier sagt [V. 11.]: „Es werden allezeit Arme unter euch sein“, wie auch alles andere Uebel. Aber dafür soll man beständig Sorge tragen, daß man diesen Dingen immer entgegen trete,\footnote{In der Erlanger: occupatur statt: occuratur.} weil sie allezeit vorkommen. Du siehst also, was es mit der Einrichtung, daß man Armuth gelobt, und mit dem ganzen Reiche des Papstes auf sich habe.

\begin{schriftzitat}
V. 7–11. Wenn deiner Brüder irgend einer arm ist etc.
\end{schriftzitat}

Das ist der Lauf der Welt, daß man, sobald ein Gesetz gegeben ist, auch ausfindig macht,
\quellseite{754}
wie man das Gesetz umgehen kann. Deshalb kommt Moses auch hier der Tücke zuvor, welche sich einstellen würde bei dem Gesetze von dem siebenten Erlaßjahr, indem der Geiz und die menschliche Hartherzigkeit so denkt, wenn etwa das Erlaßjahr nahe ist: Was soll ich thun? Soll ich leihen? Aber nach einigen Monaten wird das Erlaßjahr da sein. Dann habe ich es vergebens dahingegeben, weil keine Hoffnung da ist, es wiederzubekommen, auch nicht das Recht, es wiederzufordern. Wider diesen Geiz redet hier Moses mit wunderlichen und scharfen Worten, indem er solches ein hartes Herz nennt, ein Belialswort, eine schändliche That, ein gottloses Auge, endlich auch eine Sünde, die zu GOtt schreit, so daß du sehen kannst, daß auch Moses mit Christo übereinstimme in der Lehre vom Leihen, daß er auch unter Androhung, man solle so großer Vergehen schuldig sein, befiehlt, man solle leihen, wenngleich keine Hoffnung da sei, es wiederzubekommen, oder Recht, es wiederzufordern, so daß der Inbegriff dieser Lehre ist, man solle für die Armen sorgen in der Liebe.

\begin{schriftzitat}
V. 12–14. Wenn sich dein Bruder, ein Hebräer oder Hebräerin, dir verkauft, so soll er dir sechs Jahre dienen etc.
\end{schriftzitat}

Er erneuert (revocat) das Gesetz 2 Mos. 21,\footnote{In den Ausgaben: Ex. 22.} 2. f. [und zieht es] zu diesem Gesetze von dem Erlassen des Geliehenen, weil es auch von den Armen redet, die sich, gezwungen durch Armuth, verkauft haben und gleichsam etwas geliehen bekommen haben, indem sie sich mit ihrem eigenen Leibe zu Schuldnern gemacht haben. Er fügt aber an dieser Stelle hinzu, daß er ihn nicht leer von sich lassen soll, wenn er ihn frei gibt, und zeigt zwei Gründe dafür an, erstlich, daß sie gedenken sollen, daß sie auch Knechte waren in Egypten, zweitens, daß er [der Knecht] bei ihm ein zweifältiger Tagelöhner gewesen sei. Ich glaube, dies werde darum gesagt, weil der, welcher sich einem andern verkauft, sich zweierlei Schaden zufügt: erstlich, daß er einem andern dient, und alles für seinen Herrn erarbeitet, zweitens, daß er unterdessen das Seine versäumt, und das, was er für seinen Herrn gewinnt, für sich hätte gewinnen können. Daher würde es ganz unbillig sein, wenn er völlig leer abgeschickt würde. Deshalb sagt er [V. 18.]:

\begin{quote}
„Laß dich’s nicht schwer dünken, daß du ihn frei los gibst.“
\end{quote}

Es fügt aber Moses immer herrliche Verheißungen hinzu: daß GOtt gesegnet habe, daß er segne, und daß er segnen werde diejenigen, welche dieses thun, auf daß sie nicht daran zweifeln, es werde ihnen reichlich wiedererstattet werden, wenn sie entweder etwas gegeben oder weggeliehen haben, wie auch Christus sagt [Luc. 6, 38.]: „Gebet, so wird euch gegeben.“ Oder kommt auch das aus Sprüchwort Salomonis [Spr. 19, 17.]: „Wer sich des Armen erbarmet, der leihet dem HErrn“, und wiederum [Spr. 14, 31.]: „Wer sich des Armen erbarmet, der ehret GOtt“, und viele solche Sprüche, wie dagegen viele Drohungen, die diesen entgegengesetzt sind, wider die, welche sich nicht erbarmen, wie denn deren die Schrift voll ist. Aber dies verschlägt nichts bei der tauben Gottlosigkeit und dem Unglauben, der da meint, daß GOtt mit solchen Worten entweder scherze oder lüge, und er ist es werth, daß er der Güter dieses und des zukünftigen Lebens beraubt werde, gleichwie [1 Sam. 25, 10. f.] jener Narr Nabal zu Carmel dem David seine Hülfe versagte, aber bald, des Lebens beraubt, selbst nichts behielt.

\begin{schriftzitat}
V. 16–18. Wird er aber sprechen: Ich will nicht ausziehen von dir etc.
\end{schriftzitat}

Daß das Ohr mit einem Pfriemen an dem Pfosten in der Thür des Hauses durchbohrt ward, war ein äußerliches Zeichen beständiger Knechtschaft, wie man denn in menschlichen Dingen Zeichen haben muß, um die Personen zu unterscheiden. So gebraucht der Kaiser das Prachtkleid, das Beil, das Diadem; das Weib den Schleier; die Jungfrau die Krone; das Kind ein Knöpflein (bulla). Es war aber ein geeignetes und passendes Zeichen, daß das Ohr mit dem Pfriemen durchbohrt wurde. Der Pfriemen bedeutet, daß das Gebot seines Herrn scharf und eisern sei, nämlich, daß es durchdringe, und das Wort des Herrn bei dem Knechte wirksam sei. Das Ohr bedeutet den Gehorsam, welcher dem Gebote sich unterwirft. Daß dies aber in der Thür geschieht, bedeutet, daß ihm das Recht und die Macht (usus) gegeben wird, einzugehen und auszugehen und alles auszurichten, und daß er unter das Gesinde des Hauses aufgenommen wird. Ferner, Ohrgehänge, die aus Gold, Silber und Edelsteinen
\quellseite{755}
gemacht sind, sind ein Schmuck der Freien. Aber der Pfriemen der Sclaven ist eisern, weil der Freien Gehorsam willig und frei ist, ganz golden und köstlich, aber der der Knechte gezwungen, um Lohn, und eisern und gering.

\begin{schriftzitat}
V. 19. ff. Alle Erstgeburt, die unter deinen Rindern und Schafen geboren wird etc.
\end{schriftzitat}

Nun holt er nach und kommt wieder ganz und gar auf die Steuern, die den Priestern gegeben werden soll, nachdem die Abschmeißung über die Versorgung der Armen beendigt ist, um hinzuzufügen, was oben hätte hinzugefügt werden sollen, nämlich, daß die Erstgeburt, welche GOtte geopfert würde, so beschaffen sein müsse, daß sie durchaus ohne Fehl und Gebrechen sei. Darum sagt er hier, daß mit dem Erstling der Ochsen auch nicht geackert werden solle, und die Erstlinge der Schafe nicht beschoren werden sollen, so daß die Erstgeburt ganz unversehrt sei, und zu keinem Gebrauch oder Nutzen verwendet werde, sondern nur GOtte geopfert, und gegessen werde an dem von GOtt erwähnten Orte, wie man sich denn auch von dem Heiligen und dem, was GOtte gebührt, enthalten soll. „Du sollst essen“ [V. 20.] (sagt er), das heißt, du sollst es hingeben, daß es gegessen werde, und sollst unter denen, die es essen, versammelt sein, wie ich oben [Cap. 12] gesagt habe. So, wenn es irgend einen andern Fehl haben sollte, sollte es nicht geopfert werden, sondern ein jeglicher möchte es zu Hause in seiner Stadt essen.

\subsection{Heimliche Deutung}

Das Jahr des Erlasses der Schuld ist die ganze Zeit der Gnade, weil das ganze Reich Christi nichts Anderes ist, als fort und fort die Schulden erlassen, so daß sie erlassen werden. Daß sie aber dem Fremden nicht erlassen werden, bedeutet, daß außerhalb der Kirche Christi keine Vergebung der Sünden sei. Denn es kann denen auch nicht vergeben werden, welche nicht wollen, daß man ihnen erlasse, die sich selbst rechtfertigen und die Kirche verachten. Also muß es von diesen wiedergefordert werden, bis daß sie wiedergeben, das heißt, man muß ihnen nichts nachlassen, sondern immer fordern, damit sie anders gesinnt werden, und anders thun, bis daß sie Buße thun; dem Bruder aber, der Buße thut, soll alles erlassen werden. Denn wo der Glaube ist, da gibt es keine Sünden, die nicht vergeben und geschenkt werden könnten, wo aber der Unglaube bleibt, da gibt es keine Sünden, die nicht verdammt und gefordert werden müßten. Daß man von niemandem borge, und alle von ihm borgen [V. 6.], das heißt, daß man so gerecht lebe, daß man niemandem etwas schuldig sei, aber alle ihm schuldig seien, weil er niemanden verletzt, er aber von allen verletzt wird. Daß seine Armen sein sollen [V. 4.], und wenn sie da sind, man ihnen helfen soll, das heißt, daß man die Schwachen im Glauben und Leben tragen, lehren und ermahnen soll, und für sie beten ohne Falsch in völliger Einfalt des Herzens, und solche Leute werden allezeit unter euch sein [V. 11.]. Daß man aber den hebräischen Knecht, den Bruder, nicht ohne eine Wegzehrung von sich lassen soll, das heißt, daß man ihm nicht allein die Schuld vergeben, sondern ihm auch helfen soll, indem man noch eine Wohlthat der Liebe hinzufügt. Wiewohl dieser Knecht nach der eigentlichen geistlichen Deutung das Volk des Gesetzes ist, welches in einer harten Knechtschaft dient, so wird er doch, wenn das Erlaßjahr kommt durch das Evangelium, nicht allein freigelassen von dem Gesetze, sondern auch mit dem Worte des Evangelii verehrt, wovon er leben kann; denn er hat in einer zwiefachen Knechtschaft gedient, daß er die Werke des Gesetzes gethan, und doch nichts gewonnen hat, als daß sein Herr, das Gesetz, zugenommen hat; denn je mehr wir unter dem Gesetze Werke thun, desto mehr fordert es, und desto mächtiger wird es, und wird desto weniger erfüllt, weil nämlich der Haß gegen das Gesetz beständig zunimmt; und dabei (tamen) vernachlässigt er unterdessen auch das Seine, das heißt, er erlangt nicht eine Zuversicht des Gewissens, sondern verliert sie, während er sie vielmehr außerhalb des Gesetzes durch das Evangelium hätte erlangen können. Daß ihm aber mit dem Pfriemen sein Ohr durchbohrt wird, und er ewiglich ein Knecht bleibt, das heißt, daß er, schon frei im Geiste, dennoch sein Fleisch dem Gesetze desto stärker unterwirft und es durch das eiserne und harte Gesetz zwingt zum Gehorsam, wie Paulus sagt [1 Cor. 9, 27.]: „Ich betäube meinen Leib und zähme ihn“ etc. So bleibt er zugleich ein Knecht und ein Freier.

Daß der Erstling vom Ochsen nicht zum Ackern
\quellseite{756}
gebraucht und das erstgeborne Schaf nicht beschoren werden soll, das ist, daß wir der Gerechtigkeit des Glaubens nicht zu unserem Nutzen oder Ruhm mißbrauchen, sondern einfältiglich GOtte opfern sollen, daß er den Ruhm und die Ehre habe. Denn sie kommt nicht von unserem Werke her, ist auch nicht unter unserer Gewalt, daß wir darüber herrschen sollen, sondern sie gehört allein GOtte zu, der sie schafft und gibt. So will Paulus nicht über den Glauben der Corinther herrschen [2 Cor. 1, 24.]. Diese geistliche Deutung geht also besonders die Diener des Worts an, daß sie das Volk des Glaubens nicht ihren Gesetzen und den von ihnen verordneten Werken unterwerfen, sondern es allein GOtte opfern, daß es ihm unterthan sei und durch den Glauben regiert werde, wie Paulus sich rühmt Röm. 15, 16., daß er das Evangelium so opfere (sanctificare), daß die Heiden ein Opfer werden, GOtt angenehm, das heißt, durch das Evangelium opfert er die Heiden GOtte, daß sie nicht ihm, dem Paulus, sondern GOtte unterthan seien. Und was thun denn die Diener des Worts, wenn sie uns ihren Werken unterwerfen, als daß sie uns scheren und ihren Vortheil und Ehre suchen, nämlich die Wolle und die Frucht unserer Arbeit?

Daß kein Fehl daran sein soll, das ist, daß das Gewissen im Glauben gesund sein soll, so daß es nicht zugleich durch Werke und durch den Glauben gerechtfertigt werden und GOtte gefallen will. Darum schilt der HErr auch heftig in dem Propheten Maleachi [Cap. 1, 7. ff.] über dieses Laster, denn es ist kein Scherz mit dem Glauben und seiner Lehre. Tit. 2, 8. [Vulg.]: „Sein Wort ist heilsam und untadelig.“

\section{Das sechzehnte Capitel}

\begin{schriftzitat}
V. 1. Halte den Monat Abib, daß du Passah haltest dem HErrn.
\end{schriftzitat}

Nachdem die Stätte und die Personen beschrieben sind, und für diese gesorgt, macht er hier eine Ordnung in Bezug auf die Zeiten und Feste, besonders über die drei Hauptfeste, welche vor allen Dingen dies Volk angingen, dessen Belehrung Moses in diesem Buche auf sich genommen hat. Es sind aber diese Feste das Passah, Pfingsten und das Fest der Laubhütten, ferner das Fest der Posaunen und das Fest der Versöhnung in demselben siebenten Monate, in welchem das Fest der Laubhütten gefeiert wurde, 3 Mosi. 23, 24., nämlich das Fest der Posaunen am ersten Tage, das Fest der Versöhnung am siebenten Tage. Aber jedes dieser beiden Feste wurde nur Einen Tag gefeiert. Doch das Fest der Laubhütten wurde mitten im Monat, nämlich am fünfzehnten Tage, acht Tage lang gefeiert, an welchem alles, was männlich war, vor dem HErrn erscheinen mußte an der Stätte, die der HErr erwählt hatte, wie an dem Feste des Passah und der Pfingsten, wie er in diesem Capitel sagt.

Der Text an sich ist aber leicht, nur einige Wörter müssen erklärt werden. „Im Monate Abib“, den unser lateinischer Text [die Vulgata 2 Mos. 23, 15.] den Monat der neuen Dinge, anderswo aber [wie hier V. 1.] den Monat der neuen Früchte nennt, vielleicht, weil dann alles neu hervorkommt, und die ganze Welt grünt. Es ist aber der Monat April auch bei den Lateinern von dem Aufthun (Aprillis ab aperiendo) benannt, weil sich dann die Erde aufthut und alles neu wird, wiewohl die Hebräer den April nicht Abib, sondern Ijar\footnote{In der Erlanger und der Jenaer Ausgabe: Iier; in der Wittenberger: Eier. Beides ist falsch und es sollte Ijar gelesen werden. Vergleiche Seiffarths Berichtigungen der römischen, griechischen, persischen, ägyptischen, hebräischen Geschichte und Zeitrechnung etc.} nennen. Abib aber kommt her von Abab, welches eine kleine Ruthe oder einen zarten Zweig bedeutet, 3 Mos. 2, 14., wie diejenigen sind, wenn die Bäume im Frühling ausschlagen, an denen man, wie Christus Matth. 24, 32. anzeigt, erkennen kann, daß der Sommer nahe sei. Zu der Zeit, nämlich im April, sind sie aus Egypten gezogen; deshalb wird dieser Monat der Anfang des Jahres und der erste Monat genannt, 2 Mos. 12, 2. Andere aber machen aus dem Monat Abib den Juli. Dies wird aber widerlegt durch den Auszug der
\quellseite{757}
Kinder Israel aus Egypten, welcher zu Anfang des Jahres im Monat Abib geschah, wie er hier und 2 Mos. 34, 18. sagt.

\begin{schriftzitat}
V. 3. Sieben Tage sollst du ungesäuert Brod des Elends essen; denn mit Furcht bist du aus Egyptenland gezogen.
\end{schriftzitat}

„Brod des Elends“ nennt er es von der vergangenen Trübsal, welche sie zu der Zeit erlitten haben, da sie dieses Brod zuerst aßen, was er dadurch erklärt, daß er bald hinzufügt: „Denn mit Furcht bist du ausgezogen“ etc., das heißt, mit Angst und Zittern, wie diejenigen, welche in Nöthen sind, zu eilen pflegen und sich zu ängstigen, so daß sie so schnell als möglich fliehen. Denn das ist die Bedeutung dieses Wortes, daß es nicht bloß eilen oder furchtsam sein heißt, sondern (wie ich gesagt habe) aus Angst zu fliehen suchen, wie über David 1 Sam. 23, 26. gesagt wird, da er vom Saul auf dem Berge völlig umzingelt war, daß er verzweifelte, das ist, daß er mit ängstlichem Gemüthe darauf bedacht war zu eilen. So heißt es 2 Mos. 12, 11.: „Und sollt es essen, als die hinweg eilen“, das ist, was er hier sagt, „das Brod des Elends“, welches sie in Furcht und ängstlicher Eile aßen. Und Ps. 104, 7. [Vulg.]: „Vor der Stimme deines Donners fürchteten sie sich“, das heißt, sie hätten fliehen mögen und davoneilen. Psalm 116, 11.: „Ich sprach in meinem Zagen“, das heißt in ängstlicher Eile.

\begin{schriftzitat}
V. 8. Am siebenten Tage ist die Versammlung des HErrn, deines GOttes.
\end{schriftzitat}

Unser lateinischer Text (noster) hat das Wort „Sammlung“ (collectam) gebraucht, aber richtiger bedeutet es [\hebraeisch{עצרת}] eine Zusammenkunft (coetum), eine Versammlung (concionem), ein Zusammentreffen (synodum), eine Festversammlung (conventum), wie das Volk in dem Vorhofe (atrium) zusammenzukommen pflegt, um GOttes Wort zu hören, oder um in öffentlichen Angelegenheiten etwas zu ordnen, oder eine Wahl zu halten, wie die Volksversammlungen (comitia) der Heiden sind. So kam das Volk GOttes, nachdem sie das ganze Passahfest hindurch Gotte mit Opfern gedient und gegessen hatten an der Stätte, die GOtt erwählt hatte, am siebenten Tage an einem Orte zusammen, als, in dem Vorhofe, und daselbst wurde das Gesetz GOttes und die Gebote gelesen, welche sie angingen. Darum nennt er es eine Versammlung GOttes, des HErrn, daß sie zusammenkommen sollten, das Gesetz GOttes zu hören.

„Das Fest der Wochen“ nennt er Pfingsten von den Wochen, die gezählt werden vom Tage des Passah an, deren sieben sind, wie er sagt [V. 9.]: „Sieben Wochen sollst du zählen.“ Dies Fest zeigt er im zweiten Buche Mosis durch diese Umschreibung an [Cap. 23, 16.]: „Das Fest der ersten Ernte der Früchte, die du auf dem Felde gesäet hast“, hier aber mit einer solchen Umschreibung: „Von dem Tage, wenn man anfähet mit der Sichel in der Saat, sollst du anfangen zu zählen.“ Dies verstehe ich so: Man fängt an, die Wochen zu zählen von der Zeit des Passah, wo das Gras schon so weit emporgewachsen ist, daß man seine üppig emporwachsenden Halme (comam) schon abmähen kann. Oder wenn dies hart erscheinen sollte, so muß das Gebot des Zählens auf die Pfingstzeit bezogen werden, auf diese Weise: Du sollst anfangen zu zählen von dem Tage, das heißt, wenn es um die Zeit der Pfingsten sein wird,\footnote{Es scheint uns, daß statt fueris in den Ausgaben fuerit gelesen werden sollte. Danach haben wir übersetzt.} wenn das üppig machende Gras zuerst abgeschnitten zu werden pflegt, oder die ersten [Feld-]Früchte gemäht werden, sollst du rückwärts zählen, bis daß du durch Rechnung (computando) bis zum Passah sieben Wochen voll hast. Und dies gefällt mir besser, weil er gebietet, die Wochen zu zählen, von denen er andeutet, daß sie schon erfüllt seien, zugleich aber auch, weil es der Brauch der Hebräer ist, rückwärts zu lesen, zu schreiben und zu handeln (agere).

Dann wird es schön zusammenstimmen mit 2 Mos. 23, 16.: „Wenn du die ersten Früchte erntest“, weil um Pfingsten, besonders in jenen warmen Gegenden, die ersten Früchte da zu sein (inveniri) pflegen und geerntet werden, wie ich gesagt habe, so daß es dieselbe Zeit ist, wo die ersten Früchte geerntet werden, und wo die Sichel in die Saat schneidet (mittitur). Daher versteht er auch unter der „freiwilligen Gabe“, welche an dieser Stelle im Hebräischen Missa genannt wird, die Erstlinge der Früchte, welche sie den Priestern bringen sollten, wovon du im dritten Buche Mosis, Cap. 2 am Ende [V. 12. ff.], lesen kannst. Nachdem dies erklärt worden ist, ist alles in diesem Capitel leicht.

\quellseite{758}
Es werden solche Feste und äußerlichen Gebräuche zum Gottesdienst eingesetzt, wegen des großen Volkes, damit es, mit denselben beschäftigt, nicht durch sein eigenes Vornehmen Feste erfinde. Denn die Natur ist dazu geneigt, ja, sie hat keine Ruhe, bis daß sie Ceremonien anrichte und Gottesdienste ordne, darum ist es notwendig, daß man ihr zuvorkomme und sie behalte bei GOttes Wort, damit wir gewiß seien, daß wir mit göttlichen Dingen umgehen und Gotte gefallen. Sodann ist auch dieser Grund, daß das Volk wenigstens zwei- oder dreimal im Jahre zusammenkomme und das Gesetz GOttes höre und lerne, und in der Einigkeit des Glaubens und Lebens behalten werde.

Es sind aber drei Dinge, von denen er wollte, daß man bei diesen drei Festen daran gedächte: an Feste des Paßah sollte man gedenken des Auszugs aus Egypten; dies war das erste Fest von allen und das höchste im Anfang des Jahres, wegen des ersten und höchsten Wunders, welches ihnen erwiesen worden war, nämlich wegen der Befreiung aus Egypten. Aber am Fest der Pfingsten sollte man gedenken, daß das Gesetz am Berge Sinai empfangen worden war, da ja im dritten Monate nach dem Auszuge das Gesetz gegeben worden ist, ungefähr am fünfzigsten Tage nach dem Paßah. Am Fest der Laubhütten sollte man das Gedächtnis halten aller leiblichen Wohlthaten, die ihnen die ganzen vierzig Jahre hindurch in der Wüste erwiesen worden waren. So siehst du, daß die Feste geordnet werden, nicht der Üppigkeit, des Müßiggangs und des Vergnügens halben, sondern wegen unseres Heils und der Ehre GOttes, damit das Wort GOttes gehört werde, und man seiner Wohlthaten gedenke, auf daß wir im Glauben und der Liebe zu GOtt unterwiesen, genährt und erhalten werden.

Alle diese Feste, wiewohl noch viele andere waren, feiern wir durch die Deutung (allegoria) des Geistes in Einem Feste. Denn wir halten täglich Paßah, indem wir Christum, das Lamm GOttes, essen und opfern, das heißt, wie es Paullus 1 Cor. 11,\footnote{In den Ausgaben: „1 Cor. 10.“} 26. auslegt, verkündigen und glauben, daß er für uns geopfert sei. So haben wir alle Tage Pfingsten, indem wir das neue Gesetz, den Geist, im Herzen empfangen durch den Dienst des Wortes. Täglich feiern wir das Fest der Laubhütten, da wir lehren und dafürhalten, daß wir Fremdlinge sind in dieser Welt und weilen in den Hütten unseres Leibes, die nur eine kurze Zeit währen, indem wir wissen, daß unser πολίτευμα, das heißt, unsere Bürgerschaft oder Wandel und Wohnung im Himmel ist, wo wir unterdessen im Glauben sind und leben, und hier täglich sterben, bis daß unsere Hütte abgelegt wird etc. So sind wir fröhlich an diesen Festen vor dem HErrn in dem Süßteig der Lauterkeit und der Wahrheit, das heißt, in der Nahrung des reinen Evangelii und in dem Leben des Geistes, ohne Satzungen und menschliche Gebrechen, und opfern unsere Erstlinge des Lobes und der Danksagung.

Denn auch das, daß er gebietet, daß jeder, der männlich ist, nicht leer vor ihm erscheinen solle, war zu der Zeit deshalb verordnet, damit die Priester ernährt würden, aber zu unserer Zeit ist es durch die gottlosen Priester, um ihren Bauch zu weiden, auf die Opfergaben (offertoriis) des gemeinen Volks gezwungen worden, während es doch in Wahrheit bedeutet, daß man Lob und Dank opfere, nun aber nicht mehr vor dem HErrn auf dem Gnadenstuhl und der Bundeslade oder an irgend einer gewissen Stätte, sondern im Geist und in der Wahrheit. Das ist es, daß man nicht an einem jeglichen Orte (wie Moses hier wiederum einschärft), das Paßah opfere, sondern an der Stätte, die GOtt erwählt hätte. Denn wenn man die Meßpfaffen (sacrificulos) mit Opfergaben oder Beiträgen (pensionibus) weidet, so ist das ebenso viel, als wenn man Huren und Hurenwirthe, ja, Wölfe und wilde Thiere nährte, damit sie uns mit ihren verderblichen Lehren verschlingen. Aber den Dienern des Worts gebührt ihre Ehre, wie wir oben [Cap. 14 gegen das Ende] gesagt haben. Daher erscheint auch jetzt keiner, der männlich ist, leer vor dem Herrn, der Gläubige im Geiste dankt und Lob opfert. Das Weib, das ist, das Fleisch und die noch übrigen Schwachheiten der Sünde opfern nichts, weil in dem Munde des Sünders kein schönes Lob ist.

„Brod des Elends“ nennt er das ungesäuerte Brod, weil sie in Furcht ausgezogen sind, und er verbietet den Sauerteig. Was Sauerteig sei, lehrt Christus selbst genugsam Matth. 16, 6., wo er befiehlt, daß sie sich vor dem Sauerteig der Pharisäer hüten sollen. Darum ist das ungesäuerte Brod das gesunde und lautere Wort
\quellseite{759}
des Glaubens. Es wird ein Brod „des Elends“ genannt von der Zucht bei dem eiligen Auszuge, weil durch das Evangelium der Zorn GOttes vom Himmel offenbart wird über alle Menschen, besonders die Verkehrtheiligen, Röm. 1, 18. und 2, 16., da der Geist die Welt straft um die Sünde, um die Gerechtigkeit, und um das Gericht. Hier erschrickt das Gewissen vor dem Reiche der Sünde und des Zorns, und eilt auszugehen, indem es mit denen in der Apostelgeschichte Cap. 2, 37. spricht: „Lieben Brüder, was sollen wir thun?“ und mit Paulo [Apost. 9, 6.]: „HErr, was willst du, daß ich thun soll?“ Denn die Erkenntniß der Sünde drängt und treibt die Erwählten GOttes, daß sie da ausziehen, wo die Erstgeburt der Egypter, das heißt, die Gerechtigkeit der Werke mit dem

Vertrauen darauf und ihrem hoffärtigen Ruhme geschlagen worden ist. Von dem Fleisch des Lammes soll nichts übrig bleiben bis morgen, wo aber etwas übrigbleibt, soll es mit Feuer verbrannt werden [2 Mos. 12, 10.]. Dies wird wider die Juden gesagt, welche, nachdem das wahre Lamm, Christus, geopfert worden ist, ihn noch in der Zukunft erwarten, und es aufschieben und zögern, an den zu glauben, der gekommen ist, und die Werke des Gesetzes fahren zu lassen. Aber die Göttlichen verzehren das ganze [Lamm] durch das Feuer des Heiligen Geistes, das ist, durch das Wort und den Glauben, indem sie einfach glauben, daß alles erfüllt ist durch das wahre und ewige Lamm, nachdem der neue Morgen der Auferstehung des HErrn angebrochen ist.

\section{Das siebenzehnte Capitel}

[Cap. 16, 18.] Richter und Amtleute sollst du dir setzen in allen deinen Thoren.

Nachdem das erste Gebot von dem Glauben an GOtt ausgelegt ist, womit sowohl der inwendige als auch der äußerliche Gottesdienst gelehrt worden ist, kommt Moses jetzt auf das zweite Gebot, nämlich das Gebot der Liebe. Aber er läßt diesem eine Abhandlung vorhergehen über obrigkeitliche Personen, über das Recht des Schwertes und das Gerichtsverfahren in äußerlichen Dingen, und zwar in einer richtigen und schönen Ordnung. Denn (wie wir oben [Cap. 1, 9–12.] gesagt haben) wenn nicht Leute da sind, welche die Gesetze verwalten und vollstrecken, das Schwert führen und gebrauchen, so ist jede Gesetzgebung vergeblich, wenn sie auch noch so heilig ist, da das Schwert die Kraft und die Wirksamkeit und das eigentliche (ipsa) Leben des Gesetzes ist, welches die Bösen im Zaume hält und die Guten schützt. Hier werden wir also in einem kurzen Inbegriff darüber belehrt, daß, wo ein Gesetz gegeben wird, da auch zugleich das Schwert eingesetzt werden muß, damit das Gesetz lehre, was man thun und lassen müsse, das Schwert aber dieses vollstrecke und diejenigen strafe, welche etwas [Verbotenes] thun oder [etwas, das geboten ist,] unterlassen,

indem es die Gottlosen mit Gewalt zum Thun und Lassen zwingt. Denn ohne das Schwert ist das Gesetz unnütz und nur eine Ursache, die Uebelthaten zu vermehren. Wiederum, das Schwert ohne das Gesetz ist eine Tyrannei und ein viehisches Wüthen. Aber beides, nämlich das Gesetz und das Schwert, mit einander verbunden, ist eine schöne und feste Ordnung (politia).

Aus diesem Capitel wirst du zugleich lernen, daß öffentliche Verbrechen nicht durch den großen Haufen oder irgend welche Privatperson gestraft werden sollen, sondern durch das öffentliche Schwert von der Obrigkeit. Denn hier siehst du, daß GOtt selbst zur Handhabung seines eigenen Gesetzes Richter und Amtleute eingesetzt habe, und er fügt einige Fälle hinzu, welche (wie man sagt) vor die Schranken gehören (in terminis), durch welche er das Gerichtsverfahren beschreibt, welches bei Bestrafung von Uebelthätern beobachtet werden soll, damit niemandem Unrecht geschehe, nämlich, daß die Sache sorgfältig erforscht, gewiß erkannt und durch Zeugen überwiesen werde, und so nach dem Gesetze bestraft werde. Dies sage ich wider diese unsere neuen Propheten, die Bilderstürmer, welche, da sie neulich aus dem Himmel gefallen
\quellseite{760}
sind, alle Obrigkeit verachten, und den Pöbel aufreizen wider die, welche das Schwert aus göttlicher Gewalt führen, und das Schwert an sich reißen, und ohne Urtheil, ohne Zeugen und ohne alle Form Rechtens wüthen.

Das ist es, daß er in diesem Capitel gebietet, daß die Sachen vor die Richter gebracht und mit Zeugen bewiesen werden sollen; desgleichen, wie ein König beschaffen sein solle. So gebietet er 2 Mos. 22,\footnote{In den Ausgaben: Ex. 23.} 9., daß die Sachen vor die Götter (das ist, vor die Richter) gebracht werden sollen, und daß sie da entweder durch einen Vergleich oder durch ein Urtheil entschieden und beigelegt werden sollen. So lehrt er anderswo [5 Mos. 21, 20. f.] die Eltern, daß sie ihre widersetzlichen (rebelles) Söhne selbst vor den Richtern anklagen sollen, und daß dieselben dann gesteinigt werden sollen, und es wird [5 Mos. 22, 13. ff.] geboten, daß, wenn eine Dirne von ihrem Bräutigam nicht Jungfrau erfunden worden ist, sie gleicherweise vor Gericht gestellt werden soll, so daß du überall sehen kannst, daß durch göttliche Gewalt das Schwert und das öffentliche Gericht eingesetzt werde, um die Verbrechen zu strafen. Hiedurch wird jenen wüthenden Propheten geantwortet, daß sie auf dieses Wort GOttes dringen, in welchem er das Volk in der zweiten Person anredet, indem er sagt: Du sollst steinigen, du sollst tödten, so oder so sollst du thun, du sollst den Bösen von dir thun etc. Denn jene rühmen, daß durch dies Wort dem Pöbel ohne Unterschied das Recht des Schwertes gegeben sei, während er dies doch redet zu dem Volke, welches mit rechtmäßiger Obrigkeit und berathenden (praesidentibus) Priestern versehen und geordnet war.

Daher redet er hier auch nicht von den höchsten obrigkeitlichen Personen, als da sind die Obersten (principes), Priester und Häupter des Volks, die schon vorher durch göttliche Gewalt dem Volke von Mose gegeben worden waren, als Aaron, Josua, die zwölf Fürsten [4 Mos. 34, 18.] und die siebenzig Aeltesten [4 Mos. 11, 16.], sondern von den Richtern und Amtleuten, welche er selbst nicht einsetzt, sondern dem Volke, welches schon geordnet war, einzusetzen befiehlt. Denn die Richter sind eigentlich für die Rechtshändel da, wiewohl im Lauf der Dinge dieses Amt bisweilen wieder an die Obersten und Priester kam, wie das Buch der Richter und [die Bücher] Samuelis von den Richtern sagen. Die Amtleute aber sind eigentlich die Vollstrecker des Recht[s] und des Schwertes, wie bei uns die Scharfrichter, die Gerichtsdiener und die Profoße; ehemals wurden die angesehensten und besten Leute dazu erwählt, jetzt aber verwaltet nur irgend ein ganz verachteter und geringer Mensch dieses Amt.

Denn wenn es für einen Richter ehrbar ist, daß er mit dem Worte den Schuldigen verurtheile und verdamme, wie kann es denn eine Unehre sein, daß man denselben tödte (percutere) und das Wort des Richters im Werke vollstrecke? Wenn es unehrbar ist, einen Gottlosen mit dem Schwerte zu tödten, so würde es noch viel unehrbarer sein, ihn durch das Gericht als einen solchen zu verurtheilen und dem Schwert zu überliefern, da der Scharfrichter (praefectus) nicht richten würde, wenn es der Richter nicht geböte. Denn wir lesen, daß Salomo [1 Kön. 2, 25. ff.] dem sehr angesehenen Manne Benaja,\footnote{So die Wittenberger und die Jenaer Ausgabe. Statt: Banaiae, filio Jojadae bietet die Erlanger: Jojada Banaia.} dem Sohne Jojada’s, dies Amt übertrug, als er befahl, den Simei und den Adonia und den Joab zu tödten; und Saul trug seinem Doeg auf [1 Sam. 22, 18.], die Priester zu tödten, und David seinen Männern [2 Sam. 1, 15.], daß sie den tödten sollten, der den Saul getödtet hatte. Aber auch bei den Römern wurde dieses Amt großen und geehrten Männern übertragen, nicht weniger, als das Amt eines Richters.

Denn diese zwei Personen, den Richter und den Amtmann, verlangt die Verwaltung des Schwerts und hat sie vonnöthen. Und im Hebräischen sind durch einen schönen Anklang (allusione) \hebraeisch{שפטים} und \hebraeisch{שרים} ähnlich gemacht. Ich habe es, weil ich nichts Anderes hatte, nicht gar genau, durch „Richter und Amtleute“ übersetzt. Unter \hebraeisch{שרים} oder Amtleuten will Moses gerade diejenigen verstanden wissen, welche des Richters Diener sind, die Vollstrecker [der Todesstrafe], die Vollzieher [der Urtheile] und öffentlichen Verwalter, daher sie auch unter lateinischer Text bisweilen durch praecones [Gerichtsdiener] übersetzt, und 2 Mos. 1, 11. „Frohndvögte“ (exatores) und Werkmeister (praefectos operum) daraus macht. Christus nennt denselben schlechtweg
\quellseite{761}
einen „Diener“, Matth. 5, 25.: „Auf daß dich der Richter nicht überantworte dem Diener“ etc.

Er schreibt aber diesen Richtern und Amtleuten die Regel vor, daß sie gerecht richten sollen, das heißt, nach dem Gesetze GOttes, nicht nach ihrem Sinne. Darnach verbietet er krumme Wege (affectus), daß sie das Gesetz nicht verlassen und sich leiten und bewegen lassen durch Ansehen der Person und Geschenke. Denn diese beiden Stücke pflegen alle Gerichte zu verderben und zu beugen, daher fügt er hier auch einen Sinnspruch (gnomen) ein, indem er sagt: „Die Geschenke machen die Weisen blind und verkehren die Sachen der Gerechten.“ Das „Ansehen der Person“ begreift in sich die Furcht vor Großen, Gewaltigen und Reichen, die Liebe zu den Verwandten, die Gunst der Freunde, die Verachtung der Geringen, das Erbarmen gegen Unglückliche, die eigene Gefahr des Lebens, des guten Namens und aller Dinge. Die „Geschenke“ aber begreifen in sich den Gewinn, den Vortheil, die Hoffnung und den ganzen unerträglichen und unermesslichen Schlund des Geizes. Daher gibt Jethro 2 Mos. 18, 21. dem Moses den Rath, daß er Männer erwähle, die ohne Geiz seien, das heißt, solche seltenen Vögel, die einem schwarzen Schwan sehr ähnlich sind.

Du siehst also, was für ein Gemüth der haben muß, der das Amt eines Richters und des Schwertes ausüben soll, nämlich er muß alle Gemüthsbewegungen überwinden, die Furcht, die Liebe, die Gunst, das Erbarmen, den Geiz, die Hoffnung, den guten Leumund, das Leben und den Tod, und muß auf das einfältigste die schlichte Wahrheit und das rechte Gericht lieb haben, denn (so sagt Josaphat, der König von Juda [2 Chron. 19, 6.]) das Gericht ist des HErrn. Denn nothwendiger Weise ladet ein Richter den Haß vieler Leute auf sich, er muß durch Feindschaft (odio) in Gefahr stehen, und angefochten werden durch Erbarmen und Gemüthssucht, und, wenn er nicht allein auf GOtt sieht, kann er diese Anstöße nicht überwinden und dawider bestehen. Wie wird er aber allein auf GOtt sehen können, wenn sein Herz nicht durch den Glauben sicher befestigt ist?

Siehe deß ein Exempel an den besten Juristen, welche von Billigkeit sehr schön reden, wenn sie öffentlich lehren oder schreiben, aber sich ganz anders halten, wenn sie einzelnen Personen (privatim) Rath geben müssen, und in Widerspruch gegen sich selbst gerathen, nämlich, daß hier das Ansehen der Person oder das Erbarmen mit dem Unglück (um es aufs beste auszulegen) sie von der rechten Bahn abwendet, so daß sie sich bemühen, gleichsam zu heilen und dem Hülfslosen beizustehen, was sie zu der Zeit nicht thaten, als sie öffentlich lehrten. Und so sehen sie vor leidenschaftlicher Erregung (affectu) nicht, daß sie von der Billigkeit, die sie öffentlich gelehrt haben, abweichen. Etwas so Schwieriges, Seltenes und Hohes ist es um ein einfältiges und richtiges Auge bei einem Richter.

Dasselbe kann man an den größten Theologen sehen, als da sind Augustinus, Bernhard, und an denen, die älter sind als diese, Cyprianus, Tertullianus und ähnlichen, welche, wenn sie öffentlich die Schrift behandeln, sie lauter und rein auslegen; aber wenn ihnen Fragen vorgelegt werden, so bleiben sie selten auf der rechten Bahn, daß sie der Schrift nicht etwas nehmen und es der Sache oder der Person beilegen und die Worte GOttes mit Gewalt verdrehen sollten. Siehe sie hierauf an, wo sie gegen den Arius, Hieronymus wider den Jovinianus, Augustinus wider die Manichäer, Bernhard wider den freien Willen schreiben, und du wirst sehen, daß ich die Wahrheit sage. So gar hat keiner der Heiligen das Fleisch ganz ausgezogen oder seine Gemüthsbewegungen abgelegt, daß gar keine Hoffnung übrig ist, es sei denn, daß sowohl der Richter als der Lehrer in der Furcht GOttes handeln, und sich selbst immer verdächtig halten, sie möchten etwa das Gesetz GOttes und seine Worte nicht recht handeln. Denn Sicherheit darf in diesen Dingen nicht statthaben, wo GOttes Sache gehandelt wird, die überaus ernst ist.

\subsection{[Cap. 16, 21. 22.] Du sollst keinen Hain von Bäumen pflanzen etc}

Warum mengt denn Moses dies von dem äußerlichen Gottesdienste ein, da er schon angefangen hatte, von den Werken gegen den Nächsten zu reden? Ich antworte: Dies scheint er in der Weise einer Wiederholung zu thun, daß er in diesem gesetzlichen Volke auch das der Strafe des Schwertes unterwerfe, wo etwas wider den äußerlichen Gottesdienst gesündigt werden sollte. Dies wird dadurch bewiesen, daß er nichts vom Glauben oder von der Furcht
\quellseite{762}
sagt, sondern von dem Werke, daß eine Säule aufgerichtet, ein Hain oder ein Baum bei dem Altar des HErrn gepflanzt wird, und ein Opfer dargebracht wird, das einen Theil hat [Cap. 17, 1.]. Durch diese zwei Stücke zeigt er, als durch ein Beispiel, alle Sünde im äußeren Gottesdienste an. Sodann fügt er auch ein Exempel hinzu, daß ein Mann und ein Weib gesteinigt werden sollen, die überführt worden sind, daß sie Abgötterei getrieben haben. Durch dies Exempel wendet er sicherlich das Schwert, welches schon eingesetzt ist, auf die Sünde an, welche wider den äußerlichen Dienst GOttes begangen worden ist, und überweiset sie dem Richter und dem Diener zur Strafe. Denn das behält er sich selbst zu strafen vor, was im Herzen durch den Unglauben gesündigt wird, wie er Cap. 18, 19. von dem künftigen Propheten sagt. Siehe also, daß hier von GOtt das öffentliche Gerichtsverfahren (forum judicii) eingesetzt ist und ein Exempel desselben gegeben wird. Das erste ist, wenn ein Uebel geschieht vor den Augen des HErrn [Cap. 17, 2. ff.]. Das Uebel aber vor den Augen des HErrn, erklärt er, geschehe dann, wenn jemand mit dem Werke den Bund GOttes übertrete, indem er hingeht und anderen Göttern dient, das GOtt nicht befohlen hat. Denn er redet (wie ich gesagt habe) von der Sünde im äußerlichen Dienste GOttes, um mit einem Exempel zu zeigen, wie bei allen Verbrechen das Schwert geführt werden und das Recht vollstreckt werden solle. Er sagt aber mit einem bezeichnenden Worte „Uebels vor den Augen des HErrn“, denn es gibt vor den Augen der Menschen nichts Schöneres, Besseres und Heiligeres, als Abgötterei und Gottlosigkeit, welche durch Heuchelei und den Schein der Göttlichkeit geschmückt ist. Es ist also vor GOtt das ein böses Werk, welches ohne Glauben und Wort aus unserm eigenen Antrieb vorgenommen wird, wie gut und trefflich geschmückt (speciosum) es auch erscheinen möge. Denn das ist seine Meinung (vult), da er sagt, daß sie seinen Bund übergehen, und thun, was er nicht geboten hat, so daß du gewiß wissen kannst, wenn jemand etwas thut in dem Dienste GOttes, was er nicht geboten hat, das sei böse vor GOtt, wie er hier klar sagt, weil es nothwendiger Weise in der Zuversicht geschieht, als ob sie etwas thäten, was Gotte angenehm ist, und daß sie dadurch gerecht würden; sonst würden sie das nicht thun, was gänzlich wider den Glauben, wider seinen Bund und sein Wort ist.

Zweitens. Es ist nicht genug, daß ein böses Werk geschehen ist, sondern es muß auch angesagt und durch Gericht oder Erzählung bekannt werden, wie er hier sagt [V. 4.]: „Und es wird dir angesagt.“ Drittens. Auch dies ist nicht genugsam, sondern es soll sorgfältig nachgeforscht werden, ob es also sei, wie berichtet wird, damit man nicht einem jeglichen Angeber freventlich glaube. Diese Nachforschung aber kann nicht eine andere sein, als durch Zeugen, welche es gesehen haben und wissen, so daß ein Zeuge nicht genugsam ist, sondern „auf zweier oder dreier Zeugen Munde (sagt er V. 6.) soll alle Sache bestehen“. Und er fügt hinzu: „Auf Eines Zeugen Mund soll niemand sterben.“ Viertens. Endlich sollen diejenigen, welche überführt sind, hinausgeführt und gesteinigt werden. So siehst du, daß im öffentlichen Gerichte und bei der Auflegung öffentlicher Strafe die Verbrechen so öffentlich sein sollen, als die Strafe ist.

Aus diesem Exempel kannst du lernen, wie viele Vergehen zu tragen und zu verbergen seien, welche entweder du allein, oder andere mit dir wissen, oder welche das Gericht überall verbreitet hat. Zugleich siehst du, wie sorgfältig GOtt den Angebern und Verleumdern das Maul gestopft hat, auch denen, die ihnen zuhören, die Ohren, daß auch die Strenge des göttlichen Gesetzes hier nicht allein die wahrhaft Unschuldigen in Schutz nimmt, sondern auch die Schuldigen, so lange sie noch verborgen sind. Verborgen nenne ich sie aber, weil sie der Obrigkeit und dem Richter nicht angezeigt worden sind als solche, die mit rechten Zeugen überführt werden sollten. Und hier ist offenbar, ein wie teuflischer Pfuhl die Welt sei, wo der gute Name des Nächsten so leicht verletzt wird, und Vergehen, die entweder erlogen, oder allein durch das Gerücht und Argwohn aufgelesen sind, öffentlich weitergetragen werden, so daß das Gesetz GOttes an einer anderen Stelle wiederum befohlen hat, daß diejenigen nicht weniger strenge bestraft werden sollen, welche fremde Vergehen anklagen, die sie nicht beweisen können. Deshalb, wenn du allein jemanden sündigen siehst, so schweige; wenn du davon hörst, so halte an dich, bis daß die Sache dem Richter mit gewissen Zeugen bewiesen werden kann.

\quellseite{763}
Es hat aber Moses das Beispiel des Gerichtes lieber hergenommen von einem Vergehen im Dienste gegen GOtt, als von einem Wort oder Ehebruch oder einem anderen Verbrechen, nämlich weil dies Exempel gewiß und zuverlässig ist, und durch kein Erbarmen verziehen oder zugedeckt werden darf. Denn, gegen den Dienst GOttes sündigen, das heißt gegen den Glauben und das Wort sündigen. Hier soll man keine Rücksicht nehmen auf Liebe oder Barmherzigkeit, denn hierdurch beleidigt und verliert man GOtt, nimmt dem Worte, welches der Führer, das Licht, die Lehre und die Regel des ganzen Lebens und aller Werke ist; wenn dies verloren ist, kann kein Werk [recht] angestellt, kein Leben [recht] geordnet werden. Aber die Sünde wider die Liebe, wenn nur das Wort und die Lehre unangetastet bleibt, ist nur ein verlorenes Werk, welches nach der Regel des Wortes wiederhergestellt und wieder gutgemacht werden kann. So gar unermeßlich ist der Abstand zwischen der Sünde gegen den Glauben und das Wort, und der Sünde gegen die Liebe und das Werk. Denn die Liebe verträgt alles, duldet alles [1 Cor. 13, 7.], aber der Glaube verträgt nichts, und das Wort duldet nichts, sondern das Wort muß vollkommen rein sein, die Lehre muß allezeit durchaus gesund sein, so daß sie das Ziel sei für das Leben und für die rechte Richtung der Werke. Die Liebe aber kann schwach und unrein sein, und muß täglich vermehrt und vervollkommnet werden. Daher tödet Moses, der allerfanatümthigste Mann auf der Erde, dennoch ohne Erbarmen dreitausend Männer, welche das Kalb angebetet hatten, 2 Mos. 32,\footnote{In den Ausgaben: Exodi 34.} 28., weil sie wider das Wort, des Lebens Licht und Führer, gesündigt hatten.

\begin{schriftzitat}
V. 8–13. Wenn eine Sache vor Gericht dir zu schwer sein wird etc.
\end{schriftzitat}

Auch diesen Text haben die Papisten mit wunderlicher Mühe auf ihren Gößen gezogen, um das Papstthum zu bestätigen, darum ist es vonnöthen, daß wir ihn etwas sorgfältiger betrachten. Erstens handelt Moses hier nicht von dem Worte oder von der Lehre, oder (wie jene sagen) von den Fragen des Glaubens, von denen sie wollen, daß man sie an den Papst bringe, sondern von dem Urtheil über öffentliche und weltliche Verbrechen, wie die Worte klar sagen, von einer zweifelhaften Sache „zwischen Blut und Blut“ (das heißt, von Todtschlag), „zwischen Handel und Handel“ (das heißt, von Diebstahl, Raub, Streit, Scheltworten), „zwischen Schaden und Schaden“ (das heißt, von Schlagen, Wunden oder anderer gewalttätiger Verletzung des Leibes), wo Hieronymus, seinen Träumen folgend, „zwischen Aussatz und Aussatz“ übersetzt hat, wodurch er unseren Papisten Anlaß gegeben hat, daß sie meinten, Moses rede auch von den heiligen Geheimnissen des Glaubens, nicht bloß von weltlichen Sachen. Jetzt, da der Papist ein Diener des Geistes sein soll, und nicht des Buchstabens, und diese weltlichen Sachen fahren lassen, geht ihn dies Gesetz durchaus nicht an, wie auch kein anderes Gesetz Moses; oder wenn es jemanden anginge, so ginge es vielmehr den Kaiser oder einen König oder einen Fürsten an, so daß nach dem Exempel dieses Gesetzes die Untergebenen von dem Hofe des Oberherrn in zweifelhaften Sachen das Urtheil holen sollten.

Sodann redet dieses Gesetz von dem Orte, den GOtt erwählt hat, nämlich wo die Hütte war, das äußerliche Zeichen des gegenwärtigen GOttes. Aber im neuen Testamente ist keine solche Stätte. Denn [Joh. 4, 21.] weder auf diesem Berge, noch zu Jerusalem, sondern [R. 23.] „im Geist und in der Wahrheit“ wird GOtt angebetet; sie werden auch nicht sagen [Matth. 24, 23.]: „Siehe, hier ist Christus oder da“, denn das war dem alten und gesetzlichen Volke eigen. Daher sollten die Papisten erst beweisen, daß Rom der von GOtt erwählte Ort sei. Wann werden sie das thun? Wo ist das Wort? Wo ist das Zeichen? Es ist Eine Stätte, die allen Christen gemeinsam ist, das ist der Geist und die Wahrheit. So erklärt es Christus selbst; im Geist und in der Wahrheit (sagt er) wird GOtt angebetet und verehrt. Und gerade an der Stätte müssen die Fragen über den Glauben aufgelöst werden. An dieser Stätte ist der Papist nicht mehr denn die Magd in einer Mühle, und seine Entscheidung gilt an dieser Stätte nicht mehr als die eines Bauern, der auf dem Acker ist. Denn der Glaube gehört allen zu, und der geistliche Mensch richtet alles. Wenn nun an irgend einem Orte vollkommenere Christen sein sollten, so müssen ihnen freilich die unvollkommenen weichen, nicht um der Stätte willen, sondern wegen der Vollkommenheit des Geistes, wie in
\quellseite{764}
der Apostelgeschichte [Cap. 15, 2.] die von Antiochia ein Urtheil von den Aposteln und Aeltesten begehrten, nicht um des Ortes willen, als ob GOtt ihn erwählt hätte, da auch viele Aelteste daselbst nicht die rechte Meinung hatten, Apost. 15, 5., und die Apostel nicht dort blieben als an einem beständigen Orte, sondern um des vollkommeneren Geistes der Apostel willen.

Und, gesetzt den Fall, daß dies alles nichts beweisen sollte, und daß dies Gesetz Mosis von den Fragen, die den Glauben betreffen, und von Rom verstanden werden müßte (was unmöglich ist), so macht er doch mit Einem Worte ihr ganzes Vornehmen zu Schanden, da er hier sagt [V. 1.]: „Nach dem Gesetz, das sie dich lehren“, nämlich weil er gewollt hat, daß selbst diese Priester der Werke und des Gesetzes nicht auf ihren eigenen Geist sich verlassen und darnach ein Urtheil abgeben, sondern daß sie alles lehren sollten nach dem Gesetz, welches er vorgeschrieben hatte. Wie viel mehr müssen daher die Priester des Glaubens und des Evangelii nicht aus ihrem eigenen Geiste, sondern nach dem gewissen Worte GOttes alles entscheiden und lehren!! Aber dies haben weder der Pabst noch die Concilia jemals gethan, noch thun sie es jetzt, sondern sie beschließen alles aus der Fülle ihres vermessenen Geistes. Sie beschließen aber meistens nichts als überaus thörichte Dinge, als von Ceremonien, Pfründen, Würden und ähnlichen Narrenwerken.

Es stellt also Moses dies Gesetz für die Richter und Diener, von denen er in diesem Capitel handelt, daß die kleinen Städte sich in zweifelhaften Sachen an die Hauptstadt wenden sollten, und die Unerfahrneren an die Vollkommneren, da es gewiß war, daß auch GOtt dort die Sache lenkte, und immer solche Priester dort wären, welche das Gesetz GOttes wüßten, was GOtt keinem anderen Orte gewährt hat. Denn es war ein bürgerliches Volk, darum hat er es auch mit einer bürgerlichen Stätte und einem leiblichen Priesterthum versehen, in welchem der Geist gewiß bleiben sollte, gleichwie er uns Christum im Glauben seines Fleisches gewißlich versehen hat\footnote{Wittenberger: providet statt: providit.} zu seiner Rechten, das ist, als den, der überall gegenwärtig ist und alles erfüllt [Eph. 1, 23.]. Daher kommt es, daß der HErr vom Gnadenstuhl Antwort gab durch das Licht, das heißt, durch das Amtsschildlein (rationale [2 Mos. 28, 15.]) oder den Priester, wie Saul sich 1 Sam. 28,\footnote{In den Ausgaben: 1 Reg. penult.} befragt, daß GOtt ihm keine Antwort gegeben habe, weder durch die Propheten noch durch das Licht.

\begin{schriftzitat}
V. 14. ff. Ich will einen König über mich setzen.
\end{schriftzitat}

Weil er angefangen hatte von den obrigkeitlichen Personen, so beendigt er diese Abhandlung auch in Bezug auf die Könige und das Königreich, indem er sagt: denn dieses werde geschehen, daß die Kinder Israel Könige haben würden. Dies war ehemals dem Abraham und Isaak verheißen, da der HErr 1 Mos. 17, 6. sprach: „Es sollen Könige von dir kommen“ etc. Es kommt daher Moses hier auch dem zuvor, daß sie nicht nach ihrem Sinne und ohne das Wort GOttes einen König erwählen oder einsetzen, sondern, damit sie gewiß seien, daß sie aus GOttes Huld und Willen einen König haben, will er, daß sie durch GOttes Macht und durch sein Wort geführt werden, was wiederum etwas ganz Besonderes ist vor allen Völkern auf Erden. Denn wiewohl er alle Reiche verändert und einsetzt, und sie gibt, wem er will, so fügt er doch für dieses Volk insbesondere ein Wort hinzu, darin er ihnen seine Huld anzeigt, damit sie wüßten, daß sie das Reich GOttes hätten. Deshalb werden sie im [ersten] Buche Samuelis darüber gestraft, daß sie schwer gesündigt haben, daß sie in solcher Weise einen König begehrten, daß sie auf ihn als auf einen Menschen und ein menschliches Reich ihr Vertrauen setzen möchten, und nicht einen König in GOtte und nach seinem Worte verlangten, wie er ebendaselbst [Cap. 8, 7.] sagt: „Sie haben mich verworfen, daß ich nicht König über sie sein soll.“

Darnach, damit auch der König selbst ein gottgemäßer sein möchte, so unterweist ihn Moses, daß er allein im Vertrauen auf GOtt regieren solle, indem er wüßte, daß er allein durch sein Wort eingesetzt sei und durch dasselbe erhalten werden müsse. Deshalb befiehlt er, daß er nicht viele Rosse halten solle, das heißt, daß er nicht auf seine Streitmacht (viribus) vertrauen solle; daß er das Volk nicht wieder nach Egypten führen solle, das heißt, daß er nicht auf der Menschen Gunst und Stärke sich verlasse; daß er nicht viel Gold und Silber sammeln solle,
\quellseite{765}
das heißt, daß er nicht auf Reichthum vertraue; daß er nicht viel Weiber nehmen solle, das heißt, daß sein Glaube nicht durch die Dornen erstickt werde, wie Christus die Sorgen und Wollust dieses Lebens nennt, Luc. 8, 14. Denn das ist nicht sein Wille, daß er nicht gut versehen sei mit Rossen und Truppen, da man liest [1 Kön. 4, 26.], daß David und Salomo das Gegentheil gethan haben; so verbietet er auch nicht Bündnisse und Beistand von Menschen, da Salomo den König von Tyrus und Hiram [1 Kön. 5, 12.], David die Könige der Moabiter und der Ammoniter und den König Thoi [2 Sam. 8, 9. ff.] zu Bundesgenossen hatte. So steht auch geschrieben [1 Kön. 10, 23.], daß Salomo eine so große Menge Goldes und Silbers gehabt habe, welche ihm theils von David hinterlassen, theils von ihm selbst zusammengebracht war, als kein anderer König hatte. So hatten auch viele Könige sehr viele Weiber.

Aber darin liegt die Kraft des Gebotes, daß (wie er im Texte sagt) sein Herz nicht abgewandt werde, das heißt, daß er sich das alles nicht sammle und verschaffe, um darauf sein Vertrauen zu setzen. Denn er will, wie ein gläubiges Volk, so auch einen gläubigen König haben, der allein auf GOtt und sein Wort vertraue, darum nimmt er alles weg, worauf die Menschen ihr Vertrauen setzen. Wenn sie sich darauf nicht verlassen würden, sondern es nur gebrauchten im Glauben an GOtt, zufrieden allein mit seinem Worte, dann würde es nicht schaden, wie der beste, der reichste und mächtigste aller Könige rühmt, indem er spricht [Ps. 44, 7.]: „Ich verlasse mich nicht auf meinen Bogen, und mein Schwert kann mir nicht helfen“, und wiederum [Ps. 60, 14.]: „Mit GOtt wollen wir Thaten thun. Er wird meine Feinde untertreten“, und wiederum [Ps. 108, 11. 13.]: „Wer will mich führen in eine feste Stadt? Schaffe uns Beistand in der Noth, denn Menschenhülfe ist kein nütze“, und dergleichen viel. Er sagt nicht, daß er keinen Bogen habe, kein Schwert und keine Leute, sondern er sagt, daß er darauf nicht vertraue.

Daher gebietet hier Moses, daß der König das andere Gesetz (Deuteronomion) lesen solle, damit er GOtt zu aller Zeit fürchten lerne, nämlich, er befiehlt ihm das Wort des Glaubens, welches er in diesem Buche, wie wir gesehen haben, überaus reichlich vorträgt, damit sich sein Herz, durch den Ueberfluß, der ihm zufällt, nicht über seine Brüder erhebe, und er nicht auf seine Kräfte oder Reichthümer vermessen sei, sondern den HErrn fürchte, der nicht Gefallen hat an jemandes Beinen noch an der Stärke des Rosses [Ps. 147, 10.]. Wiederum soll er aber auch nicht verzagt werden, wenn alles verzweifelt steht, und die Feinde die Oberhand haben, das heißt, wie er hier sagt, er soll weder zur Linken noch zur Rechten weichen [V. 20.], sondern richtig und fest einhergehen in Glück und Unglück, in Stärke und Schwäche, in Ehre und Schande, indem er allein treu und stark an GOttes Wort hangt.

Die heimliche Deutung des Richters und des Dieners ist das Amt des Gesetzes, welches durch sein Urtheil die Gottlosen verdammt. Der Tod, der Diener des Gesetzes, tödtet, wie Paulus sagt [1 Cor. 15, 56.]: „Die Kraft der Sünde ist das Gesetz, der Stachel aber des Todes ist die Sünde.“ Der Richter ist rechtschaffen (integer) und muß unbeugsam sein, denn das Gesetz ist heilig, das Gebot ist gut, gerecht und heilig, und dennoch ist es ein Amt des Todes, eingesetzt, daß es tödte, daß es Zorn anrichte, denn dem Gerechten ist kein Gesetz gegeben, gleichwie das bürgerliche Schwert nicht den guten Werken zu fürchten ist, sondern den bösen, Röm. 13, 3. Daß Geschenke genommen und Personen angesehen werden, das ist, daß das Amt des Gesetzes mißbraucht wird, um die Gerechtigkeit der Werke und den Ruhm des Fleisches aufzurichten, nicht aber zur Erkenntniß der Sünde. Denn diese Gerechtigkeit der Werke macht auch die Weisen und Gerechten blind, und betrügt durch ihren Schein und ihre Zuversicht, so daß uns durch dies Sprüchwort [V. 19.] gelehrt wird, wie gefährlich das Gesetz sei, wenn es dazu gelehrt wird, daß man Werke thun soll.

Das ist es, daß er bald, nachdem er von den Richtern geredet hat, verbietet, einen Hain und Bäume und Säulen aufzurichten bei dem Altar des HErrn, das heißt, man soll die Gerechtigkeit und den Ruhm der Werke nicht aufrichten, welche GOtt haßt. Man soll auch nichts opfern, das einen Fehl hat, das heißt, sein Gesetz soll nicht mit Werken erfüllt werden. Denn das heißt das Gesetz und die Absicht des Gesetzes beflecken, welche diese ist, daß es alle Menschen schuldig mache und zur Erkenntniß ihrer Sünde bringe, Röm. 7, 7. Ja, es sollen diejenigen ohne Gnade (rigidissime) getödtet werden, welche andere
\quellseite{766}
Götter anbeten, das heißt, durch das rechte Verständnis und Erkenntniß des Gesetzes soll der Dienst (cultus) und das Vornehmen (studium) von Werken verdammt werden, wodurch man glaubt, GOtt zu dienen, während in Wahrheit ein Götze unseres Herzens, ein falscher Wahn von GOtt, verehrt wird, der, wie das Gesetz lehrt, allein durch den Glauben des Herzens geehrt wird. Daß sie durch Zeugen überführt werden sollen, bedeutet, daß man sie durch das Zeugniß der Schrift überwinden soll, wie Paulus in dem Briefe an den Titus [Cap. 1, 9.] lehrt, daß ein Bischof mächtig sein soll in der heilsamen Lehre, daß er die Widersprecher überführen könne. Ferner sollen diejenigen, welche bei dem Urtheil der Vollkommeneren im Geiste, die nach dem Gesetze GOttes urtheilen, nicht bleiben wollen, wie er hier sagt, getödtet und aus dem Mittel genommen werden, das heißt, einen ketzerischen Menschen soll man meiden und in den Bann thun, und ihn halten wie einen Heiden.

Die Einsetzung eines Königs bedeutet eigentlich Christum, der nach dem Richter, dem Gesetze, kommt. Dieser ist nicht durch unsere Wahl, sondern durch GOttes Einsetzung König in seinem Volke geworden; er ist auch nicht aus einem fremden Volke hergenommen, weil er nicht die Engel an sich genommen hat, sondern den Samen Abrahams, gleich wie mir Fleisch und Blut unter Brüder geworden, als ob er einer von uns wäre [Hebr. 2, 16. 17.]. Und er erhebt sich nicht über seine Brüder, denn, ob er wohl in göttlicher Gestalt war, erniedrigte er sich selbst, und nahm Knechtsgestalt an [Phil. 2, 6. f.]. Darum, wiewohl er Ueberfluß hat an vielen Rossen und großen Reichthümern und der Menge des Volks, das heißt, an aller Fülle des geistlichen Reichthums und geistlicher Kräfte, denn in ihm sind verborgen alle Schätze der Weisheit und der Erkenntniß [Col. 2, 3.]; sodann auch an dem Gold und Silber der Weisheit und Beredsamkeit in der Schrift; und viele Weiber hat, das heißt, durch viele verschiedene Kirchen groß und herrlich ist, so vertraut er dennoch nicht darauf, rühmt sich dessen nicht, sondern allein seines GOttes. Denn er ist der gerechte Herrscher unter den Menschen in der Furcht GOttes, wie David in seinen letzten Worten [2 Sam. 23, 3.] sagt. Er behandelt dies andere Gesetz (Deuteronomion hoc) nicht allein nach dem Buchstaben, sondern auch beständiglich nach dem Geiste, und hat ein Reich, das da lebt in der Furcht GOttes, was die Richter und das Gesetz des Todes nicht vermocht haben.

\section{Das achtzehnte Capitel}

\begin{schriftzitat}
V. 1–14. Die Priester, die Leviten des ganzen Stammes Levi, sollen nicht Theil haben etc.
\end{schriftzitat}

Auch dieses Capitel gehört zum Vorhergehenden und bildet gleichsam den Schluß der Rede von der Versorgung der Diener des Wortes, von denen er im 16. Capitel gesagt hat, nämlich, um das gänzlich zu Ende zu bringen, was er von beiderlei Obrigkeit vorgenommen hatte, sowohl in Bezug auf die heilige, die es mit dem Worte zu thun hat, als auch auf die bürgerliche, welche das Schwerdt führt. Dies thut er in solcher Weise, daß er am Ende dieses Capitels alles beschließt in dem künftigen Propheten, welcher erweckt werden soll, der ein neues Reich und ein neues Predigtamt anrichten werde, und alles abschaffen, was er bisher gesagt hat. Denn auf diesen [Propheten] hat Moses ganz und gar sein Absehen, auf diesen richtet er sein Amt, der das Ende aller Dinge ist. Daher macht er in diesem Capitel einige Zusätze zu dem, was er oben gesagt hat, wie es denn Mosis Weise ist, daß er wiederholt, um etwas hinzuzufügen, und so das Ganze zu vollenden.

Das Erste ist, daß die Priester nicht Theil haben am Lande mit den Stämmen Israel, sondern von den Opfern des HErrn ernährt werden sollen. Er zeigt aber an, was ihnen von den Opfern zukommen solle, nämlich die Schulter, beide Backen und der Wanst von einem jeglichen geopferten Schafe und Ochsen, desgleichen die Erstlinge des Getreides, des Weins, des Oels, der Wolle. Hier geschieht der Zehnten nicht Erwähnung, welche die Leviten empfangen sollten, auch nicht der Erstgeburt, weil
\quellseite{767}
er oben von diesen geredet hat, und hier wiederholt er nur, um hinzuzufügen, was sie von den Opfern haben sollten. Ueber die Schulter und die Brust lies im dritten Buche Mosis die ersten drei Capitel\footnote{per tria prima capitula. Das Wort prima ist wohl durch ein Versehen hineingekommen, denn in den ersten drei Capiteln findet sich davon nichts. Doch wird in drei nach einander folgenden Capiteln der Schulter und der Brust Erwähnung gethan, nämlich Cap. 7, 30–34, Cap. 8, 29. und Cap. 9, 20. f.} und im zweiten Buche Mosis Cap. 29 [V. 27.]. „Bauch“ (ventriculum) nennt er hier entweder den Magen oder den Unterleib, nämlich den untersten Theil des Bauches. Ich bin nicht ganz gewiß, ob die Priester von jedem Opfer, sei es nun ein Sündopfer oder ein Brandopfer oder ein ganzes Opfer, diese Theile nehmen sollten. Ich möchte es gern von allen Opfern behaupten, ausgenommen vom Brandopfer, wo, wie es scheint, das Ganze verbrannt worden ist, wiewohl auch dies mich nicht völlig überzeugt (satis permovet). Doch ich überlasse dies den Leuten, die mehr Muße haben [als ich].

Zweitens, was er hinzufügt von dem fremden Leviten, daß man ihn aufnehmen soll in die Ordnung des Dienstes an dem Orte der Hütte, das gehört auch zu der Versorgung der Leviten, von denen er ebenfalls oben [Cap. 14, 29.] gesagt hat, und hier am Ende wiederholt, um dieses hinzuzufügen, nach seiner Weise, wie ich gesagt habe. Aber man muß glauben, daß dies Gebot hart gewesen sei für diejenigen, welche in der Hütte dienten, daß sie gezwungen wurden, ihre Theile kleiner zu machen und in so viele zu zertheilen, besonders wenn sie geizig und ungläubig waren. Aber GOtt hat dies geboten, damit der Glaube und die Liebe geübt werde, daß sie lernen möchten, gern das Ihre mitzutheilen, in der gewissen Zuversicht, daß GOtt sie ernähren werde. Aber es ist auch kein Zweifel, daß dies Gesetz so gehalten worden ist, wie sonst alle anderen, so daß wegen der Geringsfügigkeit des Antheils wohl kaum wenige oder auch gar keine Leviten versucht haben, sich an den Ort der Hütte zu begeben: doch wenn es etliche hätten thun wollen aus Liebe zum Wort und um das Gesetz zu lernen, so mochten sie hier ihre Versorgung haben.

Drittens, und dies ist der letzte Zusatz zu allem, was oben von der äußeren Verehrung GOttes gesagt worden ist, nämlich, daß sie sich hüten sollten vor den Greueln der Heiden, welche er hier der Ordnung nach anführt, und neun aufzählt, welche alle wider den rechten und reinen Gottesdienst streiten, der durch sein Wort eingesetzt ist. Diese unterscheiden wir nun, wiewohl wir für das Hebräische nicht in allen Stücken Gewähr leisten können, doch so gut als möglich mit glaubwürdigen Muthmaßungen auf diese Weise: Die ersten, welche Moses \hebraeisch{קסם} nennt, Weissager, Wahrsager sagen wir; und dies Wort wird nicht selten in gutem Sinne genommen, wie auch der Name Prophet in bösem Sinne. Weissager sind nun die, welche Zukünftiges vorhersagen, oder die Schrift und die Worte GOttes auslegen, aber nach ihrem Sinne. Die zweiten \hebraeisch{מעוננים}, in Bezug auf welche fast alle übereinstimmen, auch die Juden, daß es Tagewähler sind. Die dritten \hebraeisch{מנחשים} nennen wir: die auf Vogelgeschrei achten (augures); und wiewohl sie davon den Namen haben, daß sie auf die Vögel achten, so bedeutet dies Wort doch auch 4 Mos. 23,\footnote{In den Ausgaben: Num. 24.} 23. das Achtgeben auf andere Dinge, als: „Es ist kein Zauberer in Jakob“, und 1 Mos. 44, 5.: „Ist's nicht das, damit mein Herr zu weissagen pflegt?“ Die vierten \hebraeisch{מכשפים}, diese sind, wie mir meinen, eigentlich Zauberer und Hexen, 2 Mos. 22, 18.: „Die Zauberinnen sollst du nicht leben lassen.“ Die fünften \hebraeisch{חוברים} sind Beschwörer, was jener Psalm [Ps. 58, 6.] bezeugt, welcher sagt: „Welcher nicht hört die Stimme des Beschwörers, der wohl beschwören kann.“ Die sechsten \hebraeisch{אוב}; das ist der Wahrsagergeist (Python), jener Geist, von dem in der Apostelgeschichte [Cap. 16, 16.] erzählt wird, daß er ein solcher ist, der geheime und verborgene Dinge offenbart. Diese nennt man bei uns auf deutsch kluge Männer und kluge Frauen. Die siebenten \hebraeisch{ידעונים}\footnote{In den Ausgaben: jedoeni, doch hat Luther auch hier wohl den Plural gesetzt, wie in den andern Fällen; darauf weist auch der Context hin.}; diese sind die Wahrsager, welche Zeichen auslegen, oder ihre Vermuthungen auf dieselben gründen. Die achten, die die Todten fragen, sind diejenigen, welche die Geister zu Rathe ziehen, welche bei Nacht umgehen, von denen man meint, es seien bisweilen Seelen Abgeschiedener. Die neunten, welche Moses zuletzt setzt, die ihre Söhne oder Töchter durchs Feuer gehen lassen, das heißt, die ihres Samens dem Götzen Moloch geben, wie er anderswo
\quellseite{768}
[3 Mos. 18. 21.] sagt. Dies alles, weil es menschliche Erfindungen sind ohne das Wort GOttes, soll verdammt sein in dem Volke GOttes, dessen Leben und Werke durch das gewisse Wort GOttes geleitet werden muß.

\begin{schriftzitat}
V. 15. Einen Propheten, wie mich, wird der HErr, dein GOtt, dir erwecken, aus dir und aus deinen Brüdern, dem sollt ihr gehorchen.
\end{schriftzitat}

Dies ist die Hauptstelle dieses ganzen Buches und eine ganz ausdrückliche Weissagung von Christo, dem neuen Lehrer, daher auch die Apostel diese Stelle stark anziehen; und Moses setzt diesen Spruch hier überaus angemessen an das Ende, wo er von dem Priesterthum, von dem Königreiche, von den obrigkeitlichen Personen und allem Gottesdienste ausgebeugt hatte, um zu zeigen, daß später ein anderes Priesterthum sein werde, ein anderes Königreich, ein anderer Gottesdienst, ein anderes Wort, durch welches der ganze Moses abgethan würde. Und hier beschreibt Moses klärlich, daß es ein Ende mit ihm nehmen werde, und tritt seine Meisterschaft dem zukünftigen Propheten ab. Deshalb wollen wir seine Worte recht sorgfältig ansehen.

Erstlich ist es nothwendig, daß dieser Prophet ein neues Wort bringe, welches Moses nicht gelehrt habe, denn es verheißt hier GOtt, daß er seine Worte in seinem Mund geben werde. Wenn aber hier nicht ein anderes Wort wäre, so hätte nicht verheissen werden sollen, daß es durch diesen Propheten gebracht werden solle, sondern es wäre genügend gewesen zu sagen: „Der soll dein Mund sein“, wie von Aaron im zweiten Buch Mose Cap. 4.\footnote{In den Ausgaben: Exod. 3.} 16. gesagt wird, damit angezeigt würde, daß jener Prophet die Worte Mosis und sein geschriebenes Gesetz lehren werde. Denn da er sagt: Den sollst du hören, welcher erweckt werden wird, gleich wie mich, lehrt er, daß sein Wort, und das Wort dieses Propheten ganz verschieden sind, und dies bestätigt er dadurch, daß er sagt, daß das Volk auf dem Berge Sinai einen solchen Propheten begehrt habe, daß der mit ihnen reden sollte, da sie schon durch Mosen das ganze Gesetz gehört hatten. Aber es kann kein anderes Wort über das Wort Mosis hinaus sein, als das Evangelium, da alles, was zum Lehren des Gesetzes gehört, durch Mosis auf das vollkommenste und reichlichste ausgerichtet ist, so daß durchaus nichts hinzugefügt werden kann. Denn was könnte zu den heiligen zehn Geboten hinzugefügt werden (um alles Uebrigen zu geschweigen)? Was kann Erhabeneres gelehrt werden, als daß man an GOtt von ganzem Herzen glaube, ihm vertraue, ihm liebe, ihn fürchte, daß man GOtt nicht versuche etc.? Ferner, was kann Gerechteres und Heiligeres gelehrt werden als das, was Mosis von dem äußerlichen Gottesdienste, von der Obrigkeit, von der Liebe gegen den Nächsten verordnet? Hier haben also die Juden keinen Grund dafür, daß sie schwätzen, dieser Prophet werde ein solcher sein, der Mosen auslege. Moses erklärt sich selbst in diesem Buche in solcher Weise, daß ein anderer nicht vonnöthen ist, ein anderer auch nicht Einen Buchstaben oder Titel hinzufügen kann, dadurch er klarer oder vollkommener werde. Da nun, weil über die vollkommene Lehre des Gesetzes hinaus kein anderes Wort sein kann, als das Wort der Gnade, so folgt, daß dieser Prophet nicht ein Lehrer des Gesetzes sein werde, sondern ein Diener der Gnade. So zwingt dieser Text die Juden klärlich, von diesem Propheten etwas Anderes zu erwarten, als sie in Mose haben.

Zweitens. Wenn dieser neue Prophet nicht ein anderes Wort brächte, so hätte Moses ihn nicht mit sich vergleichen müssen, indem er sagt: „Der HErr wird ihn erwecken, wie mich.“ Denn alle andern Propheten, welche Mosen gelehrt haben, und nicht ein anderes Wort aufgebracht haben (suscitaverunt), waren nicht wie Moses, oder Mosis Gleichen, sondern geringer als Moses, nämlich Diener des Wortes Mosis, welche das lehrten, was Moses vorgeschrieben hatte, so daß sie in allem diesem nicht irgend einen anderen Propheten oder sich selbst, sondern Mosen selbst und seine Worte hörten. Denn Moses redet in ihnen, und er legt seine Worte in ihren Mund, und sie sind sein Mund dem Volke gegenüber. Aber er unterstellt sich diesen nicht, daß er diesen Propheten sich unterwerfe, und seine Worte in dessen Mund lege, sondern er erklärt, daß der HErr selbst, durch sich selbst, Seine Worte in dessen Mund geben werde, und er sagt, daß Er ihm gleich sein werde im Amte und in dem Gehorsam, den man seinen Worten leisten sollte (audientia), wodurch er ihn sicherlich von dem Gehorsam gegen sich selbst [nämlich gegen Moses
\quellseite{769}
] ausnimmt, und ihn allen Propheten vorzieht, welche Mosen lehren.

Daß er ihn aber von dem Gehorsam gegen Mofes ausnimmt und ihn allen Propheten vorzieht, welche Mosen lehren, das heißt fürwahr, daß das Amt des Gesetzes aufhören und ein neues aufgerichtet werden soll, da kein Mensch von dem Dienste des Gesetzes frei ist, sondern alle dem Gesetze unterworfen sind. Deshalb ist es vonnöthen, daß dieser Prophet den Mosen gleich sei (nach der Macht zu lehren und zu befehlen, denn das will er ausdrücken, da er sagt: „Wie mich“), daß er höher sei als Mofes und Größeres lehren werde. Denn wenn er nicht größer wäre, so würde Mofes ihm nicht den Gehorsam und das Ansehen übergeben. Wenn er aber nicht Größeres lehrte, so wäre er nicht größer. Denn hier redet er nicht von der Gleichheit der persönlichen Würde zwischen Mofes und jenem Propheten, sondern von der Gleichheit der Gewalt oder des Amtes. Denn hier handelt er nicht von dem Leben, den Sitten oder den Werken Mosis oder dieses Propheten, sondern von der Lehre, wie der Text genugsam ausweist, da ein Prophet vom Lehren [seinen Namen hat] und um des Lehrens willen kommt, und hier befohlen wird, daß man ihm gehorchen solle (audientia erga eum).

Wenn also von der Lehre beider gehandelt wird, so wird aus der Vergleichung der Lehre beider leicht offenbar werden, was er nothwendiger Weise predigen müsse. Denn Mofes ist ein Diener des Gesetzes, der Sünde, des Todes, weil er lehrt und bringt zu thun, und durch die Strahlen des Gesetzes alle schuldig macht des Todes, straffällig wegen der Sünde; er fordert, aber er gibt nicht, was er fordert. Da nun aber dieser Prophet den Mosen vorfindet, der dieses lehrt, und er selbst neben ihm zum Lehrer eingesetzt wird, so muß sein Wort etwas Anderes lehren; er kann aber nichts Anderes lehren als Sünde, Zorn und Tod, wenn er nicht Gerechtigkeit, Gnade und Leben lehrt. Deshalb muß er nothwendiger Weise ein Lehrer des Lebens, der Gnade und der Gerechtigkeit sein, gleichwie Mofes ein Lehrer der Sünde, des Zorns und des Todes ist. Es soll aber beides in gleicher Weise gehört werden, wie es von GOtt erweckt ist, weil durch das Gesetz alle gedemüthigt werden müssen, und durch das Evangelium alle erhöht werden. Sie sind gleich an göttlichem Ansehen, aber ungleich nach der Frucht ihres Amtes und einander völlig entgegengesetzt, da jener die Sünde und den Zorn, welche Mofes durch sein Amt erregt, in seinem Amte durch die Gerechtigkeit und die Gnade aufhebt (evacuat). Dieser Prophet fordert also nichts, sondern scheut das, was Mofes fordert.

Wir haben daher aus dieser Stelle die beiden Aemter des Wortes, welche nothwendig sind zur Seligkeit des menschlichen Geschlechtes, das Amt des Gesetzes und das Amt des Evangelii, das eine zum Tode, das andere zum Leben, und wie sie zwar gleich sind, wenn man auf die Gewalt sieht, aber überaus unähnlich, wenn man die Frucht betrachtet. Denn das Amt Mosis ist ein zeitliches, welches schließlich ein Ende haben soll durch die Ankunft des Amtes Christi, wie er hier sagt: „Dem sollt ihr gehorchen.“ Aber das Amt Christi wird durch kein anderes zu Ende kommen, da es die ewige Gerechtigkeit bringt und die Sünde abthut, wie es Dan. 9, 24. heißt. Darum wird hier das levitische Priesterthum ganz und gar aufgehoben und ihm ein Ende gemacht, weil es eingesetzt war, um Mosen zu lehren. Wenn nun dem Priesterthum ein Ende gemacht wird, so hört auch das Gesetz auf, wie der Verfasser des Briefes an die Hebräer [Cap. 7, 12.] sagt: „Wo das Priesterthum verändert wird, da muß auch das Gesetz verändert werden.“ So kann hier kein anderer der Prophet sein als Christus selbst.

Aus allem diesem folgt, wie ganz ungehörige, ja, verderbliche Lehrer die sind im neuen Testamente, welche mit Gesetzen und Werken die Gewissen quälen, da diese Abweisung von Christo dies Amt ganz vernichtet und aufhebt. Ferner sind die noch viel schädlicher, die durch ihre Satzungen und menschlichen Gesetze die ganze Welt beschweren. Denn wenn das Amt dieses neuen Propheten das Amt\footnote{Der Text scheint uns hier mangelhaft zu sein. Entweder ist ein zweites ministerium einzufügen (dies haben wir gethan), oder statt legis ist legem zu lesen.} des Gesetzes Mosis nicht leidet, welches doch göttlich ist, wie wird er in seinem Reiche die Gesetze der Menschen leiden? so daß zu sehen kommt, daß durch diesen einen Text das ganze unordentliche Wesen (chaos) der papistischen Tyrannei mit seinen Mönchen umgestoßen wird.

Doch hier wirst du sagen: Aber man findet überall im Evangelio und in den Episteln der
\quellseite{770}
Apostel Gebote, also wird entweder dieser Prophet nicht unser Christus sein, oder seine Lehre wird sich in nichts von dem Gesetze Mosis unterscheiden. Darauf ist kurz zu antworten: Die Gebote des neuen Testamentes werden gegeben für die Gerechtfertigten und die neuen Menschen im Geiste. Denn da wird nichts gelehrt oder geboten, was nicht allein die Gläubigen angeht, die alles freiwillig thun, nicht aus Noth oder wider ihren Willen. Das Gesetz aber wird gegeben für den alten Menschen, der todt ist in Sünden, daß es ihn dränge und ihm seine Sünde zeige. Denn das heißt in Wahrheit und eigentlich das Gesetz lehren. Deshalb findet es den Menschen nicht allein unwillig, sondern auch unvermögend, das zu thun, was das Gesetz fordert, wie es hier im Texte sagt, daß das Volk sich am Tage der Versammlung geweigert habe, und die Stimme des Gesetzes nicht hören konnte; darum habe es einen anderen Lehrer gefordert, welcher das Wort zu ihnen redete, das sie ertragen könnten.

Daher liegt das Verständnis dieser Sache in der Erkenntniß und dem rechten Unterschiede des Gesetzes und des Evangelii, daß du wissest, des Gesetzes Lehre sei nur diejenige, welche gebietet, was die Gottlosen und verlorenen Menschen thun sollen, wie es 1 Tim. 1, 9. heißt: „Denn Gerechten ist kein Gesetz gegeben, sondern den Gottlosen“ etc. Wo aber gottselige Leute sind, da ist das Gesetz schon abgethan, welches nur damit umging,\footnote{Wir sind der Lesart der Wittenberger gefolgt: egit statt: agit in der Erlanger.} daß es die Gottlosen demüthigte durch die Erkenntniß ihrer Sünde und ihres Unvermögens. Aber die Lehre des Evangelii ist diejenige, welche lehrt, woher du das empfangen kannst, wodurch du das Gesetz erfüllen mögest, und in dieser Hinsicht nichts gebietet, auch den Geist nicht drängt, der von selbst, willig geworden durch den Glauben, läuft. Es [das Evangelium] fügt aber einige Gesetze hinzu, aber nur wegen der Ueberreste des alten Menschen im Fleische, damit dies getödtet werde, welches noch nicht gerecht geworden ist; von diesen Gesetzen ist jedoch der Geist seinerseits frei, der allein an dem Glauben ein Genüge hat. Hiervon haben wir anderswo reichlich geredet.

Nun wollen wir die Worte ansehen. „Einen Propheten aus deinem Volke und aus deinen Brüdern, wie mich, wird dir der HErr erwecken, den sollt du hören.“ Hier weissagt er, daß Christus wahrer Mensch sein werde, und von dem Geblüte der Juden, weil das Heil von den Juden kommt, Joh. 4, 22. Bisher ist aber aus diesem Volke niemand aufgetreten, der ein anderes Wort gelehrt hätte als Moses, und ein neues Amt eingesetzt hätte, als allein dieser unser Christus. So viele Propheten auch vor ihm gewesen sind, haben sie doch alle Mosis gehalten und gelehrt. Dieser aber hat nicht allein die Juden, sondern auch alle Heiden in der ganzen Welt von Mose befreit und das neue Wort des Evangelii gegeben. Da er aus den Juden herstammt, wie sowohl Heiden als Juden bezeugen, so wird durch diese Thatsache bewiesen, daß er der ist, von dem Moses hier redet, und daß diese Weissagung in ihm erfüllt ist.

\begin{schriftzitat}
V. 16. Wie du denn von dem HErrn, deinem Gott, gebeten hast zu Horeb etc.
\end{schriftzitat}

Im zweiten Buche Mosis Cap. 20 liest man nichts von dieser Bitte, warum redet denn Moses so? Darauf ist zu antworten: Moses schließt aus den Worten des Volkes, daß diese Bitte in ihnen gewesen sei, da sie sagen: Ich will die Stimme des HErrn, meines GOttes, ferner nicht mehr hören, und dieses überaus große Feuer nicht mehr sehen, damit ich nicht sterbe [2 Mos. 20, 19. 5 Mos. 5, 25.]. Aus diesen Worten, sage ich, nicht wie sie lauten, sondern nach dem, was sie meinen, entnimmt er dies, denn er sieht auf ihre innerste Herzensmeinung und auf den Hauptpunkt ihres Begehrens. Denn sie waren durch den Schrecken vor der Stimme so bewegt, daß sie mit dem größten Verlangen ihres Herzens ein sanfteres und lieblicheres Wort ersehnten. Denn der ihnen durch die Stimme des Gesetzes angedrohte Tod zwang sie, nach dem Leben und dem Wort der Seligkeit zu seufzen, und in solcher Angst wünschten sie mehr, als daß sie hätten mögen sollen, um ein anderes Wort zu bitten. Aber GOtt, der die Tiefen des Herzens erforscht, hört und antwortet nicht nach den auswendig tönenden Worten, sondern nach dem inwendig glühenden Verlangen, gleichwie Christus im Evangelio [Luc. 19, 2. ff.] den Zachäus ansah, nicht nach dem, was er auswendig redete, sondern nach dem Verlangen seines Herzens. Denn auch Zachäus würde nicht gewagt haben, auswendig darum zu bitten, daß
\quellseite{771}
Christus in seinem Hause einkehren sollte, und doch wünschte er innerlich nichts mehr. Daher nahm er ihn auch mit Freuden auf, da sein Begehren erfüllt wurde, welches er vorher kaum empfand.

Es zeigt also Moses an dieser Stelle, mit welcher Brunst und Hitze des Herzens man nach dem Worte der Gnade Verlangen tragen müsse, da es ja die Gabe GOttes ist, von der insonderheit das Wort des Paulus verstanden wird, da er im Briefe an die Ephesier Cap. 3, 20. sagt, daß GOtt uns gebe und thue über alles, das wir bitten oder verstehen. Denn er haßt die trägen und kalten Beter (petitores), welche hoffen, daß sie es durch viele Worte erlangen, aber er will unaussprechliches Seufzen haben, welches diejenigen sicherlich nicht haben können, welche ihre Tage in eiteln Leben zubringen, und im Lande derer wohnen, die ein sündiges Leben führen, sondern diejenigen, welche durch die Empfindung des Todes und den Biß der Sünde zur Verzweiflung gebracht werden, so daß auch sie mit dem Volke Israel sich weigern, ferner die Stimme GOttes zu hören und das überaus große Feuer zu sehen, das heißt, die Kraft des Gesetzes und den Stachel der Sünde, das heißt, den Tod und GOttes Zorn, der sie gleichsam wie ein Feuer bereits verzehrt, so daß sie das Wort des Lebens zu hören begehren.

Und beachte, wie schön Moses sein Amt dem künftigen Propheten abtritt. Denn er sagt, daß das Volk darum gebeten habe, daß ein anderer Prophet erweckt werden möchte, während sie doch 2 Mos. 20, 19. baten, daß Moses selbst mit ihnen reden sollte, aber nichts davon sagten, daß außer Mose ein anderer erweckt werden solle. Durch diese Rede zeigt Moses an, daß das Volk vor Angst nicht gewußt habe, was es bitten solle, oder wie es bitten solle. Denn es war thöricht, darum zu bitten, daß sie Moses hören möchten, da Moses nichts Anderes lehren und hören lassen (sonare) durfte, als was sie selbst auf dem Berge hörten, das ist, das Gesetz. Denn so pflegen erschreckte und angstvolle Leute zu bitten, was sie nicht wissen, wie Paulus Röm. 8, 26. sagt: „Wir wissen nicht, was wir beten sollen, oder wie wir beten sollen, sondern der Heilige Geist vertritt die Heiligen“ etc. Deshalb legt Moses ihre Bitte hier auch aus nach ihrem Verlangen, als wenn er sagen wollte: Du hast thöricht gebeten, daß ich mit dir reden möchte, da du dich fürchtetest, jene Stimme auf dem Berge zu hören, weil dies beständiglich die Stimme meines Amtes ist, aber GOtt, der dein Herz ansieht, hat gesehen, daß du nicht mich, sondern einen anderen erbeten, der lieblicher redete als ich, und da er dies dein Gebet erhört, wird er dir einen solchen Propheten geben, wie du ihn haben willst.

Aber was für einen willst du? Freilich einen solchen, der diejenigen, welche durch meine Stimme zerschlagen sind, das heißt, durch das Gesetz der Sünde und des Zorns, heile, der die gesund mache, welche zerbrochenes Herzens sind, der die Erschreckten und Zerknirschten los mache und tröste. Mit welcher Arznei? Freilich durch das Wort des Lebens und der Seligkeit, so daß die Stimme derer gehört werde, die da Frieden verkündigen, Gutes predigen auf den Bergen [Jes. 52, 7.], daß der da sei, von dem nachher Maleachi [Cap. 3, 1.] vorhergesagt hat: „Und bald wird kommen zu seinem heiligen Tempel der Engel des Bundes, deß ihr begehret“, und nach ihm Jesaias [Cap. 61, 1. 2.]: „Der Geist des HErrn ist über mir, darum hat er mich gesalbt, und hat mich gesandt, den Elenden zu predigen, zu heilen, die zerbrochenes Herzens sind, zu predigen ein gnädiges Jahr des HErrn.“ Um einen solchen (sagt er) hast du zu Horeb gebeten, da du mich in thörichter Weise erbäteist; einen solchen wird dir der HErr an meiner Statt geben. Siehe du zu, daß du ihn hörest anstatt meiner. Diesem werde ich weichen, weil du mich nicht ertragen kannst, und darum gebeten hast, daß ich zum Schweigen gebracht und abgethan werde.

\begin{schriftzitat}
V. 18. Ich will meine Worte in seinen Mund geben.
\end{schriftzitat}

Auch dies beweiset klar, daß jener Prophet anderes lehren werde, da er ja solche Worte noch nicht in Mosis Mund gegeben, ihm dieselben auch nicht geboten hat, sondern er verheißt, daß er sie künftighin geben und gebieten werde. Deshalb muß dies Wort ein anderes sein als das Wort des Gesetzes, welches er schon in Mosis Mund überaus reichlich und aufs vollkommenste gegeben und geboten hatte, so sehr gewaltig, daß es durch seine Kraft die schon Erschreckten in den Tod und Verzweiflung jagte. Denn vergeblich würde er verheißen, daß es noch gegeben und geboten werden solle, wenn er es schon in den Mund Mosis gegeben und geboten hätte. Da nun der künftige Prophet etwas Anderes lehren
\quellseite{772}
sol, so ist vonnöthen, daß sein Wort nicht ein Wort des Gesetzes sei, welches tödtet und schreckt, sondern welches lebendig macht und tröstet. So sehen wir, daß durch diese Stelle die Propheten unterwiesen worden sind, wo sie von einem anderen Bunde, einem anderen Worte, einem zukünftigen Priesterthum weissagen, wie Jeremias [Cap. 31, 31. f.]: „Siehe, es kommt die Zeit, da will ich einen Bund machen, nicht wie der Bund gewesen ist“ etc., und Ps. 110, 4.: „Du bist ein Priester ewiglich“, und Jes. 2, 3.: „Das Wort des HErrn wird ausgehen von Zion“ etc.

\subsection[§. 19. Und wer meine Worte nicht hören wird]{§. 19. Und wer meine Worte nicht hören wird etc., von dem will ich’s fordern}

Durch dies Wort schafft er die ganze bürgerliche Obrigkeit ab, die bisher in Mose gelehrt worden ist, denn er hatte das Schwert und die Richter eingesetzt, um das Gesetz Mosis zu erhalten und zu verwalten. Aber hier, wo er ein neues Wort einsetzt, ordnet er kein Schwert, keinen Richter, sondern er droht, er selbst werde es sein, der da räche, wie er auch im Propheten Micha sagt, Cap. 5,\footnote{In den Ausgaben: Cap. 6.} 14.: „Ich will Rache üben mit Grimm und Zorn an allen Heiden, so nicht gehorchen wollen.“ Denn da dies die Predigt zum Leben und zur Seligkeit ist, soll niemand zu derselben gezwungen werden durch menschliche Gewalt, da GOtt ein williges Herz verlangt, welches glaubt. Da dies aber durch menschliche Kraft weder geleistet noch erzwungen werden kann, so muß man die allein GOttes Gericht überlassen, welche nicht glauben wollen, weil er selbst sie schon verdammt hat, da er sagt [Joh. 3, 18.]: „Wer nicht glaubt, der ist schon gerichtet“, und hier: „Von dem will ich’s fordern.“

Hieraus siehst du, wie toll die Bischöfe und Fürsten heutzutage sind, welche die Ketzer oder Gottlosen mit Gewalt zum Glauben zwingen, und diejenigen, welche sich nicht zwingen lassen wollen, tödten und verbrennen. Sie maßen sich sicherlich das Gericht an, welches GOtt für sich selbst in Anspruch nimmt, und zwingen, wo niemand gezwungen werden kann, und fordern, was weder sie, noch irgend ein Mensch leisten kann. Aber dadurch, daß sie dies thun, zeigen sie, wie sie Christum und sein Wort gar nicht verstehen, noch auch sehen, was Moses hier lehre. Denn da sie aus dem christlichen Wesen eine Art äußerlicher Larve machen, daß sie meinen, diejenigen glauben, welche aus Furcht vor der Strafe oder Schande äußerlich den Namen Christi bekennen, und nicht Acht darauf haben, daß diese Sache allein dem Geiste angehört und ganz frei ist von Menschen: so werden sie alsobald gezwungen, dieselbe dem Schwerte und Gericht zu unterwerfen, da ja jedes äußerliche Gesetz und äußerlicher Gottesdienst (religio) mit Nothwendigkeit weltlicher (profanis) Richter bedarf, wie wir auch an dem Gesetz Mosis gesehen haben. So ist es geschehen, daß sie aus dem Worte ein Gesetz, aus dem Glauben Werke, aus der Wahrheit einen äußeren Schein gemacht haben, darum sind sie gezwungen gewesen, endlich auch das Schwert anzuwenden.

Mit derselben Tollheit, aber unter einem anderen Schein, rasen unsere neuen Propheten, welche, nachdem sie, von eitler Ehre besessen, bessere und höhere Dinge lehren wollten als das Wort des Lebens, welches uns Christus zu unserer Zeit von neuem wiederhergestellt hat, durch den Geist des Teufels dahin gebracht werden, daß sie es verlachen, daß sie von uns über den Glauben und die Liebe belehrt werden sollten. Sie haben aber anstatt des Wortes eine neue Art von Geist erfunden, neue, nichtige Wörter, mit denen sie das gemeine Volk bethören. Und damit wir ihr Gift meiden mögen, da sie zu unserer Zeit ihr Wesen treiben (currunt), so wollen wir kurz ihre Narrheiten berühren, um die Gottlosen zu erinnern, daß sie sich ja nicht durch ihre prächtigen Worte verführen lassen. Um diese Dinge leichter verständlich zu machen, wollen wir zuerst unsere Lehre (nostra) vortragen. Denn so lehren und halten wir, und zwar mit Wahrheit:

Zuerst und vor allen Dingen muß der Prophet gehört werden, von dem Moses hier sagt, das heißt, unser Heil fängt nicht an mit irgend einem Werke, das unser ist, sondern von dem Hören des Wortes des Lebens, da ja auch das Gesetz um deswillen vorhergeht, damit es demüthige und die Sünde zeige, und uns lehre, daß wir aus unseren Kräften nichts vermögen, Röm. 4, 15.: „Das Gesetz richtet nur Zorn an“, in einem solchen Maße, daß das Volk das Gesetz nicht allein nicht thut, sondern es ferner nicht einmal hören will, da es sie tödtet, und deswegen um ein anderes Wort bittet, welches
\quellseite{773}
lebendig mache, ehe wir irgend etwas thun, wie du hier gesehen hast.

Wenn man aber das Wort hört, so wird der Heilige Geist gegeben, welcher das Herz reinigt durch den Glauben [Apost. 15, 9.]. Röm. 10, 17.: „Der Glaube kommt aus der Predigt“, zwar nicht bei allen, welche hören, sondern bei welchen GOtt will. Denn der Wind bläset, wo er will, Joh. 3, 8., nicht, wo wir wollen. Wenn wir aber den Heiligen Geist empfangen haben durch den Glauben, so sind wir alsbald eben durch diesen Glauben gerecht, ohne irgend ein Werk welches wir thun, allein durch GOttes Gabe, und empfinden Frieden und ein gutes Gewissen gegen GOtt, Röm. 5, 1., und rufen fröhlich und sicher: Abba, lieber Vater. Nun ist nicht mehr Furcht oder ein knechtisches Fliehen vor dem Zorne GOttes, sondern ein kindlicher Zugang zu GOtt durch die Gnade, in welcher wir stehen und uns rühmen, Röm. 5, 2.

Wenn aber das Herz in solcher Weise inwendig durch den Glauben im Geiste gerechtfertigt und befriedet ist, so geht man dann bald nach außen auf verschiedene Weise. Erstens geschieht mit dem Munde das Bekenntniß zur Seligkeit, und es wird auswendig das Wort gepredigt, durch welches wir inwendig selig geworden sind, damit auch andere durch uns selig werden, gleichwie auch wir durch die, welche eher waren als wir, damit das Reich Christi gemehrt werde. Hier erhebt sich die Verfolgung, Kreuz und Trübsal um des Wortes willen in der Welt, welche das Licht und das Leben haßt, weil sie die Finsterniß liebt. Dadurch geschieht es, daß der Glaube bewährt und vollendet wird zur vollkommenen Hoffnung, die uns unserer Seligkeit gewiß macht. Zweitens wird unser alter Mensch angegriffen, in welchem die Ueberbleibsel der Sünde sind, welche durch ihr Gesetz in den Gliedern wider den Geist streiten, daß wir nicht thun, was der Geist will. Hier wird die Tödtung des Fleisches vorgenommen. Drittens kommen als Früchte gute Werke hervor als Zeugen des Glaubens und des Geistes gegen den Nächsten in der Liebe, Freundlichkeit, Frieden, Gütigkeit etc., Gal. 5, 22. Dies ist der kurze Inbegriff und die rechte Ordnung des christlichen Lebens, welche vom Worte abhängt und anfängt, und sich erstreckt bis zu den Werken der Liebe gegen den Nächsten.

Diese Ordnung lassen diese neuen Propheten anstehen und verfahren so: Erstlich ergreifen sie die Tödtung. Diese theilen sie in viele Stufen oder Pfunde (pondera),\footnote{Vgl. Tischreden, Cap. 37, § 1. Walch, St. Louiser Ausg., Bd. XXII, 1010.} welche sie mit verschiedenen Namen benennen, als, eine besteht in der Entgräbung, eine andere in der Vermunderung, eine andere in der Disciplin, eine andere in der Langenweile, eine andere in der Entleerung der Creatur, und niemandem gestehen sie die Gerechtigkeit zu, bis daß er durch diese Stufen fortgeschritten ist zur vollkommenen Tödtung; dann werden sie mit dem Heiligen Geiste erfüllt, daß sie alles vermögen, alles wissen, alles wollen. Denn nicht aus der Predigt des Wortes, das sie prahlerisch verachten, sondern aus der Tödtung kommt ihnen der Geist. Denn sie ärgern sich daran, daß sie sehen, daß die Gerechten und Gottseligen noch mit Ueberbleibseln der Sünde zu schaffen haben. Darum nennen sie dieselben Heuchler und [solche Leute], welche den Glauben vorgeben und stehen aus den Silben der Schrift, nicht schöpfen aus dem Geiste. Sie aber sind ohne alle Ueberbleibsel der Sünde, darum haben sie den wahren Glauben aus dem Geiste.

Du siehst also hier ganz offenbar den Teufel, welcher sie blind macht, daß sie von der Tödtung anfangen, welche doch dem Worte und dem Glauben und dem Geiste folgen muß, und mit dem Geiste aufhören, welcher allem vorausgehen muß, und einen ganz neuen Schein (larvam) von Werken unter dem Namen „Tödtung“ herbeibringen. Unterdessen verachten sie das Wort selbst, so daß sie nichts als einen äußerlichen Schein aufrichten und auf eine neue Weise lehren, daß durch unsere Kräfte die Gerechtigkeit zuwege gebracht werde, nämlich eine solche Gerechtigkeit, welche gänzlich ohne Ueberbleibsel der Sünde ist, was doch in diesem Leben unmöglich ist, Röm. 7. Doch dabei sind diese Leute, die sich des Geistes rühmen und der völligen (purissimae) Tödtung und der reinsten Gerechtigkeit, nicht allein voll von allen Lastern, sondern auch besessen bis zu einer unversöhnlichen Raserei. Denn nicht leicht wirst du jemals größere Hoffahrt sehen können, größere Ueberhebung, Vermessenheit, Anmaßung, Neid, Haß, Zorn, Ungeduld, Geiz, Kleinmüthigkeit (wenn die Sache nicht gut vorwärts geht), Kühnheit
\quellseite{774}
(wenn es glüht), darnach auch Unzuverlässigkeit.\footnote{Es ist hier mit der Wittenberger Ausgabe infidelitatem zu lesen statt infelicitatem in der Jenaer und Erlanger. Letzteres paßt nicht in den Zusammenhang.} Ja freilich, so tödten sie sich, so entgräben sie sich (onerantur crassities suas)! Wahrlich, durch diese Sünde werden sie vielmehr vom Teufel beschwert, als von GOtt geehrt.\footnote{onerantur — honorantur, ein Wortspiel.}

Sodann, weil sie diesen verwirrten und unverdauten Schein der äußerlichen Tödtung und der Gerechtigkeit hoch aufwerfen, sind sie gezwungen, auch das Schwert einzusetzen, um dies auszuüben und zu erhalten, wie es denn der Brauch ist bei allen äußerlichen Gesetzen. Und gleichwie ihre Tödtung und ihre Gerechtigkeit verwirrt und verkehrt ist, so ist bei ihnen auch nichts als ein verwirrtes Durcheinander (confusum chaos) im Gebrauch des Schwerts. Denn nicht der rechtmäßigen Obrigkeit, sondern dem gemeinen Pöbelhaufen wider die rechtmäßige Obrigkeit übergeben sie es, um alle Gottlosen (das heißt, die es nicht mit ihrer Weisheit halten) zu tödten, so daß ihre Secte nicht sowohl wegen ihrer Thorheit und Unverständigkeit, als wegen ihres offenbaren Wüthens und ihrer Raserei zu verdammen ist, nämlich der Teufel ist es, der sich auf diese Weise an dem wieder aufstehenden und wieder blühenden Worte rächt.

So wollen wir nun diese Propheten fahren lassen, die durch ihre offenbare Raserei genug und übergenug widerlegt sind, und bleiben in der reinen Lehre und Erkenntniß Christi, indem wir wissen, daß die Summe unserer Religion sei der Glaube im Geiste, der durch das Wort geschenkt ist; sodann, daß noch Sünde übrig bleibe im Fleische, welche ausgefegt werden muß, und daß nicht eine Tödtung erfordert werde, welche vollkommen ausgeführt sei, sondern die vollkommen ausgeführt werden soll. Deshalb müsse man auch diejenigen in Liebe tragen, nicht aber verdammen, welche noch schwach sind im Glauben und Leben, Röm. 15, 1. Dies ist das höchste Lob und das erste Gesetz unter den Christen, die schon gerecht sind, wie Paulus sagt [Gal. 6, 2.]: „Einer trage des andern Last, so werdet ihr das Gesetz Christi erfüllen.“ Darum müssen wir uns mit dem höchsten Fleiße hüten, daß wir keinem Werke die Gerechtigkeit zuschreiben, sondern dieselbe allein dem Glauben im Herzen überlassen nach dem Worte Röm. 10, 10.: „So man von Herzen glaubt, so wird man gerecht“; wiederum, daß wir nicht einen schläfrigen Glauben nähren, der die Tödtung vernachlässige, das heißt, wir sollen die Tödtung so ausüben, daß wir weder meinen, daß, wenn diese geschehen ist, dadurch die Gerechtigkeit oder der Geist erlangt werde, wie diese Propheten rathen, noch auch, wenn sie unterlassen ist, weniger der Geist und die Gerechtigkeit durch das Wort gebracht werde.

Diese Reinheit der Erkenntniß Christi ist freilich gar zart und gering (rara) auch bei heiligen und geistlichen Leuten, nämlich so sehr schadet der Schein der Werke und die Bosheit des ungetödteten Fleisches, daß auch die besten Leute sich mehr und eher dazu bringen lassen, daß sie ihr Fleisch tödten und Werke thun, als daß sie auf das Wort dringen und den Glauben und Geist erlangen, als ob durch die Werke und die Tödtung der Glaube und Geist käme, während im Gegentheil die Tödtung und die Werke sowohl aus dem Geist als aus dem Glauben kommen, der Geist aber und der Glaube nur aus dem Worte des Lebens kommen. Denn es ist auch dieses eine grobe Bosheit des Fleisches, welche mit allem Fleiße getödtet werden muß, daß sie die Werke und die Tödtung so groß macht und hochschätzt, daß sie die Kraft des Wortes nicht sieht, welche mächtig ist durch die Macht GOttes, die Gläubigen selig zu machen [Röm. 1, 16.], und häufig große Männer betrügt. Dies sei nun genug von den falschen Propheten, daß sich ein gottloser Mensch vor ihnen hüte und das unverfälschte (sanum) und reine Wort GOttes bewahre. Es folgt:

\begin{schriftzitat}
V. 20. Wenn ein Prophet vermessen ist etc.
\end{schriftzitat}

„Es muß ja Aergerniß kommen“ [Matth. 18, 7.], und bei dem wahren und unverfälschten Worte GOttes mengen sich immer falsche Lehrer ein, wie Petrus und Paulus und Christus bezeugen: „Es werden falsche Lehrer unter euch sein, wie auch unter dem Volke falsche Propheten waren“, 2 Petr. 2,\footnote{In den Ausgaben: 1 Pet. 2.} 1., und: „Nach meinem Abschied werden unter euch kommen greuliche Wölfe“, Apost. 20, 29. Denn der Teufel schläft nicht, sondern wenn die Leute schlafen, säet er Unkraut zwischen den guten Samen. (Es erinnert uns also auch Moses der Gefahr,
\quellseite{775}
daß wir gewiß seien, es werde geschehen, daß neben dem Worte der Wahrheit gottlose Leute, entbrannt von der eitlen Ehre Lehrer zu sein [Jac. 3, 1.], dafür angehen sein wollen, daß sie Neues und Besseres lehren, wie auch unsere Propheten thun, von denen eben geredet ist. Daher sagt er hier: „Welcher vermessen ist zu reden in meinem Namen, das ich ihm nicht geboten habe zu reden“, nämlich, daß nur Frevel und Vermessenheit sie treibt zu reden, ohne GOttes Befehl.

Er zeigt aber zwei Arten von Propheten an, welche wider die rechte Lehre aufstehen, die eine kommt im Namen des HErrn, die andere im Namen fremder Götter. Die im Namen des HErrn kommen, sind die, welche das Wort GOttes vorwenden und die heilige Schrift gebrauchen, weil GOtt in derselben redet und GOttes Name in derselben ist, da er sie ja selbst gegeben hat. Diese sollten sich am meisten einstellen (erant futuri) zur Zeit Christi unter den Juden, wo sie sich der heiligen Schrift als des Wortes GOttes bedienten wider das Evangelium, für die Gerechtigkeit des Gesetzes, und nach ihnen unter den Ketzern. Die im Namen anderer Götter kommen, sind die, welche ohne Schrift und GOttes Wort kommen, aber ihren eigenen Geist vorgeben, wie die Manichäer und nach ihnen der Papst und alle, die mit menschlichen Satzungen das Volk verderben, unter welche auch unsere Propheten zu rechnen sind, welche uns fort und fort vom Himmel herunter fallen, und wollen, daß man ihrem Geiste glaube, auch wenn er ohne die Schrift rede.

\begin{schriftzitat}
V. 21. Ob du aber in deinem Herzen sagen würdest etc.
\end{schriftzitat}

Ich habe oben beim dreizehnten Capitel gesagt, wie diese zwei Dinge mit einander stimmen, daß er dort gesagt hat: auch wenn das Zeichen komme, welches ein falscher Prophet vorhergesagt hat, solle man ihm nicht glauben, hier aber sagt: wenn es nicht komme, so sei es ein Zeichen, daß er falsch sei, nämlich, es seien zwei Arten von Lehren; die eine, die schon angenommen sei und durch göttliche Gewalt oder Wunder bestätigt. Wider diese Art soll man keine Wunder zulassen, selbst nicht einmal einen Engel vom Himmel. Denn durch diese Wunder werden wir versucht, wie er oben [Cap. 13, 3.] gesagt hat, wie an jenem Propheten gezeigt wird 1 Kön. 13, 1. ff., welcher wider den Altar zu Bethel sogar durch zwei Wunder bestätigt war, und doch, da er von einem andern Propheten versucht ward, zu Falle kam und, von einem Löwen geschlagen, starb. So hat Paulus vorherverkündigt [2 Thess. 2, 9. ff.] gleichwie Christus [Matth. 24, 24.], daß Zeichen und lügenhafte Wunder wider die Wahrheit kommen würden, auf daß die Ungläubigen verführt würden. So hat der Teufel, um den Glauben zu zerstören, viele Zeichen gethan bei den Wallfahrten (visitationibus) zu den Heiligen, mit Geistererscheinungen und dergleichen.

Die andere [Art der Lehre] ist die, welche erst noch angenommen werden soll. Hier muß man nicht glauben, wenn nicht Zeichen geschehen, weil GOtt nie ein neues Wort geredet hat, ohne es durch Zeichen zu bestätigen, wie er das Gesetz Mosis mit den größten Zeichen bestätigt hat, das Evangelium auch mit überaus großen und fast beständigen, da es zuerst angenommen ward. Deshalb muß man seine Zustimmung aufschieben, wenn eine neue Lehre kommt, bis daß Zeichen geschehen. Wenn wir das thun, so ist GOtt getreu, der nicht zuläßt, daß sie Zeichen thun. Wenn wir es aber nicht thun, so gestattet er mit Recht, daß Zeichen geschehen, damit wir betrogen werden und verderben.

So habe ich meinen rasenden Propheten gethan; da sie rühmten, ihr Geist sei so groß, wie wir wissen, daß auch der Apostel Geist nicht gewesen sei, nämlich, daß sie mit GOtt selbst, gleichwie Moses, täglich von Angesicht zu Angesicht redeten, habe ich sie geheißen, sie sollten Wunder thun, sonst wären sie des Glaubens nicht werth. Sie haben es freilich versprochen, aber ich sehe oder höre bis jetzt noch nichts. Ja, trotziglich habe ich gesagt und sage es noch: Mein GOtt wird dafür sorgen, daß euer Gott keine Zeichen und Wunder thue.\footnote{Vgl. die Tischreden des Cordatus No. 125. Walch, St. Louiser Ausgabe, Bd. XXII, 1822.}

Diese Wunder betrügen aber unerfahrene Leute, weil sie nicht Acht darauf haben, oder weil sie nicht wissen, wie groß die Gewalt und die Kraft des Teufels sei, mit wie großem Blendwerk er es zuwege bringen könne, daß du meinst, du seiest blind auf einem Auge, verletzt am Schenkel oder einem andern Gliede, oder am ganzen Leibe krank, während doch da in Wahrheit
\quellseite{776}
nichts vorgeht; darnach aber, wenn du etwas gelobt, oder irgend einen Heiligen anrufst, läßt er ab, so daß die Gesundheit zurückkehrt. Dann meint man, daß dies durch das Verdienst der Heiligen und aus göttlicher Kraft geschehe. So kann er einen, der im Wasser [anscheinend] erstickt ist, einige Stunden lebendig erhalten ohne wahrnehmbaren Athem und Empfindung, da ihm nichts leichter ist, als die Sinne zu täuschen.

In allen diesen Dingen muß man nicht sowohl auf die Wunder sehen, als auf die Ursache derselben. Wenn es sich um Werke und äußeren Schein handelt, daß diese bewährt werden sollen, so ist es der Teufel. Aber in Sachen des Worts und des Glaubens spielt er nicht in solcher Weise mit Wundern, weil er das Wort des Glaubens haßt. Wenn nun auch etwas Wahres geschieht, so daß es eine rechte Krankheit und eine rechte Gesundheit ist, soll man deshalb nicht größeren Glauben schenken, weil der Tausendkünstler die Kräfte des Leibes kennt und die Eigenschaften der Kräuter und aller heilsamen und aller schädlichen Dinge, welche er nach seinem Belieben bald anwendet, bald fernhält, wie er durch die Hexen thut, welche die Leiber derer verletzen, welche sie hassen, sie bisweilen auch tödten, wiederum das Verletzte auch wieder gesund machen etc.

Dies sei genug von den Propheten. Die heimliche Deutung davon, daß die Leviten unter den Kindern Israel nichts besitzen, ist, daß die Diener des Worts sich nicht in Händel der Nahrung flechten sollen, und sich aller Dinge enthalten, auf daß sie dem HErrn gefallen, der ihr Erbe ist. Das heißt, das Wort ist ihr Besitzthum und ihr Reichthum, mit dem sie GOtt dienen und das Volk erhalten, daß sie das Land der Verheißung, das ist der Seelen Seligkeit, erwerben. Auch das Recht der Priester an den Opfern, nämlich die Schulter, die Backen, der Wanst und die Erstlinge, bedeutet dasselbe, daß sie in der Uebung (usu) und im Dienste des Worts seien, so daß sie mit dem Arme kräftig seien im Werke, und mit den Backen die Gottlosen beißen und zermalmen, mit Strafen, mit Schelten, mit inständigen Bitten. Mit dem Wanst aber sollen sie sich der Schwachen annehmen und sie tragen; die Erstlinge aber sollen sie gebrauchen, indem sie lehren, daß man alles der Gnade GOttes zuschreiben müsse, nicht unsern Kräften.

Die Wahrhager aber und das ganze Gelichter von Greueln, welche Moses hier aufzählt, bedeuten die gottlosen Lehrer, welche aus ihrem Kopfe weissagen, was GOtt ihnen nicht geboten hat, und die Gewissen beladen mit gewissen Tagen, gewissen Werken, gewissen Stätten, gewissen Personen. Deshalb täuschen sie die Gewissen mit trügerischen und scheinbaren Meinungen und nichtigen Worten, nicht anders als die Zauberer und Beschwörer die Augen verblenden mit ihrem Blendwerk, so daß sie meinen, an etwas zu empfinden oder zu sehen, wo sie nichts empfinden noch sehen. Solche Leute nennt Paulus Eph. 4, 14. πανούργους\footnote{Das ist: listige, boshafte Leute. An der angeführten Stelle steht: ἐν πανουργίᾳ.} und sie werden recht eigentlich durch diese Täuscher und Zauberer abgebildet. Sie lassen auch ihre Kinder durchs Feuer gehen, das heißt, sie entzünden und verderben das Volk durch einen falschen Eifer in der Gottseligkeit, wie dies leicht angewendet und verstanden werden kann.

\section{Das neunzehnte Capitel}

\begin{schriftzitat}
V. 1–4. Wenn der HErr, dein GOtt, die Völker ausgerottet hat etc.
\end{schriftzitat}

Nachdem die Bestimmungen über den Gottesdienst und die Obrigkeit vollendet sind, fährt Moses fort mit dem zweiten Theil des Gesetzes, nämlich der Liebe, welcher in der zweiten Tafel enthalten ist. Daher handelt dies ganze Capitel vom Todtschlag und vom Geize, wie ein zufälliger Todtschlag verziehen, ein absichtlicher aber bestraft werden solle. Darnach ordnet er durch ein allgemeines Urtheil, daß jedes Unrecht durch gleiche Strafe (talionis lege) gerächt werden soll. Denn es fordert die Liebe, daß für den öffentlichen Frieden gesorgt werde, aber wenn die
\quellseite{777}
Bösen nicht im Zaume gehalten werden, so kann der Friede nicht bestehen. Daher kommt es der Liebe zu, die Gesetze streng und ohne Erbarmen zu verwalten, zur Rache über die Bösen, damit die Guten und Gotteligen erhalten werden. Wiederum ist es Sache der Liebe, daß die Strenge des Gesetzes durch Billigkeit gelindert werde, damit der Barmherzigkeit erlange, der zwar einem bösen und schädlichen Menschen ähnlich ist, aber dennoch weder böse noch schädlich ist, weil er es wider seinen Willen und ungern gethan hat. Daher wird hier befohlen, drei Städte auszusondern zur Zuflucht für die Leute, welche unabsichtlich einen Todtschlag begangen haben.

\begin{schriftzitat}
V. 4–7. Wenn jemand seinen Nächsten schlägt, nicht vorsätzlich x.
\end{schriftzitat}

Er erklärt mit einem Exempel, wie er den nicht vorsätzlichen Todtschlag verstanden wissen wolle, welches handgreiflich (crasso) genug ist, nämlich, wenn jemand mit etwas Anderem umgeht, als, wenn er Holz haut, und ihm das Eisen vom Stiele führe und den Nächsten erschlüge. Hier siehst du gar schön, wie GOtt die ἐπιεικειαν oder Billigkeit in den Gesetzen erfordere, wie er will, daß man nicht allein nach den Thaten richte, sondern nach dem Willen. Denn dies Exempel muß ganz allgemein auf alle Gesetze gezogen werden, so daß bei allem, womit gesündigt sein mag, vielmehr auf das Herz als auf die Hand des Sündigenden gesehen werde, wie hier Moses ausdrücklich erklärt, daß derjenige des Todtschlags nicht schuldig sei, welcher nicht aus Haß, sondern zufällig jemand getödtet habe, ohne seinen Willen. Aber diese Gesinnung (animus) muß doch mit gewissen Zeichen und Umständen bewiesen werden, damit nicht ein jeglicher nach einer solchen Begebenheit sagen möge: Ich habe es nicht mit Willen gethan. Deshalb setzt Moses das Exempel von dem Eisen, das vom Stiele fährt, und gibt klärlich den Umstand an, welcher beweiset, daß der Todtschläger mit etwas Anderem beschäftigt gewesen sei und das Werk des Todtschlags wider seinen Willen gethan habe.

So auch, wenn jemand aus Unvorsichtigkeit einem anderen ein Auge ausschlüge, entweder mit der Faust oder mit einem Steine, aber durchaus etwas Anderem umging, als daß er ihm das Auge ausschlagen wollte, so ist er dieses Unrechts nicht schuldig. Ebenso, wenn er einen anderen an irgend einem Gliede verletzte oder ihm einen Schaden zufügte, wie ich Augenzeuge gewesen bin (vidi factum), daß ein Bruder den andern mit dem Schwerte durchbohrte und tödtete, da er den wilden Eber, der seinen Bruder unter sich hatte, erstechen wollte. Denn so steht es in diesem Leben, daß jeder von seinem Nächsten Fährlichkeit erwarten muß. Denn wer unter Leuten (societate communi) leben will, und Frieden, Sicherheit, Schutz genießen, und sich aller Rechte, Vortheile und Güter bedienen, der muß auch die Gefahren, Schäden und Nachtheile tragen, und alles Ungefähr, wie es auch immer geschehen mag. So sagen die Juristen: Ein Bürger ist dem andern einen Brand schuldig, das heißt, wenn jemand ein Bürger sein will, so muß er mit seinem Nachbar in Gefahr der Feuersbrunst oder irgend eines anderen Schadens stehen.

\begin{schriftzitat}
V. 8. 9. Und so der HErr, dein GOtt, deine Grenze weitern wird x. So du anders alle diese Gebote halten wirst x.
\end{schriftzitat}

Du siehst, daß es sich mit den gesetzlichen Verheißungen für das zeitliche Reich ganz anders verhalte als mit den Verheißungen der Gnade oder den geistlichen Verheißungen. Denn die gesetzlichen Verheißungen gründen sich auf den Willen und das Verdienst der Menschen, wie er hier sagt: „So du anders halten wirst.“ Darum sind sie nicht feststehend noch immerwährend, weil keine Verdienste der Menschen Bestand haben (permanent), gleichwie auch die Menschen selbst nicht, mögen sie auch noch so heilig sein. Aber die geistlichen Verheißungen gründen sich ohne Ansehen der Verdienste allein auf die ewige Barmherzigkeit GOttes, die aus Gnaden umsonst gegeben wird, denn „es liegt nicht an jemandes Wollen oder Laufen, sondern an GOttes Erbarmen“ [Röm. 9, 16.]. Daher herrscht in jenen [gesetzlichen Verheißungen] der freie Wille, ja, er verderbt und richtet alles zu Grunde. In diesen aber herrscht und erhält die Gnade und der Geist.

\begin{schriftzitat}
V. 10. Auf daß nicht unschuldig Blut vergossen werde.
\end{schriftzitat}

Du siehst, daß ein Todtschlag, der nicht vorsätzlich geschieht, von GOtt selbst für einen solchen erklärt wird, der keine Schuld habe (innoxium), wie er auch kurz zuvor [V. 6.] gesagt hat: „Daß
\quellseite{778}
er die Seele dessen nicht schlage, an dem kein Urtheil des Todes ist, weil er keinen Haß vorhin zu ihm getragen hat.“ Siehe das Herz, welches des Hasses nicht schuldig ist, macht, daß auch die That nicht die Schuld des Todtschlags habe. Darum will er auch viele Städte zur Zuflucht haben, weil sich in einem großen Volke viele Fährlichkeiten zutragen, die einer von dem andern erleidet.

Und es ist zu merken, daß das Wort, da er sagt [V. 6. Vulg.]: „Damit nicht etwa, der dem Getödteten am nächsten steht, ihn verfolge, und schlage die Seele dessen, der des Todes nicht schuldig ist“, durchaus nicht so verstanden werden darf, als ob jener Rächer, sofern er ein Privatmann ist (privata manu), den Todtschläger töten solle oder dürfe, wie der Text für ausführliche Leute zu lauten scheint, sondern daß, indem er sich der Sache des Getödteten annimmt, die Richter durch gesetzmäßige Zeugen die Sache erledigen sollen. Weil dies auf sein Betreiben und Dringen geschieht, so wird gesagt, daß er es thue, nach dem Worte [der Juristen]: Was andere, durch uns veranlaßt, thun,\footnote{In der ersten Ausgabe und in der Wittenberger: Quod per nos alii faciunt, was wir angenommen haben. In der Erlanger und in der Jenaer: Quod nos per alios facimus.} davon wird mit Recht gesagt, daß wir es gethan haben. Denn keinem Menschen ist es gestattet, sich selbst zu rächen, sondern: „Mein ist die Rache“, spricht GOtt [5 Mos. 32, 35.]. Deshalb hat er Richter und Obrigkeit eingesetzt.

\begin{schriftzitat}
V. 11–13. Wenn aber jemand Haß trägt wider seinen Nächsten etc.
\end{schriftzitat}

Auch hier siehst du, daß die Rache der Obrigkeit befohlen werde, denn er sagt [V. 12.]: „Es sollen die Aeltesten in seiner Stadt hinschicken, und ihn von dannen holen lassen“, und dennoch, gleich als ob der Rächer selbst es thun sollte, fügt er gleich darauf hinzu: „Und sollen ihn in die Hände des Nächsten oder des Rächers geben“, nicht als ob er selbst, als Privatperson, ihn mit dem Schwerde schlagen solle, sondern er soll ihn vor Gericht seiner Schuld überführen und zum Tode bringen. Ja, gleich als wenn das ganze Volk dies thun sollte, sagt er [V. 13. Vulg.]: „Du sollst dich sein nicht erbarmen, und sollst das unschuldige Blut aus Israel thun.“ Es bleibt also nur übrig, daß jene Rechte und die öffentlichen Vollstreckungen der Rechte zwar dem Volke befohlen werden, aber einem solchen Volke, welches mit obrigkeitlichen Personen und Richtern versehen ist, welche, wenn der Rächer anklagt und die Zeugen es bewähren, nach dem Gesetze mit gerechtem Urtheil richten und tödten. Er will aber, daß sie mit einem vorsätzlichen Todtschläger kein Erbarmen haben sollen, weil er aus Bosheit gesündigt und den öffentlichen Frieden gestört hat. Und durch dies Beispiel sieht du, daß die Freiheiten der Oerter, die bei uns anstatt der Zufluchtsstädte sind, nur den Unschuldigen nützen sollen, den Schuldigen aber nichts nützen. Denn so gebietet er auch 2 Mos. 21,\footnote{Wittenberger und Jenaer: Exod. 23. Erlanger: Exo. 25.} 14., daß man einen vorsätzlichen Todtschläger ergreifen und sogar vom Altar des HErrn wegreißen soll und ihn tödten. Heutzutage aber dienen jene Freistätten (immunitates) ganz und gar ohne allen Unterschied und meistens nur den Schuldigen zu einem Zufluchtsorte, und beschirmen öffentliches Unrecht, vielleicht nach dem Exempel des Romulus, des Vaters der Römer, welcher zuerst den Räubern Rom zu einer Schützerin (patronam) gemacht hat durch die Freistätte.

\begin{schriftzitat}
V. 14. Du sollst deines Nächsten Grenze nicht zurück treiben.
\end{schriftzitat}

Eine treffliche Stelle wider den Geiz und Tyrannei, um dessenwillen nicht selten auch Todtschlag geschieht, um die Grenzen auch einnehmen zu können. Deshalb kommt es dem Amte der Obrigkeit zu, dafür zu sorgen, daß ein Bürger vor dem andern sicher sei hinsichtlich seines Erbes. Besonders in diesem Volke Israel war dies nothwendig, wo er wollte, daß die Erbgüter nicht mit einander vermengt werden sollten.

\begin{schriftzitat}
V. 15–18. Es soll kein einzelner Zeuge wider jemand auftreten etc.
\end{schriftzitat}

So hat er auch oben im 17. Capitel, V. 6., gesagt: „Auf Eines Zeugen Mund soll niemand sterben.“ Dies sind schöne Einwürfe (gnomae), würdig, daß sie häufig in aller Munde wären, denn auch Christus führt diesen Spruch mehr als Ein Mal an [Matth. 18, 16. Joh. 8, 17.]: „Alle Sache soll auf zweier oder dreier Zeugen Mund bestehen.“ Wie aber, wenn beide Zeugen die Unwahrheit (falsum) sagen, was
\quellseite{779}
soll dann der Richter thun? Darauf ist zu antworten: Wenn er weiß, daß sie die Unwahrheit sagen, so soll er nicht richten; wenn er es aber nicht weiß, soll er richten nach der Aussage der Zeugen. Denn dies ermahnet die Ordnung des Gerichts, welche von GOtt eingesetzt ist, welcher befiehlt, daß man zweien Zeugen glaube. GOtt wird schon zusehen, wie er jenen Zeugen Gleiches wiedervergelte, und sie mögen zusehen, wie sie es vor Gotte verantworten, daß sie seiner Ordnung und seines Ansehens gemißbraucht haben zum Verderben des Nächsten. Solche falschen Zeugen traten auf wider Naboth, den Jesreeliten [1 Kön. 21, 13.], wider Christus [Matth. 26, 60. f.] und Stephanus [Apost. 6, 13.]. Und wer auf solche Weise umkommt, der muß sein Unglück als ein solches ansehen, wie das ist, davon wir oben gesagt haben, daß es ein Bürger von dem andern zu erwarten habe, als da sind Feuersbrunst, Pest, ansteckende Krankheit, Krieg, Wunden und allerlei andere Beschädigungen. Denn gleichwie man von dem Nachbar die Gefahr der Feuersbrunst leiden muß, so auch falsches Zeugniß und alle anderen Beschwerlichkeiten.

\begin{schriftzitat}
V. 19. So sollt ihr ihm thun, wie er gedachte seinem Bruder zu thun.
\end{schriftzitat}

Hier hast du einen offenbaren und klaren Text, daß der Versuch (conatus) bestraft werden solle, wenn es auch nicht zur Ausführung der That kommt, weil er befiehlt, den Zeugen zu tödten, und ohne Barmherzigkeit Seele um Seele von dem zu nehmen, welcher wider eine Seele falsches Zeugniß abgelegt hat. Denn um deswillen fügt er hier alsbald das Gesetz der Wiedervergeltung hinzu (legem talionis), daß man dem, der falsches Zeugniß redet, thun solle, wie er selbst hat thun wollen, so daß man auch hier sagen kann: Es gibt kein billigeres Gesetz, als daß diejenigen, welche Anschläge machen, jemanden zum Tode zu bringen (necis artifices), durch ihren eigenen Anschlag umkommen. Denn in solcher Weise wird „Zahn um Zahn, Auge um Auge“ etc. nicht allein von der vollendeten That, sondern auch von dem Unternehmen des Versuchs verstanden.

Aber hier ist eine Frage übrig, welche den unschuldigen und unvorsätzlichen Todtschläger betrifft. Wenn er unschuldig ist, warum wird er einer solchen Strafe unterworfen, daß er gezwungen wird, außerhalb seiner Stadt in der Stadt der Zuflucht als ein Verbannter zu leben, bis daß der Hohepriester stirbt? Ferner: Wenn er etwa außerhalb der Grenze der Freistadt gefunden und getödtet wird, so wird bestimmt, daß der Bluträcher [wenn er ihn erschlägt] des Blutes nicht schuldig sei, wie im vierten Buch Mosis Cap. 35, 27. geschrieben steht? Darauf ist zu antworten: Gleichwie derjenige, welcher von ungefähr getödtet ist, zugleich mit den Seinen gezwungen ist, dies Uebel von seinen Nachbaren und denen, mit welchen er in bürgerlicher Weise zusammenwohnet, zu erleiden, wie ich gesagt habe, so muß auch jener nach demselben Rechte seine Verbannung tragen als ein Uebel, welches ihm dadurch zustößt, daß er Nachbaren hat (ex vicinitate). Das Gesetz trifft aber hierüber eine so strenge Bestimmung, um von dem vorsätzlichen Todtschlage desto mehr abzuschrecken, und um den Schaden, wiewohl er nicht mit Willen zugefügt worden, aber doch den Freunden schwer zu tragen ist (molestam), auf irgend welche Weise auszugleichen und wieder gutzumachen, gleichsam\footnote{Statt vel ut in der Jenaer und der Erlanger haben wir mit der Wittenberger velut angenommen.} durch einen dem entgegengesetzten Schaden.

\subsection{Heimliche Deutung}

Den Nächsten tödten bedeutet, daß man mit dem Worte der Lehre seine Seele verderbe und das Gewissen verführe. Wenn dies ohne Vorsatz geschieht, das heißt, wenn etwa das Wort nicht nach seinem rechten Sinne gehandelt wird, wie die Väter bisweilen zu thun pflegen, da sie aus dem Felsen (petra), Christo, den Petrus gemacht haben, oder die Werke allzusehr gelehrt und die Worte des Glaubens gar unvorsichtig (imprudentius) auf den äußerlichen Schein gezogen haben, als, da sie Regeln, wie man leben solle, gewisse Ordnungen des Fastens, der Feste und anderer Ceremonien aufgerichtet haben, haben sie, soviel die That anbelangt, zwar die Seelen getödtet, aber sie haben das durchaus nicht vorgehabt, daß sie tödten wollten, sondern das Geschick und der Schade (casus) der menschlichen Natur, welche blind und unvollkommen im Glauben ist, hat sie so zu Fall gebracht. Diesen wird eine Freistadt angewiesen, die Verheißung des Heils, nämlich das Wort: „Vergib
\quellseite{780}
uns unsere Schuld“, und es wird ihre Schwachheit getragen, bis daß der Priester stirbt, das heißt, bis daß sie in völligerem Glauben erkennen, daß Christus es sei, der durch seinen Tod alle Sünden vernichtet und uns frei macht. Denn durch dieses Bewußtsein werden wir alle erhalten, daß wir wissen, daß alle Sünden, mögen sie nun in Unwissenheit oder aus Schwachheit begangen sein, in Christo erlassen werden, wenn wir zu ihm fliehen und bitten, daß sie uns vergeben werden mögen.

Wer aber in glühendem Hasse wissentlich jemanden tödtet, der wird ohne Barmherzigkeit getödtet. Dies sind die Ketzer, welche auf Satans Antrieb damit umgehen, daß sie die Seelen verführen, sie in ihre Secten locken, von Christo und dem reinen Worte wegreißen, ja, sich ergötzen und rühmen wegen der Menge derer, die verloren gehen und ihnen anhangen. Deshalb verdienen sie keine Gnade, haben auch keine Vergebung der Sünde, da sie durch ihr eigenes Urtheil verdammt sind, Tit. 3, 11., und eine Sünde zum Tode begehen, welche, wie sie nicht getragen wird, so auch nicht vergeben wird. Denn es ist eine Lästerung wider den Heiligen Geist, für welche man nicht bitten soll, wie Johannes [1. Ep. 5, 16.] sagt.

Nicht, daß man nicht bitten solle für irgend eine Sünde, sondern daß man GOtt nicht bitten soll, er möge sie nicht für eine Sünde halten oder als eine Gerechtigkeit annehmen, wie Moses im vierten Buche Cap. 16, 15. bittet, daß GOtt sich nicht wenden möge zu ihrem Opfer, das heißt, daß er es ihnen als eine Sünde zurechnen möge, wie es denn in Wahrheit Sünde ist, weil sie es nicht [für eine Sünde] rechnen, sondern rühmen als eine Gerechtigkeit. Sonst wird für alle Sünde gebeten, die wir erkennen, und so wird sie nicht zugerechnet, wenn sie auch andauern sollte, wie denn die Ueberbleibsel der Sünde im Fleische sind, Gal. 5, 17. ff. Denn das ist die Meinung Johannis, da er nicht will, daß man für eine Sünde zum Tode bitte, das heißt, er will nicht, daß sie für eine Sünde gehalten werde, die nicht zugerechnet werde, oder für keine Sünde, wie er doch will, daß man von allen Sünden halte, die wir erkennen und bekennen. Denn wenn man betet für eine solche Sünde, daß sie nicht für Sünde gehalten werde, das heißt bitten, daß ihre Ketzerei und gottloses Wesen GOtte nicht mißfallen möge, oder daß er ein GOtt sei, dem gottlos Wesen gefalle, wider den 5. Psalm, V. 5.

Daß man die Grenzen nicht zurücktreiben soll, welche die Vorigen gesetzt haben [V. 14.], das heißt, daß man nichts hinzuthun soll zu der von den Aposteln überlieferten Lehre, als ob man dadurch den Gewissensangelegenheiten besser rathen könne. Und diese Stelle werfen die Sophisten und die Bischöfe für sich hoch auf, indem sie für ihre Satzungen und Gewohnheiten schreien: man solle die Grenzen nicht verrücken, welche die Väter gesetzt haben. Unter Vätern verstehen sie aber ihre Bischöfe und Lehrer, aber nicht die Apostel. So begründen (stabiliunt) sie ihre Fabeln mit einer bildlichen Rede, indem sie dabei nicht sehen, wie sie selbst die ersten sind, und zwar allein von allen (auch wenn die geistliche Deutung Bestand hätte), welche die Grenzen verrücken, nicht nur des Glaubens und des Geistes, welche durch die Vorigen, die Apostel und Christus, im Evangelio gesetzt sind, sondern auch ihre eigenen, die von ihren Vorigen und von ihnen selbst gesetzt sind, da das ihr einziges Bestreben ist, daß sie Gesetze mit Gesetzen ändern und häufen, und, wie jener sagt, Gesetze festhalten und für Geld wieder aufheben, so daß man sieht, daß sie die Gesetzgebung anstatt eines Spiels und der Würfel haben, um damit in den Gewissen der Menschen ihr Spiel zu treiben; und doch halten sie anderen entgegen, sie sollen die Grenzen nicht verrücken, welche die Vorfahren gesetzt haben.

Das Gesetz der Wiedervergeltung: „Auge um Auge“ etc., bedeutet, daß die falschen Lehrer von GOtt gestraft werden, gleichwie sie selbst Schaden gethan haben, so daß sie getödtet werden, wenn sie getödtet haben; wenn sie verletzt haben und die Leute schwach gemacht im Glauben oder Werken und das reine Wort nicht gelehrt haben, so sollen sie nach demselben Maße auch selbst wieder von GOtt empfangen, wie Paulus im Briefe an die Römer Cap. 14,\footnote{In den Ausgaben: Roma. 15.} 15. sagt, daß sie wider Christum sündigen, indem sie die Schwachen verderben, um derer willen Christus gestorben ist. So sind die falschen Zeugen, welche wider den Nächsten Zeugniß ablegen, diejenigen, welche ein Gewissen der Sünde und des Todes machen, wo es nicht ist, so daß sie die Seelen tödten, welche nicht sterben, und
\quellseite{781}
die lebendig machen, welche nicht leben. Diese nun soll man sorgfältig erforschen [V. 18.], und ihnen thun, wie sie gedachten zu thun, das heißt, durch das Wort der Schrift soll man sie überwinden und verdammen, das heißt, durch das Wort sollen sie zur ewigen Verdammniß verurtheilt werden, wenn sie nicht Buße thun, und es soll ihnen bewiesen werden, daß das falsch sei, was sie gelehrt haben. So werden sie durch den Geist des Mundes Christi getödtet und verdammt, auf daß sie leben, wenn sie wollen, oder ewiglich verderben sammt ihrer Lehre.

\section{Das zwanzigste Capitel}

\begin{schriftzitat}
V. 1. ff. Wenn du in den Krieg zeuchst wider deine Feinde etc.
\end{schriftzitat}

Moses bleibt noch bei der Sache des Todtschlags bis zum 23. Capitel. Deshalb, damit alles in diesem Volke durch ein gewisses Wort geleitet würde, gibt er ihnen auch ein Gesetz für den Krieg, welches wahrlich sehr schön ist und ihnen immer den herrlichsten Sieg verschafft hätte (victoriosissimam), wenn sie es gehalten hätten. Denn so oft die Juden geschlagen worden sind, haben sie dies erlitten, weil sie dies Gesetz nicht gehalten haben, da es unmöglich war, daß sie hätten besiegt werden können, wenn sie Krieg führten nach diesem Gesetze. Zuerst nimmt er ihnen das Vertrauen auf ihre Waffen und Macht, damit sie nicht meinen möchten, sie seien es, die da siegten und triumphirten, und legt eine unüberwindliche Grundlage und beständige Stärke des Sieges, nämlich ein Wort der herrlichen Verheißung und des Glaubens, indem er spricht: „Fürchte dich nicht vor ihnen, denn der HErr, dein GOtt, ist mit dir.“ Danach befiehlt er, daß dies Wort vor der Schlacht von dem Priester gepredigt und eingeschärft werden solle, indem er sprechen solle: „Israel, höre zu. Ihr gehet heute in den Streit wider eure Feinde; euer Herz verzage nicht, fürchtet euch nicht, erschrecket nicht und lasset euch nicht grauen vor ihnen; denn der HErr, euer GOtt, gehet mit euch, daß er für euch streite mit euren Feinden, euch zu helfen.“ Lieber, was ist es Wunder, wenn ein Volk, das mit solchen Verheißungen gewappnet ist, auch ohne Waffen nicht allein seine Feinde, sondern sogar die ganze Welt besiege und schlüge? Welche Mauern, welche Waffen, welche Macht, welche Menge könnte denen Widerstand leisten, für welche der HErr streitet? Sind denn nicht die Egypter im rothen Meere dies inne geworden, da sie sprachen [2 Mos. 14, 25.]: „Lasset uns fliehen von Israel; der HErr streitet für sie wider uns“?

Daraus ist das trotzige Rühmen hergekommen [Ps. 44, 7.]: „Ich verlasse mich nicht auf meinen Bogen, und mein Schwert kann mir nicht helfen“, und wiederum [Ps. 147, 10.]: „Der HErr hat nicht Gefallen an jemandes Beinen, noch an der Stärke des Rosses“, und [Ps. 33, 17. 16.]: „Rosse helfen nicht, und ein Gewaltiger wird nicht errettet durch seine große Kraft“, und [Ps. 144, 1. 10.]: „Gelobet sei der HErr, der meine Hände lehret streiten, und meine Fäuste kriegen, denn er gibt den Königen Sieg.“ So sagt Jonathan, der Sohn Sauls [1 Sam. 14, 6.]: „Es ist dem HErrn nicht schwer, durch viel oder wenig helfen.“ Daher wird gelehrt, daß die Kriege dieses Volkes im Glauben an das Wort GOttes geführt werden sollten. Denn durch den Glauben erschlug David den Goliath und erlangte alle seine Siege. So oft sie aber ohne den Glauben allein mit ihrer großen Zahl und Macht Krieg geführt haben, sind sie geschlagen worden.

Zweitens. Die so mit dem Worte GOttes Gewappneten bekleidet er auch mit der Larve der leiblichen Rüstung, und stellt das Heer in Schlachtordnung auf, und heißt die Neulinge und die Verzagten heimgehen, und richtet eine auserwählte Mannschaft zu, und ordnet alles nicht anders, als ob sie im Vertrauen auf ihre Macht und Waffen den Krieg führen sollten, nach der Weise der anderen Völker. Das thut er, um sich unter der äußerlichen Larve zu verbergen, und heimlich, durch den Glauben an sein Wort, den Sieg zu bereiten, damit es nämlich nicht schiene, als ob Israel GOtt versuchte, indem es bloß und unbewaffnet kämpfte. Jetzt
\quellseite{782}
aber lernt es, daß es zwar Waffen führen und Kriegsrüstung gebrauchen soll, aber nicht darauf vertrauen. Denn er will nicht, daß sie durch die Waffen oder durch die Rüstung, sondern unter den Waffen und der Rüstung siegen sollen, damit so der Glaube eine Statt habe, der mit dem zu thun hat, daß man nicht siehet [Hebr. 11, 1.], und ein Versehen Gottes vermieden werde, welches das Ansehen läßt, was GOtt äußerlich zum Gebrauch im Kriege gegeben hat. Daß GOtt sich so verberge (de hac larva), davon haben wir auch oben gesagt. Denn David hat nicht gesagt: Ich will ohne Schwert Krieg führen, und ohne meinen Bogen wird mir geholfen werden, sondern: Ich will zwar mit Schwert und Bogen Krieg führen, aber nicht darauf meine Hoffnung setzen. Ich will nicht ohne Brod und Wein leben, aber ich will die Hoffnung des Lebens nicht auf das Brod und den Wein setzen. So steht es in allen anderen Dingen.

Drittens, damit die Juden nicht meinen möchten, sie wären nur dazu geboren, daß sie die Menschen tödten sollten, und nicht vielmehr dazu, daß sie die Gottlosen und Unruhigen züchtigen, schreibt er ein Gesetz vor, daß sie den Feinden zuerst Frieden anbieten sollen, was auch etliche Heiden gehalten haben, wie der Dichter\footnote{Virgil, Aeneis, lib. VI, v. 853.} von den Römern sagt:

\begin{quote}
Parcere subjectis et debellare superbos.
\end{quote}

\begin{quote}
[Der Unterworfenen schone, die Stolzen mußt du bekriegen.]
\end{quote}

Und Augustus hat den Frieden immer dem Kriege vorgezogen, und sagte, daß man nur aus dringender Noth Krieg führen solle, da ja der Krieg, sagte er, einem goldenen Angelhaken gleich sei, wo die Kosten größer sind als der Ertrag, und doch, wenn er verloren ist, durch Einen Fang nicht ausgeglichen werden kann. Und wir sehen fast in allen Kriegen durch den Ausgang, welchen sie haben, daß GOtt als Richter gegenwärtig sei, so daß diejenigen, welche den Frieden verlassen und zuerst Krieg anfangen, immer besiegt werden. So fiel auch der König Josias [2 Chron. 35, 22. f.], da er den Frieden, der ihm angeboten wurde, verachtete, und Krieg gegen den König von Egypten anfing. So sind die Weltreiche, welche durch Krieg zuwege gebracht sind, durch Krieg untergegangen, und selbst das letzte Reich der Römer, welches alles unter sich gezwungen hat, ist ein Raub aller Völker geworden.

Viertens. Dieses Gesetz des Krieges wird gegeben wider die auswärtigen Völker, daß sie, wenn diejenigen, welche sich widersetzen, unterworfen sind, alles, was männlich ist, tödten sollen, aber der Weiber, der Kinder und des Viehes und der ganzen Beute der Stadt verschonen, und unter sich austheilen. Aber die Cananiter und Amoriter selbst, in deren Lande sie wohnen würden, gebietet der HErr gänzlich zu verbannen, damit nichts übrig bleibe, davon sie Gottlosigkeit und Greuel lernen könnten.

Nach dem heimlichen Verstande werden hier die in der Sache des Glaubens zu führenden Kriege bedeutet. Hier ist vor allen Dingen vonnöthen, daß wir wissen, GOtt sei unter uns und kämpfe für uns; in dieser Erkenntniß können wir sicher und beherzt die Kriege des HErrn führen. Darnach aber [ist vonnöthen zu wissen], daß unsere Kräfte, das heißt, Weisheit, Erkenntniß, Verstand, sodann auch unsere Waffen, das ist, das Schwert des Geistes und die Schrift, nur ein äußerliches Bemühen (operam) seien, unter welchem GOtt in uns mächtig ist, und streitet und siegt, wie Paulus sagt [2 Cor. 10, 4.]: „Unsere Waffen sind nicht fleischlich, sondern mächtig vor GOtt“, damit wir uns nicht rühmen, als ob wir es seien, die den Sieg errungen haben, sondern GOtt allein uns überall zum Siege verhilft, der in uns allen alles wirkt.

Wir haben drei Arten von Feinden: erstlich, die den angebotenen Frieden annehmen, und sich uns unterwerfen; zweitens, die Widerstand leisten. Diese beiden bezeichnet Moses als auswärtige Völker, welche fern vom Lande der Verheißung gelegen sind. Die dritte Art ist innerlich und in der Heimath, als da sind die Cananiter, Amoriter etc. Diese bedeuten dreierlei Ketzer. Die ersten sind nicht hartnäckig, geben mit leichter Mühe der Ermahnung Raum, und machen als Unterworfene Frieden und dienen in Liebe. Die zweiten sind diejenigen, welche Widerstand leisten und sich mit den Waffen der Schrift vertheidigen. Von diesen soll man tödten, was männlich ist, die übrigen soll man zur Beute nehmen, das heißt, durch das Wort des Geistes sollen die obersten und mächtigsten unter ihnen
\quellseite{783}
verdammt werden, und, wenn sie hartnäckig sind, soll man sie, gleich als wären sie todt, auf immer fahren lassen und für Heiden und Zöllner halten; die übrigen unter ihnen aber soll man annehmen, und ausheilen, damit unsere Liebe an ihnen zu schaffen habe (distribuendi in opera caritatis). Aber die einheimischen und nahe gelegenen soll man gänzlich vertilgen mit allem. Diese sind es, die ganz verhärtet und verstockt sind, unverbesserliche Leute. Diese soll man durch das Wort in den Bann thun, und von uns absondern.

Denn hier schlägt nicht ein äußerliches Schwert die Leiber, sondern das Wort GOttes die Gewissen. Wenn aber die Gewissen überwunden sind, so ist genug überwunden, denn es ist der hochfahrende Sinn zerstört, welcher sich wider die Erkenntniß GOttes erhob [2 Cor. 10, 5.]. Das Uebrige, das noch im Fleische gelebet wird, muß getragen werden, wiewohl es schwach ist. Denn nachdem die Gottlosigkeit des Herzens getödtet ist, schadet die Schwachheit des Fleisches nicht, da die Sünde in den Gliedern nun nicht mehr zugerechnet wird, nachdem das Haupt der Schlange abgehauen ist, das heißt, der Sinn und Wahn der Gottlosigkeit. Aber wo das Haupt und der Wahn nicht weicht, da muß auch das, was man im Fleische auf das beste lebt, zugleich mit der Gottlosigkeit des Herzens getödtet werden, das heißt, in den Bann gethan werden, weil ihnen keine Vergebung, keine Barmherzigkeit, keine Gnade übrig ist für die Sünde, die noch übrig ist im Fleisch, da sie die Gottlosigkeit des Herzens festhalten; sodann ist auch alles verdammt, was sie, wie sie sich dünken lassen, Gutes thun. Diese sind die Cananiter und die nahe gelegenen Amoriter.

Frieden anbieten, das bedeutet, daß man das Evangelium des Friedens und die Lehre des ruhigen und fröhlichen Gewissens, das heißt, des reinen Glaubens, den gottlosen und unruhigen Lehrern und Peinigern der Gewissen anbiete, die sich mit Werken abmühen.

Daß man die Bäume, welche Früchte tragen, nicht abhauen soll, bedeutet, daß man, wenn sie rechte Aussprüche der Schrift haben, diese nicht leugnen solle. Denn es ist noch nie irgend eine Ketzerei gewesen, die nicht auch irgend etwas Wahres gesagt habe. Deshalb muß man das Wahre um des Falschen willen nicht leugnen. Die nicht Frucht tragenden Bäume soll man abhauen und zu Bollwerken verwenden, das bedeutet, daß man die Schriftstellen, welche sie verkehrt anführen und übel entgegenhalten, wieder auf diejenigen zurücktreiben soll, welche sie so gebrauchen, und sie widerlegen durch ihre eigene unnütze Lüge.

Daß ein neuer Ehegemahl, und der ein Haus baut, und der einen Weinberg pflanzt, nicht in den Krieg ausgesendet werden soll, bedeutet, daß ein Neuling nicht in das Amt des Wortes gesetzt werden soll (wie Paulus [1 Tim. 3, 6. f.] lehrt), „auf daß er sich nicht aufblase, und nicht dem Lästerer in den Strick falle“. So auch nicht ein Furchtsamer, das ist, der schwach ist im Glauben, sondern, wie er den Titus [Cap. 1, 9.] ermahnt, ein solcher, der halte ob dem Worte der reinen Lehre, und mächtig sei in dem Worte, zu ermahnen und zu überwinden die Widersprecher, damit er nicht im Kriege entweder umkomme oder fliehe, das heißt, entweder durch Gottlosigkeit überwältigt werde oder den gottlosen Lehrern nachgebe. Denn was er hier vom Einweihen des Hauses und vom Gemeinmachen des Weinberges und vom Heimholen des Weibes sagt, bedeutet, daß ein Streiter des Worts durch die Anwendung (usu) geübt sein muß, so daß er die Lehre nicht allein vom Hören, sondern auch durch die Erfahrung kenne. Denn nach hebräischer Weise wird ein Haus dann eingeweiht, wenn es durch den Gebrauch und Bewohnen in Besitz genommen wird, und der Weinberg wird gemein gemacht, wenn sich jemand desselben, dadurch daß er ihn gebraucht und Nutzen davon hat, bemächtigt. Die Braut wird heimgeholt, wenn man zusammen mit ihr lebt.

\quellseite{784}
\section{Das einundzwanzigste Capitel}

\begin{schriftzitat}
V. 1. ff. Wenn man einen Erschlagenen findet im Lande, das dir etc.
\end{schriftzitat}

Moses bleibt in diesem Capitel noch bei den Sachen, welche Krieg und Blut betreffen. Zuerst setzt er eine große Feierlichkeit (ceremoniam) ein, welche stattfinden soll, wenn ein Erschlagener im Felde gefunden wird, und man den Thäter nicht weiß: wie man die Schuld abthun soll. Dieses Uebel gehört auch zu dem, was bürgerlich ist, und wovon wir gesagt haben, daß man es leiden muß, wenn man mit Nachbarn zusammenleben will (jure vicinitatis). Warum sollten sonst diejenigen am meisten mit der Schuld belastet werden, welche in der nächsten Nachbarschaft der Stelle des Todtschlags (occisi) wohnen, da es doch möglich ist, daß er von ferner Wohnenden oder von Fremden getödtet wäre? Dieses Gesetz wird aber, nach meinem Dafürhalten, deshalb gegeben, damit das rohe Volk des Gesetzes durch dies äußere Gepränge der Traurigkeit und der Buße\footnote{Wir haben mit der Jenaer poenitentiae angenommen statt poenitentia in der Wittenberger und in der Erlanger.} abgeschreckt und vom Todtschlage zurückgehalten würde, gleichwie man die Kinder durch äußere Sachen, die man ihnen vor Augen führt (spectaculis), und Schreckbilder vom Bösen abzuschrecken pflegt.

Zweitens gibt er ihnen ein Gesetz über die gefangenen Weiber. Und hier siehst du, wie viel das Gesetz den jüdischen Kriegsleuten zuließ, daß sie nicht allein mehrere Weiber haben durften, sondern auch, wenn sie Lust dazu hatten (si amor jusserit), eine Heidin und Kriegsgefangene. Denn da er oben von den Kriegsleuten gehandelt hat, welche nicht neue Eheleute sind [Cap. 20, 1–4.], so ist es offenbar, daß fast alle, welche Kriegsdienste leisten sollten, wenigstens länger als ein Jahr Eheleute gewesen sind, und diesen verheiratheten Kriegsleuten wird gestattet, auch eine heidnische Frau mit nach Hause zu nehmen, und, was noch mehr ist, darnach, wenn sie ihm mißfällt, gibt er [Moses] ihm das Recht, daß er sie von sich lasse, aber als eine Freie, so daß sie einen andern heirathen, aber nicht verkauft oder geschändet werden (prostitui) konnte. Denn das ziemt sich nicht für die bürgerliche Ehrbarkeit, daß man eine Gedemüthigte verkaufe oder schände.

Aber was ist das? Stand es denn den Juden frei, die Weiber zu schänden? Es hat der Herr nicht das Wort „schänden“ (prostituendi), sondern das Wörtlein \hebraeisch{עברי}, welches verkaufen oder verdingen (locare) bedeutet, so daß mir der Sinn zu sein scheint, sie solle nicht verkauft werden für Geld, oder auf irgend eine andere Weise um Gewinnes willen weggegeben werden, an wen der Mann will, nämlich an einen Knecht oder einen Freund, sondern sie soll ganz frei ausgelassen werden, so daß sie vielmehr, frei von ihm, heirathen kann, wen sie will.

Kurz, auch dieses Gesetz gehört theilweise zu den Dingen, die Moses den Juden wegen ihrer Här tigkeit nachgelassen hat, wie Christus von dem Scheidebriefe lehrt Matth. 19, 8. Denn es ist unbillig, daß ein Weib ohne ihre Schuld und allein aus Muthwillen von ihrem Manne verstoßen werde (dimitti); dennoch wird den harten und ungeschlachten (barbaris) Leuten diese Unbilligkeit zugelassen, damit sie nicht eine größere Unthat (malum) begingen. Es war aber bei diesem Volke die Vielweiberei gebräuchlich, sowohl nach dem Exempel der Väter als nach dem Recht des Gesetzes.\footnote{Der letzte Satz fehlt in der Wittenberger.}

Es wird aber dieser Gebrauch, daß sie ihre Nägel beschneide, und die Kleider ablege, darinnen sie gefangen ist, und sitze und ihren Vater und ihre Mutter einen Monat lang betraure, eingesetzt, damit die Härte der That gelindert werde, nämlich, daß eine Heidin in dem Volke GOttes zur Ehe genommen wird, so daß sie durch diese Reinigungen gleichsam geheiligt werde und mit dem heiligen Volke vermischt, nachdem sie ihres Volkes und ihrer Freundschaft vergessen hat und nach dem Fleische eine neue Creatur geworden ist. Denn der rohe Pöbel muß durch äußerliche Ceremonien und Gepränge zur Ehrbarkeit und Ehrfurcht geleitet und angehalten werden, da er sich durch den Geist nicht leiten läßt.

\quellseite{785}
Drittens gibt er ein Gesetz über zwei Weiber und ihre Nachkommenschaft, damit der Mann nicht nach seiner Zuneigung und Liebe dem Sohne der Frau, die er am liebsten hat, die Erstgeburt gebe, der nicht der Erstgeborne ist. Denn die Erstgeburt war die edelste und höchste Ehre der Kinder, die nicht durch den Willen eines Menschen geschah, wie die Ehe mit einer Kriegsgefangenen, sondern durch den himmlischen Segen, wobei kein Anziehen der Person stattfindet. Darum darf sie nicht nach menschlichem Willen [von einem auf den andern] übertragen werden, oder nach Willkür des Liebenden oder der Geliebten verändert werden, so daß das ein allgemein gültiger Satz (gnome) ist, den dies Gesetz enthält: Was GOtt als ein Geschenk gibt, darf der Mensch nicht wegnehmen. Dies Gesetz war gar hart, da es sowohl die Feindseligkeit (vindictam) gegen die verhaßte Frau als auch die Gunst gegen die geliebte verhindert. Denn Feindseligkeit und Gunst abzulegen, dazu gehört eine nicht geringe Tugend, ja, es ist der Natur unmöglich. Darum wird es durch das Gesetz und mit Gewalt erzwungen, was die Bosheit weder will noch kann. Du siehst also, daß die Vielweiberei durch das Gesetz gestattet werde.

Viertens [gibt er] das Gesetz von den ungehorsamen Söhnen, welche mit dem Tode bestraft werden sollen. Der beste und gerechteste Gesetzgeber ist Moses, welcher nicht bloß über geringere Vergehen das Todesurtheil ausspricht, sondern viel mehr auch über größere. Unter Menschen steht es so mit den Gesetzen und Bräuchen, daß sie Mord und Diebstahl mit dem Tode bestrafen, den Ehebruch aber selten mit Todesstrafe. Ferner, den Ungehorsam und das schändliche Leben der Kinder bestrafen sie nicht mit dem Tode, viel weniger Vergehungen wider heilige Dinge (sacrilegia), und Gottlosigkeit und Lästerungen wider GOtt und GOttes Wort. Hier aber gebietet Moses so streng, daß die Widerchristlichkeit der Kinder bestraft werden soll, daß er sogar befiehlt, daß die Eltern selbst die ersten sein sollen, die diese Todesstrafe veranlassen, da sie selbst ihre eigenen Söhne vor Gericht bringen sollen, sie anklagen und wider sie zeugen. So hoch schlägt GOtt den Gehorsam und die Ehrerbietung gegen die Eltern an, und wollte GOtt, daß auch wir diesem Gebote nachfolgten, damit der zügellosen und unverschämten Jugend mehr Furcht und Scheu beigebracht würde, welche durch so jämmerliche Beispiele und so freche Sitten überall zu Grunde geht.

Fünftens und letztens verordnet er, daß die Uebelthäter, welche mit dem Tode am Galgen bestraft sind, an dem Tage abgenommen werden sollen, und daß man sie die Nacht nicht am Galgen lassen solle. Er fügt den Grund hinzu: „Denn ein Gehenkter ist verflucht vor GOtt, auf daß das Land nicht verunreinigt werde.“ (Er sagt in thätiger Form (active): Er ist eine Verfluchung GOttes (maledictio Dei), das heißt, er wird bei GOtt und vor GOtt und von GOtt für einen Verfluchten gehalten.“ Nun aber soll man das aus den Augen thun, wovon man weiß, daß es Gott so mißfalle, daß es auch verflucht ist vor ihm, damit es nicht den Schein gewinne, wenn du es so hängen läßt, als ob du Lust daran habest, daß das Land mit dem Fluche GOttes erfüllt und so verunreinigt werde. Dies alles wird verordnet, um den Todtschlag verhaßt zu machen, ebenso wie jener gottesdienstliche Gebrauch (religio) mit dem Leichnam dessen, der als ein Erschlagener gefunden wurde, damit sie vom Töten abgeschreckt würden, in einem solchen Grade, daß sie auch den, der des Todes würdig war, gezwungen durch die Nöthigung des Gesetzes, selbst aber nur ungern tödten möchten, und den Getödteten als ein greuliches Bild (spectrum) des Todtschlags wegnehmen sollten, damit nicht durch die Gewöhnung, den Leichnam anzuschauen und zu dulden, das Grauen vor dem Tödten gering werde und erfolgte.

\subsection{Heimliche Deutung}

Eben [Cap. 19] ist gesagt worden, daß „das Töden“ eine zwiefache heimliche Deutung habe. Denn es wird entweder die Gottlosigkeit oder die Gottlosigkeit getödtet; beides ist ein geistliches Töddten. Ersteres wird bedeutet durch den Todtschlag als die Schuld. Letzteres wird bedeutet durch das Schwerdt des rächenden Gerichts als die Strafe. So wird der in übler Weise getödtet, welcher durch das Wort der Gottlosigkeit von GOtt und dem Worte des Lebens entfremdet wird; es wird aber in rechter Weise getödtet, der durch das Wort des Geistes zur Strafe geschlagen wird, und geschreckt wird durch den künftigen Zorn GOttes, damit er Buße thue und lebe. Wenn er nicht so lebendig gemacht werden
\quellseite{786}
will, so wird er in zwiefachem Tode ein Todter bleiben, sowohl in seiner Gottlosigkeit als auch in dem zukünftigen Zorn über seine Gottlosigkeit, mit dem ihn GOtt tödten wird. Wir aber tödten geistlicher Weise, indem wir schrecken. Daher bedeutet der gefundene Leichnam entweder einen getödteten Bruder oder irgendwelche getödteten Leute (populum), die wir von falschem Wahne angesteckt finden, dessen Urheber oder Lehrer nicht bekannt ist. Hier müssen alle, welche die Nächten sind, hinzulaufen, und, erschrocken in gottseeligem Gewissen, diese Schuld auf sich nehmen, indem sie fürchten, es möchte durch ihre Schuld geschehen sein, oder es möchte der unter uns sein, der es gethan hat, besonders die Aeltesten, das heißt, die Diener des Worts, und nicht anders handeln, um dies Aergerniß aufzuheben, als wenn sie es selbst angerichtet hätten.

Sie sollen deshalb eine junge Kuh von den Kindern nehmen, die noch nicht am Joche gezogen hat, das ist Christus in seinem Fleische, welches er von unserem menschlichen Geschlechte angenommen hat. Er wird aber im Geiste genommen, das heißt, gepredigt und geglaubt. Er ist aber niemals unter dem Joche gewesen, das heißt, er war ohne Sünde, dem Gesetze nichts schuldig, sondern nach seinem Willen für uns geopfert. Sie sollen sie in ein rauhes und felsichtes Thal führen, welches niemals beackert ist, auch nicht Samen empfangen hat, und sollen ihr daselbst den Hals abhauen, das heißt, sie sollen predigen, daß er in der Welt und für die Welt gelitten habe, welche unfruchtbar, rauh und undankbar ist gegen GOtt, nie zuvor das Wort GOttes gehört oder empfangen hat, sondern nur prangte in Gerechtigkeit und Weisheit des Fleisches. Dann sollen, während die Leviten dabeistehen, die Aeltesten zu dem Erschlagenen hinzutreten, und ihre Hände über die junge Kuh waschen, und beten, daß ihnen dies Blut nicht zugerechnet werde, das heißt, über Christum, der so gepredigt und geglaubt ist, sollen die Diener des Worts unschuldig werden, und\footnote{et fehlt in der Jenaer.} handeln [als ob es ihre eigene Sache wäre],\footnote{Eine solche Ergänzung: ac si sua res agatur, scheint uns durch den Zusammenhang und durch den Parallelismus am Ende dieses Absatzes geboten.} und beten, daß ihnen dies Aergerniß verziehen werde. So wird durch Christum alles vergeben, sowohl die

rechten und erkannten Sünden als auch die unbekannten und die, von denen wir aus Furchtsamkeit den Verdacht hegen, daß sie Sünden seien. Ein treffliches Beispiel für diese geistliche Deutung gibt dir Paulus in der ganzen Epistel an die Galater, wo er, da er auf dem Acker der Kirche den Leichnam eines Erschlagenen gefunden hatte, und nicht wußte, wer es gethan hätte, als der Nächste hinzuläuft und handelt, als wenn es seine eigene Sache wäre. Siehe daselbst nach.

Das gefangene Weib, die zur Ehefrau genommen wird, meint St. Hieronymus, sei die Beredsamkeit. Denn in solcher Weise halten sie die heimlichen Deutungen für eine Art Spiel. Wir sagen, daß die heimlichen Deutungen auf den Glauben angewendet werden müssen. Deshalb ist dies Weib entweder die Synagoge, oder die Kirche aus den Heiden. Wir verstehen darunter (facimus) die Synagoge, welche schon war durch die Gaben, die Paulus Röm. 9, 4. 5. aufzählt, und sagt: „Welchen gehört das Gesetz und der Bund, die Herrlichkeit, die Väter“ etc., und völlig geeignet zur Liebe und Ehe, das heißt, Christi Wort zu hören. Das Haar (caesaries) war ihr glänzendes Priestertum; dieses ist abgescheren und abgethan. Das Kleid, darin sie gefangen ist, der äußerliche Schein der Gerechtigkeit, mußte auch abgelegt werden, damit die wahre Gerechtigkeit des Kreuzes eingeführt werde, welche alles, was zuvor Gewinn war, für Dreck achtet [Phil. 3, 7. 8.]. Sie wird auch die Nägel beschneiden, das heißt, sie wirft auch die Härte des Gesetzes, welche die Gewissen zerfleischt, verlieren. Sodann soll sie Vater und Mutter betrauern, das heißt, sie soll erkennen und bekennen, daß die Sünde eine Erbsünde sei und in der Natur, welche durch die Gerechtigkeit des Gesetzes nicht aufgehoben werden kann, sondern vielmehr vermehrt wird. Und so soll sie verheirathet werden (maritetur), das heißt, das Wort GOttes empfangen und glauben, und das Eheweib Christi sein.

Bis hieher hat die geistliche Deutung von der Synagoge für Einen Theil Bestand, wie es denn der Brauch der Schrift ist, von dem Ganzen zu reden in der Figur der Synedoche, da sie nur einen Theil meint. Nun folgt von dem andern Theile, der verstoßen wird, nämlich es ist dasselbe Weib, welches angenommen und verstoßen wird, denn es ist ein und dasselbe Volk, zum
\quellseite{787}
Theil angenommen, zum Theil verlassen. Wenn dir nun nachher das Weib mißfällt, das heißt in diesem gottlosen Theile, so sollst du sie frei lassen, das heißt, Christus hat die Juden verlassen und unterwirft sie auch nicht mehr dem Gesetze, übergibt sie auch nicht einem anderen Meister, sondern sie [die Synagoge] kann frei hinweggehen, wie es sie gelüstet, und dem Meister anhangen, den sie will. So hat der HErr Jes. 5, 2. ff. vorherverkündigt, sein Weinberg werde ihn verlassen. Nun kannst du diese heimliche Deutung ganz allgemein auf alle Werkheiligen ziehen, von denen einige durch das Wort gefangen werden, welche das Haar abscheren, die Nägel beschneiden und allen Ruhm der früheren Gerechtigkeit verdammen, und erkennen, daß sie Sünder sind. Wiederum, etliche werden als Verstockte in ihrem Sinne gelassen.

Zwei Weiber Eines Mannes sind die zwiefache und doch einige Kirche aus den Juden und Heiden. Die Synagoge ist die feindselige, denn sie tödtet die Propheten und kreuzigt Christum etc.; die Kirche der Heiden ist die Geliebte, denn sie nimmt das Wort mit Freuden an. Dennoch hat die Synagoge den Erstgebornen, weil aus ihr Christus kommt, die Apostel, das Wort, und nicht aus den Heiden, denn das Heil kommt von den Juden“, Joh. 4, 22. Deshalb soll man auch heutzutage die Juden nicht verachten, weil von ihnen, nicht von uns, alle Herrlichkeit herkommt, wie Paulus thut Röm. 9, 3. ff. Denn sie sind die ersten Christen gewesen, und „ihnen ist verheißen und vertraut, was Gott geredet hat“ [Röm. 3, 2.].

Der Sohn, welcher widerpenstig und ein Trunkenbold ist, bedeutet einen Lehrer in der Kirche, der erstlich durch das Wort gezeugt ist, darnach aber verderbt durch menschliche Lehren und Träume, und ihnen hartnäckig anhängt. Dieser soll vor dem Gericht der Kirche angeklagt und überführt werden, und mit Steinen, das heißt, mit harten Sprüchen der Schrift getödtet werden, das heißt, überwunden werden, damit er lebe. Wenn er nun nicht leben will, so muß das harte Urtheil des Todes und des Zornes GOttes angewendet werden, damit er so getödtet werde und lebe.

Der wegen einer Uebelthat Gehängte, der wegen des Fluches GOttes abgenommen werden soll, damit das Land nicht verunreinigt werde, bedeutet, daß ein Gottloser, der durch das Wort verdammt ist, und nicht Buße thut, ganz hinweggethan werden soll, und als ein Fluch vor GOtt abgesondert und für einen Bann gehalten werden soll, damit er nicht, wenn er unter uns geduldet würde, sehr viele von uns anstecke und verderbe, wie es heißt 1 Cor. 5,\footnote{In den Ausgaben: 1 Cor. 15.} 13. 6.: „Thut von euch selbst hinaus, wer da böse ist“, denn „ein wenig Sauerteig versäuert den ganzen Teig“.

Aber was sollen mir zu Paulus sagen, der im Briefe an die Galater Cap. 3, 13. diese Stelle von Christo versteht, indem er sagt: „Christus ward ein Fluch für uns, denn es stehet geschrieben: Verflucht ist jedermann, der am Holze hängt.“ Hieronymus quält sich an der Stelle wunderlich ab, da er sagt, Moses redet von einem gehängten Uebelthäter, Christus aber habe keine Sünde gethan. Aber er sieht nicht, daß Paulus auf das einfältigste redet, und Moses, der ganz allgemein redet, richtig und dem Buchstaben nach, ohne bildliche Rede, von Christo versteht. Denn wenn jeder Gehängte ein Fluch ist vor Gott,\footnote{Erlanger: Deo statt: Dei in den anderen Ausgaben.} wie Moses lehrt, so ist auch Christus ein Fluch vor GOtt, und wenn jeder Gehängte ein Uebelthäter ist, so ist auch Christus ein Uebelthäter. Deshalb hätte hier nicht davon gehandelt werden sollen, wie Christus ohne Sünde gewesen sei, sondern, wie er Sünde gehabt habe; desgleichen, nicht davon gehandelt werden sollen, wie er der Gebenedeite und diesem Fluche nicht unterworfen gewesen sei, sondern, wie er demselben unterworfen gewesen sei und habe sein können.

Hier löst Paulus mit Einem Worte die ganze Schwierigkeit, indem er sagt: „Er ward ein Fluch für uns“; für uns, sage ich, nicht für sich. Durch dies Wort lehrt er klarlich, daß jemand zweierlei Sünde habe und daß ein Gehängter ein Fluch vor GOtt werde; einmal für sich und die eigene Sünde; derartig sind alle anderen; das andere Mal für andere und für fremde Sünde; ein solcher ist einzig und allein Christus, welcher sagt, Ps. 69, 10.: „Die Schmach derer, die dich schmähen, fallen auf mich“, und Jes. 53, 4.: „Er trug unsere Sünden“, und wiederum [V. 8.]: „Um die Sünden meines Volks habe ich ihn geplagt“, und [V. 10.]: „Der HErr wollte ihn also zerschlagen um unseiner
\quellseite{788}
Sünden willen“ etc. Mag nun jemand für sich oder für einen anderen gehängt werden, so bleibt der einfache Verstand des Gesetzes, daß jeder für die Sünde Gehängete ein Fluch vor GOtt sei. Deshalb ist Christus, indem er unsere Sünde trug, für uns in Wahrheit ein Fluch geworden, wie der Text dem Buchstaben nach lautet. Wie er beschnitten worden ist und unter das ganze Gesetz gethan für uns, während er für sich frei war von jedem Gesetze, so ist der, „der von keiner Sünde wußte, für uns zur Sünde gemacht, auf daß wir würden in ihm die Gerechtigkeit, die vor GOtt gilt“, 2 Cor. 5,\footnote{In den Ausgaben: 2 Cor. 6.} 21. So ist der ein Fluch geworden, der den Fluch nicht kannte, auf daß wir in ihm ein Segen GOttes wären. Daß aber dieser Fluch vor GOtt vom Holze genommen werde, damit das Land nicht verunreinigt werde, bedeutet, daß man predige und glaube das Evangelium des Kreuzes, nämlich, daß das Kreuz sei Herrlichkeit, der Tod Leben, die Sünde Gerechtigkeit, der Fluch Segen, die Verdammniß Seligkeit, oder wie Paulus [Gal. 6, 14.] sagt, daß wir uns rühmen des Kreuzes unseres HErrn, daß uns die Welt gekreuzigt sei, und wir der Welt, nach dem Worte [Matth. 5, 10.]: „Selig sind, die verfolgt werden“ etc. Sonst, wenn dies nicht gepredigt und geglaubt wird, gibt es nichts Abscheulicheres, nichts Greulicheres als Kreuz und Tod, und das ganze Land wird durch diese Meinung und dies Aergerniß des Kreuzes verunreinigt, denn „den Unreinen ist nichts rein, den Reinen ist alles rein“ [Tit. 1, 15.]. Dies [daß alles rein ist] geschieht durch das Wort, jenes [daß nichts rein ist] ohne das Wort.

\section{Das zweiundzwanzigste Capitel}

\begin{schriftzitat}
V. 1. ff. Wenn du deines Bruders Ochsen oder Schaf siehest irre gehen, so sollst du dich nicht entziehen von ihnen etc.
\end{schriftzitat}

Bis hieher hat Moses fast nur insoweit die Liebe gelehrt, als sie zum vierten und fünften Gebote gehört, nämlich von dem Gehorsam gegen die Eltern und vom Meiden des Todtschlags, durch den die Person des Nächsten beschädigt wird. Nun geht er weiter zu den übrigen, damit er lehre, dem Nächsten Liebe zu erweisen in Betreff äußerlicher Güter, und behandelt ingemein das sechste und siebente Gebot, in denen Ehebruch und Diebstahl verboten wird, aber Keuschheit und Wohlwollen geboten, welche der Liebe zustehen, und Früchte des Glaubens sind, Gal. 5, 22. Er gibt aber in diesem Capitel verschiedene Gesetze, wie man das Wohlwollen und die Keuschheit beweisen möge.

Erstlich befiehlt er, daß man einen Gegenstand, den der Nächste verloren hat, sei es nun ein Ochse, der irre geht, oder ein Kleid oder irgend eine andere Sache, die verloren ist, an sich nehme, bewahre und wiedergebe, und nicht vernachlässige als etwas Fremdes, sondern Sorge dafür trage, als wäre es unser Eigenthum. Solcher Weise, wenn jemand den Esel seines Bruders fallen sieht, so soll er ihm aufhelfen, als ob es sein eigener wäre. Du siehst, daß durch dies Gebot dasselbe gelehrt werde, was Paulus Phil. 2, 4. sagt: „Sehet auf das, daß des andern ist.“ Dies ist durchaus der Liebe eigen, 1 Cor. 13, 5.: „Die Liebe suchet nicht das Ihre“, nur daß das Evangelium, indem es das Gesetz völliger erklärt, lehrt, daß dies auch den Feinden gethan werden soll, während Moses dem rohen Volke des Gesetzes nachläßt, die Feinde zu hassen.

Zweitens. [V. 5.:] „Ein Weib soll nicht die Waffen eines Mannes tragen, und ein Mann soll nicht Weiberkleider anthun.“ Hier wird nicht verboten, daß, um Gefahr zu vermeiden, oder aus Scherz beim Spiel, oder um die Feinde zu täuschen, ein Weib nicht möge die Waffen eines Mannes tragen, oder der Mann ein Weiberkleid anthun, sondern, daß dieses im Ernste und im gewöhnlichen Leben und Brauche nicht geschehe, damit die gebührliche Ehrbarkeit und Würde in beiden Geschlechtern erhalten werde, da es schandbar ist, wenn ein Mann Weiberkleider trägt, und nicht anständig, wenn ein Weib Mannswaffen führt.

\quellseite{789}
Er scheint aber durch dies Gesetz irgend ein Volk zu strafen, in welchem diese Weise stattgefunden haben mag, wie die Schriftsteller insgemein von den Egypten schreiben, daß die Weiber Mannsarbeit, die Männer aber Frauenarbeit thun, die Weiber Geschäfte betreiben und die Männer ernähren, die Männer das Haus besorgen, und viele solche Dinge in einer verkehrten und entgegengesetzten Weise, als es bei anderen Völkern Brauch ist, so daß man dies Gesetz als ein allgemeines verstehen muß, nämlich, daß ein Weib weibliche Verrichtungen, ein Mann Mannsgeschäfte besorge, und überhaupt ein jeglicher, mit seinem Schicksal zufrieden, sein Handwerk treibe und seiner Gaben brauche, Röm. 12, 7. 8., der Schuster bei seinem Leisten bleibe, und der gemeine Mann nicht Rathsherr sei, sondern die Mannigfaltigkeit und Verschiedenheit der Personen bestehe zu einem guten Regimente. Wenn dies verkehrt wird, soll es so schandbar sein, als wenn nach einer verkehrten Ordnung der Dinge ein Weib die Waffen eines Mannes trüge, und ein Mann Weiberkleider anhäte. Gott will, daß sein Volk sich fern halte von den Sitten dieser Völker, da er an denselben einen Greuel hat.

Drittens. Wenn auf dem Wege ein Nest gefunden wird, kann man die Jungen nehmen, soll aber die Mutter, die darauf sitzt, fliegen lassen. Was lehrt doch dies Gesetz anders, als daß sie auch durch wohlwollende Behandlung der Thiere Sanftmuth und Wohlwollen lernen sollen? Sonst würde man es für ein thörichtes Gesetz ansehen müssen, da es über eine so gar unbedeutende Sache nicht allein ein Gebot gibt, sondern auch denen, die es halten, Wohlergehen und langes Leben verheißt, indem er sagt [5. 7.]: „Auf daß dir’s wohl gehe, und lange lebest“ etc. Auch dies kann durch eine allgemeine und gleichsam sprüchwörtliche bildliche Redensweise (figura) darauf gezogen werden, daß man sich hüte vor aller übermäßigen Ausnutzung (usus) des Nächsten, wie wir im Deutschen sagen: Ein williges Pferd soll man nicht zu sehr treiben. So sollen wir der Freunde und Nächsten in solcher Weise gebrauchen, daß wir ihnen etwas übrig lassen, damit wir ihrer auch zu anderer Zeit und bei anderer Gelegenheit gebrauchen können. So soll nun dies ein Sprüchwort für die Hebräer sein: Wenn man die Jungen genommen hat, so soll man die Mutter fliegen lassen, damit es ihnen wohlgehe. Sonst, wenn sie ihre Freunde gänzlich ausgesogen haben, so muß es ihnen endlich übel gehen, so daß dies Gesetz gleichsam wider das zuerst gegebene ist, gegen die, welche sich zwar von anderen gern Wohlthaten und Liebe erweisen lassen, selbst aber keine erzeigen, als ob das des HErrn Wille sei, daß diejenigen, welche immer reichlich geben, nie aufhören sollten [zu geben], aber die, welche empfangen, nie satt werden sollten [zu nehmen], und daß durch die Arbeit jener der Müßiggang dieser ernährt werden sollte [2 Cor. 8, 13.], was im gesellschaftlichen (sociali) Leben etwas ganz Unbilliges ist.

Viertens. Wenn ein Haus erbaut ist, so soll auf dem Dache rings herum eine Mauer gemacht werden, damit nicht jemand herabfalle und eine Blutschuld auf das Haus bringe. Er redet nach der Sitte des Volkes, welches die Häuser so baute, daß die Dächer oben flach waren, gleichwie die Straßen, die mit Steinen gepflastert sind. Daher lesen wir auch im Evangelio [Luc. 5, 19.], daß sie das Dach aufbrachen und den Gichtbrüchigen durch die Ziegel hernieder ließen, und Christus sagt, daß das Evangelium auf den Dächern gepredigt werden solle [Matth. 10, 27.]. Denn da kam man zusammen, spielte und aß, gleichwie bei uns in den Hausfluren. Es kann auch dies ein sprichwörtliches und allgemeines Gesetz sein, daß man in der öffentlichen Gesellschaft so baue, und sich im Verkehr so verhalte, daß man dem andern keine Gefahr, Nachtheil oder Schaden verursache, sondern vielmehr ihm Schutz und Nutzen verschaffe, damit du nicht durch eines anderen Unglück deinen Nutzen suchest, wie man spricht, daß Eigennutz Städte und Reiche zerstöre.\footnote{Siehe Walch, St. Louiser Ausg., Bd. VIII, Anhang B., Sprüchwörter s. v. Reichhart.}

Fünftens. Es soll der Weinberg nicht mit mancherlei Samen besäet werden, es soll auch nicht zugleich mit einem Ochsen und Esel geackert werden, und niemand soll ein Kleid von Wolle und Leinen zugleich gemengt gebrauchen. Dies Gesetz kann dem Buchstaben nach einen Grund der Ehrbarkeit im Auge haben, aber ich halte dafür, daß auch dies eine sprichwörtliche und allgemeine Lehre sei, durch welche die Juden belehrt werden, daß sie einträchtig, einerlei Sinnes, einfältig, aufrichtig sein sollen, ohne Secten und Anrichten (studiis) von Parteiungen. Denn
\quellseite{790}
ein Haus, das in sich selbst uneins ist, wird müste, wie das Evangelium [Luc. 11, 17.] sagt, und durch Eintracht werden kleine Güter groß, durch Zwietracht zerfallen die größten Güter, sagt Sallust. Daher sagen etliche Leute, daß civitas [Stadt] seinen Namen gleichsam von civium unitas [Einigkeit der Bürger] habe.

Er fügt aber die Ursache hinzu, warum mancherlei Same verböten sei, nämlich, „daß du nicht in deiner Fülle heiligest solchen verschiedenen Samen neben dem Einkommen des Weinberges“.\footnote{Hier ist in der Erlanger Ausgabe falsch interpungirt. Das Punctum vor semen ist zu tilgen und nach vineae ein Punkt zu setzen.} „Fülle“ nennt er [den Ertrag] der härteren Früchte, wie er „Thränen“ [den Ertrag] der saftigeren (liquidiorum) Früchte nennt, 2 Mos. 22,\footnote{In den Ausgaben: Exod. 23.} 29.: „Deine Fülle und Thränen sollst du nicht verziehen“ zu opfern.\footnote{Vergleiche Luther’s Randglosse zu dieser Stelle. Malch, St. Louiser Ausg., VIII, 1682.} Also Gott, welcher schlichtes Wesen (simplicitatem) und Aufrichtigkeit liebt, will nicht, daß man ihm Erstlinge und Zehnten von gemengtem Samen opfere. So will er, daß auch seine Bürger schlicht und aufrichtig einerlei Sinnes seien und [einträchtige] Sitten bewahren, daß sie nicht ackern sollen mit einem Ochsen und Esel, nicht ein aus gemengtem Stoffe verfertigtes Kleid gebrauchen, und alles ungleich werde.

Sechstens gebietet er, daß sie Käpplein\footnote{funiculos in fimbris = Schnüre, die zu Fransen gezupft sind.} machen an den vier Fittigen des Mantels. Dies wird den Juden zu tragen auferlegt, als ein besonderes Zeichen, wodurch sie sich von anderen Völkern als durch eine Losung und ein Erkennungszeichen unterscheiden sollten.\footnote{Wir haben mit der Wittenberger discernantur angenommen statt discernuntur in der Jenaer und in der Erlanger.} Nicht, daß der Gesetzgeber dies allein im Auge gehabt habe, sondern es sollte ein Wahrzeichen (Symbolum) des Bekenntnisses sein, wodurch sie bekennen sollten, daß sie das Volk GOttes seien, und daraus der Anlaß entstände, daß sie gehaßt würden und Verfolgung erlitten, damit GOtt durch diesen Spruch (proverbio) des Gesetzes seine Bürger zur Geduld heranbilde, und sie, geübt durch die Widersacher, desto stärker, einträchtiger und einmüthiger würden.

Dies ist gesagt über das Wohlwollen und die Liebe gegen den Nächsten in leiblichen Dingen kraft des siebenten Gebots: „Du sollst nicht stehlen.“ Nun wird er handeln von der Keuschheit kraft des sechsten Gebots. Moses kümmert sich in diesem Capitel nicht um die Reihenfolge der Gebote, gleichwie er im vorigen eher vom Todtschlage gesagt hat, als von dem Ungehorsam der Kinder, denn im gemeinen Leben geschieht und ereignet sich das nicht nach der Ordnung, was sonst wohl der Reihe nach geschrieben werden kann; dies gibt Moses dadurch zu verstehen, daß er diese Ordnung der Gebote vernachlässigt.

Siebentens [gibt er ein Gebot] von einem Weibe, welche von ihrem Manne nicht Jungfrau befunden worden ist, und von ihm in ein böses Geschrei gebracht, entweder mit Wahrheit oder aus Haß etc. Es ist ungewiß, was das für eine Weise (ritus) gewesen sei, um öffentlich die Jungfrauschaft zu erkennen oder zu beweisen. Denn ich sehe, daß sonst alle darin übereinkommen, daß man privatum beim ersten Beilager erfahren könne, ob sie geschwächt oder ungeschwächt sei, weil der erste Beischlaf unangenehm und beschwerlich sei, da das Jungfernhäutchen verletzt wird, und ein Anzeichen davon auf dem Kleide zurückbleibt durch das Blut des verletzten Häutchens, wiewohl Augustinus auch dies ungewiß macht, welcher sagt, daß eben dieses Häutchen bisweilen von den Frauen (matronis) verletzt werde, oder auch von den Mädchen selbst, indem sie es mit dem Finger untersuchen. Aber wir wollen zugeben, es sei gewiß; wie kann es öffentlich beweisen werden? Wird der Mann angehen, die Beiwohnung sei weder unangenehm noch schwer gewesen? Wenn nun das Weib dies leugnet? Wer soll sie des Bluts vorzeigen, welches am Kleide nachgeblieben ist? Wer würde es dulden, daß diese schandbaren Dinge (pudenda) im öffentlichen Gerichte gezeigt werden? Vielmehr will ich sagen, daß ich diese Sitte nicht kenne. Ich möchte wohl denen beistimmen [die da sagen], daß der Vater und die Mutter des Mädchens, wenn sie die Tochter weggeben, zuvor durch ehrbare Frauen ihre Jungfrauschaft erforschen, welche nachher Zeugen sein können vor Gericht, daß sie Jungfrau gewesen sei, und daß es dies sei, was hier Moses sagt, daß der Vater das Kleid ausbreite, indem er sagt: „Siehe, hier ist die Jungfrauschaft meiner Tochter
\quellseite{791}
“, das heißt, durch diese Frauen thue ich dar, daß sie Jungfrau gewesen ist, daß zum Zeichen breite ich das Kleid aus, dessen Ausbreiten bedeuten soll, daß die Sache offenbar und bewiesen sei, wie es bei diesem Volke Sitte war, das Kleid auszuschütteln, wenn sie drohten, wie Nehemia that, Cap. 5,\footnote{In den Ausgaben: „7.“} 13., und das Kleid zu zerreißen, wenn sie zürnten, so auch das Kleid auszubreiten, wenn sie ausdrücken wollten, eine Sache sei klar, deutlich und licht; wiewohl auch diese Sitte, eine Braut vor der Hochzeit durch Frauen unterrichten zu lassen, gar schändbar ist. Doch da wir etwas Anderes nicht haben, wollen wir es unterdessen zulassen.

Auch dies Gesetz dient zur allgemeinen Anwendung wider Verleumder und Ausprenger von bösen Gerüchten. Du siehst, wie schwer diese Uebelthat hier bestraft wird; erstlich, daß er gezüchtiget wird, nämlich mit den vierzig Streichen weniger eins, darnach wird er um hundert Seckel Silbers gebüßt, die dem Vater des Mädchens gegeben werden sollen, endlich wird er auch gezwungen, sie selbst, die er so gehaßt und selbst übel berüchtigt hat, ewiglich zu behalten zur Strafe, während derjenige, welcher ein Mädchen einfach geschändet hatte, dem Vater des Mädchens nur fünfzig Seckel geben, und sie behalten sollte, wie er hernach sagen wird.

Ferner, das Weib, welches nicht Jungfrau erfunden ist, steinigen sie. Durch diesen Beweisgrund wird erkannt, daß dies Gesetz rede von einer Braut, die Hurerei treibt, was ein rechter Ehebruch ist. Denn ein Mädchen, welches schlechte Hurerei treibt, wird ihrem Schänder zur Ehe gegeben, wie ich gesagt habe. Oder wenn es nicht von einer Braut verstanden wird, so wird die Hurerei des Mädchens so strenge bestraft, weil sie heimlich und williglich gehurt hat, und ihren künftigen Ehemann betrogen, und das Haus ihres Vaters entehrt hat. Unter dem Mädchen aber, die ihrem Schänder gegeben werden soll, wird eine solche verstanden, die mit Gewalt geschändet ist. Und dies ist mir das Wahrscheinlichere.

Achtens setzt er eine dreifache Art des Ehebruchs. Erstlich, wenn ein Mann bei eines anderen Weibe schläft, so sollen sie beide sterben. Zweitens, wenn jemand bei eines anderen Braut schläft, aber in der Stadt, so sollen sie beide sterben. Unter Stadt wird hier jeder Ort verstanden, wo das Mädchen sich durch Schreien hätte schützen können, es aber nicht gethan hat, als wenn sie in einem Walde oder auf dem Felde war und wußte, daß Leute nahe wären, welche sie hören konnten. Denn eine Stadt wird hier nur des Beispiels halben gesetzt, weil in derselben Leute zur Hand sind, um zu helfen. Die also nicht schreit, zeigt an, daß sie sich mit ihrem Willen hat schänden lassen. Drittens, wenn jemand bei der Braut eines anderen schläft, aber auf dem Felde, da soll allein der Ehebrecher sterben. Und hier wird das Feld auch nur Beispiels halber genannt, daß darunter ein jeglicher Ort verstanden werde, wo ein einzelnes Mädchen sich durch Geschrei nicht helfen kann, als, wenn jemand sie im Hause oder in der Vorhalle oder in der Kammer allein ergriffe, wie Amnon seine Schwester Thamar; als, wenn jemand die Braut jemandes im Bette durch ein gezogenes Schwert zwänge, nicht zu schreien. Diese ist freilich nicht auf dem Felde, und dennoch ist sie mehr als allein; sie würde schreien, wagt es aber nicht aus Furcht vor dem Tode; darum soll die Billigkeit des Gesetzes es auslegen, als ob sie geschrieen hätte, wie der Text sagt.

Zum letzten ist das Gesetz: Wer ein Mädchen, die nicht Braut ist, schändet, der soll ihrem Vater fünfzig Seckel geben, und soll sie zum Weibe haben. Im zweiten Buche Mosis Cap. 22, 16. nennt er diese fünfzig Seckel eine „Morgengabe“, die er der Jungfrau geben soll; die stellt er hier fest auf so viele Seckel, welche nach unserer Münze ungefähr zwölf Gulden ausmachen, denn ein Seckel ist der Ort eines Guldens.\footnote{quartale aurei — ein Ortgulden, ein Viertelgulden.} Er fügt aber dort hinzu, daß er sie zwar zum Weibe haben soll, [aber nur] wenn der Vater des Mädchens will; wenn er es nicht will, muß er nichtsdestoweniger diese Morgengabe bezahlen. Wollte doch GOtt, daß dieses Gesetz auch bei uns als Exempel gebraucht würde, um des elenden weiblichen Geschlechtes willen.\footnote{Der letzte Satz fehlt in der Wittenberger.}

\subsection{Heimliche Beratung}

Das Thier des Bruders, welches irre geht, oder die verlorene Sache, die man wiedergeben und sich derselben annehmen soll, als wäre es die eigene, bedeutet, daß wir das suchen sollen,
\quellseite{792}
maß Christi ist, wie er selbst auch gethan hat (sei es nun ein Esel, das ist, ein Langsamer, Unschlachtiger und Schwacher im Glauben oder Leben, oder irgend etwas Anderes, als da sind Irrthümer und Mängel,)\footnote{Diese Klammern sind von uns gesetzt.} im Worte, in Ermahnen, im Wandel.

Ein Weib trägt nicht die Waffen eines Mannes, und ein Mann nicht Weiberkleider, wenn der Glaube nicht verkehrt wird durch Werke, das heißt, wenn den Werken das nicht zugeschrieben wird, was dem Glauben gebührt, und wiederum, wie diejenigen, welche besessen sind von Satzungen der Menschen, welche auf Menschen als ihre Lehrer vertrauen, denen man nur durch die Werke der Liebe dienen muß, was sie nicht thun. Dagegen Gotte, dem sie vertrauen sollen, dienen sie mit selbsterwählten Werken und Diensten, und vertrauen nicht auf sein Wort. Beides aber ist ein Greuel, denn einen Menschen soll man nicht glauben, sondern mit dem Werke dienen, Gotte aber soll man glauben, nicht ihm dienen mit dem Werke. So wird der Mann ein Weib, wenn man den Menschen glaubt; das Weib wird ein Mann, wenn man Gotte mit Werken dient, davon anderswo gesagt ist.

Die Vogelmutter wird frei gelassen, und die Jungen werden auf dem Wege gefangen, das heißt, die Synagoge wird verlassen, welche durch das Gesetz die Apostel und die ersten Gläubigen unter sich hatte, die Christus aufgenommen hat, in diesem Leben sein Wort unter den Heiden auszubreiten. Das Haus, erbaut mit einer Mauer rings um das Dach, ist die Kirche Christi, welche er durch sein Wort so befestigt hat, daß niemandem eine Gefahr oder ein Aergerniß zum Schaden der Seele aus ihr erwachsen kann, sondern daß sie in jeder Hinsicht sicher und getrost gebraucht werden kann.

Es wird der Weinberg nicht mit mancherlei Samen besät, das heißt, in der Kirche ist nicht zugleich die Lehre des Glaubens und der Werke, sondern das reine und einfache Wort des Glaubens. So wird auch nicht zugleich mit einem Ochsen und Esel geackert, das heißt, es ist in der Kirche nicht ein Lehrer, der zugleich den Glauben und die Werke lehre, sondern ein aufrichtiger und treuer Diener in der Wahrheit. So gebrauchen wir auch nicht ein Kleid, welches genennet ist von Wolle und Leinen, indem wir so leben, daß der Glaube unser Leben sei, nicht die Heuchelei der Werke, weil das Gewissen das nicht zugleich ergreifen kann, daß du dich auf den Glauben und auf die Werke gründest, wiewohl das Werk auswendig für den Leib und für den äußerlichen Menschen nothwendig ist.

Es soll also der mancherlei Same getrennt werden, der Esel, das Leinen, gesondert auf sein Feld, und sie sollen nicht vermenget werden, das heißt, es soll so gelehrt, geglaubt, gelebt werden, daß der Glaube allein das Gewissen regiere ohne die Werke, die Werke aber das Fleisch brücken ohne den Glauben, damit so inwendig im Herzen der Geist und die Freiheit sei, das Gesetz aber und das Joch der Werke auswendig am Leibe, daß das Fleisch nicht Freiheit habe, sondern Knechtschaft, der Geist aber Freiheit, nicht Knechtschaft, auch nicht nach beiden Seiten hin entweder [allein] Knechtschaft oder [allein] Freiheit.

Die Zäplein an den vier Zipfeln des Mantels bedeuten das Bekenntniß des Glaubens im äußerlichen Wandel durch die ganze Welt, welches in den vier Weltgegenden ausgebreitet werden soll. An diesem Zeichen sollen die Christen erkannt werden, und den Heiden eine Thorheit, den Juden ein Aergerniß sein.

Eine Ehebrecherin bedeutet immer eine Seele, die Christo untreu wird (fornicantem), indem sie das Wort verläßt und auf Werke vertraut. Darum wird durch das Wort des Geistes sowohl sie selbst getödtet, als auch ihr Mann, das heißt, der Lehrer des Irrthums. Ferner werden hier einige mit ihrem Willen geschändet, und schreien nicht, obgleich sie in der Stadt sind. Diese sind es, die, des Wortes überdrüssig geworden, sich zu den Fabeln kehren, und sich aus freien Stücken Meister holen, da ihnen die Ohren darnach jücken Neues zu hören, fragen auch nicht die Gläubigen, die da sind, um Rath, daß sie befreit werden möchten. Andere sind, die allein auf dem Felde geschändet werden wider ihren Willen. Diese sind vielmehr nicht Ehebrecherinnen. Dazu gehören die, welche inwendig an aufrichtigem Glauben stehen, aber gezwungen werden, auswendig gottlose Lehrer zu hören und unter bösen Leuten zu leben, wie Lot in Sodom, und gottlose Leute unter gottlosen Tyrannen.

Die Jungfrau aber, die nicht vertraut ist,
\quellseite{793}
ist die Synagoge oder irgend eine andere Gemeinschaft (congregatio) ohne das Wort, welche aber doch einen schönen Schein und Geistlichkeit (religiosa) hat in ihrer eigenen Gerechtigkeit. Deren Jungfrauschaft nimmt (violat) Christus durch sein Wort und behält sie als Eheweib, wenn es der Vater will, das heißt, wenn sie ihre väterlichen Satzungen verläßt und verlassen wird. Wenn sie dieselben aber nicht verläßt, so gibt er ihr eine Morgengabe, das heißt, zeitliche Vortheile für jene Gerechtigkeit und das Hören des Worts. Denn man muß eine zwiefache Jungfrau annehmen, wenn die heimliche Deutung Bestand haben soll, eine Jungfrau des Glaubens und der Gnade, und eine Jungfrau des Gesetzes und der Werke. Beide sind in ihrem Gottesdienste (religione) keusch und heilig. Die Schändung der letzteren ist Christo angenehm, die der ersteren ist ein Ehebruch und eine Todsünde. So ist der Mann, welcher die Frau übel berüchtigt, daß sie nicht Jungfrau sei, derselbe Christus, welcher die durch Werke gleisende (religiosam) Jungfrauschaft der Synagoge anklagt durch das Evangelium, daß sie nicht rein im Glauben sei. Aber die Eltern, welche ihm starken Widerstand leisten, bedeuten den Kampf, durch welchen die Gerechtigkeit der Werke mit dem Zeugniß der Schrift und dem ausgebreiteten Kleide vertheidigt wird. Und Christus wird gezüchtigt mit Streichen, das heißt, mit Lästerungen, und wird gebüßt, endlich wird er gezwungen sie zu behalten, weil in seinem Reiche schön scheinende Heuchler sein müssen, welche er ernähren muß, und ihre [der Heuchler] Lehrer unter seinem Namen, die sich nicht überwinden lassen wollen. Wenn sie aber überführt wird, daß sie im Glauben nicht Jungfrau sei (non esse virgo fidei), so wird sie getödtet, das heißt, durch das Wort GOttes getödtet, damit sie lebe im Geiste, nachdem sie den Ruhm der gesetzlichen Jungfrauschaft verloren hat.

\section{Das dreiundzwanzigste Capitel}

Cap. 22, 30. Niemand soll seines Vaters Weib nehmen 2c.\footnote{Im Hebräischen ist dies der erste Vers des 23. Capitels.}

Mosis fährt fort, die Gebote der Liebe zu lehren, und bringt noch mit verschiedenen Gesetzen auf das sechste und siebente Gebot hinsichtlich der Keuschheit und des Friedens und der öffentlichen Ehrbarkeit, bis zum 27. Capitel. Im dritten Buch Mosis Cap. 18, 6. ff. werden die Personen aufgezählt, welche zur Ehe zu nehmen nicht recht, und verboten ist; hier wird nur Eine Person genannt, nämlich die Stiefmutter, welche man nicht nehmen soll, wenn der Vater todt ist, viel weniger, wenn er lebt. Dadurch wird angezeigt, daß man viel weniger nähere Blutsverwandte nehmen soll, als die Schwester, die Tochter, die Mutterschwester (amita) etc. Dennoch glaube ich, daß die Ehe im zweiten Grade nicht selten nachgelassen worden sei, da auch an vielen anderen Orten nicht eine so große Strenge des Gesetzes ist.

Es ist aber wahrscheinlich, daß jener Corinth-er, welchen Paulus 1 Cor. 5, 1. verdammt, da er sah, daß die Freiheit vom Gesetze Mosis gelehrt wurde, dieser Freiheit mißbraucht habe, und wider dieses Gesetz das Weib seines Vaters genommen habe. Daher bringt auch Paulus dort nicht so sehr darauf, daß er wider das Gesetz gehandelt habe, als wider die Natur, indem er sagt, es sei eine solche Hurerei, die sich auch bei den Heiden nicht finde, die von dem Naturgesetz geleitet werden, weil die Natur lehrt, daß man an einer solchen Ehe ein Greuel habe, so daß es vergeblich ist, wenn sich jemand vermessen würde, Freiheit von dem Gesetze vorzuwenden.

Zweitens, ein Beschnittener, dem die Hoden ausgeschnitten oder zerstoßen sind, soll nicht in die Gemeine des HErrn kommen. Hier redet er nicht von der Stätte des Gottesdienstes, sondern von der Gemeine selbst, und in die Gemeine kommen oder eingehen ist etwas Anderes, als in der Gemeine sein und leben. Denn wie sollten doch die vom Hören des Gesetzes und vom Gottesdienste abgehalten werden, welche nicht
\quellseite{794}
das Vermögen haben zu zeugen? da das Gesetz insgemein allen gegeben ist, welche vom Samen Israels sind, und auch die Heiden nicht vom Gesetze und Gottesdienste ausschließt, vielmehr nennt man diese Leute Judengenossen (proselytos), und gar oft wird befohlen, daß man sich der Fremdlinge annehme.

Es ist also „eintreten in der Gemeinde“, ein Amt oder öffentlichen Dienst (magistratu) verwalten, denn er sagt nicht: in die Gemeinde, sondern: in der Gemeinde eintreten, wie auch Christus Joh. 10, 1. diejenigen, welche nicht zur Thür hinein gehen in den Schafstall, Wölfe nennt, das heißt, gottlose Lehrer. Es mögen daher in der Gemeine Verschnittene sein, aber sie sollen keine obrigkeitlichen Aemter bekleiden, wegen dieses körperlichen Gebrechens, damit nicht das Amt geringgeachtet oder geschmäht werde; sodann auch, weil solche Leute insgemein ein feiges Herz haben und weibisch sind. Das obrigkeitliche Amt aber in einem ehrbaren Gemeinwesen erfordert einen Mann, wie Bias sagt: Regiment weist aus, was einer für ein Mann ist.\footnote{Magistratus ostendit virum. Vgl. St. Louiser Ausg., Bd. VII, 1374.}

Drittens, ein Hurenkind,\footnote{Im Text: mamzer; in der Vulgata: mamzer; bei Corbatus No. 539, Schlaginhaufen No. 95 und Bindseil I, 374: manser. Hebräisch: \hebraeisch{ממזר}. Vgl. St. Louiser Ausgabe, Bd. XXII, 1628 und Einleitung, S. 51.} das ist ein von einer Hure Geborener, soll nicht in die Gemeine kommen, das heißt, er soll nicht tauglich erachtet werden, ein Amt oder obrigkeitliche Stellung inne zu haben, wegen der Verachtung und Schande seiner Herkunft und Geburt, auch nach dem zehnten Glied, das heißt, ewiglich, weil auch dies sich nicht geziemt in einem ehrbaren Gemeinwesen. Sodann, weil es ein Sprichwort ist, daß Hurenkinder selten gute Leute sind, und immer ein großes Gebrechen an sich haben. Dies Gesetz wird auch bei vielen Völkern gehalten, besonders in etlichen Reichsstädten (municipiis), so daß sie dieselben auch nicht zu den Handwerken zulassen, vielleicht weil sie ehemals ihre Nichtswürdigkeit erfahren haben. Sodann ist es billig, daß man denen, die in Ehren und ehelich geboren sind, in dem äußerlichen Gepränge und Schein der Welt mehr Ehre erweise, als Unehelichen, damit nicht die Ehe der Hurerei gleich geachtet werde, und die Ehre, die einem Unwürdigen erwiesen wird, Gelegenheit gebe, die Hurerei gering zu achten und die Ehe verächtlich zu halten.

Viertens. Die Ammoniter und Moabiter sollen nicht in die Gemeinde kommen, auch nach dem zehnten Gliede, das heißt, für immer. Diesen Heiden wird nicht ein körperliches Gebrechen oder ein Mangel der Geburt vorgeworfen, sondern ihre Gottlosigkeit und ihr Haß gegen die Kinder Israel, daß sie ihnen nicht allein nicht zuvorkamen mit Wasser und Brod auf dem Wege etc., sondern daß sie auch wünschten, daß sie ganz und gar vernichtet würden, da sie begehrten, daß sie durch Bileam verflucht werden sollten, wodurch sie ihren heidnischen\footnote{gentile odium, das heißt, ihren Haß als Heiden wider das Volk der Juden. Heutzutage würden wir es „Nationalhaß“ nennen.} und unauslöschlichen Haß zeigten, von dem sie nie ablassen würden, wie auch nachher die Historien bezeugen. Aber daß man einen solchen Menschen zum Führer oder zu einer obrigkeitlichen Person mache, von dem man weiß, daß er aus natürlichem Haß dein Verderben suche, ist nicht nur unpassend und ungereimt, sondern auch thöricht und frevelhaft und ein Versuchen GOttes. Aber Ehen mit ihnen einzugehen und andere Verträge, und Verkehr mit ihnen zu haben, wird hier nicht verboten. Denn auch Ruth, das moabitische Weib [Ruth 1, 4.], war ein heiliges und gotteliges Weib. Denn er sagt hier von den Männern, nicht von den Weibern, so daß ich glaube, daß es nicht verboten gewesen sei, weder die Tochter oder Enkelin eines Hurenkindes, noch die angenommene oder vor der Verschneidung eines Verschnittenen gezeugte Tochter zur Ehe zu nehmen.

„Du sollst“, sagt er [nach der Vulg. V. 6.], „mit ihnen keinen Frieden machen, sollst ihnen auch nicht Gutes wünschen dein Lebenlang ewiglich.“ Was ist dies? Verbietet er denn, daß man den Elenden oder Feinden oder Fremdlingen von diesen Völkern nicht wohlthun soll? Hat nicht David den König der Ammoniter Hanon trösten lassen [2 Sam. 10, 1.]? Hat er nicht Frieden gehalten mit seinem Vater? Ich antworte: Man muß auf den Gegenstand sehen, von welchem der Gesetzgeber in der Gegenwart redet. Er redet aber von den Feinden, die von heidnischem und unauslöschlichem Haß glühen, daß dieselben nicht in ein obrigkeitliches Amt gesetzt werden sollen. Es wird verboten (sage ich), solchen Häßern und übelwollenden Leuten den öffentlichen Frieden und das gemeine Wohl in
\quellseite{795}
die Hand zu geben (exhiberi), auf daß ihr Haß nicht für gerecht erklärt, und sie nicht in obrigkeitliche Aemter angenommen werden, damit nicht ihr Volk (natio) einmal die Oberhand erhalte und die Herrschaft gewinne. Aber er verbietet nicht, Privatpersonen Privatwohlthaten zu erweisen, gebietet es vielmehr.

Fünftens. Die Edomiter und die Egypter läßt er nach dem dritten Gliede in die Gemeine des HErrn kommen, denn der Edomiter (sagt er) ist dein Bruder, und der Egypter ist dein Gastfreund gewesen. Deshalb, wiewohl du aus ihnen keinen König nehmen kannst, magst du ihnen doch obrigkeitliche Aemter und höhere Stellen (praefecturas) übertragen, aber nur denen, die nach dem dritten Gliede geboren sind, so daß sie schon ihre Art verloren haben und in das Geblüt Israels übergegangen sind. Das dritte Glied muß man aber von der Zeit an rechnen, da sie zum Judenthum bekehrt und Judengenossen geworden sind; was sollte sonst das dritte Glied sein? Und so siehst du, daß den Ammonitern und Moabitern und den Hurenkindern in Israel die Ehe nicht versagt war, da er sagt, daß sie auch nicht nach dem zehnten Gliede, nämlich, nachdem sie Judengenossen geworden waren, in die Gemeinde kommen sollten. Diese Glieder wären nicht vorhanden, wenn sie nicht Eheleute unter den Israeliten und Judengenossen wären.

Sechstens befiehlt er, daß sie, wenn sie aus dem Lager gehen wider die Feinde, sich hüten sollen vor allem Bösen. Ein dunkler Ausspruch, der erstlich so verstanden werden kann, daß sie, wenn sie wider die Feinde kämpfen wollen, zuerst zusehen sollen, daß sie selbst gut und ohne Schuld seien, damit sie nicht von einem gerechten Gericht GOttes betroffen würden, indem sie selbst einen Balken im Auge haben, und suchen den Splitter aus dem Auge eines anderen zu ziehen, wie es ihnen mit den Philistern widerfuhr, 1 Sam. 4,\footnote{In den Ausgaben: 1 Reg. 5.} 2., und mit den Gibeoniten, Richt. 20.\footnote{In den Jenaer und der Erlanger: Jud. 17.} 21., und anderswo. Denn auch die herrliche Verheißung Cap. 20, 14., da er zusagt, er werde im Kriege mit ihnen sein, sollte sie nicht so sicher machen, daß sie meinen möchten, sie könnten ungestraft böse sein, oder GOtt wolle mit den Bösen sein, oder sie könnten dem Gerichte Gottes entgehen, wenn sie dasselbe thäten, was sie richteten und bekämpften.

Zweitens [kann dieser Ausspruch verstanden werden] auf diese Weise: Wenn du kämpfen willst, so gib Acht, daß du nicht eine böse Sache willst, sondern eine gute und gerechte, denn die Verheißung GOttes erstreckt sich nicht auch auf böse und frevelhaft herbeigezogene Sachen, um Krieg zu führen, da er sagt, er wolle für sie streiten [Cap. 20, 4.].

Aber der folgende Text scheint zu erzwingen, daß es von einer gewissen Kriegszucht (religione militari) und züchtigem Verhalten im Lager während des Krieges verstanden werden müsse, da er selbst auslegt, was das Böse sei; erstlich, wenn jemandem des Nachts etwas widerfahren ist, soll er aus dem Lager hinausgehen, und zurückkehren, nachdem er sich bei Sonnenuntergang gebadet hat; zweitens sollen sie zur Nothdurft aus dem Lager gehen und sich ehrbarlich verhalten aus Ehrerbietung gegen GOtt, der im Lager gegenwärtig ist, damit GOtt, wenn sie dies verachten und unehrerbietig sind, nicht ihnen die Feinde, sondern vielmehr sie den Feinden übergebe. Dies ist also das Böse, vor dem sie sich hüten sollen, wenn sie in den Kampf gehen wollen.

Diese beiden Dinge sind gar klein und gering, ja, natürlich, warum bringt er denn so sehr darauf und macht sie groß? Aus einem gar guten und gerechten Grunde, denn der Krieg ist an sich schon etwas sehr Stürmisches, der nichts weniger leiden will als irgend ein Gesetz oder Ehrbarkeit. Denn siehe die Kriegsleute (milites) an, ob irgend welche wilden Thiere grausamer, zügelloser, unbändiger, frecher und schamloser seien, sowohl in Thaten als Worten, als sie. Denn nichts, was ihnen beliebt, ist ihnen dort verwehrt, wenn du auf Sitten und Ehrbarkeit siehst. Es will also der HErr, daß sein Volk mit der größten Ehrerbietung, mit der größten Ehrbarkeit, mit der größten Zucht Krieg führe. Darum macht er auch Bestimmungen über diese kleinen und natürlichen Dinge, damit sie sich durch Uebung in diesen Dingen und nach deren Exempel viel mehr hüten, und einen Greuel haben möchten an Thaten und Worten, in denen sich eine wollüstige und gröbere Unehrerbietigkeit, Schamlosigkeit und Zuchtlosigkeit kundgibt in anderen Dingen, die nicht nothwendig und nicht natürlich sind. Denn wer, weil ihm des Nachts
\quellseite{796}
etwas widerfahren ist, den HErrn scheut und aus dem Lager geht, wie sollte der bleiben, wenn er ein böses Gewissen hätte wegen Hurerei, Diebstahls, Fluchens oder einer anderen Sünde? Und wer den natürlichen Gang mit solcher Ehrerbietung selbst vor dem Lager verbirgt, wie sollte der sich unterstehen, im Lager irgend etwas in schamloser Weise zu thun oder zu reden? Daher treibt er das rohe und hartnäckige Volk durch dieses Gesetz zur Furcht GOttes beim Kriegführen, damit es keine Vermessenheit und Unverschämtheit ablege, und mit Furcht im Glauben an GOtt streite.

\subsection{Siebentens}

Sie sollen einen entlaufenen Knecht seinem Herrn nicht ausliefern, sondern ihn bei sich wohnen lassen, wo es ihm gefällt, und ihn nicht schinden. Auch dies ist dunkel, denn man darf das Eigenthum (rem) eines anderen nicht behalten. Ein Knecht aber ist das Eigenthum des Nächsten, ebenso wie ein Ochse oder Esel, den man auch wiedergeben muß, wenn man ihn gefunden hat. Es muß dies also verstanden werden von einem grausamen Herrn, der den Knecht sucht, um ihn zu tödten (ad mortem). Denn es deutet auf Grausamkeit hin, daß er verbietet, ihn zu „überantworten“ (tradi), aber daß man jemanden in die Hände eines anderen überantwortet, hat insgemein die Bedeutung, daß er zum Tode überantwortet werde. Er will also, daß den Knechten das Leben erhalten werde, und ihnen Nahrung gegeben werde, damit in einem wohlgeordneten Gemeinwesen den Herren nicht alles freistehe. Und um sein Leben zu erhalten, durfte er entfliehen und frei werden, da ihm auch für ein ausgeschlagenes Auge die Freiheit gegeben werden mußte, 2 Mos. 21,\footnote{In den Ausgaben: Exo. 22.} 26. Denn so gebührt es sich für ein gutes Gemeinwesen, daß man den Knechten Leben und Nahrung lasse.

\subsection{Achtens}

„Es soll keine Hure sein unter den Töchtern Israel, und kein Hurer unter den Söhnen Israel.“ Er redet von Hurenwirthen und -wirthinnen, oder öffentlichen und schamlosen Huren, welche sich vom Hurengewerbe ernähren, weil auch diese öffentliche Schamlosigkeit und Zügellosigkeit böser Leute sich nicht ziemt für ein wohlgeordnetes Gemeinwesen, welches durch ehrbare Ehen und rechtmäßige (liberali) Nachkommenschaft wachsen muß. Läßt denn nun Moses Huren aus anderen Völkern zu? Das scheint der Text zu verstehen zu geben.\footnote{Statt der Worte: Sic videtur textus innuere, was sich in der Jenaer und in der Erlanger findet, auch in der Straßburger Ausgabe von 1525, hat die Wittenberger die Worte: Non credo [Ich glaube es nicht]. Daß die letztere Lesart Luthers Sinne entspricht, daran ist kein Zweifel. Vielleicht ist hier im ersten Druck etwas (durch Versehen) ausgefallen.}

\subsection{Neuntens}

Du sollst weder Hurenlohn noch Hundegeld opfern. Ein wunderliches Gesetz! Freilich der erste Theil, von dem Hurenlohn, wird, wie man leicht erkennt, deshalb gegeben, damit öffentliche Schandbarkeit verabscheut werde. Denn GOtt haßt räuberische Brandopfer [Jes. 61, 8.], so auch alles, was mit Sünde und Gottlosigkeit erworben ist, da er will, daß man ihm mit Ehrerbietung diene. Darum soll das nicht in das Haus des HErrn gebracht werden, sei es als Opfer oder zum Bau, woran Gemeinheit, und Schande oder Schimpf haftet, damit nicht GOtt auch gering geachtet werde, der es duldet, daß ihm mit so schlechten und schändlichen Dingen gedient werde, wie der große Haufe leicht diejenigen zu verachten pflegt, von denen sie sehen, daß sie an verächtlichen Dingen Gefallen haben. Denn er will durchaus, daß der rohe Haufe durch einen äußeren Schein (larva) erfaßt, und in Furcht und Ehrerbietung gehalten werde.

Und dies ist der Grund, warum auch das Geld, mit dem ein Hund gekauft wird, nicht geopfert werden soll, denn ein Hund wird in der Schrift für das verächtlichste Thier gehalten, so daß er auch zu Sprüchwörtern dient, wie Abner sagt [2 Sam. 3, 8.]: „Bin ich denn ein Hundskopf?“ nämlich nichts Verächtlicheres hätte er anführen können. Wer also mit dem Erlös (pretio) eines Hundes GOtt ehrt, der verachtet GOtt und hält ihn gering, als wenn er das für Mist gelöste Geld Gott opfern wollte. Diese Zeichen machen kund, eine wie geringe Meinung das Herz von GOtt habe.

Es könnte alles, was bisher gesagt ist, zu dem Vorigen gehören, da er lehrte, daß diejenigen, welche kämpfen wollen, sich vor Bösem hüten sollten, nämlich, daß hierdurch ein ehrerbietiges Verhalten im Lager erfordert werde, und daß er durch den Brauch und das Exempel dieser beiden schändlichen und verächtlichen Dinge insgemein allen schändlichen und niedrigen Gewinn
\quellseite{797}
und Geld dem Volke zu einem Greuel machen wolle. Denn es wird das Geringere und Kleinere verboten, damit erkannt werde, daß Großes und Bedeutenderes noch mehr verboten sei. Denn wer nicht wagt, Hundegeld oder Hurenlohn zu opfern, von denen das eine ohne Sünde, das andere zwar mit Sünde, aber doch ohne dem Nächsten Schaden zu thun (sine injuria), erworben ist, wie sollte der Gotte so gering achten, daß er wagen sollte, das zu opfern, was er durch Diebstahl, Raub, List, Betrug, Ungerechtigkeit an sich gebracht hat?

Zehntens. Daß man an einem Bruder nicht wuchern solle, sondern an einem Fremden, von diesem Gesetze ist oben [Cap. 15, 16.] genug gesagt worden, wo wir die Verheißung behandelt haben,\footnote{In der Erlanger: tractamus statt: tractavimus in den anderen Ausgaben.} daß die Juden Wucher treiben würden. Daraus wird leicht erkannt werden, wie den Juden der Wucher an Heiden und Fremden erlaubt gewesen sei, nicht durch ihr Verdienst oder Recht, sondern durch den Willen GOttes.

Elftens gebietet er, daß ein Gelübde gehalten werden solle, sonst wird es als Sünde angerechnet werden. Ueber Gelübde haben wir anderswo überflüssig geredet, jetzt wollen wir nur kurz sagen, daß zwei Stücke bei einem Gelübde sein müssen, damit es gottseelig sei und [so beschaffen, daß es] gehalten werden müsse. Erstlich, daß es nicht wider die Gottesfurcht sei, oder daß es nicht ein böses Gelübde sei. Denn in solchen Zusagen muß man sein Versprechen nicht halten, als, wenn du ein gutes Werk gelobst, um dadurch gerecht zu werden und Sünden zu tilgen, so ist es ein gottloses und böses Gelübde, weil es wider den Glauben an GOtt ist, der allein gerecht macht und die Sünden wegnimmt, Röm. 1, 16. [Cap. 3, 28.] Derartig sind die Gelübde aller Ordensleute heutzutage in der ganzen Welt. Wenn du aber etwas solcher Art gelobst, allein zur Ehre GOttes und frei, oder dem Nächsten zu Nutz, so ist es gottseelig und muß gehalten werden. Zweitens, wenn du ein Gelübde thust, in Dingen, die nicht dein sind, oder sonst unmöglich, als, wenn du gelobtest, dem Kaiser ein Königreich zu geben, während du doch ein Bettler wärest, oder die Sonne zu verfinstern. Derartig ist das Gelübde der Keuschheit, weil Keuschheit in keines Menschen Gewalt steht, sondern allein Gottes.

Zwölftens. Im Weinberg und in der Saat des Nächsten mag man die Früchte essen, aber von dort nichts hinweggetragen. Dies ist ein Gesetz über bürgerliche Freundschaft, welches anzeigt, daß einer dem andern förderlich und dienstlich sein soll, gleichsam als wollte er lehren, daß unter Freunden alles gemeinsam sei, doch in solcher Weise, daß der andere nicht beschwert werde, oder sich dessen, was sein ist, ein anderer bemächtige und es an sich reiße.

\subsection{Heimliche Wartung}

Daß man des Vaters Weib nicht nehmen soll, bedeutet, daß man die Synagoge und jegliche Secte fahren lassen soll, welche uns durch das Gesetz und durch Werke erzogen hat, und daß sie nicht verbunden werden soll mit dem Wort des Lebens. Dies hat auch Christus gethan, und wir sollen es ihm nach thun, weil der Glaube es nicht leidet, daß die Gerechtigkeit der Werke bleibe und mit ihm verbunden werde. Die Synagoge aber hatte unserm Vater, das heißt, das Gesetz, welches Paulus Röm. 7, 2. den Mann nennt etc.

Ein Eunuch, oder der, welchem die Hoden ausgeschnitten oder zerstoßen sind, ist der, welcher von der heiligen Schrift keine Kenntniß hat, oder eine verletzte und durch die verkehrte Meinung von den Werken zerstoßene. Ein solcher muß nicht in das Amt des Wortes gesetzt werden, da er untauglich ist, da er Hoden haben sollte, das ist, das Gesetz und die Propheten, welche Zeugniß ablegen von der Gerechtigkeit des Glaubens, Röm. 3, 21., daß er mächtig sei in der heiligen Schrift, Tit. 1, 9.

Das Hurenkind ist noch viel ärger; es ist geboren von einer Hure [, das ist ein solcher Mensch], welcher ganz und gar in weltlichen Satzungen der Menschen und in fleischlichen Meinungen steckt, das heißt, er wird noch viel weniger jemals zum Dienste des Worts geschickt sein, sondern er ist auch eine Schmach für die Kirche, und nimmt die Ehre weg, die wir von Christo haben. Was sind aber jetzt die Bischöfe in der ganzen Welt anders als Hurenkinder? Sie sind nicht bloß Verschnittene.

\quellseite{798}
Die Moabiter und Ammoniter sind der geistlichen Deutung nach recht eigentlich die Juden, und nach ihnen die Ketzer. Denn beide sind die allerbittersten Feinde der Kirche, welche der von der Welt ausgehenden und leidenden Kirche hätten mit Wasser und Brod entgegen kommen sollen, um ihr zu helfen und sie zu trösten, als ihre Hausgenossen und Gastfreunde. Sie stehen ihr aber nicht allein nicht bei, sondern sie setzen auch noch den Bileam, falsche Lehrer, auf, welche mit lästerlichen Lehren den rechten Weg der Kirche verfluchen. So wünschen sie mit dem glühendsten Hasse, daß die Kirche vertilgt werde nach Leib und Seele. Darum sind sie niemals, fort und fort nicht, ins Predigtamt zu setzen.

Dazu stimmen auch die Namen. Ammon heißt so viel als: ein trauriges Volk, weil solche Werklehrer traurige und betrübte Gewissen machen. Moab heißt so viel als: vom Vater, denn sie sind äußerlich stolz, und rühmen sich, daß sie allein aus GOtt sind und daß ihre Dinge recht und wahr seien, wie denn Moab in den Propheten gestraft wird als stolz und anmaßend, und Ammon als bitter und grausam. Dazu reimt sich auch ihre Geburt, daß sie vom Vater herkommen aus blutschänderischer Verbindung. Der Vater ist das Gesetz, aus welchem sie durch unrechten Gebrauch geboren werden, weil sie das Gesetz so nehmen, als erfordere es Werke, nicht daß es die Kraft der Sünde sei, wie Paulus es nennt [1 Cor. 15, 56.]: „Die Kraft der Sünde ist das Gesetz“; jene aber würden sagen: Die Kraft der Gerechtigkeit ist das Gesetz. Bileam heißt so viel als: ihr Verschlinger, weil die gottlosen Lehrer immer ärger werden und viele verschlingen; dazu verschlingt der Bauch auch ihr Hab und Gut, und sie verderben die Leiber durch das schwere Joch der Werke.

Solche Dinge leiden wir auch heutzutage. Diejenigen, welche uns beistehen sollten, wo wir über dem Evangelio leiden, gehen selbst davon und verlassen uns nicht allein, sondern verfluchen uns noch obendrein durch ihre lästerlichen Lehrer. Aber GOtt erhört sie nicht, ja, er wandelt ihren Fluch in Segen, daß wir desto mehr zunehmen durch ihr Wüthen. Darum muß man nicht Frieden mit ihnen machen, noch ihnen Gutes wünschen, nämlich so, daß ihr Wüthen Gedeihen habe und glücklich hinausgehe, sondern man muß dawider beten, daß sie fallen von ihrem Vornehmen und ausgestoßen werden, V. 5, 11. Sie werden aber Söhne „Beor“ genannt, weil dies einen Thoren bedeutet. Denn sie reden in der That Thörichtes, und widersprechen sich selbst, wie auch Bileam that, da er segnete, wider seinen Wunsch [4 Mos. 23, 20.]. Denn es ist kein Ketzer, der von verständigen Männern nicht in seinen Lügen ergriffen werde, daß er wider sich selbst rede.

Die Edomiter sind Brüder, aber schwach im Glauben, und fallen bisweilen in einen Wahn. Diese sollen nach dem dritten Gliede zugelassen werden, das heißt, wenn sie gewachsen sind und mit uns stärker im Glauben geworden sind. So sind die Egypter unsere Gastfreunde, welche schwach sind im Leben. Keine von beiden sind wegen ihrer Schwachheit für Greuel zu achten, sondern man soll sie annehmen, so daß sie endlich auch lehren mögen.

Im Kampfe wider die Feinde soll man sich vor allem Bösen hüten, so daß, wenn jemandem Nachts etwas widerfahren ist, er aus dem Lager hinausgehen soll etc., das heißt, wenn im Streit und in der Aufsechtung aus Furcht vor dem Tode ihm das Wort entfallen ist, daß er dasselbe entweder verleugnet, oder nicht recht handelt, der soll sich demüthigen und sich von der Kirche absondern, bis daß die Sonne untergehe und die Auferstehung ein Ende habe; dann bade er sich, und nachdem er seine Sünde bekannt hat, kehre er zu seiner Pflicht zurück. Denn GOttes Wort muß mit Zucht und Ehrerbietung gehandelt werden, und es muß wider die Feinde gekämpft werden, damit wir nicht durch Vermessenheit und Mangel an Ehrerbietung das Wort verlieren, und, wo GOtt uns verläßt, überwunden werden.

So soll auch der Auswurf der Nothdurft außerhalb des Lagers mit einem Schäuflein eingegraben und verborgen werden, das heißt, die natürlichen Gebrechen, welche im sündlichen Fleische noch übrig sind, soll man tödten und durch das Wort des Glaubens verbergen, damit GOtt sie nicht sehe. Sie werden aber verborgen, wenn wir sie in unserem Leibe nicht herrschen lassen. Solcher Weise muß also ein Diener des Worts sowohl die Lende des Gemüths als auch die Lenden des Fleisches umgürtet haben, daß er gesund sei im Glauben des Wortes, und heilig, nachdem das Fleisch getödtet ist.

Der Knecht soll seinem Herrn nicht überantmordet
\quellseite{799}
werden; das Fleisch soll auf solche Weise gedämpft werden (mortificanda), daß es nicht getödtet werde (occidatur). Denn das Gesetz, der Herr des Fleisches, verdammt es ganz und gar und will es todt haben. Aber man muß Rücksicht auf das Fleisch nehmen, daß es lebe, jedoch nicht, daß es sündige, wiewohl es ohne böse Lust nicht leben kann, aber es soll nicht regieren. Denn damit wir Gottes wirken mögen und den Glauben lehren, so ist es vonnöthen, daß das Fleisch lebe, wiewohl es nicht ohne Gebrechen leben kann, welche durch das Schäflein verborgen werden. So lehrt auch Paulus die Colosser [Cap. 2, 23.], daß dem Fleische, als einem schwachen Gefäß, seine Ehre gebühre.

Die Hure und der Hurer bedeuten die gottlosen Heiligen in heuchlerischem Gottesdienst, wie es denn dem Buchstaben nach scheint, daß im Volke der Juden gewisse Secten im Gottesdienste gewesen seien, welche \hebraeisch{קדשים}\footnote{In den lateinischen Ausgaben: codesim.} genannt wurden, wie wir von Josaphat, dem Könige Juda, lesen, der sie von dem Tempel hinwegtat [2 Kön. 23, 7.]. Es bedeutet aber \hebraeisch{קדשים} und \hebraeisch{קדשה} Hurer und Huren, wiewohl es nach dem Klang (allusione) des Wortes auch heilige Männer und heilige Weiber bedeutet. Denn \hebraeisch{קדוש} heißt heilig, Jes. 6, 3.: „Heilig, heilig, heilig.“ Deshalb glaube ich, daß durch \hebraeisch{קדשים} die Verehrung des Priapus bezeichnet werde, in welcher sie Unzucht trieben etc., und es sind (wie ich gesagt habe) die Heiligen und Gerechten des Satans unter einem schön scheinenden Gottesdienste.

Der Hurenlohn scheint die zeitliche Herrlichkeit des gottlosen Gottesdienstes zu sein, von dem ich geredet habe. Diese Herrlichkeit soll nicht geopfert werden, das heißt, gepredigt werden, als ob sie Gotte gefalle oder uns gezieme, da der Ruhm und der Lohn unserer Gerechtigkeit ein gutes Gewissen gegen GOtt ist, und auswendig Kreuz und Schande vor der Welt. Beides fliehen jene Hurer und suchen Ruhm vor der Welt, aber inwendig sind sie schlecht und schandbar wie die Hunde, und darum ist ein solcher Lohn Gotte ein Greuel, das ist, ein verworfenes und böses Gewissen.

Der Bruder soll nicht mit Wucher bedrückt werden, sondern der Fremde, das heißt, dem Gerechten ist kein Gesetz gegeben, welches Werke erfordert, sondern dem Ungerechten, damit die Gottlosen durch die Erkenntniß der Sünde gedemüthigt werden. Denn das Amt des Gesetzes ist, daß es Wucher fordere, mehr als wir vermögen, bis daß wir nach dem Glauben seufzen und, nachdem wir Brüder geworden sind, vom Wucher befreit werden, Ps. 72, 12. ff.

Die Gelübde lassen nicht leicht eine heimliche Deutung zu, weil dies ein sittliches Werk ist; doch kann gesagt werden, es sei überhaupt nur ein Gelübde, welches wir Gotte schuldig sind, nämlich daß wir glauben, ihn preisen und loben, daß er unser GOtt ist. Dies Bekenntniß und Lob wird bedeutet durch alle äußerlichen Gelübde, welche Gotte geleistet werden. Ueber dies Gelübde sagt der 50. Psalm, V. 14.: „Opfere GOtt Dank, und bezahle dem Höchsten deine Gelübde“, und der 116. Psalm, V. 14.: „Ich will meine Gelübde dem HErrn bezahlen, vor allem seinem Volk“, so daß „Gelübde“ ebendaselbe ist als das Opfer des Lobes und des Dankes; denn dies Eine sind wir GOtt schuldig; wir haben es versprochen und können es halten.

Von dem Weinberge des Nächsten essen und nichts hinwegtragen bedeutet, daß sich ein jeglicher Gerechte der Güter des anderen Gerechten im Geiste bemächtige. Dies ist die Gemeinschaft der Heiligen, durch die es geschieht, daß alles allen gemeinsam ist, und doch durch den Gebrauch und gegenseitigen Nutzen dem anderen nichts genommen oder gemindert wird, sondern ganz gelassen wird, ja, durch den Gebrauch und die Gemeinschaft zunimmt. So viel fehlt daran, daß von derselben etwas weggetragen oder weggenommen werden könnte.

\quellseite{800}
\section{Das vierundzwanzigste Capitel}

\begin{schriftzitat}
V. 1. ff. Wenn jemand ein Weib nimmt und ehelicht sie etc.
\end{schriftzitat}

Daß das Gesetz der Ehescheidung ein bloß weltliches sei, und nachgelassen um der Härtigkeit des gemeinen Volkes willen, beweist Christus genugsam Matth. 19, 8. Daher lesen wir auch kein Exempel dieses Gesetzes, welches uns heilige Männer gelassen hätten. Denn auch Abraham ließ seine Hagar erst dann von sich, als GOtt es ihm von neuem gebot [1 Mos. 21, 10–12.]. In einem Gemeinwesen aber muß vieles nachgelassen werden um der Harten und Unbeugsamen willen, damit nicht größeres Uebel geschehe.

Darum muß man unter den Gesetzen einen Unterschied machen, daß etliche seien, die Gutes befehlen, andere, die Böses nachlassen. Unter denen, die Gutes befehlen, ist das erste das Gesetz des Glaubens, welches auch über die Liebe herrscht, in solcher Weise, daß, wenn es der Glaube gegen GOtt erfordern sollte, dem Nächsten die Liebe zu versagen ist, weil GOtt, der durch den Glauben verehrt wird, einem Menschen vorgezogen werden muß, dem durch die Liebe gedient wird. Nach dem Glauben ist es die Liebe, welche alle Gesetze lenkt, sowohl ceremonielle als auch weltliche; doch der Glaube ist ausgenommen.

Zweitens. Ein junger Ehemann soll ein Jahr lang frei sein von allen gemeinen Lasten. Dies ist ein Gesetz der Billigkeit und des bürgerlichen Nutzens. Denn es ist billig, daß man dem Neuvermählten (sponso) mit seinem Weibe ein Jahr gestatte. Es ist nützlich, damit das Gemeinwesen wachse durch Nachkommenschaft und Haushaltungen (familiis), was nicht geschehen würde, wenn der Ehemann sofort nach der Hochzeit dem Kriege und gemeinen Gefahren ausgesetzt würde, und es würden nicht leicht Männer gefunden werden, die sich verheirathen würden. Daß dies aber ein allgemeines Gesetz sei, und auf andere Fälle ausgedehnt werden müsse, erhellt genugsam aus Cap. 20, 5–7., wo ein Bräutigam, und der ein Haus gebaut, und der einen Weinberg gepflanzt hat, für frei erklärt wird vom Kriegsdienste. In gleichem Sinne wird dies auch verstanden von der Uebernahme irgend eines anderen neuen Amtes, wodurch ein Hausstand gegründet wird, so daß dies Gesetz als Sprüchwort dienen könnte von der Freiheit aller, die neu eintreten (neophytorum), nämlich: Einem Bräutigam muß man ein Jahr geben etc.

Drittens. „Du sollst nicht zum Pfande nehmen den untersten und obersten Mühlstein.“ Dies Gesetz lehrt ebenfalls in einem allgemeinen Sprüchworte, daß man um Schulden willen niemandem sein Handwerk, von dem er sich nährt und erhalten wird, legen soll, wie bei uns grausamen Leuten (barbaros) solches Wüthen sich findet, daß wir die Schuldner entweder einkerkern, bis sie den letzten Heller bezahlen, oder ihnen ihr Handwerk oder ihre Arbeit verbieten, während er hier sagt, es sei genug, daß er seine Seele zum Pfande gegeben habe, das heißt, selbst ganz und gar Schuldner sei, und er durch dieses Handwerk das erwerben muß, womit er seine Schuld bezahle. Es ist also grausam und unbillig, ihm sein Handwerk zu hemmen, nämlich, daß er gezwungen wird, die einfache Schuld mit zwiefachem Schaden zu bezahlen; der eine ist, daß seine Mühle unterdessen müßig steht, der andere, daß er anderswoher neue Schulden auf sich ladet, oder das Seine verkauft, um zu bezahlen.

Viertens. Wer einen aus Israel gestohlen hat, daß er ihn entweder verkaufe oder schände, das heißt (wie mir oben [Cap. 21, 14.] dasselbe Wort ausgelegt haben), um Gewinnes willen verdingt, oder einem anderen überantwortet, der soll getödtet werden. Dies Gesetz bezieht sich auf Menschendiebe, Mörder und Räuber, und diejenigen, welche Menschen fangen, um Geld zu erpressen, wie bei uns die Geier und Wölfe der Kaufleute sind, die man Edelleute nennt.

Fünftens. Vor dem Aussatz soll man sich hüten nach der Lehre der Priester, der Leviten; nur sollen sie lehren, wie ich (sagt er) ihnen geboten habe.\footnote{Nach der Vulgata: juxta id, quod praecepi eis.} Das heißt, die Aussätzigen sollen von dem gemeinen Verkehr der Bürger abgesondert
\quellseite{801}
werden, weil der Aussatz ansteckend ist, und ein Gemeinwesen ein reines und kräftiges Volk vonnöthen hat. Er gibt das Beispiel der Miriam, der Schwester Mosis, welche auf Befehl des HErrn ausgeschlossen worden ist [4 Mos. 12, 14.]. Wenn nun ihrer nicht verschont worden ist, daß sie im Lager hätte bleiben dürfen, so soll keines Menschen verschont werden. Andere wollen es so verstehen: Habe Acht, daß du nicht eine solche Sünde begehest, daß du mit dem Aussatz geschlagen werdest, wie Miriam damit geschlagen wurde, oder wenn du damit geschlagen wirst, daß du thuest, wie es sich für einen Aussätzigen gebührt gegenüber den Leviten.

Sechstens. Wer eine Schuld einfordert, der soll nicht in das Haus eingehen, nämlich um zu nehmen, was ihm beliebt, denn das ist Gewalthätigkeit (wie nun unsere grausamen Wucherer\footnote{danistae von δανειστής. Die Wittenberger und die Jenaer Ausgabe scheinen dies Wort für einen Eigennamen, gebildet von „Dan“, gehalten zu haben, denn sie schreiben es mit einem großen Anfangsbuchstaben.} mit den Schuldnern verfahren, indem sie ihnen durch richterliche Gewalt rauben, was sie wollen), sondern das, was der Schuldner nach seinem Vermögen herausbringt, das soll er nehmen. Sodann soll es wiedergegeben werden vor Sonnenuntergang, wenn er arm sein und dessen bedürfen sollte. Es läßt dies Gesetz freilich zu, daß Seine wiederzufordern, was das Evangelium verbietet,\footnote{Die Worte: quod evangelium prohibet fehlen in der Wittenberger.} wiewohl es der Gewalthätigkeit wehrt. Aber in einem bürgerlichen Volke sind solche Gesetze, die etwas nachlassen, nothwendig, um größeres Uebel zu vermeiden. Denn wenn das Wiederfordern nicht stattfände, so würde die Nichtswürdigkeit müßiger Leute Nahrung finden, und das Vermögen der Gottseeligen und Guten würde erschöpft werden. Wiewohl der Gottseelige im Glauben beides leidet, so darf dies doch das Gesetz und die Obrigkeit nicht dulden, gleichwie der Gottseelige auch alles andere Uebel trägt, und das Schwert ist dennoch ein Rächer wider die Bösen, welches auch\footnote{Statt tamen in den Ausgaben mag wohl etiam zu lesen sein, was wir angenommen haben.} ein Gottseeliger für andere gebrauchen kann, wie ich anderswo gesagt habe.

Siebentens.\footnote{Hier hat die Wittenberger wieder: Sexto und darnach im folgenden falsche Zählung.} Man soll dem Armen seinen Lohn nicht vorbehalten, damit er nicht zum HErrn schreie etc. Auch dieses und das eben vorhergehende Gesetz zeigt genugsam, wie GOtt sich der Armen und Betrübten annimmt, so daß alle, die in der Welt verlassen und unterdrückt sind, sich in GOtt trösten können. Zugleich lehren sie, daß in der Noth den Armen auch die Schulden erlassen werden sollen, und daß man ihnen in einem wohlgeordneten Gemeinwesen zu Hülfe komme, so daß ein Einschreiten (actio) wegen Schuldforderungen nur stattfindet wider die Müßiggänger und Leichtfertigen.

Achtens. „Die Väter sollen nicht für die Kinder sterben“ etc. Das heißt, die Bösen sollen so gestraft werden, daß nicht die Guten auch zugleich verletzt werden, wie im Evangelio verboten wird, das Unkraut auszuraufen, um des Weizens willen, damit dieser nicht verletzt werde. Denn wenn ein Schuldiger nicht bestraft werden kann, ohne daß zugleich ein Unschuldiger verletzt werde, so ist der Böse vielmehr zu dulden, als zugleich der Gute zu verletzen. Dies Gesetz ist für die obrigkeitlichen Personen nothwendig, damit sie nicht in einer stürmischen Erregung des Gemüths die Strenge des Gesetzes sollten lassen, und durch Strafen mehr sündigen als der, welcher die Strafe verdient. Derartig ist es, wenn wegen Eines Bösewichts ein ganzes Gemeinwesen in Krieg verwickelt wird, und ähnliche Dinge.

Neuntens. „Du sollst das Recht des Fremdlings und des Waisen nicht beugen, und sollst der Wittwe nicht das Kleid zum Pfande nehmen.“ Alle diese Gesetze, wie du siehst, werden geordnet um der Liebe und des gegenseitigen Wohlwollens willen. Daß das Recht des Waisen nicht gebeugt werde, das schließt zugleich ein, daß einem Waisen oder einem Fremdling kein Schade zugefügt werden soll. Und das Kleid der Wittwe bedeutet auch ganz allgemein alles, was sie zum Leben nöthig hat. Höher als alle Schulden und Rechte steht das Gesetz, daß dem Armen die nothwendige Nahrung und Kleidung gebührt, die er nicht fahren lassen soll, damit der Gläubiger bezahlt werde, sondern der Gläubiger soll nachlassen und geben, damit der Arme leben könne. Und er fügt die Ursache hinzu [3. 18.]: „Du sollst gedenken, daß du Knecht in Egypten gewesen bist“, als ob er sagen wollte: Wie du wolltest, daß man
\quellseite{802}
dir thäte. Denn wie du ein Knecht gewesen bist, so kannst du wiederum zum Knechte gemacht werden. Also lerne an deinem eigenen Exempel, was du ihnen schuldig seiest und thun sollest.

Zehntens. Das, was übrig bleibt auf dem Acker, im Weinberge, im Oelgarten, soll man nicht sammeln, sondern den Armen lassen. Dies bestätiget, was schon gesagt ist, daß man den Armen nicht bloß Nahrung und Kleidung lassen soll, sondern ihnen auch noch von unserem Ueberfluß dazugeben, damit sie in ihrer Noth an unseren Gütern theilnehmen, indem wir wiederum daran gedenken, was wir wollten, daß uns in der Noth geschehe, wie er hier sagt: „Gedenke, daß auch du ein Knecht gewesen bist in Egypten“, und es wiederum geschehen mag, daß du dienstbar sein mußt. Ich glaube, daß Moses hier gar reichlich und klar das siebente Gebot gehandelt hat: „Du sollst nicht stehlen“, wo verboten wird, Schaden zu thun an äußeren Gütern, und geboten wird, auf den Nutzen des Nächsten zu sehen.

\subsection{Heimliche Deutung}

Die verstoßene Frau ist die Synagoge und alle Secten der Werkheiligen. Die „Unlust“ oder die Garstigkeit ist die schön scheinende Gerechtigkeit der Werke, die voll ist von Gottlosigkeit und Här tigkeit des Herzens. Der Scheidebrief selbst ist das Zeugniß der Schrift, welchen Christus ihr in die Hand gibt, welchen sie auch wider sich selbst trägt. Der zweite Mann ist das Gesetz. Wenn dieser gestorben ist, das heißt, nachdem das Gesetz abgethan ist, so nimmt er [der erste Mann, Christus] sie auch so nicht wieder an, wenn sie nicht glaubt; denn die Synagoge kommt nie wieder zu ihrer ersten Herrlichkeit, wird auch nie wieder von ihrem ersten Manne angenommen, nachdem sie den anderen geheirathet hat, wenngleich er stirbt, weil Christus die Gerechtigkeit der Werke nicht neben dem Glauben an ihn zuläßt.

Der neue Ehemann ist ein Neuling im Glauben, der nicht in das öffentliche Predigtamt gesetzt werden soll, wie Paulus [1 Tim. 3, 6.] sagt. Es kann auch Christum bedeuten, der die neuen Gläubigen freundlich regiert und sie nicht alsbald dem Kreuze und der Verfolgung aussetzt, bis daß sie stark werden. Er ist also mit ihnen ein Jahr lang fröhlich, das heißt, er schenkt ihnen in dem eben erst angenommenen Worte Frieden und Freude, weil die Hochzeitsleute nicht Leid tragen, so lange der Bräutigam bei ihnen ist, Matth. 9,\footnote{In der Jenaer und der Erlanger: Matth. 12.} 15.

Man sagt richtig, daß der untere und der obere Mühlstein Hoffnung und Furcht bedeuten, oder das Gesetz und das Evangelium. Denn das Gesetz, das Wort des Zorns, ist der obere Mühlstein; es erhält die Furcht, und demüthigt durch Erkenntniß der Sünde. Das Evangelium, das Wort der Gnade, erhält durch den Glauben das Gewissen, daß es nicht in Aengsten stehe. Beiderlei Predigt ist nothwendig. In einem allgemeinen Bilde aber bedeutet es, daß nach beiden Seiten hin nicht etwas zu viel geschehe, so daß du die Sünde weder allzu hart bestrafest, noch auch den Sünder allzu liebreich behandelst, ihn weder zu sehr tröstest, noch zu sehr schreckest. Denn die Schuld bedeutet die Sünde, um derentwillen der untere Mühlstein weggenommen wird, wenn allzu sehr gestraft wird; wiederum wird der obere weggenommen, wenn allzugroße Nachsicht geübt wird.

Einen Bürger stehlen bedeutet, durch falsche Lehre (verbo) vom Glauben abwenden und ihn überantworten den Werken und der äußerlichen Gerechtigkeit, welche Lohn sucht, und reiche Liebe macht, das heißt, gottlose Lehrer, Joh. 10, 8.: „Alle, die vor mir gekommen sind, die sind Diebe gewesen“ etc. Er wird aber getödtet durch das Evangelium, damit er lebe, oder es wird ihm der ewige Tod angekündigt.

Es ist ganz klar, daß der Aussatz die Seuche der geistlichen Gottlosigkeit bedeute, darum muß man sich sorgfältig davor hüten nach dem Urtheil und Amt der Leviten. Denn ein wenig Sauerteig versäuert den ganzen Teig [Gal. 5, 9.].

Der die Schuld wiederfordert, empfängt nicht was er will, sondern was der Schuldner ihm heraus trägt, das heißt, der Sünder soll nicht gebüßt werden (mulctandus) nach unserer Willkür, sondern nach dem, was er ertragen kann, und wenn er ganz arm und schwach ist, so muß ihm alles wiedergegeben und nachgelassen werden, um der Schwachheit des Fleisches willen, nach dem Worte (Gal. 6,\footnote{In den Ausgaben: Gal. 5.} 1.): „So ein Mensch etwa von einem Fehler übereilt würde“ etc. „Vor Sonnenuntergang“, das heißt, ehe er verzwei-
\quellseite{803}
Seele oder durch Traurigkeit verzehret werde, wie Paulus [2 Cor. 2, 7.] redet.

Du sollst dem Arbeiter sofort seinen Lohn geben, das heißt, denen, die bedrückt sind durch die Arbeit und ermüdet durch die Last des Gesetzes und gequält durch ein böses Gewissen, sollst du das Wort des Lebens verkündigen und sie ermahnen zum Glauben, sonst ruft er zum Himmel, weil die Bedrückung der Gewissen die größte Tyrannei ist; dieses [Geschrei]\footnote{Zu quam wird etwa vocem clamantium zu ergänzen sein.} erhört GOtt, der versöhnt ist, bald wider die Tyrannen der Seelen, es sei denn, daß kein Zorn da sei, in welchem er die ganze Welt straft.

Der Vater soll nicht für den Sohn sterben, das heißt, wiewohl wir alle in der Sünde Adams sterben, so wird doch diese Sünde erlassen, nachdem das Evangelium offenbar ist. Und wer danach stirbt, der stirbt nicht mehr in fremder, sondern in seiner eigenen Sünde, wie auch Hesek. 18, 2. 4. beweiset, daß es geschehen werde, daß ein jeglicher in seiner eigenen Sünde sterbe. Denn nun werden die Zähne der Kinder nicht mehr stumpf, weil die Väter Heerlinge gegessen haben, weil es jetzt im Geiste zugeht, so daß, wer da stirbt, wegen seines Unglaubens stirbt.

Daß das Uebrige der Ernte und der Weinlese den Armen gelassen werde, bedeutet, da wir durch den Glauben an Christum alles erlangen, daß wir unser Uebriges den Sündern dienen lassen sollen, damit auch sie mit uns leben durch den Glauben und das Wort, welches sie von uns haben. Und wir müssen nicht so im Fleische leben, als ob wir allein glauben und selig werden wollten, sondern wir sollen wissen, daß wir deshalb leben, damit wir anderen auch zum Leben und zur Seligkeit nütze seien.

\section{Das fünfundzwanzigste Capitel}

Erstlich wird ein Gesetz über die Strafen der Uebelthäter gegeben, daß diejenigen, welche nicht mit dem Tode, oder um Geld gestraft werden sollen, sondern mit Schlägen, so gestraft werden nach dem Maß ihrer Missethat, nur soll dies nicht über vierzig [Streiche] hinausgehen. Die Ursache wird angegeben: auf daß dein Bruder nicht verachtet und scheußlich vor euch werde. Er will also, daß die Strafen aufgelegt werden nicht nach unserem Muthwillen oder aus Nachgier, sondern nach einer gemäßigten Regel und mit Rücksichtnahme auf die Ehrbarkeit. Das gebührt sich für ein rechtschaffenes Gemeinwesen.

Zweitens. „Du sollst dem Ochsen, der da drischet, nicht das Maul verbinden“ [, dieses wird geboten], damit sie, geübt durch gütiges Verhalten gegen die Thiere, desto wohlwollender würden gegen die Menschen. Es ist aber ein sprüchwörtlicher Ausspruch, den Paulus 1 Cor. 9, 9. ff. reichlich auslegt, so daß er sagt: „Sorget GOtt für die Ochsen?“ als ob er sagen wollte: Wiewohl GOtt für die Ochsen sorgt, so läßt er dies doch nicht um der Ochsen willen schreiben, da sie nicht lesen können, so daß die Meinung des Paulus ist: Dieser Spruch wird nicht bloß von den Ochsen verstanden, sondern ingemein von allen Arbeitern, daß sie von ihrer Arbeit leben sollen, wie auch Christus sagt [Luc. 10, 7.]: „Ein Arbeiter ist seines Lohnes werth.“ Darum ist derjenige, welcher einem Arbeiter nicht seinen Lebensunterhalt gibt, oder ihm nicht auch allerlei Gutes mittheilt, Gal. (6,\footnote{In den Ausgaben: Gal. 5.} 6., der, welcher dem Ochsen, der da drischt, das Maul verbindet.

Drittens. Das Gesetz, daß man die hinterlassene Frau des Bruders nehmen und Samen erwecken soll für den verstorbenen Bruder, ist aus einem sehr guten Grunde gegeben worden. Erstlich [ist der Grund], welchen der Text anführt, damit die Geschlechter nicht untergingen, sondern gemehret würden, was dazu dient, ein Gemeinwesen zu heben und stärker zu machen. Zweitens, daß GOtt auf diese Weise die Witwen und das elende [weibliche] Geschlecht versorgt, daß man sie aufnehmen und ernähren sollte, da ein alleinstehendes (per se) Weib ein schwaches und elendes Geschöpf ist, und noch viel mehr eine Wittwe, welche zugleich verlassen und verachtet ist.

\quellseite{804}
Er zwingt aber zu dieser Liebesverweisung durch eine überaus große Schmach, daß er ein Barfüßer genannt werden, und man vor ihm ausspieen soll: Pfui dich an! daß er werth sein soll, daß ihn alle verachten und über ihn auf die Erde speien und sprechen: Du bist ein Pfui-dich-an, weil er das Gemeinwesen, in welchem er lebt und dessen Rechte er genießt, nicht erbauen noch mehren will. Der unbeschuhte Fuß aber soll ein Zeichen der Schande sein und eine Ursache ewiger Schmach, daß er werth sei, auch barfüßig zu sein, das heißt, ohne Geschlecht (familia) und Untergebene, die durch den Schuh bedeutet werden, da er sich selbst durch diese Eine That, daß er das Geschlecht seines Bruders nicht erbauen will, zu einem Barfüßer macht, so daß das Zeichen der That ähnlich ist, mit der er sündigt.

Viertens. Dem Weibe soll ohne Barmherzigkeit die Hand abgehauen werden, welche, um ihren Mann zu erretten, einen andern Mann, der mit ihm streitet, bei der Scham ergreift. Fürwahr ein thörichtes Gesetz, wie es lautet. Es will aber dies besagen, daß ein Weib, weil es kleinmüthig ist, da Schaden zu thun pflegt, wo wir am meisten verletzt werden. Zum Beispiel, eine solche Zunge, die bereit (levis) ist, übel zu reden und zu schmähen, verletzt leicht den guten Ruf, da sie vor allem das angreift, dessen wir uns zu schämen haben. Dies bedeutet das leibliche Ergreifen der Scham, wodurch ebenfalls die Person des Mannes sehr verletzt und vergewaltigt wird (cogitur), und dennoch ist es eine weibliche Gewaltthätigkeit, welche sogar mehr ausrichtet (vincit) als die Stärke eines Mannes; um der Schandbarkeit und Schamlosigkeit zu geschweigen, welche dies Gesetz an einem Weibe auch verdammt, wiewohl dieselbe in diesem Falle dem Weibe und dem Manne nützlich und nothwendig ist, so daß die allgemeine Meinung dieses Gesetzes diese sein könnte, man solle nicht Böses thun, damit Gutes daraus komme.

Fünftens. Es muß in einem Gemeinwesen rechtes Gewicht und rechtes Maß erhalten werden, damit der Arme und der Nächste nicht betrogen werde. Dies hat ebenfalls allgemeine Geltung für jeden Umsatz (commutationes) in allen Händeln (contractuum), daß der Verkäufer rechte und gleiche (aequa) Waare gebe für das Geld des Käufers. Denn hier keimt die Habgucht unglaubliche Nichtswürdigkeiten und Kunstgriffe im Verändern, Verfälschen, Herausputzen und Herrichten der Waaren, so daß es nicht ein geringer Theil der Sorge sein muß, daß die Obrigkeit hierin für das Gemeinwesen wohl Acht habe.

Endlich gebietet er, daß die Amalekiter gänzlich vertilgt werden sollen unter dem Himmel wegen der grausamen Schandthat, daß sie die auf dem Wege Ermüdeten nicht allein nicht erquickten, sondern sogar mit schändlicher Hinterlist die Hintersten tödteten.

Daß er aber befiehlt, die Amalekiter zu vertilgen, thut er nicht, damit sie selbst sich rächen sollen, sondern GOtt rächt sich durch die Hände Israels. Sonst, wenn er gewollt hätte, daß sie ihre eigene Rache suchen sollten, so hätte er es nicht so lange Zeit verschoben, da die schon lange gestorben waren, welche verletzt worden waren, sondern er hätte ihnen sofort zu derselben Zeit gestattet, sich zu rächen. Deshalb wird uns hier gelehrt, daß erst das eine rechte Rache ist, welche nicht nach unserem Muthwillen, sondern auf GOttes Befehl und um GOttes willen ausgeführt wird, so daß wir allein Werkzeuge der Rache GOttes seien, und uns nichts von der Rache anmaßen, wie du hier siehst.

\subsection{Heimliche Deutung}

Das genau bestimmte Maß der Schläge, damit der Bruder nicht scheußlich umhergehe, bedeutet, daß das Fleisch so gezüchtigt werden soll, und die Sünde so gestraft werden, daß seine Laster bedrückt werden, nicht daß es selbst zu Grunde gerichtet werde. Die Zahl „vierzig“ bedeutet auch anderswo alle Stufen und den Inbegriff der Tödtung des alten Menschen in der Schrift.

Daß man dem Ochsen, der da drischet, nicht das Maul verbinden soll, kann außer der Gleichnissrede (metaphoram), von der wir gesagt haben, daß dadurch bedeutet werde, es gebühre den Dienern des Worts der Lebensunterhalt von ihrem Amte, auch noch nach der heimlichen Deutung (allegorice) bedeuten, daß einem Lehrer nicht verboten werden solle, das Wort frei zu predigen, durch welches er Stoppeln und Heu verzehrt, das heißt, die gottlosen Leute bekehrt und Christo einverleibt.

Die hinterlassene Frau des Bruders ist die Synagoge oder die Kirche, welche der gestorbene
\quellseite{805}
Christus in der Welt als eine Unfruchtbare zurückließ, weil er von ihr die Kinder des Gesetzes nicht empfing. Wir nehmen sie daher auf und schwängern sie mit dem Worte, damit sie nicht uns oder unserem Namen, sondern ihm und seinem Namen Kinder gebäre, so daß, wiewohl wir die Leute lehren, wir sie doch nicht anders als im Namen Christi lehren, damit die Kinder, die durch das Wort gezeugt sind, nicht Pauliner, Apolliner oder Petriner, sondern allein Christen heißen, auf daß wir sein Haus bauen, nicht das unsere.

Wer aber dies nicht thun will, den soll die Kirche verklagen als einen solchen, der für sich, nicht für Christum lehren wolle, nicht Gotte, sondern Menschen gefällig sein, die Jünger nicht zu Christo, sondern zu sich selbst ziehen. Daher soll sie ihm öffentlich den Schuh ausziehen, das heißt, ihn der Befugniß zu lehren berauben, und die Erklärung abgeben, daß sein Wort nicht „ein Treiben des Evangeliums“ [Eph. 6, 15.] sei, wie Paulus sagt, auf daß alle wissen, daß er nicht richtig wandele, gestiefelt in der evangelischen Wahrheit, wie Paulus dem Petrus that. Ferner soll sie vor ihm auspeien, damit sie ihn durch den Schimpf des Ketzer namens verächtlich mache, daß er ein Barfüßer genannt werde, das heißt ein Ketzer, der nicht in den Schuhen des Evangelii einhergeht.

Die beiden Männer, die mit einander had ern, sind das Gesetz und das Evangelium, oder vielmehr die Lehrer und Diener des Gesetzes der Werke und des Evangeliums der Gnade. Diese streiten beständig über Werke und Gerechtigkeit. Das Weib, welches die Scham des anderen Mannes ergreift, ist das weise Volk, welches diesen Streit beilegen will dadurch, daß es verbietet zu lehren, und [gebietet] auf einerlei Weise zu lehren, wodurch es den Dienst des Wortes hindert. Deren Hand soll ohne Barmherzigkeit abgehauen werden, das heißt, ihr Werk und Vornehmen soll ganz und gar verworfen werden, weil dieser Streit in Ewigkeit nicht geendig[t] wird, es sei denn, daß GOtt Richter sei. Zorn und Streit soll zwar abgethan werden, wo es sich um die Sache der Liebe handelt, aber in Sachen des Glaubens kann und muß er nie aufgehoben werden, weil Rotten sein müssen. [1 Cor. 11, 19.].

Daß zweierlei Gewicht verboten und gleiches Gewicht geboten wird, das heißt, daß in der Kirche ein aufrichtiges Wesen (simplicitatem) erhalten werde, so daß wir einerlei Sinn und Meinung haben, dessen Paulus oft Erwähnung thut [1 Cor. 13, 11.].

Daß die Amalekiter getödtet werden sollen, bedeutet, daß durch das Wort des Geistes die ermürgt werden sollen, welche die Schwachen im Glauben und die auf dem Wege Gottes Ermat teten und die durch viele Trübsale Ermüdeten nicht allein nicht trösten, sondern obenein noch den Schmerz der Wunden hinzufügen, und den Verstorbenen tödten. Dessen gibt dir der 109. Psalm, V. 16., ein Exempel, indem er sagt: „Darum, daß er so gar keine Barmherzigkeit hatte“ etc., darum [Ps. 69, 29.] sollen sie „mit den Gerechten nicht angeschrieben werden“ etc. Denn es soll ihnen verkündigt werden, daß sie des ewigen Todes werth seien, wie sie denn in der That sind.

\section{Das sechsundzwanzigste Capitel}

\begin{schriftzitat}
V. 1. ff. Wenn du ins Land kommst etc.
\end{schriftzitat}

Nachdem die Erklärung der heiligen zehn Gebote beendigt ist, schließt er in diesem Capitel alle Predigten, die über die Gebote gehalten worden sind, indem er gleichsam als einen Folgerungssatz oder anstatt eines Schlusses der Rede die Gebräuche beim Opfern der Erstlinge und Zehnten hinzufügt, und so den Bund zwischen GOtt und den Kindern Israel bestätigt.

Du siehst aber, wie schön er lehrt, die Worte zu setzen [V. 3. ff.] beim Opfern der Erstlinge, auf daß wir wissen, dieses Opfer sei nichts Anderes als ein Lob und Danksagung für die Erlösung von den Feinden und für den Besitz des Landes, nicht anders, als wenn ein Zinsmann durch seinen Zins seinem Herrn bekennt, daß er seinen Acker durch die Gnade desselben habe. Es waren aber die Erstlinge die [Früchte],
\quellseite{806}
welche sie am Pfingstfeste den Priestern opferten, von denen im 16. Capitel gesagt ist.

So macht er es auch mit den Zehnten im dritten Jahr, indem er lehrt, daß sie den Leviten, den Waisen, den Wittwen und den Fremdlingen gegeben werden sollen, mit dem Bekenntniß, daß sie das Werk der Liebe erfüllt haben. „Ich habe gebracht“ (spricht er [V. 13.]), nicht zum Opfer, sondern an einen besonderen Ort, an welchem in jeder Stadt solche Zehnten für die Armen gesammelt wurden, gleichsam in einen Sack (bursam) oder einen gemeinen Kasten.

„Ich habe deine Gebote nicht übergangen (nämlich, so viel dies Werk anbelangt; sonst, wer erfüllt auch nur Ein Gebot völlig?), habe auch nicht davon gegessen in meinem Leide (das heißt, ich habe es nicht traurig, sondern fröhlich gethan, denn einen fröhlichen Geber hat GOtt lieb), und habe nicht davon gethan in Unreinigkeit“ (das heißt, nicht mit Sünden oder Unreinigkeit, sondern in Heiligkeit habe ich es abgesondert). Denn das Opfer der Unreinen ist ebensowenig [Gotte] gefällig als das Opfer der Leidtragenden. „Ich habe auch nicht zu den Todten davon gegeben“, denn dadurch würde der auch verunreinigt werden, welcher dieses thut, und so würde die Gabe der Zehnten auch verunreinigt sein.

\begin{schriftzitat}
V. 16. Heutiges Tages gebeut dir der HErr, dein Gott etc.
\end{schriftzitat}

Dies sind Worte des bestätigten Bundes, daß sie die vorhergenannten Gebote halten sollen, und er fügt die Verheißung hinzu, daß, wenn sie die Gebote halten würden, der HErr sie hoch machen werde an Ruhm, Preis und Ehre über alle Völker, was er auch gethan hat, so lange sie den Bund hielten. Und dies ist das Ende der heiligen zehn Gebote und [anderer] Gebote. Was nun in den folgenden Capiteln gehandelt wird, sind etliche Verordnungen, welche nach Gelegenheit der Zeit gehalten werden sollen, Ermahnungen, Aufträge, Glücke und Segnungen, um das Volk zu erinnern, daß es die Gebote GOttes halte.

Die heimliche Deutung der Erstlinge ist oben [Cap. 15] genugsam gehandelt worden, daß dieselben bedeuten das Bekenntniß des Glaubens und die Danksagung der Gerechtigkeit des Geistes, da wir zugleich bekennen, wie großem Uebel wir unterworfen gewesen seien, aus dem der HErr uns befreit hat, und wie große Güter es sind, die wir durch den Glauben empfangen haben. Das Bringen der Zehnten aber bedeutet, daß wir durch die Liebe ganz und gar in den Dienst des Nächsten dahingegeben werden, wie Paulus sagt [Gal. 5, 13.]: „Durch die Liebe diene einer dem andern.“ Dies geschieht aber nicht, wenn du nicht zuvor als ein durch den Glauben Gerechtfertigter alle Gebote GOttes erfüllt hast und gerecht bist. Ferner: fröhlich, nicht traurig oder ungern; endlich ohne Unreinigkeit, das heißt, daß du nicht begehrest, daß man dir Gleiches wiedervergeltet, sondern aus einfältigem und reinem Herzen. Du sollst nichts zu den Todten geben, das heißt, es soll kein unfruchtbares Werk in der Liebe geschehen, sondern es soll alles lebendig, heilig, umsonst, fröhlich sein und GOtt wohlgefällig.

\section{Das siebenundzwanzigste Capitel}

\begin{schriftzitat}
V. 1. ff. Und Mose gebot sammt den Aeltesten Israel dem Volk etc.
\end{schriftzitat}

Was er hier geboten hat von dem Altar, der auf dem Berge Ebal aus unbehanenen Steinen errichtet werden sollte, und von dem Gesetze, welches darauf geschrieben werden sollte, ist ein zeitliches Gebot, welches nur einmal erfüllt werden sollte. Josua erfüllte es auch, nachdem er die Stadt Ai genommen hatte, Jos. 8, 30. ff., indem daselbst vor allem Volke die Segnungen und die Flüche gelesen wurden und die Brandopfer und Dankopfer dargebracht, wie hier Moses sagt, freilich, nun durch dies äußere Gepränge (wie ich gesagt habe) das rohe Volk zu reizen, daß es des Gesetzes eingedenk sei, da es durch äußere Ceremonien mehr bewegt wird als durch bloße, nackte Worte. Dieser Altar gereichte später zu großem Aergewiß
\quellseite{807}
, daß das Volk dort opferte und GOtt diente, aufgeblasen durch das Beispiel der Väter, welche dort nur dies Eine Mal geopfert hatten, und zwar auf Anordnung Mosis, der befohlen hatte, daß dies geschehen sollte. Denn auch das Weib von Samaria hielt dies Joh. 4, 20. Christo entgegen und sprach: „Unsere Väter haben auf diesem Berge angebetet.“ Aber sie sahen nicht, daß dies zeitweilige Gebot nicht zu einem immerwährenden Exempel gemacht werden müßte, da sie anderswo den ausdrücklichen Befehl hatten, 2 Mos. 29,\footnote{In den Ausgaben: Exodi 19.} 10. f., daß sie an der Stätte opfern sollten, die der Herr erwähnt hätte, was Moses auch so oft in diesem Buche [Cap. 12, 14. 16, 2.] wiederholt. Die Stätte aber, welche er erwähnt hatte, war immer gewiß, nämlich, wo die Bundeslade war, über welcher er verheißen hatte unter ihnen zu wohnen, so daß wir sehen, daß wir nicht auf das Exempel der Väter oder ein vergangenes Wort GOttes, welches mit einem vergangenen Ereigniß (facto) schon aufgehört hat, vertrauen und uns darauf gründen sollen, sondern allein auf das gegenwärtige Wort GOttes, welches eine gegenwärtige Thatsache (facto) anbetrifft.

Das ist in diesem Capitel zu beachten, daß Moses zu den Segnungen die angenehmeren (potiores) Stämme erwähnt hat, nämlich Simeon, Levi, Juda, Isaschar, Joseph (welcher zwei Stämme in sich begreift, Ephraim und Manasse) und Benjamin, und so sind hier sieben Stämme, alle Kinder der Freien, der Rahel und der Lea; zu den Flüchen aber sechs Stämme, die Kinder der Mägde, der Silpa und der Bilha, mit den beiden Söhnen Rubens und Sebulon von der Freien, Lea. Vielleicht weil der Segen ehrenvoller und würdiger ist, als der Fluch, deshalb mußte er mit den vornehmeren Stämmen geziert werden.

Wiederum, hier zählt er keine Segnungen auf, sondern nur Flüche, da er doch sieben Stämme zum Segen verordnet, und in Josua geschrieben steht, daß sie auch die Segnungen [Jos. 8, 33. f.] ausgerufen haben. Ich meine, dies sei um deswillen geschehen, daß Moses anzeigen will, er sei ein Diener der Sünde und des Todes durch das Gesetz, welches alle unter dem Fluche hält, Gal. 3, 10. Und schließlich ist dieses Volk des Gesetzes auch durch den Fluch zu Grunde gegangen. Josua aber fügt den Segen hinzu, damit er Christum bedeute, der allen den Segen bringt, welchen [begehren] von dem Fluche des Gesetzes [erledigt zu werden und] nach ihm seufzen. Und Juda wird hier um Christi willen auch unter die Zahl derer verordnet, welche da stehen sollen, um zu segnen, weil er von Juda kommen sollte.

Auch darauf hat Moses wohl Acht gehabt, daß er da, wo er vom Segen sagte, hinzugefügt hat [V. 12.]: „das Volk“, indem er sagt: „zu segnen das Volk“, aber da, wo er von dem Fluche sagte, „das Volk“ ausgelassen hat, indem er spricht [V. 13.]: „zu fluchen“; sodann, daß er eher vom Segen geboten hat als vom Fluche, daß zuerst der Segen da sei, als das, was vortrefflicher, besser und wünschenswerther für das Volk ist, der Fluch aber verhaßt, so daß er ihn lieber hätte ansehen lassen, damit er das Volk nicht träfe.

Es ist aber dieser steinerne Altar, auf den das andere Gesetz (Deuteronomium) geschrieben wird, nach der heimlichen Deutung das Volk des Gesetzes selbst, auf welches als auf ein hartes Herz, welches den Geist noch nicht hat, der [das Herz] erweicht, das Gesetz der Werke geschrieben wird. Ferner sind die unbehaüenen Steine die Israeliten selbst, welche durch das Gesetz zwar äußerlich dem Ansehen nach geordnet werden, und mit Kalk ordentlich zu einem Altar geglättet werden, aber an sich bleiben sie innerlich unbearbeitet, und haben ein Herz, das unwillig ist gegen das Gesetz.

Die beiden Berge aber bedeuten, daß das Volk in zwei Theile geschieden ist, nämlich, daß unter ihnen solche Leute sind, die im Geiste das Gesetz halten; diese sind auf dem Berge Ebal und ererben den Segen; die andern aber, welche [das Gesetz] nur durch Werke ohne den Geist zum Ausdruck bringen (exprimunt), die sind auf dem Berge Grisim und liegen unter dem Fluch, Gal. 3, 10.

Hier aber muß man nicht an dem Apostel Paulus vorübergehen, welcher ein wunderbarer Ausleger dieser Stelle ist [Gal. 3, 10.]: „Verflucht sei jedermann, der nicht bleibt in alle dem, das geschrieben stehet in dem Buche des Gesetzes, daß er es thue.“ Denn dies legt er in bejahender Weise (affirmative) aus, auf diese Weise, da er sagt: „Denn die mit des Gesetzes Werken umgehen, die sind unter dem Fluch.“

\quellseite{808}
Denn es steht geschrieben: „Verflucht sei jedermann, der nicht bleibt“ \&c. Denn was kann Wunderlicheres gesagt werden, als daß von dem, der das Gesetz mit Werken ausübt, gesagt wird, er bleibe nicht in den Worten des Gesetzes oder beweise dieselben nicht [in seinem Leben] (confirmare)? Wie kann es sein, daß ebenderselbe, welcher es [das Gesetz] thut, die Worte des Gesetzes nicht thut? Soll man denn etwa sagen, daß der bei seiner Schwester, Schwiegermutter, Stiefmutter, beim Vieh liege, der nicht bei diesen liegt? Und soll er denn so verflucht sein, weil er bei ihnen liegt, da er doch nicht bei ihnen liegt? Und soll er denn so verflucht sein, weil er bei ihnen liegt, da er doch nicht bei ihnen liegt? Denn es gilt dieselbe Weise zu reden für irgend einen Theil des Gesetzes, wie für das ganze Gesetz, so daß der gleicherweise verflucht ist, welcher, indem er Eines thut, es nicht thut, wie der, welcher, indem er das Ganze thut, es nicht thut.

Wiewohl ich hierüber in meiner [kürzeren] Auslegung des Galaterbriefs\footnote{In Galatis meis. Damit ist die „kürzere Auslegung der Epistel an die Galater“ gemeint, welche sich in unserer Ausgabe Bd. VIII, 1352 ff. findet.} genug gesagt habe, so muß ich es doch, weil es die Gelegenheit fordert, kurz wiederholen. Paulus liest Mosen im Geiste, wie er Röm. 7, 14. sagt: „Das Gesetz ist geistlich, ich aber bin fleischlich“ \&c. Darum, weil das Fleisch den Geist nicht erfüllt, ist es unmöglich, daß irgend ein Theil des Gesetzes erfüllt werde, viel weniger das ganze, wie er Röm. 8, 3. sagt: „Denn das dem Gesetz unmöglich war (sintemal es durch das Fleisch geschwächet ward).“ Darum, wenn du auf den äußeren Schein (larvam) der Werke siehst, so sind viele, welche das Gesetz erfüllen, meinest in vielen Stücken, und so scheinen sie nicht unter dem Fluch zu sein, wie die Heuchler. Wenn du aber auf den Geist siehst, der das Gesetz lieb hat, so ist niemand, der auch nur einen Buchstaben oder Tüttel erfüllt. Und so sind alle unter dem Fluch, und diejenigen am meisten, welche sich auswendig mit Werken, ohne den Geist inwendig, im Gesetze üben. Denn diese thun die Werke, entweder gezwungen durch die Strafe oder gereizt durch den Nutzen, in beständigem Haß gegen das Gesetz, und so erfüllen sie es nicht, während sie es erfüllen, thun es nicht, während sie es thun. Denn sie thun es allein mit der Hand und mit dem Werke, unterlassen es aber und thun es nicht mit dem Herzen. Aber GOtt siehet das Herz an [1 Sam. 16, 7.], und nach demselben richtet er durch sein geistliches Gesetz.

So ist es wahr, wer allein Fleisch ohne den Geist ist, der liegt bei seiner Schwiegermutter, Schwester und beim Vieh, und thut alles, was verboten ist, und unterläßt alle Gebote, während er nicht bei ihnen liegt, und während er nicht Verbotenes thut, und die Gebote nicht unterläßt, denn, wenn du auf das äußerliche Werk siehst, so ist er heilig, aber wenn du auf das Herz siehst, so ist kein Werk so schändlich, daß er es nicht thun würde, wenn Scham, Furcht und Strafe nicht entgegenstände, so daß Augustinus recht gesagt hat: Es ist keine Sünde, die ein Mensch gethan hat, welche ein anderer Mensch nicht auch thun könnte, denn wir sind alle von Einem Teige, Ein Fleisch, haben alle dieselbe Bosheit, da wir so durch die erste Sünde Adams verderbt und verflucht sind. Daraus folgt denn, daß die Juden die Flüche im 28. Capitel sogar dann am meisten verdienet haben, wenn sie durch die scheinbarlichsten Werke ohne den Geist glänzten, wie man aus allen Propheten ersehen kann, welche solche ihre Heuchelei durchhecheln und sie mit solchen Flüchen schrecken.

\section{Das achtundzwanzigste Capitel}

\begin{schriftzitat}
V. 1. ff. Und wenn du der Stimme des HErrn, deines GOttes, gehorchen wirst, daß du haltest \&c.
\end{schriftzitat}

Wie dies Capitel das längste ist, so ist es auch das leichteste, denn es enthält weiter nichts als einen großen Haufen Segnungen und Flüche, wodurch er das vorhergehende Capitel weiter ausgeführt und erklärt. Darum haben wir nichts, was die Historie anbetrifft, das wir hier auslegen könnten,\footnote{In der Erlanger: commendemur statt: commentemur.} nur daß wir sagen, man müsse das recht von Herzensgrunde (affectu) betrachten,
\quellseite{809}
wie erschrecklich es sei, unter dem Gesetze zu sein. Denn diese Flüche oder solche, die diesen gleichen, trägt das Gewissen, welches der Sünde dient, das in allen Dingen Qual und Schande empfindet, denn „die Gottlosen haben nicht Frieden“ [Jes. 56, 21.], und kein Ort ist sicher, keine Stunde ruhig (secura), keine Creatur zuverlässig. Wiederum, diese und dergleichen Segnungen genießt ein reines und gottseliges Gewissen, dem alles lieblich ist und günstig (applaudunt).

Daher beschreibt er auch nach allen Flüchen, indem er sie gleichsam in eine Hauptsumme zusammenfaßt, das böse Gewissen, indem er sagt [V. 65.]: „Der HErr wird dir ein bebendes Herz geben und verschmachtete Augen und eine durch Traurigkeit verdorrte Seele.“ Denn so macht das Bewußtsein der Sünde seiner Natur nach furchtsame Leute, so daß sie auch durch das Rauschen eines fliegenden Blattes erschreckt werden, 3 Mos. 26, 36., und, wie der Dichter sagt: Er erbleicht vor jeglichem Blitze. Sodann macht es das Gesicht traurig und bleich; denn gleichwie ein fröhliches Herz ein fröhlich Angesicht macht (wie Salomo [Sprüchw. 15, 13.] bezeugt), so gräbt ein trauriger Geist Runzeln in die Stirne, so daß der Dichter wiederum den glücklich nennt, der wegen keiner Schuld erbleichen muß. Ein solches Gesicht nennt Moses hier „verschmachtete Augen“, weil die Augen der vorzüglichste Theil des Gesichts sind, in welchen Frohsinn oder Traurigkeit am meisten zu Tage tritt. Die Seele aber, die durch Traurigkeit verdorrt ist, ist das ganze Leben in den äußerlichen Sinnen, traurig, niedergeschlagen, mürrisch und verdrießlich, so daß der Mensch keine Lust hat zu leben, und dennoch sich scheut zu sterben, wie folgt:

\begin{schriftzitat}
V. 66. f. Und dein Leben wird vor dir schweben. Nacht und Tag wirst du dich fürchten, und deines Lebens nicht sicher sein. Des Morgens wirst du sagen: Ach, daß ich den Abend erleben möchte! Des Abends wirst du sagen: Ach, daß ich den Morgen erleben möchte!
\end{schriftzitat}

Ich habe noch keine Stelle gesehen, welche den Jammer eines bösen Gewissens klarer ausspräche, mit so passenden und geeigneten Worten und Aussprüchen. Denn so steht das Herz, welches GOtt beleidigt hat, das heißt, welches gequält wird durch das Bewußtsein der Sünde. Es fürchtet sich auch da (wie der Psalmist [ille Ps. 53, 6.] sagt), wo alles sicher ist. Jesaias aber hat es alles in einem Worte zusammengefaßt und sagt [Jes. 57, 20.]: „Die Gottlosen sind wie ein ungestüm Meer, das nicht stille sein kann.“ So treiben den Gottlosen die Stürme der Gedanken, der Furcht, des Mißtrauens, der Verzweiflung, so daß er in Wahrheit einem ungestümen Meere verglichen werden kann.

Man muß aber nicht meinen, daß er wolle, das ganze Volk solle allen diesen Flüchen, und jeder einem jeglichen Fluche unterworfen sein, gleichwie es auch nicht so zu verstehen ist, daß das ganze Volk mit allen Segnungen geehrt werden solle, nämlich, daß jeder einen jeglichen Segen erlange, sondern es ist genug, daß dies alles so erfüllt werde, daß man einen Theil für das Ganze nimmt (per synecdochen). Während zwar den Gottseligen alle Segnungen geschenkt werden, so geschieht dies doch in solcher Weise, daß sie vertheilt werden (distributim): der eine bekommt diese, ein anderer jene, das ganze Volk und alle bekommen etliche; desgleichen etliche [Segnungen] zu immer anderer und anderer Zeit: so müssen auch die Flüche vertheilt werden auf Personen, Oerter, Zeiten und eine gewisse Anzahl (numerum), so daß etliche allen zuteil werden (als Gefangenschaft, Verwüstung, Dienstbarkeit), etliche eine Zeitlang (als Trübsal unter den Feinden), etliche einigen Leuten (als Pestilenz, Krankheiten, Vertreibung).\footnote{Diese letzte Klammer ist von uns gesetzt der Gleichmäßigkeit halber.} Und dies ist die Ursache, weshalb die Gottlosen ihre Flüche nicht genugsam empfinden, wiederum, die Gottseligen vielen Flüchen unterworfen zu sein scheinen, während wir nicht an einem jeglichen alle einzelnen Segnungen sehen.

Nach der heimlichen Deutung werden durch diese leiblichen Flüche alle die Schäden und das Verderben bedeutet, welches Menschensatzungen und Verdrehen der Schrift bei den Gottlosen hervorbringt. Da ist Mangel, Hunger, Durst, Fieber, Hitze (ardor), Schwären, mancherlei Seuchen, Krebsgeier (devoratio) der Kinder, Gefangenschaft, Dienstbarkeit, Schande, Schmach, verächtliche Arbeit und Verlust aller Güter. Aber während dies alles geschieht und ausgeführt wird, wird es nicht empfunden werden vor dem Wüthen und der Raserei des Gemüthes, das bezaubert ist durch jene Greuel. Es wird
\quellseite{810}
aber hernach empfunden werden. Und hier ist nicht eine Vertheilung der Uebel nach Orten, Personen, Zeiten und Anzahl, sondern eine jegliche Seele leidet alle diese Uebel zugleich und immerdar.

Wiederum jene leiblichen Segnungen bedeuten die himmlischen Segnungen, mit denen wir in Christo gesegnet werden. Hier ist Ueberfluß, Vorrath, Reichthum, Gesundheit, Stärke, Heil, Sieg, Ehre, Herrlichkeit und die Fülle aller Güter, und auch diese sind nicht nur theilweise (per synecdochen) vertheilt auf Personen, Stätten, Zeiten und eine Anzahl, sondern alles ist für alle immer und überall vorhanden.

\section{Das neunundzwanzigste Capitel}

\begin{schriftzitat}
V. 1. Dies sind die Worte des Bundes, den der HErr Mose geboten hat, zu machen mit den Kindern Israel, in der Moabiter Lande, zum andermal, nachdem er denselben etc.
\end{schriftzitat}

Dies Capitel mit den zwei folgenden ist gleichsam die letzte Predigt, die Moses dem Volke gehalten hat, und begreift Ermahnungen und Erinnerungen und anderes der Art in sich, um ihnen das Gesetz aus Herz zu legen. Und er sagt, daß dieser Bund gemacht sei „in der Moabiter Lande, zum andermal, nachdem er denselben mit ihnen gemacht hatte auf dem Berge Horeb“, das heißt, er ist erneuert worden um derer willen, die nicht auf Horeb waren und inzwischen geboren waren, da ihre Väter in der Wüste starben, damit sie nicht meinen sollten, das Gesetz, welches den Vätern gegeben und von denselben nicht gehalten worden war, gehe sie nicht an. Deshalb mußte der Gesetzgeber Moses so lange leben, daß er auch dem anderen Geschlechte (generationi) das Gesetz erneuerte.

\begin{schriftzitat}
V. 2. ff. Ihr habt gesehen alles, was der HErr gethan hat etc.
\end{schriftzitat}

Ihr, sagt er, und eure Augen haben gesehen, da er zu den Nachkommen derer redet, die solches gesehen haben, sie selbst es aber nicht gesehen haben, da sie in der Wüste geboren waren, und alle Väter gestorben waren, ausgenommen Josua und Caleb, an deren Statt sie aufgenommen waren. Es sei denn, daß dies so verstanden werde, daß dies von den Nachkommen gesagt werde, die fast schon erwachsen waren, als sie aus Egypten zogen, und unter den Vätern, die gestorben waren, diejenigen [verstanden werden], welche im vollen Mannesalter standen, und, wie der Text sagt, die Kriegsleute, welche vom zwanzigsten Jahre an und darüber gezählt worden waren, 4 Mos. 1,\footnote{In den Ausgaben: Nume. 2.} 3., daß in solcher Weise Moses die anredet, welche beim Auszuge aus Egypten oder zur Zeit der Zählung neunzehn Jahre und darunter waren, so könnte das Wort Moses nicht bestehen. Oder man müßte sagen, daß er nach der Redefigur der Synedoche zu denen rede, die dem Volke angehören, wiewohl dies schon gestorben war, bis auf zwei etc., welche diese Zeichen gesehen hatten. Aber die erste Meinung gefällt mir besser.

Moses redet dies aber in der Weise eines Vorwurfs und gebraucht diese ihre frühere Härtigkeit dazu, die Ermahnung schärfer zu machen, als wenn er sagte: Sehet zu, daß ihr künftig hin die Gebote GOttes besser haltet, als ihr sie bisher gehalten habt. Denn wiewohl ihr dies alles gesehen habt, so seid ihr dennoch um nichts verständiger und gehorsamer geworden, „denn (sagt er) der HErr hat euch bis auf diesen heutigen Tag noch nicht gegeben ein Herz, das verständig wäre, Augen, die da sähen, und Ohren, die da höreten“. Darum ist zu fürchten, daß er vielleicht dieses euch gleicher Weise auch forthin nicht gebe. Deshalb habt ihr mit Furcht und Sorgfalt darauf aus zu sein, daß ihr durch Demuth diese Gnade verdienet, daß ihr ein verständiges Herz, hörende Ohren und sehende Augen erlanget.

Hier siehst du, daß es nicht genug ist, daß Zeichen erwiesen werden und das Wort gelehrt werde, wenn nicht auch inwendig der Geist da ist, welcher lehrt. So hat sich Pharao mit den Seinen nichts gekehrt an die Zeichen Moses, und die Pharisäer lästerten sogar die Zeichen Christi. So groß ist die Härtigkeit des menschlichen
\quellseite{811}
Herzens, daß es sich durch keine Zeichen und Wunder bewegen läßt, durch keine Worte sich weisen läßt, durch keine Drohungen sich schrecken läßt.

Und es ist ganz klar, daß aus dieser Stelle die Weissagung des Jesaias Cap. 6, 10. genommen ist, und ähnliche, welche die Evangelisten, da auch dieselbe Härigkeit [des jüdischen Volks] sie bewegte,\footnote{Wittenberger: momente statt: movente.} so wiederholen [Matth. 13, 14. Marc. 4, 12. Luc. 8, 10. Joh. 12, 40.], indem er sagt:\footnote{In der Form, wie hier, findet sich der Spruch nur bei Jesaias, aber nicht bei den Evangelisten. Deshalb kann es wohl kaum zweifelhaft sein, daß statt dieses in den Ausgaben dicentis (bezogen auf Esaiae) gelesen werden muß, wonach wir übersetzt haben.} „Blende ihre Augen, und laß ihre Ohren dicke sein“ etc. Und offenbar zeigt Moses durch dieses Wort an, daß diejenigen, welche mit des Gesetzes Werken umgehen, vor anderen solche Leute sein würden, daß sie die Gebote GOttes nicht halten, da er gewiß ist, daß nicht allen diese Gnade geschenkt werde. Deshalb redet er so, um sie zu demüthigen und zu schrecken, wenn etwa solche Leute da wären, welche erkennen, daß es nicht in ihren Kräften stehe, das Gesetz zu halten, und gedemüthigt GOtt um seine Gnade bitten möchten, wie geistliche Leute thun; die vermessenen Heuchler aber thun es nicht.

Darnach wiederholt er die Wohlthaten, daß vierzig Jahre lang in der Wüste ihre Kleider nicht veraltet sind und ihre Schuhe nicht veraltet, daß sie kein Brod gegessen und keinen Wein und berauschend Getränk getrunken haben, wodurch der HErr ihnen zeigte und sie wissen sollten, daß er ihr GOtt wäre, nämlich, daß er sie nähren und kleiden könne ohne menschliches Zuthun (opere), wie er oben Cap. 8, 3. gesagt hat: „Der Mensch lebt nicht vom Brod allein, sondern von einem jeglichen Worte Gottes“ etc.

Darnach erzählt er den Sieg, den sie erlangt haben über die beiden Könige der Amoriter, und daß das Land unter sie ausgetheilt sei, und schließt: „So haltet nun“ etc., „auf daß ihr weislich handeln möget in alle eurem Thun“, das heißt, damit ihr weislich und glücklich alles hinausführen möget, was ihr euch vorgenommen habt. Denn dies drückt das hebräische Wort [\hebraeisch{השכיל}] aus: „daß ihr weislich handeln möget“, gleichwie von David geschrieben steht [1 Sam. 18, 30.], daß er klüglicher handelte als alle Knechte Sauls, wie auch zu Josua gesagt wird [Jos. 1, 7.]: „Auf daß du weislich handeln mögest“, das heißt, glücklich und gedeihlich (prosper), und nicht eine Thorheit begehest; dies wird aber geschehen, wenn sie das Gesetz GOttes halten werden.

Darnach bestätigt er den Bund, indem er mit Namen alle Stände (ordines) aufzählt, die Obersten, die Amtleute, die Aeltesten, die Stämme, die Kinder, die Weiber, die Fremdlinge, bis auf die Holzhauer und Wasserschöpfer, damit alle und jeder einzelne, sowohl die Anwesenden als auch die Abwesenden wüßten, daß dieser Bund sie angehe. Moses hat sicherlich in einem so großen Volke diese Worte nicht allein reden können, denn wie sollten ihn alle hören? Es ist also nothwendig, daß er dies in einer vorgeschriebenen Predigt durch die Priester und Leviten an verschiedenen Orten in dem ganzen Volke (per populum) habe predigen lassen.

Er schlägt aber die heimliche Bosheit des Herzens und kommt ihr zuvor, durch welche viele sich innerlich trösten, wenn sie die Drohungen und Flüche GOttes hören, und sprechen: Es wird so schlimm nicht sein, und so sicher verharren in ihrer Gottlosigkeit. Dies thun am meisten, betrogen durch den Schein der Gottseligkeit, die Werktreiber und die Abgöttischen, welche glauben, daß das sie nicht angehe, was das Gesetz droht, wie man bei allen Propheten sehen kann. Denn wie die Gottseligkeit ihrer Natur nach die Furcht GOttes mit sich bringt, so die Gottlosigkeit und die Heuchelei Sicherheit. Wider diese Sicherheit der Gottlosen geht nun Moses sehr scharf vor, bis zum Ende dieses Capitels, wie dies auch alle Propheten thun, und sagt zuvor, daß ihnen am meisten die Flüche und Verwüstungen drohen, so daß sie, umgekehrt (subversi) wie Sodom und Gomorra, das Gespräch (fabula) und das Gespött der ganzen Welt werden.

Und Moses zeigt an dieser Stelle, daß seine Predigten besonders wider die Heuchler und Werktheiligen donnern, die sich vermessen, ohne den Geist das Gesetz zu erfüllen, so daß aus dieser Stelle die Propheten einen großen Theil ihrer Predigten wider dieselben Heuchler und Abgöttischen genommen haben. Daher nennt er [V. 18.] einen solchen Menschen eine Wurzel, die Galle und Wermuth trägt, indem er mit
\quellseite{812}
bitteren Worten und Sachen die bittere und verderbliche Heuchelei angreift. Dieser Worte gebraucht auch die Epistel an die Hebräer [Hebr. 6, 8.]\footnote{Es sind nicht die Worte „Galle und Wermuth“ im Briefe an die Hebräer gebraucht, sondern: „Dornen und Disteln“. Daher ist „dieser Worte“ durch „bitterer Worte“ auszulegen.} und fast auch Petrus in der Apostelgeschichte, Cap. 8, 23., wider Simon, den Zauberer. Es ist also ein solcher Werkslehrer eine Wurzel, die Galle und Wermuth trägt, denn er bringt durch seine Lehre und Exempel bittere Früchte, das heißt, die Werke des Gesetzes ohne den Geist, welche durch falschen Schein auswendig süß sind, inwendig aber das Gewissen beunruhigen und verbittern. Das heißt, sie stellen GOtt dar als bitter und erzürnt, da er ja durch solche Gerechtigkeit gelästert wird, indem unseren Kräften das zugeschrieben wird, was seinem Geiste zukommt.

„So fährt (sagt er) die Trunkene mit der Durstigen dahin.“ „Die Trunkene“ nennt er einen Lehrer, der voll eigener Meinungen ist im Gesetze GOttes. Daher wird in den Propheten, besonders in Jesaia, viel von dieser Trunkenheit und Speien und Wein gesagt [Jes. 28, 1. Jer. 25, 27. 48, 26.]. „Die Durstige“ aber ist der Zuhörer, welcher begierig ist, die Gottseligkeit kennen zu lernen, und ohne allen Unterschied alles annimmt und trinkt.

Am Ende fügt er hinzu [V. 29.]: „Das Geheimniß des HErrn ist uns offenbaret, daß wir thun sollen alle Worte dieses Gesetzes.“ Auch dies gehört zur Schärfe der Ermahnung, als ob er sagen wollte: Die Geheimnisse, welche GOtt über uns denkt, daß er den Gottlosen fluche und die Gottseligen segne, und alles was er will, daß wir thun sollen, werden uns, und nicht anderen Völkern, offenbart, aber nicht vergeblich. Als ob er sagen wollte: Darum fürchtet euch. Derjenige, welcher solches offenbart, offenbart es um deswillen, daß er euch schrecke, damit ihr gewiß seiet, alles das werde geschehen, wenn ihr nicht alle Stücke dieses Gesetzes haltet.

\section{Das dreißigste Capitel}

\begin{schriftzitat}
V. 1. ff. Wenn nun über dich kommt dies alles etc.
\end{schriftzitat}

In diesem Capitel knüpft Moses die Hoffnung der Buße an und die Verheißung der Vergebung der Sünden, wenn sie Buße thun, damit diejenigen, welche durch die Unterweisung der Trübsal wieder zurecht kämen, etwas hätten, dazu sie ihre Zuflucht nehmen könnten, nämlich diese Verheißungen, wie der König Manasse sich zu denselben wendet. Denn wie den harten und sicheren Gottlosen ein hartes und heftiges Schelten, Fluch und Drohung vonnöthen ist, wie wir im vorigen Capitel gesehen haben, so ist den Zerschlagenen und Gedemüthigten und Erschreckten eine liebliche und heilsame Verheißung und Trost vonnöthen. Und Moses hält ihnen in diesem Capitel fürwahr gar herrliche Verheißungen der Gnade vor, so daß er sagt [V. 5.], daß ihnen alles Gute werde erwiesen werden, daß ihrer viel mehr werden würde als zuvor, [V. 9.:] daß sich der HErr über ihnen freuen werde, so daß der Text für den Leser leicht ist; alles dies ist für ein betrübtes und verzweifeltes Gewissen nothwendig.

\begin{schriftzitat}
V. 11. ff. Denn das Gebot, das ich dir heute gebiete, ist dir nicht verborgen etc.
\end{schriftzitat}

Diese Stelle ziehen die Sophisten darauf, daß sie sagen, es sei uns möglich, die Gebote GOttes zu erfüllen, ja, es sei uns leicht. Denn diese blinden Leser fallen so hinein, daß sie das, was Moses sagt: „Es ist dir nicht verborgen“ (Non est supra te = es ist dir nicht zu hoch) so verstehen: Es geht nicht über deine Kräfte, ist dir auch nicht unmöglich, und bedenken nicht, daß, wenn dieses wahr wäre, man vergebens nach der Gnade GOttes trachte. Denn was ist es vonnöthen, daß mir die Gnade GOttes beistehe in dem, was mir möglich ist und nicht über meine Kräfte? Wenn sie den gebräuchlichen Text gehabt hätten, so würden sie sich noch viel mehr herausgenommen haben und aufgeblasen gewesen sein. Denn der Hebräer sagt so: „Dieses Gebot, welches ich dir heute gebiete, ist dir nicht wunderbar
\quellseite{813}
noch schwer.“ Denn \hebraeisch{נפלאות} bedeutet bisweilen das Wunderbare und Schwere, was unsere Kräfte übersteigt, und dazu wir nicht zu gelangen vermögen. Doch wir wollen nun zu Moses kommen.

Moses will an dieser Stelle die Weise zeigen, wie das Gebot GOttes erfüllt werde, und sagt erstlich: „Es ist dir nicht verborgen, noch zu ferne“, das heißt, du meinst, das Gebot könne alsbald erfüllt werden, wenn du es hörst; darum, wenn es nicht gehört würde, so möchtest du vielleicht sagen: Ich würde es thun, wenn es da wäre. Jetzt, da nicht da ist, was ich thun soll, wer wird es mir vom Himmel oder von jenfeit des Meeres herholen, daß ich es höre und thue? Aber es wird nicht erfüllt durch Hören und durch ein Werk, vielmehr wenn das Wort in deinem Munde und deinem Herzen sein wird. Im Munde, sage ich, und im Herzen muß es sein, wenn du es erfüllen willst, nicht im Ohre oder in der Hand. Darum sind im Texte diese beiden Stellen zu beachten, die eine [V. 12.]: „Daß wir’s hören und thun“, die andere [V. 14.]: „In deinem Munde und in deinem Herzen.“ In diesen beiden Stellen, sage ich, liegt die Bedeutung (vis) und das Verständnis dieses Textes.

Denn es sind zweierlei Erfüller des Gebotes; die ersten, von denen er sagt: „Daß wir’s hören und thun.“ Diese, verlassen auf ihre Kräfte, sagen und rühmen sich alsbald, daß sie es thun, wenn sie es gehört haben. Diese machen nicht ihre Kräfte geltend (causantur), sondern das Hören, und sagen: Wenn wir es nur hören, so werden wir es auch alsbald thun. Diese verwirft Moses, und sagt von den anderen: „Das Wort ist in deinem Munde und in deinem Herzen“, das heißt, es fordert das Herz, nicht die Hand. Es ist aber nicht in dem Herzen, wenn es nicht durch den Geist GOttes hineingeschrieben wird. Denn in dem Herzen ist, was ich will, was ich liebe, was ich gern habe, was mir gefällt. Daß also das Gesetz im Herzen sei, das heißt, daß es geliebt werde, was unseren Kräften unmöglich ist. Diesem ist der 37. Psalm nachgefolgt und sagt [V. 31.]: „Das Gesetz seines GOttes ist in seinem Herzen“, und der erste Psalm, V. 2.: „Er hat Lust zum Gesetz des HErrn.“ Wenn aber das Gesetz im Herzen geliebt wird, so wird es in Wahrheit erfüllt, denn das Gesetz fordert ein Herz (affectum), welches es liebe und Gefallen an ihm habe; so wird ihm genuggethan und so wird es erfüllt, wenn man es liebt.

So ist nun Moses klar: „Das Wort ist in deinem Munde und in deinem Herzen, daß du es thust“, als ob er sagen wollte: Du wirst es nicht thun, es sei denn in deinem Munde und in deinem Herzen.

Warum setzt er aber zuerst: „In deinem Munde“, ehe er sagt: „In deinem Herzen“, da es doch eher im Herzen ist als im Munde, weil es eher geliebt als gelehrt wird? Dies thut er, um die Art und Weise anzuzeigen, wie das Gebot GOttes erfüllt wird, daß dies geschehen werde durch das Wort des Evangelii, welches zuerst mit dem Munde gepredigt, darnach durch die Predigt (auditu) geglaubt wird. So lenkt also Moses durch diesen Text das Volk hin zu einem anderen künftigen Worte, welches, mit dem Herzen aufgenommen, Liebe gegen sein Gebot hervorbringt, als wenn er sagen wollte: Mein Gebot wirst du nicht erfüllen, wenn du es hörst, sondern wenn du es von Herzen liebst. Dies wirst du nicht thun, wenn das Wort nicht mit dem Munde gepredigt und mit dem Herzen geglaubt wird. So wird mein Gebot nicht schwer und nicht ferne sein.

Es bleibt noch übrig, daß wir jetzt den Paulus Röm. 10, 6–8. ansehen, der diese Stelle anders anführt, als Moses sie hat, indem er sagt: „Die Gerechtigkeit aus dem Glauben spricht also: Sprich nicht in deinem Herzen: Wer will hinauf gen Himmel fahren? (Das ist nicht anders, denn Christum herabholen.) Oder, wer will hinab in die Tiefe fahren? (Das ist nicht anders, denn Christum von den Todten holen.) Aber was sagt die Schrift? Das Wort ist dir nahe in deinem Munde und in deinem Herzen. Dies ist das Wort vom Glauben, das wir predigen“ etc. Hier sagen wir erstlich, daß Paulus absichtlich den Moses nicht habe wörtlich anführen wollen, wenigstens im ersten Theile, sondern in einem reichen (abundante) Geiste aus Mose die Gelegenheit genommen habe, wider die Werkheiligen gleichsam einen neuen und eigenen Text zu verfassen. Deshalb hat er auch das Stück ausgelassen, welches zweimal bei Moses gesetzt ist [V. 12. und 13.]: „Und uns holen, daß wir’s hören und thun.“ Ferner sagt er nicht, daß so geschrieben stehe, sondern er sagt, daß die Gerechtigkeit aus dem Glauben so spreche. Nachher aber, wo er Moses wörtlich anführt, sagt er: „Aber was sagt die Schrift?“
\quellseite{814}
um anzuzeigen, daß es die Worte seines Geistes und nicht Mosis seien, welche er der Gerechtigkeit aus dem Glauben zuschreibt, wiewohl eine Aehnlichkeit (analogia) unter ihnen ist.

Es stimmen also im einfältigen Verstande Moses und Paulus überein im letzten Theile, denn Paulus will, daß durch die Worte: „Das Wort ist dir nahe“ etc. gelehrt werde, daß man mit dem Herzen glaubt zur Gerechtigkeit durch das Wort des Glaubens. Ebendieselbe will Moses, da er das Urtheil ausspricht (sentit), sein Gebot werde erfüllt, wenn das Wort im Herzen sei, das heißt, wenn es geliebt wird, was durch den Glauben geschieht. Aber das, was Paulus aus Anlaß der Worte Mosis spielt: „Sprich nicht in deinem Herzen: Wer will hinauf gen Himmel fahren? Das ist nicht anders, denn Christum herabholen“ etc., das donnert er wider die Werktheiligen. Denn da diese mit ihren Werken und aus ihren Kräften das Gebot erfüllen, so leugnen sie, daß Christus gen Himmel gefahren sei, und niedergefahren sei zur Hölle, das heißt, sie leugnen, daß er für uns gestorben und auferstanden sei, damit wir durch das Wort des Glaubens gerechtfertigt würden und das Gesetz erfüllten. Denn er sieht auf ihr Herz, da er sagt: „Sprich nicht in deinem Herzen: Wer will hinauf gen Himmel fahren?“ als ob er sagen wollte: Der Gottlose glaubt nicht, daß Christus gen Himmel gefahren sei, sondern er zieht ihn vom Himmel herab, da er nicht glaubt, daß durch seine Auferstehung der Glaube geschenkt werde, welcher das Gesetz erfüllt.

\section{Das einunddreißigste Capitel}

\begin{schriftzitat}
V. 1. ff. Und Mose ging hin, und redete diese Worte etc.
\end{schriftzitat}

„Er ging hin“, nämlich, nachdem die Predigten vom Gesetze, die Ermahnung, Drohung und die Weise, das Gesetz zu erfüllen, beendigt waren, und alles vollbracht, was sein Amt an dem Volke erforderte. Nachdem er aber zurückgekehrt war, nämlich an einem andern Tage, redete er, was in diesem Capitel geschrieben steht. Es wird aber geschrieben, wie er sich zum Tode bereitet habe, Abschied genommen und verheißen habe, daß die Heiden vertilgt werden sollten, wie er den Josua an seine Statt als Führer eingesetzt habe, wie er befohlen habe, daß das Gesetz geschrieben werden und in jedem siebenten Jahr gelesen werden sollte, und ein Lied zum Gedächtniß gemacht habe. Und so beschließt Moses seine Predigt in diesem Capitel. Denn das 32. Capitel enthält dieses Lied, das 33. Capitel die Segenssprüche; das 34. Capitel kann Moses nicht geschrieben haben.

Erstlich nun sagt er, er sei hundert und zwanzig Jahre alt, und könne nicht mehr aus und eingehen, das heißt, sich der Angelegenheiten des Volkes annehmen und sein Amt verwalten. Aber auch dies würde ihn nicht ganz verhindert haben, wenn der HErr nicht gesagt hätte, daß er nicht über den Jordan gehen solle. Er verkündigt ihnen daher, daß der HErr ihnen ein Führer sein werde wider die Heiden, die vertilgt werden sollten; unter dem HErrn aber werde Josua ihr Führer sein an seiner Statt, und er fügt herrliche Verheißungen hinzu, indem er sagt [V. 6.]: „Seid getrost und unverzagt, fürchtet euch nicht, und lasset euch nicht vor ihnen grauen; denn der HErr, dein GOtt, wird selber mit dir wandeln, und wird die Hand nicht abthun, noch dich verlassen.“ Dasselbe wiederholt und redet er [V. 7.] zu Josua, da dieser vor das Volk gerufen worden war, und setzt ihn zum Führer ein, damit in solcher Weise sowohl das Volk als auch der Führer sich im Glauben gründen möchten auf das Wort der Verheißung GOttes, ihres obersten Führers.

Nachdem nun der Führer eingesetzt und ihm sammt dem Volke Muth eingesprochen worden ist, ordnet er in Bezug auf das Gesetz an, daß es in jedem siebenten Jahre öffentlich gelesen werden solle durch die Leviten an der Stätte, die der HErr erwählet hätte. Darauf ruft der HErr [V. 14.] Moses und Josua, und befiehlt ihnen, daß sie ein Lied machen sollen zum Zeugniß der künftigen Uebertretung und Strafe, denn es war gewiß, daß sie später in Uebertretungen gerathen
\quellseite{815}
würden, weil sie nicht alle den Geist hatten; da sie, wie Moses sich beklagt [V. 27.], während er noch lebte, dem HErrn ungehorsam waren, wie viel mehr nach seinem Tode! Endlich wird befohlen [V. 26.], daß das Buch des Gesetzes in die Seite der Lade des Bundes gelegt werde, zum Zeugniß wider sie. Und nachdem so [V. 28.] die Aeltesten und Amtleute versammelt worden waren, redete er das Lied.

Es ist oben [Cap. 3] die heimliche Deutung des sterbenden Moses und der Einsetzung Josua’s berührt worden, darum muß es hier kurz wiederholt werden. Moses ist das Amt des Gesetzes, welches nicht zur Vollkommenheit führt, das heißt, zur Gerechtigkeit, sondern es zeigt die Sünde, und erfordert die Gnade, welche es nicht verleiht, darum stirbt er und läßt ab diesseits des Jordans in dem Lande der Moabiter, das heißt, in der Gerechtigkeit der Werke. Diesem folgt Josua als Führer, das heißt, das Amt der Gnade, und führt das Volk hinein in das Land der Verheißung, das heißt, zu der wahren Gerechtigkeit in Christo, indem die Israeliten trockenen Fußes durch den Jordan gehen, das heißt, indem sowohl die Sünde als auch der Tod zurückweicht und der Gnade Statt gibt.

\section{Das zweiunddreißigste Capitel}

Das Lied Mosis ist voll Scheltens und Vorwürfe wegen der so vielen und so großen Wohlthaten, die dem undankbaren und bösen Volke von GOtt erwiesen worden sind. Und er zeigt deutlich an, daß es geschehen werde, daß sie GOtt verlassen, sich zu fremden Göttern kehren, und die vorhergenannten Flüche über sie kommen werden. Und hierin thut Moses ein seiner würdiges Werk (officium), das heißt, das des Amts des Gesetzes würdig ist, nämlich, daß er beißt, anklagt, schilt, droht, flucht, und durch und durch nichts als Zorn anzeigt, und doch wohl weiß (sentit), daß das Volk dadurch um nichts besser gemacht werde, so daß man hier gleichsam als in einem Spiegel auch die Natur des Gesetzes sehen kann, daß es Zorn anrichte und unter dem Fluche festhalte.

\begin{schriftzitat}
V. 1. Merket auf, ihr Himmel etc.
\end{schriftzitat}

Himmel und Erde, das heißt, alle Creaturen ruft er als Zeugen an. Und es geschieht auch, daß, wenn die Gottlosen heimgesucht werden, sie meinen, daß alle Creaturen feindselig gegen sie seien. Dies thut das Bewußtsein der Sünde, die ihnen durch das Gesetz angezeigt ist. So stimmt die ganze Creatur mit dem Gesetze überein, als ein Zeuge desselben.

\begin{schriftzitat}
V. 2. Meine Lehre triefe wie der Regen etc.
\end{schriftzitat}

Er wünscht aber, daß seine Lehre träufeln und fließen möge wie der Regen und der Thau, wie die Tropfen und Platzregen über Kraut und Gras, das heißt, daß sie kräftig und fruchtbar sei, damit er nicht vergeblich singe und einem Tauben eine Geschichte erzähle, da er von großen und nothwendigen Dingen singt, das heißt, daß man dem wahren Gotte dienen solle, und von der Gefahr derer, die GOtt verlassen und sich zu den Götzen wenden.

Tropfen, Regen und Thau setzt er den Wolken ohne Wasser entgegen, als da sind menschliche Lehren, welche nichts nützen. Darum will er, daß sein Wort ein Regen sei, nicht eine leere Wolke. Kraut aber und Gras setzt er den Kieseln, dem Sande und ähnlichen unfruchtbaren Dingen entgegen, auf welche es vergeblich regnet, wiewohl der Regen heilsam ist, das heißt, er wünscht Zuhörer, bei denen sein Wort Frucht bringe. Denn es ist beides erforderlich, daß die Lehre heilsam sei, und der Zuhörer sich lehren lasse (docilis).

\begin{schriftzitat}
V. 3. Denn ich will anrufen den Namen des HErrn.
\end{schriftzitat}

Das heißt, ich will davon fügen, daß man nur Einen GOtt anrufen und ihm dienen soll, wider die Götzendiener, die in meinem Volke sein werden.

\subsection{Gebt unserm GOtt allein die Ehre}

Das heißt, ehret ja keinen andern Gott, schreibet auch nicht irgend einem irgend welche Herrlichkeit oder Gewalt zu, sondern allein unserem
\quellseite{816}
Götte. Denn er allein hat Herrlichkeit, Größe, Kraft. Alle anderen Götter sind eitel, lügenhaft und nichts.

\begin{schriftzitat}
V. 4. Er ist ein Fels. Seine Werke sind unsträflich.
\end{schriftzitat}

Er nennt GOtt einen Felsen, weil er getreu und zuverlässig ist denen, die auf ihn trauen. Seine Werke nennt er heilsame, unsträfliche, vollkommene Werke, und redet von den Werken, die er in uns und allen Gottseeligen wirkt, die ihn in der Wahrheit anbeten, und setzt dieselben den Werken der Eitelkeit entgegen, mit welchen wir, da wir sie aus uns selbst erdichtet haben, entweder GOtt selbst oder anderen Göttern dienen. Diese [Werke] sind nicht unsträflich und vollkommen, sondern verworfen und verdammt.

Es ist dasselbe, was da folget: „Alles, was er thut, das ist recht“, das heißt, das ganze Leben seiner Gläubigen ist recht und wahrhaftig, geht einher im Worte und Gebot, nicht auf eigenen Wegen und Meinungen. So „ist GOtt treu, und kein Böses an ihm, gerecht und fromm ist er“. Dies alles sagt er von GOtt, sofern man ihm dient und an ihn glaubt, nicht allein von GOtt nach seinem Wesen, als wenn er wider falsche Götter und die, welche ihnen dienen, sagte: Wer diesem GOtt dient, der ist gewiß, der hat einen treuen GOtt, mit welchem er Frieden hat. In diesem Dienste ist nichts Böses, sondern nur Gerechtigkeit und rechtschaffenes Wesen, und dies alles ist treu, gewiß und wohl gegründet. Dagegen die Abgöttischen haben nichts Sicheres, nichts Treues oder Gewisses, sondern sie wanken, da ihr Gewissen immer unstät und ungewiß ist, sodann auch alles, was sie leben und vornehmen, ist nur Bosheit und Schlechtigkeit, so sehr sie sich auch anders erscheinen mögen.

\begin{schriftzitat}
V. 5. Die verkehrte und böse Art fällt von ihm ab; sie sind Schandflecke und nicht seine Kinder.
\end{schriftzitat}

So sagt er von ihnen aus, daß es mit ihnen dahin kommen werde, daß sie nicht bleiben würden in der heilsamen Lehre und dem rechten Gottesdienste, sondern, wie es mit einer bösen, ehebrecherischen Art (wie auch Christus [Matth. 12, 39.] sie nennt) zu geschehen pflegt, sie werden verderbt, so daß sie nun nicht mehr Kinder sind, wegen ihrer Schändlichkeit und ihres verdorbenen Gottesdienstes, den sie selbst erfunden haben, sondern Bastarte, ja, Widersacher.

\begin{schriftzitat}
V. 6. Dankest du also dem HErrn, deinem GOtt, du toll und thöricht Volk? Ist er nicht dein Vater und dein HErr? Ist’s nicht er allein, der dich gemacht und bereitet hat?
\end{schriftzitat}

Hier fängt er an, dem undankbaren Volke Vorwürfe zu machen, und siehe, wie glühende und überaus heftige Worte er gebraucht. Für dies alles (sagt er), daß GOtt gerecht und fromm ist, seine Werke unsträflich sind etc., und daß er dein Vater ist, was ihm hoch zu stehen thun, daß er dich erlangte; desgleichen, daß er dich gemacht hat, da du kein Volk warst, und dich zu einem fertigen Volke zugerichtet hat, das heißt, bereitet und gemacht, damit alles wohl hinausginge, für alles dies (sage ich) verderbst und verkehrst du alles, so daß ihr nun nicht mehr Kinder eines solchen Vaters seid, sondern eine verkehrte Art, welche fremden Göttern und ihrem eigenen Gelüsten nachgeht.

\begin{schriftzitat}
V. 7. Gedenke der vorigen Zeit bis daher, und betrachte, was er gethan hat an den alten Vätern. Frage deinen Vater, der wird dir’s verkündigen, deine Aeltesten, die werden dir’s sagen.
\end{schriftzitat}

Er beginnt, der Ordnung nach die Wohlthaten aufzuzählen, welche über das undankbare Volk ausgeschüttet und an demselben verloren gewesen sind. Deinen Vater hat er in der Einzahl gebraucht, anstatt der Mehrzahl, als ob er sagen wollte: Wenn du auf die alten Väter siehst, als Abraham, Isaak und Jakob und ihre Nachkommen, so wirst du erkennen, wie viel Gutes er dir in jenen gethan hat, wie sehr er für dich Sorge getragen hat.

\begin{schriftzitat}
V. 8. Da der Allerhöchste die Völker zertheilete, und zerstreuete der Menschen Kinder, da setzte er die Grenzen der Völker nach der Zahl der Kinder Israel.
\end{schriftzitat}

Die erste Wohlthat ist, daß er, obgleich er GOtt ist, und aller Völker GOtt, da er sie ja zerstreuete, 1 Mos. 11, 8., und den ganzen Erdkreis unter sie austheilte, einem jeglichen Volke seinen Theil und seine Grenze, wie Paulus in der Apostelgeschichte, Cap. 17, 26., diesem Verse folgend, lehrt, sich dennoch nicht alle zu dem Volke erwählte, dessen GOtt er insbesondere wäre durch wahren und rechten Gottesdienst, sondern er setzte Grenzen, daß allein da, wo die Kinder Israel wären, sein Volk wäre; wo aber
\quellseite{817}
die Grenze der Kinder Israel aufhörte, da sollte das Volk GOttes aufhören, und die Grenzen der Heiden anfangen.

\begin{schriftzitat}
V. 9. Denn des HErrn Theil ist sein Volk, Jakob ist die Schnur seines Erbes.
\end{schriftzitat}

Das heißt, sein Volk erstreckt sich nicht weiter als Jakob oder Israel ist. Denn durch diese Worte wird beschrieben, wie dies Volk insonderheit GOttes Volk ist.

\begin{schriftzitat}
V. 10. Er fand ihn in der Wüste, in der dürren Einöde, da es heulet. Er führte ihn und gab ihm das Geleit. Er behütete ihn wie seinen Augapfel.
\end{schriftzitat}

Das heißt, er machte auf dem Berge Sinai einen Bund mit ihm, da fand er ihn und nahm ihn an, darnach hat er ihn in der Wüste umhergeführt und unterwiesen und behütet, daß ihm nie etwas mangelte. Dies ist nun die zweite herrliche Wohlthat.

\begin{schriftzitat}
V. 11. Wie ein Adler aufführet seine Jungen und über ihnen schwebet. Er breitete seine Fittige aus, und nahm ihn, und trug sie auf seinen Flügeln.
\end{schriftzitat}

Mit diesen Worten zeigt er an, wie er sich ihrer angenommen habe in der Wüste, ihre Sitten getragen, und sie versucht und ihnen wohlgethan habe, damit sie fliegen lernten, das ist, auf ihn vertrauen, wie wir Cap. 8 gesehen haben.

\begin{schriftzitat}
V. 12. Der HErr allein leitete ihn, und war kein fremder Gott mit ihm.
\end{schriftzitat}

Das heißt, unter Mose ging in der Wüste der rechte Gottesdienst im Schwange. GOtt selbst ging vor ihnen her, sowohl mit dem Worte als auch in der äußerlichen Leitung der Wolke und des Feuers.

\begin{schriftzitat}
V. 13. Er ließ ihn hoch herfahren auf Erden, und nährete ihn mit den Früchten des Feldes, und ließ ihn Honig saugen aus den Felsen, und Oel aus den harten Steinen.
\end{schriftzitat}

Das ist die dritte Wohlthat, durch welche der Reichthum beschrieben wird, der ihnen im Lande Canaan geschenkt war, nämlich, daß sie in einem hohen und trefflichen Lande wohnten, und durch die Früchte des Feldes reichlich ernährt wurden, das heißt, er gab ihnen, zu wohnen in einem guten und fruchtbaren Lande. Sodann sagt er in dichterischer Weise und in bildlicher Rede, es werde Honig aus den Felsen und Oel aus den Steinen gesogen, nämlich, daß Land sei so gut, daß da, wo nicht Acker ist, sondern Fels und Stein, dennoch Oel sei und die Bienen Honig tragen, da ja auch die Oelbäume gern an steinichten und unfruchtbaren Orten wachsen, und die Bienen überall Honig tragen, auch wo kein Acker ist.

\begin{schriftzitat}
V. 14. Butter von den Kühen, und Milch von den Schafen, sammt dem Fett von den Lämmern, und feiste Widder und Böcke mit fetten Nieren, und Weizen, und tränkte ihn mit Traubenblut.
\end{schriftzitat}

Auch hier beschreibt er in dichterischer Weise, wie es in Liedern zu geschehen pflegt, den Reichthum, daß sie mit Butter und Milch und fetten Lämmern genähret werden sollten, und mit gemästeten Widdern, und mit Böcken mit fetten Nieren, das heißt, die auch wohl genähret und gemästet waren.

(Eine andere Uebersetzung\footnote{Nämlich die Vulgata, welche bietet: Et arietum filiorum Basan.} hat so übersetzt: „Und mit Widdern der Kinder Basan“, indem der Uebersetzer nicht die hebräische Weise zu reden kannte, daß „Kind Basan“, das heißt, ein Kind der Fettigkeit, dasselbe sei als fett, wie ein Kind des Jahres, das heißt, ein Jährling; ein Kind des Friedens, das heißt, ein Friedfertiger; ein Kind des Todes, das heißt, dem Tode verfallen. So nennt er den Wein „Traubenblut“, weil in jenen Ländern insgemein rother Wein wächst.

\begin{schriftzitat}
V. 15. Da er aber fett und satt ward, ward er geil. Er ist fett und dick und stark worden und hat den GOtt fahren lassen, der ihn gemacht hat. Er hat den Fels seines Heils gering geachtet.
\end{schriftzitat}

Das heißt, den, der ihn gemacht und ihn aus allem Uebel errettet hat, auf den er hätte vertrauen sollen, den hat er, durch Reichthum aufgeblasen und satt geworden, vernachlässigt, indem er anderen Göttern anhing, wie folgt:

\begin{schriftzitat}
V. 16. Und hat ihn zu Eifer gereizet durch Fremde, durch die Greuel hat er ihn erzürnet.
\end{schriftzitat}

Die ihm hiefür hätten dankbar sein sollen, wendeten sich und reizten ihn vielmehr, indem
\quellseite{818}
sie ihm Böses für Gutes vergälten, wie wir in den Historien sehen, wie sie mit fremden Göttern gehurt haben. Ein gesättigter Bauch erträgt nicht die Gottseligkeit, denn er steht in Sicherheit und kümmert sich nicht um GOtt.

\begin{schriftzitat}
V. 17. Sie haben den Feldteufeln geopfert, und nicht ihrem GOtt; den Göttern, die sie nicht kannten, den neuen, die zuvor nicht gewesen sind, die eure Väter nicht geehrt haben.
\end{schriftzitat}

Das heißt, den Götzen im Felde, nämlich in Hainen, in Thälern, auf Bergen, von denen ihnen im Gesetz nichts gesagt war, von denen sie auch nichts gehört hatten, sondern nachdem ihnen die Ohren jückten; so wie jemand irgend einen neuen Gottesdienst aufbrachte, so strömten sie aufs stärkste hinzu, und verließen GOtt, den sie von ihren Vätern gelernt hatten.

\begin{schriftzitat}
V. 18. Deinen Fels, der dich gezeuget hat, hast du aus der Acht gelassen, und hast vergessen GOttes, der dich gemacht hat.
\end{schriftzitat}

Das ist ein bitteres Schelten und Vorwurf der Undankbarkeit, daß sie für so viel Gutes so viel Böses vergälten. Darum folgt jetzt, mit welcher Strafe sie für eine so große Schuld geschlagen werden sollen.

\begin{schriftzitat}
V. 19. Und da es der HErr sahe, ward er zornig über seine Söhne und Töchter.
\end{schriftzitat}

Der Zorn ist um so erschrecklicher, weil die Söhne und Töchter einen so großen Vater mit solchen und so großen Sünden beleidigten.

\begin{schriftzitat}
V. 20. Und er sprach: Ich will mein Antlitz vor ihnen verbergen, will sehen, was ihnen zuletzt widerfahren wird, denn es ist eine verkehrte Art, es sind untreue Kinder.
\end{schriftzitat}

„Daß das Antlitz GOttes verborgen wird“, das heißt, daß sie blind werden, daß sie ihn weder erkennen noch finden, so sehr sie ihn auch suchen und rufen. Die Bedeutung hiervon ist bekannt aus dem Gegentheil [Ps. 80, 4.]: „Laß leuchten dein Antlitz, so genesen wir.“ Er nennt sie „Untreue“, das heißt, die nicht in Wahrheit und von Herzen Kinder sind, sondern äußerlich heucheln mit einem erdichteten Schein (hypocrisi) der Werke, und rühmen, daß sie Kinder seien.

\begin{schriftzitat}
V. 21. Sie haben mich gereizet an dem, das nicht GOtt ist; mit ihrer Abgötterei haben sie mich erzürnet. Und ich will sie wieder reizen an dem, das nicht ein Volk ist; an einem närrischen Volk will ich sie erzürnen.
\end{schriftzitat}

Das heißt, wie sie mich verlassen und einen anderen Gott angenommen haben, was mir nicht zu leiden stand, so will ich sie wiederum verwerfen, und ein anderes Volk annehmen, welches sie auch sehr quälen wird, wie es geschehen ist nach dem Zeugniß des Apostels Röm. 11, 25., da die Heiden angenommen sind durch das Evangelium. Daher zürnen die Juden bis auf den heutigen Tag in unversöhnlicher Weise, daß wir sagen, sie seien nicht das Volk GOttes, und behaupten, wir seien das Volk GOttes, laut dieses Verses.

\begin{schriftzitat}
V. 22. Denn das Feuer ist angegangen durch meinen Zorn, und wird brennen bis in die unterste Hölle, und wird verzehren das Land mit seinem Gewächs, und wird anzünden die Grundfeste der Berge.
\end{schriftzitat}

„Das Feuer“ bedeutet hier nicht allein Feuer, sondern unter dem Feuer [wird bedeutet] die ganze Verwüstung, wie das Feuer des Königs zu Babylon und der Römer war, wo das Land, wiewohl nicht das ganze Land, verbrannt und verwüstet wurde, wie in Jerusalem und an anderen Orten, wo auch nicht ein Stein auf dem andern gelassen wurde. Dies bedeutet, daß er sagt, es brenne bis in die Hölle, das heißt, es wird auf der Erde nichts übrig sein, weder an Gebäuden noch an Gewächs, sondern auch der Boden selbst und die Berge, wo die Gebäude standen, werden verbrannt werden, gleichsam als wollte das Feuer von oben herab durch die Erde bis in die Hölle hinein brennen, wie wir sehen an den Orten, die durch Feuer verwüstet sind.

\begin{schriftzitat}
V. 23. Ich will alles Unglück über sie häufen, ich will alle meine Pfeile in sie schießen.
\end{schriftzitat}

Die Pfeile bedeuten hier die Plagen und das Auslassen des göttlichen Zorns. Und er selbst beschreibt nun im folgenden Verse diese Pfeile und Übel, indem er sagt:

\begin{schriftzitat}
V. 24. Vor Hunger sollen sie verschmachten, und verzehret werden vom Fieber und jähem Tode. Ich will der Thiere Zähne unter sie schicken, und der Schlangen Gift.
\end{schriftzitat}

\quellseite{819}
Hier nennt er vier oder drei Pfeile, den Hunger, das Fieber, die Pestilenz, die wilden Thiere, darnach fügt er im folgenden Verse das Schwert hinzu, und die Propheten wiederholen öfter diese Plagen. Unsere Uebersetzung [die Vulgata], die da sagt: Cum furore trahentium super terram, hat nicht beachtet, daß dies eine dichterische Umschreibung ist, wie sie in Liedern stattzuhaben pflegt, und daß „die sich im Staube entlang schleppen“ (trahentes in pulvere) die Schlangen genannt werden.

\begin{schriftzitat}
V. 25. Auswendig wird sie das Schwert berauben, und inwendig das Schrecken; beide Jünglinge und Jungfrauen, die Säuglinge mit dem grauen Manne.
\end{schriftzitat}

Das ist, es wird ihnen aller Schutz genommen, und alles, darauf sie vertrauen können.

\begin{schriftzitat}
V. 26. Ich will sagen: Wo sind sie? Ich werde ihr Gedächtnis aufheben unter den Menschen.
\end{schriftzitat}

Das heißt, nachdem alles, was sie haben, zerstört ist, und zu nichts gemacht, wird man ihrer spotten: Wo ist jetzt euer Königreich? Wo ist euer Priesterthum? Nicht als ob nichts oder keine Juden übrig wären, sondern daß sie keinen gewissen Ort, Reich oder Priesterthum haben werden, wie alle anderen Völker haben. Und „daß ihr Gedächtnis aufgehoben werde“, das heißt, daß nichts mehr an ihnen gerühmt werde, wie vorher, daß sie etwas seien oder vermöchten, nach der Weise anderer Völker, und wie sie zuvor auch waren.

\begin{schriftzitat}
V. 27. Wenn ich nicht den Zorn der Feinde scheuete, daß nicht ihre Feinde stolz würden, und möchten sagen: Unsere Macht ist hoch, und der HErr hat nicht solches alles gethan.
\end{schriftzitat}

Dies sagt er, weil er bisweilen diesen Zorn aufgehoben hat und auch die Unwürdigen, und die es nicht verdienten, befreit hat; aber dies hat er gethan um seines Namens willen, damit er nicht von den Abgöttischen verlästert würde. Er hebt also in diesem Verse das Rühmen der Juden auf und kommt demselben zuvor, da sie vielleicht sagen möchten: Wir sind nicht so böse, wie Moses singt, er singt auch nicht von uns, denn, wie die Erfahrung bezeuget, so hat GOtt uns oft wunderbarlich befreit. Diesen antwortet er: Das ist nicht durch euer Verdienst geschehen, die ihr alles, was zuvor gesagt ist, immer verdientet, sondern weil ich euch zum Volke angenommen hatte, so bin ich gezwungen, für meinen Namen Sorge zu tragen an euch, die ihr dessen überaus unwürdig seid. Aber dies werde ich nicht ewiglich thun, es wird dennoch endlich euer verdientes Ende kommen.

\begin{schriftzitat}
V. 28. Denn es ist ein Volk, da kein Rath innen ist, und ist kein Verstand in ihnen.
\end{schriftzitat}

Das heißt, es kümmert sich um nichts, weder um Zorn noch um Gnade, hat auch nicht Acht, wie ihre Sache vor mir stehe; vermessen bilden sie sich ein, ich sei ihnen gnädig, während ich erzürnt bin. Darum sind sie nicht darauf bedacht, daß sie Rath für sich schaffen und sich bessern.

\begin{schriftzitat}
V. 29. O daß sie weise wären, und vernähmen solches, daß sie verstünden, was ihnen hernach begegnen wird!
\end{schriftzitat}

Dies alles wird darum gesagt, weil sie ihren gegenwärtigen Zustand und ihre künftige Strafe nicht bedenken, gleich als ob sie sicher wären, daß sie gut wären und das Beste hoffen, während doch weitaus das Gegentheil da ist und bevorsteht.

\begin{schriftzitat}
V. 30. Wie geht es zu, daß Einer wird ihrer tausend jagen, und zween werden zehntausend flüchtig machen? Ist’s nicht also, daß sie ihr Fels verkauft hat, und der HErr hat sie übergeben?
\end{schriftzitat}

Das heißt, sowohl durch die babylonische als auch durch die römische Gefangenschaft sollten sie aus der Erfahrung erkennen, daß sie von GOtt verlassen seien, da Ein Heide tausend Juden in die Flucht treibt, während zuvor das Gegentheil geschah, daß Einer tausend Heiden in die Flucht schlug, 3 Mos. 26, 8.

\begin{schriftzitat}
V. 31. Denn unser Fels ist nicht wie ihr Fels, deß sind unsere Feinde selbst Richter.
\end{schriftzitat}

Dies verstehe so, daß gesagt werde, daß die anderen Götter und Verehrer der Götter selbst zugestanden haben, der GOtt Israels sei ein solcher GOtt, daß, wenn er es nicht dahingab, niemand sein Volk Israel überwinden konnte, wenn er es aber dahingab, sie niemals Stand halten konnten, wie Bileam deshalb dem Könige der Moabiter rieth, daß er das Volk zur Sünde verleiten sollte, und dann würde er siegen. Hier ist also Bileam einer der Zeugen,
\quellseite{820}
von denen Moses hier sagt, der herrliche Dinge von GOtt gesagt hat. Zeugen sind auch die Egypten, welche sagten, 2 Mos. 14, 25.: „Lasset uns fliehen von Israel, der HErr streitet für sie.“ Zeugen sind auch die Philister, 1 Sam. 5, 7., welche sich fürchteten vor der Bundeslade GOttes. Es haben also keine Götter ihren Völkern so beigestanden, so ihre Macht erzeigt etc. Denn es ist kein anderes Volk so hoch erhoben, wie oben Cap. 2, 25. gesagt ist.

\begin{schriftzitat}
V. 32. Denn ihr Weinstock ist des Weinstocks zu Sodom, und von dem Acker Gomorra; ihre Trauben sind Galle, sie haben bittere Beeren.
\end{schriftzitat}

Das heißt, anstatt der süßen Frucht tragen sie die überaus bitteren und argen Früchte der Abgötterei, und sind werth, daß sie so untergehen wie Sodom und Gomorra. Denn er sagt alles mit heftiger Entrüstung wider die Gottlosigkeit und Abgötterei, welche wüthet unter dem Namen GOttes und unter dem Volke GOttes.

\begin{schriftzitat}
V. 33. Ihr Wein ist Drachengift und wüthiger Ottern Galle.
\end{schriftzitat}

Das heißt, ihre Lehre und ihr Leben ist so schädlich und vergiftet, daß sie sehr viele unheilbar vergiftet, weil er vergebens durch so viele Propheten bearbeitet wird etc. Es bleibt ein vergifteter Weinstock.

\begin{schriftzitat}
V. 34. Ist solches nicht bei mir verborgen, und versiegelt in meinen Schätzen?
\end{schriftzitat}

Das heißt, alle diese Gottlosigkeit weiß ich und halte sie verborgen, wiewohl ich äußerlich gezwungen werde, diesen Gottlosen wohlzuthun, wegen meines Namens. Sie aber sehen nicht, was bei mir über sie beschlossen ist, deshalb sind sie sicher. Darum schreckt er sie mit dieser ernsten Drohung, damit sie darauf sehen, nicht, wie viel Gutes sie als Unwürdige empfangen, sondern wie viel Böses sie verdient haben, damit nicht jemand glaube, er habe um deswillen einen gnädigen GOtt, weil er ihn mit Gütern überschütte, vielmehr sich dann am meisten fürchte vor dem, was bei ihm verborgen ist, und versiegelt in seinen Schätzen.

\begin{schriftzitat}
V. 35. Die Rache ist mein, ich will vergelten. Zu seiner Zeit soll ihr Fuß gleiten; denn die Zeit ihres Unglücks ist nahe und ihr Künftiges eilet herzu.
\end{schriftzitat}

Dies sagt er, weil endlich die Larve der Guthaten aufhören wird, unter welcher sie sicher sind, aber die Strafe wird hervorkommen, welche so lange um seines Namens willen aufgeschoben ist, nämlich, daß sie ganz und gar unter den Römern zu nichts gemacht werden sollen, was sie schon längst und immer wohl verdient hatten. Denn so gleitet ihr Fuß, daß ihr Königreich nicht mehr besteht. Und so siehst du, daß Moses klarlich weissage, es werde dermaleinst das Ende dieses Volkes kommen, das heißt, seines Königreichs.

\begin{schriftzitat}
V. 36. Denn der HErr wird sein Volk richten, und über seine Knechte wird er sich erbarmen. Denn er wird ansehen, daß ihre Macht dahin ist, und beide das Verschlossene und Verlassene weg ist.
\end{schriftzitat}

Das heißt, nicht einmal seines Volkes wird er verschonen, damit sie sich nicht rühmen, daß sie sein Volk seien. Er wird aber verhöhnt werden, ohne Ansehen der Person, allein mit denen, die ihm dienen, das heißt, mit den übriggebliebenen des Volks, die an Christo besteht sind. Sonst wird ihre Macht bei allen dahin sein, so daß auch das Verschlossene und Verlassene weg ist, das heißt, nichts vom Königreiche wird übrig sein im Lande Juda, indem alle überall unter die Heiden zerstreut sind.

\begin{schriftzitat}
V. 37. Und man wird sagen: Wo sind ihre Götter? ihr Fels, darauf sie traueten?
\end{schriftzitat}

So höhnt sie sowohl das Gewissen als auch der Verfolger. Dann werden sie sehen, was sie unter ihrer Sicherheit gewesen seien und gethan haben, wer ihr Gott gewesen sei, und wem sie gedient haben, wie folgt:

\begin{schriftzitat}
V. 38. Von welcher Opfer sie Fett aßen, und tranken den Wein ihres Trankopfers. Lasset sie aufstehen und euch helfen, und euch schützen.
\end{schriftzitat}

Während sie doch meinten, daß sie zu der Zeit von den Opfern des HErrn äßen, aber jetzt sehen sie zu spät ihre Abgötterei (vanitatem).

\begin{schriftzitat}
V. 39. Sehet ihr nun, daß ich’s allein bin, und ist kein GOtt neben mir? Ich kann tödten und lebendig machen, ich kann schlagen und kann heilen, und ist niemand, der aus meiner Hand errettet.
\end{schriftzitat}

So lehrt die Erfahrung, was sie zuvor unter der Sicherheit nicht gelernt haben.

\quellseite{821}
\begin{schriftzitat}
V. 40. Denn ich will meine Hand in den Himmel heben, und will sagen: Ich lebe ewiglich.
\end{schriftzitat}

Dies wird gesagt in der Weise, die ein Schwörender hat, oder jemand, der sich rühmt, daß außer GOtt sich niemand (nihil) rühmen könne weder von seinem Sein (esse) noch von seinem Leben (vivere), sondern er allein sagt: Ich bin, ich lebe; nämlich, damit er gefürchtet und unsere Vermessenheit vernichtet werde.

\begin{schriftzitat}
V. 41. Wenn ich den Blitz meines Schwerds wetzen werde, und meine Hand zur Strafe greifen wird; so will ich mich wieder rächen an meinen Feinden, und denen, die mich hassen, vergelten.
\end{schriftzitat}

Das heißt, schrecklich ist meine Rache, weil sie durchdringt (und nicht aufgehalten werden kann, wie der Blitz) durch meine Feinde, mögen sie nun Juden sein oder Heiden.

\begin{schriftzitat}
V. 42. Ich will meine Pfeile mit Blut trunken machen, und mein Schwert soll Fleisch fressen, über dem Blut der Erschlagenen und über dem Gefängniß, und über dem entblößten Haupt des Feindes.
\end{schriftzitat}

Diese Rache will ich vergelten (sagt er), daß meine Pfeile voll Bluts sein sollen, und mein Schwert soll nicht verwunden, sondern die Leiber fressen: so viele werden getödtet, so viele gefangen worden sein. Sodann wird auch das Haupt der Feinde bloß sein, das heißt, sie werden des Königreichs beraubt werden, daß sie ohne die Herrlichkeit des Königreichs durch die ganze Welt zerstreut sein sollen.

\begin{schriftzitat}
V. 43. Jauchzet alle, die ihr sein Volk seid.
\end{schriftzitat}

Das heißt, mit den gläubigen Uebrigegebliebenen werden die Heiden im Glauben verbunden werden.

\begin{quote}
Denn er wird das Blut seiner Knechte rächen.
\end{quote}

Nämlich an denen, welche die Propheten und Apostel tödten.

\begin{quote}
Und wird gnädig sein dem Lande seines Volks.
\end{quote}

Nämlich [des Volks], welches übrig ist von Israel, das sich bekehrt hat; diesem wird er ein neues Reich bereiten unter Christo, Amen.

\begin{schriftzitat}
V. 48. Und der HErr redete mit Mose etc.
\end{schriftzitat}

Schließlich beschreibt er das Gebot des HErrn vom Tode Mosis, wie er ihm das ganze Land Canaan gezeigt habe, aber ihm nicht gestattete hineinzugehen, wegen der Sünde des Unglaubens, die er begangen hatte (admissum) am Hadermasser. Das Gebirge Abarim ist dasselbe als das Gebirge Arabien's, auf dessen Spitze der Berg Nebo ist, auf der Spitze [des Gebirges] Pisga [5 Mos. 34, 1.], auf welchen Moses hinaufstieg und starb.

\section{Das dreiunddreißigste Capitel}

\begin{schriftzitat}
V. 1. 2. Dies ist der Segen, damit Mose, der Mann GOttes, die Kinder Israel vor seinem Tode segnete, und sprach: Der HErr ist von Sinai kommen, und ist ihnen aufgegangen von Seir; er ist hervorgebrochen von dem Berge Paran, und ist kommen mit viel tausend Heiligen.
\end{schriftzitat}

Dies verstehe ich so, daß es die herrliche und prächtige (insignem) Ankunft des HErrn im Lande der Verheißung bedeute. Denn dahin ist er gekommen, indem er anfing von dem Berge Sinai, wo er den Bund des Gesetzes gab, und ist aufgegangen von Seir, das heißt, er ging durch das Land Seir bis nach Canaan, und so kam er gleichwie die aufgehende Sonne, um mit einem neuen Reiche jenes Land zu erleuchten. Es ist dasselbe, daß er hervorgebrochen ist von dem Berge Paran, weil er dort durchzog in den Grenzen Edom. Er kam aber mit viel tausend Heiligen, das heißt, mit dem Volke Israel, deren zwar viele Tausende waren, wiewohl nicht jeder einzelne heilig; dennoch waren tausende von Heiligen da, weil viele unter ihnen heilig waren.

\begin{quote}
Zu seiner rechten Hand ist ein feuriges Gesetz an sie.
\end{quote}

Ob du dies nun in der That von dem Gesetze Mosis verstehen willst, mit dem er sie in seiner Rechten regiert, daß es im Geist brenne
\quellseite{822}
und geliebt werde, nämlich von dem Heiligen, oder [ob du es verstehen willst] von der Feuersäule, mit der er sie führte, darüber streite ich nicht, wiewohl mir der erstere Verstand besser gefällt, weil es eine außerordentliche Wohlthat war, daß diesem Volke das Gesetz, das Wort GOttes, anvertraut wurde.

\begin{schriftzitat}
V. 3. Wie hat er die Leute so lieb!
\end{schriftzitat}

Nämlich, da er so bei ihnen wohnt, mit ihnen wandert und gleichsam vertraulichen Umgang hat in diesem ganzen Zug durch die Wüste etc.

\begin{quote}
Alle seine Heiligen sind in seiner Hand.
\end{quote}

Das heißt „seine Heiligen“, nämlich des Volkes oder in dem Volke; „sind in seiner Hand“, das heißt, du vertheidigst sie, denn so liebst du die Leute.

\begin{quote}
Sie werden sich setzen zu deinen Füßen, und werden lernen von deinen Worten.
\end{quote}

Das heißt, auch dies ist ein Unterpfand seiner Liebe, daß er das Amt des Wortes in diesem Volke eingesetzt hat, wo die Zuhörer zu seinen Füßen sitzen und sein Wort hören können, wie in der Synagoge und im Tempel oder in der Stiftshütte.

\begin{schriftzitat}
V. 4. 5. Mose hat uns das Gesetz geboten, dem Erbe der Gemeine Jakob etc.
\end{schriftzitat}

Hier redet Moses von sich, was er für ein Mann gewesen sei nach seinem Amte. Erstens, daß er selbst das Gesetz gebracht habe, von dessen Amte er schon geredet hat. Zweitens [V. 5.], daß er in der vollen Macht (plenitudine) eines Königs war, wenn die Fürsten des Volkes mit den Stämmen Israels versammelt waren. Das heißt, er war anstatt eines Königs und aus GOttes Macht und an der Statt GOttes unter ihnen, so daß er, über das geschriebene Gesetz hinaus, sie auch durch Aussprüche (oraculo) der lebendigen Stimme regierte aus dem Munde GOttes.

\begin{schriftzitat}
V. 6. Ruben lebe und sterbe nicht, und sein Pöbel sei gering.
\end{schriftzitat}

Derselbe Segen ist es, den Jakob gibt 1 Mos. 49, 4. [Vulg.]: „Du sollst nicht wachsen“, nämlich, daß er nicht ganz und gar untergehen soll wegen der Sünde, die er an seinem Vater begangen hat.

\begin{schriftzitat}
V. 7. Dies ist der Segen Juda: HErr, erhöre die Stimme Juda.
\end{schriftzitat}

Dies wird verstanden von den Königen Juda, die oft den Sieg erlangt haben, mehr durch den Glauben und Gebet als durch das Schwert.

\begin{quote}
Und bringe ihn zu seinem Volke [Vulg.].
\end{quote}

Das heißt, zur Herrschaft über Israel und Juda.

\begin{quote}
Laß seine Macht groß werden, und ihm müsse wider seine Feinde geholfen werden.
\end{quote}

Das heißt, er möge stark werden im Glauben und befreit werden von Feinden; denn dieses Königreich hat viele Feinde erleiden müssen.

\begin{schriftzitat}
V. 8. Und zu Levi sprach er:
\end{schriftzitat}

\begin{quote}
Dein Recht und dein Licht bleibe bei deinem heiligen Mann, den du versucht hast zu Massa, da ihr hadertet am Haderwasser.
\end{quote}

Er will sagen: Das Priesterthum und dein Amt sei glücklich, sowohl mit Beten als Lehren, wie es mit Mose in der Wüste war, den du versucht hast. Denn „Recht und Licht“ ist das, was Moses in das Amtsschildlein (hosen) legte, 2 Mos. 28, 29. f., wodurch man GOttes Antwort (oraculum) empfing in Sachen, die man vornehmen wollte.

\begin{schriftzitat}
V. 9. 10.
\end{schriftzitat}

\begin{quote}
Wer zu seinem Vater und zu seiner Mutter spricht: Ich sehe ihn nicht; und zu seinem Bruder: Ich kenne ihn nicht; und zu seinem Sohn: Ich weiß nicht; die halten deine Rede und bewahren deinen Bund. Die werden Jakob deine Rechte lehren und den Israel dein Gesetz; die werden Räuchwerk vor deine Nase legen, und ganze Opfer auf deinen Altar.
\end{quote}

Das heißt, es werden viele Heuchler sein, sowohl mit Werken als auch mit Worten, weil ihr Herz nicht rechtschaffen ist. Die aber zum Theil selbst der Welt abgestorben sind, und allein Gotte leben und ihm dienen, die thun und lehren recht. Solche Leute fordert das Gesetz, deren Räuchwerk und Opfer ist angenehm. Er zeigt also an, wie die Diener des Wortes beschaffen sein sollen, nämlich der Welt abgestorben. Dies wünscht Moses, indem er sagt:

\begin{schriftzitat}
V. 11. HErr, segne sein Vermögen, und laß dir gefallen die Werke seiner Hände.
\end{schriftzitat}

\quellseite{823}
Das heißt, was er zu schaffen hat im Dienste des Worts, das laß glücklich hinausgehen, daß es dir gefalle und Frucht schaffe bei dem Volke.

Zerschlage den Rücken derer, die sich wider ihn auflehnen, und derer, die ihn hassen, daß sie nicht aufkommen.

Das heißt, die falschen Propheten und Könige und Fürsten und Heuchler, welche seinem Vermögen mit ihren Abgöttereien (vanitatibus) sich entgegenstellen, laß kein Gedeihen haben.

\begin{schriftzitat}
V. 12. Benjamin, der Geliebte des HErrn, wird sicher bei Ihm wohnen; allezeit wird Er ihn schützen, und er wird zwischen seinen Schultern wohnen.
\end{schriftzitat}

Vielleicht sagt Moses dies, weil der Stamm Benjamin zwischen den Stämmen Juda und Ephraim wohnte, die beide königlich und mächtig waren, und den Tempel in seinem Stamme hatte.

\begin{schriftzitat}
V. 13. Und zu Joseph sprach er: Sein Land liegt im Segen des HErrn. Da sind edle Früchte vom Himmel, vom Thau, und von der Tiefe, die unten liegt.
\end{schriftzitat}

Das ist, daß er einen gnädigen Himmel habe, wodurch ihm alles überaus reichlich zuwächst, so daß Himmel, Erde, Meere und alles ihm heilsam ist.

\begin{schriftzitat}
V. 14. Da sind edle Früchte von der Sonne, und edle reife Früchte der Monden.
\end{schriftzitat}

Das heißt, daß Sonne und Mond günstig sind für die Früchte der Erde.

\begin{schriftzitat}
V. 15. Und von den hohen Bergen gegen Morgen und von den Hügeln für und für.
\end{schriftzitat}

Das heißt, daß sowohl die ersten als auch die letzten Berge, die vorn und hinten gelegen sind, alle Ueberfluß haben, wiewohl nach meiner Meinung dies alles bildlich geredet zu sein scheint, so daß du verstehen mußt, daß das Reich Israel in Blüthe stehe, da der Himmel günstig ist, und der Thau, und die Tiefe, die Sonne, der Mond, die alten Berge, die Hügel der Welt, das heißt, da ihm GOtt und die Menschen alles geben, Könige, Fürsten, die Großen, die Reichen, die Völker und das gemeine Volk; nämlich diese verschiedenen Stände zeigt er durch diese verschiedenen Namen an. Denn von den Früchten des Landes sagt er alsbald [V. 16.]: „Und edle Früchte von der Erde und was drinnen ist“, das heißt, alles, was die Erde hervorbringt, und was Thiere und Menschen zeugen.

\begin{schriftzitat}
V. 16. Die Gnade deß, der in dem Busch wohnte, komme auf das Haupt Joseph, und auf den Scheitel des Nasir unter seinen Brüdern.
\end{schriftzitat}

Das heißt, die Gnade Christi sei über ihnen, daß ein Ueberfluß da sei an Propheten. Das ist auch geschehen, denn das Reich Israel hatte viele Propheten.

\begin{schriftzitat}
V. 17. Seine Herrlichkeit ist wie ein erstgeborner Ochse.
\end{schriftzitat}

Das heißt, es ist ein ansehnliches Reich an Macht und Männern.

Und seine Hörner sind wie Einhörners Hörner.

Das heißt, Könige und Fürsten und Herzoge, die überaus stark sind zum Kriege.

Mit denselben wird er die Völker stoßen zu Hauf, bis an des Landes Ende.

Nämlich des Landes Israel, indem er sie austreibt.

\begin{schriftzitat}
V. 18. Das sind die Tausende Ephraim und die Tausende Manasse. Sebulon, freue dich deines Auszugs; aber Jssachar, freue dich deiner Hütten.
\end{schriftzitat}

Er bedient sich entgegengesetzter Anspielungen. Denn Sebulon hat seinen Namen vom Wohnen und Bleiben, und von dem sagt er, daß er sich seines Auszugs freue. Wiederum Jssachar hat seinen Namen vom Löhne, der durch das Ausziehen gemehret wird, und von diesem sagt er, daß er sich freue, zu Hause zu bleiben. Er will also sagen: Glückselig seiest du, Sebulon, wenn du ausziehest; glückselig seiest du, Jssachar, wenn du nicht ausziehest, sondern bleibst.

\begin{schriftzitat}
V. 19. Sie werden die Völker auf den Berg rufen, und daselbst opfern Opfer der Gerechtigkeit.
\end{schriftzitat}

Dieses scheint mir zu der Historie der Debora zu gehören, wo Barak, nachdem ein Heer auf dem Berge Thabor gesammelt worden war, den Sieg über Sisera davongetragen hat [Richt. 4, 4. ff.].

Denn sie werden die Menge des Meeres saugen, und die versenkten Schätze im Sande.

\quellseite{824}
Das heißt, sie werden Ueberfluß haben an Waaren, die das Meer liefert, und an Metallen des Landes.

\begin{schriftzitat}
V. 20. Und zu Gad sprach er: Gad sei gesegnet, der Raummacher; er liegt wie ein Löwe und raubet den Arm und die Scheitel.
\end{schriftzitat}

Dies scheint mir der König Jehu erfüllt zu haben, welcher Raum machte für die Gottesfurcht, nachdem die Baalsdiener vertilgt waren, und durchbrach wie ein Löwe, nachdem in tapferer Weise auch zwei Könige geschlagen waren, gleichsam Arm und Scheitel [2 Kön., Cap. 9 und 10].

\begin{schriftzitat}
V. 21. Und er sahe, daß ihm ein Haupt gegeben war, ein Lehrer, der verborgen ist.
\end{schriftzitat}

Ich glaube, daß dies von den guten Propheten gesagt sei, welche unter Jehu in seinem Reiche sicher waren, welche vorher unter den Baalsdienern verfolgt waren, so daß der Singular „ein Lehrer“ für den Plural steht. „Ein Lehrer“ (portio doctoris) steht also für viele Lehrer (collectio doctorum).

Welcher kam mit den Obersten des Volks und verschaffte die Gerechtigkeit des HErrn, und seine Rechte an Israel.

Denn er führte die Obersten und das Volk von jenen Abgöttereien ab, daher ist er auch von GOtt gelobt worden.

\begin{schriftzitat}
V. 22. Und zu Dan sprach er: Dan ein junger Löwe.
\end{schriftzitat}

Denn sie haben tapfer gekämpft wider Laïs oder Lesem [Richt. 18, 27.].

Er wird fließen von Basan.

Das heißt, sie werden grünen von der Stadt Dan bis gen Basan.

\begin{schriftzitat}
V. 23. Und zu Naphtali sprach er: Naphtali wird genug haben, was er begehret, und wird voll Segens des HErrn sein.
\end{schriftzitat}

Denn das Land seines Stammes ist fruchtbar, wie in Galiläa.

Gegen Abend und Mittag wird sein Besitz sein.

„Mittag“ nämlich des galiläischen Meeres.

\begin{schriftzitat}
V. 24. Und zu Asser sprach er: Asser sei gesegnet mit Söhnen.
\end{schriftzitat}

Nämlich, er habe viele Nachkommenschaft.

Er sei angenehm seinen Brüdern.

Das heißt, er habe viel Gunst.

Und tunke seinen Fuß in Oel.

Das heißt, er wird viel Oel haben.

\begin{schriftzitat}
V. 25. Eisen und Erz sei an seinen Schuhen.
\end{schriftzitat}

Das ist,\footnote{Statt ideo in der Erlanger ist mit der Wittenberger und der Jenaer id est zu lesen.} er habe viel Erz und Eisen.

Dein Alter sei wie deine Jugend.

Das heißt, er möge hervorragen durch Kräfte des Leibes.

Du siehst, daß hier der Stamm Simeon ausgelassen werde und nicht gesegnet. Aus welcher Ursache, wer kann das wissen? Denn daß die Juden vorgeben, es sei geschehen wegen des Verkaufens Josephs, und wegen der Hurerei mit BaalPeor, daran liegt mir nichts. Hieronymus hält dafür, daß Judas Ischarioth, der aus diesem Stamme kommen würde, dies verdient habe; aber man lasse Fabeln Fabeln bleiben. Aber es ist gewiß, daß Simeon kein eigenes Loos hatte, wie die anderen Stämme, sondern er hat es gehabt im Stamme Juda [Jos. 19, 1.]. Darum muß er sich begnügen lassen an dem allgemeinen Segen, welcher folgt.

\begin{schriftzitat}
V. 26. Es ist kein GOtt, als der GOtt des Gerechten. Der im Himmel sitzt, der sei deine Hülfe, und deß Herrlichkeit in Wolken ist.
\end{schriftzitat}

Das heißt, der seine Majestät zeigt durch den Donner.

\begin{schriftzitat}
V. 27. Das ist die Wohnung GOttes von Anfang, und unten in der Welt sind seine Arme.
\end{schriftzitat}

Das heißt, wiewohl er im Himmel wohnt, so herrschen doch seine Arme in der Welt, daß er die Seinen befreie, wie folgt.

Er wird vor dir her deinen Feind austreiben und sagen: Sei vertilget.

Nämlich allein durch das Wort wird er ihn niederwerfen.

\begin{schriftzitat}
V. 28. Israel wird sicher alleine wohnen, und das Auge Jakobs wird sein auf dem Lande, da Korn und Most ist.
\end{schriftzitat}

\quellseite{825}
Das heißt, er wird sich ergötzen in einem fruchtbaren Lande.

Dazu sein Himmel wird von Thau triefen. Nämlich, damit er das Land fruchtbar mache.

\begin{schriftzitat}
V. 29. Wohl dir, Israel, wer ist dir gleich? O Volk, daß du durch den HErrn selig wirst.
\end{schriftzitat}

Nicht durch Menschenkräfte, wie andere Völker.

Der deiner Hülfe Schild und das Schwert deines Sieges ist. Deinen Feinden wird’s fehlen; aber du wirst auf ihrer Höhe einhertreten.

Als wollte er sagen: Mit deinen Kräften wärest du deinen Feinden nicht gewachsen, aber der ist mächtig, der dein Schwert ist.

Nun vergleiche das Lied Mosis mit diesem Segen, und du wirst entgegengesetzte Dinge sehen. Denn dort verdammt er das ganze Volk, hier segnet er jeden Einzelnen, aber um der Verheißung willen, die den Vätern geschehen ist. So wird wegen weniger Guten dem ganzen Volke der Segen gegeben, und dennoch wird schließlich um der Bösen willen das Volk mit dem Königreiche zerstört, nachdem die Guten von ihnen abgesondert sind.

\section{Das vierunddreißigste Capitel}

\begin{schriftzitat}
V. 1. Und Mose ging von dem Gefilde der Moabiter auf den Berg Nebo etc.
\end{schriftzitat}

Dies Capitel hat Moses nicht geschrieben, sondern Josua oder Eleasar, es sei denn, du wollest sagen, er habe seinen Tod, da er ihn ja vorherwußte, auf diese Weise beschrieben. Daß aber kein Mensch sein Grab habe finden können [V. 6.], davon sagt man mit Recht, daß es um deswillen geschehen sei, damit die Juden daraus nicht einen Abgott machten. Aber es bedeutet mehr, daß, da das Licht des Evangelii gekommen ist, das Amt des Gesetzes so abgethan werde, daß die Juden es weder finden noch jemals wieder herstellen können.

Daß aber geschrieben wird, daß Moses, da er im hundertundzwanzigsten Jahre seines Alters starb, so kräftig gewesen sei, daß [V. 7.] seine Augen nicht dunkel worden und sein Aussehen (forma faciei) nicht verfallen war, das wird als ein großes Wunder erzählt, da von den übrigen Vätern, als Isaac, Jakob, Eli und anderen geschrieben wird, daß sie wegen ihres hohen Alters dunkle Augen gehabt hätten, und noch heutzutage diejenigen, welche im Alter etwas vorgerückt sind, abnehmen an Augen und Aussehen. Darum\footnote{Statt idem in den Ausgaben haben wir ideo angenommen.} ist es ein Wunder, daß die Gestalt seines Angesichts so blühend war, daß es ihm weder an Farbe noch an Schönheit (species) mangelte, da es die Natur doch mit sich bringt, daß Greise ungestalt werden durch Runzeln und bleiches Aussehen. An Moses ist dies aber noch viel mehr ein Wunder, da er ja in seinem ganzen Leben mit so vielen Mühsalen, Gefahren und Geschäften überhäuft war, daß es zu verwundern ist, daß er nicht vor anderen erschöpft und ungestalt gewesen, oder eher gestorben ist. Unsere Uebersezung [die Vulgata] hat: „und seine Zähne sind nicht bewegt worden“, nämlich getäuscht durch die Aehnlichkeit des Ausdrucks, da Frischheit (vigor [\hebraeisch{לֵחַ}]) und Kinnbacken (maxilla [\hebraeisch{לֶחִי}]) im Hebräischen fast mit denselben Buchstaben geschrieben werden.

Es bedeutet aber die Dauerhaftigkeit (perseverantia) des Sehens und der Kraft bei Moses nach der heimlichen Deutung, daß die Kraft des Gesetzes durch keine Länge der Zeit, durch keine Menge der Werke abnimmt, sondern die schuldigen Gewissen immerdar drängt und in Unruhe erhält, bis daß es stirbt, das heißt, bis es durch Christum abgethan wird, durch Einsetzung des neuen Amts der Gnade. Es war aber, als Christus ankam, das alte Gesetz auch seinem Ende nahe, nach dem Worte in dem Briefe an die Hebräer [Cap. 8, 13.]: „Was aber alt und überjahret ist, das ist nahe bei seinem Ende.“

Schließlich wird Moses [V. 10. ff.] gepriesen wegen seiner Größe, daß nach ihm kein Prophet aufgestanden sei, der ihm gleich gewesen wäre in so großen Zeichen, die durch ihn in
